\documentclass[11pt]{article}

\usepackage[margin=1in]{geometry}
\usepackage{amsmath,amssymb,mathtools}
\usepackage{booktabs,longtable,array}
\usepackage{textcomp}
\usepackage[hidelinks]{hyperref}
\numberwithin{equation}{section}

\newcommand{\NS}{\textsc{Newspeak}}
\newcommand{\kw}[1]{\ensuremath{\mathsf{#1}}}
\newcommand{\sel}{s}
\newcommand{\undef}{\mathord{\uparrow}}
\newcommand{\nil}{\kw{nil}}
\newcommand{\Top}{\kw{Top}}
\newcommand{\Object}{\kw{Object}}
\newcommand{\Public}{\kw{public}}
\newcommand{\Protected}{\kw{protected}}
\newcommand{\Private}{\kw{private}}
\newcommand{\dom}{\mathop{\mathrm{dom}}}
\newcommand{\classof}{\mathop{\mathrm{classOf}}}
\newcommand{\super}{\mathop{\mathrm{super}}}
\newcommand{\mixin}{\mathop{\mathrm{mixin}}}
\newcommand{\decl}{\mathop{\mathrm{decl}}}
\newcommand{\access}{\mathop{\mathrm{access}}}
\newcommand{\selector}{\mathop{\mathrm{selector}}}
\newcommand{\methods}{\mathop{\mathrm{methods}}}
\newcommand{\lookup}{\mathop{\mathrm{lookup}}}
\newcommand{\definingClass}{\mathop{\mathrm{definingClass}}}
\newcommand{\kws}[1]{\ensuremath{\mathsf{#1\mathord{:}}}}
\newcommand{\selsym}[1]{\ensuremath{\mathord{\#}\text{\texttt{#1}}}}
\newcommand{\lit}[1]{\ensuremath{\mathsf{lit}(#1)}}
\newcommand{\dnu}{\kws{doesNotUnderstand}}
\newcommand{\PastFuture}{\ensuremath{\text{\texttt{Past\textasciigrave Future}}}}
\newcommand{\rulelabel}[1]{\textsc{#1}}

\title{A Formal Dynamic Semantics for Newspeak\\\large Design draft 0.14: Machine-checked companion}
\author{}
\date{\today}

\begin{document}
\maketitle

\begin{abstract}
This document develops a complete dynamic semantics for \NS.  ``Complete''
means that surface conveniences, nested classes, mixin application, all access
levels and send forms, initialization, closures and non-local return, actors,
eventual references, and causally connected reflective program change must
either have rules or appear as explicit outstanding obligations.  This
draft fixes the semantic domains and the access-controlled method-dispatch
kernel, the complete sequential machine, its lifting to isolated actors with
asynchronous sends and E-ordered delivery, and the mirror-capability layer for
transactional program, object, activation, continuation, and complete actor-
stack change.  Reified per-frame control, paused actors, atomic stack
installation, and ordinary-VM resumption provide the machinery required by an
integrated debugger.  The reflection rules make VM-cohort propagation and the ordering of installs with
active turns and in-flight messages explicit.  An optional final layer makes
garbage collection, \kw{WeakArray}, \kw{WeakMap}, and their observable
interaction with exact-instance and mixin-application enumeration explicit
without changing the preceding transition rules.  A source-front-end layer
now defines PEG recognition, AST construction, stable source identities,
static elaboration, and a metacircular compilation pipeline whose parser,
elaborator, compiler, and Simulator are ordinary reflectively changeable
Newspeak objects.  A Lean companion formalizes every layer, proves the
determinism, preservation, equivalence, fixed-point, transaction, and
refinement obligations summarized in the final section, and contains no
admitted theorem or additional axiom.  The remaining compiler/Simulator work
is implementation verification: a concrete installed pair must supply the
mandatory artifact-admission and frame-representation certificate already
specified here.
This draft intentionally does not replace the language specification yet.
\end{abstract}

\clearpage
\begingroup
\footnotesize
\tableofcontents
\endgroup
\clearpage

\section{Authority, scope, and method}

The normative source is working version 0.109 of the \NS{} Programming
Language Draft Specification (henceforth \emph{the Specification})
\cite{newspeak-spec-0.109}.
Section 3.3 of \emph{Modules as Objects in Newspeak} \cite{bracha10a} is
explanatory prior art.  If the sources disagree, the Specification governs.

The surface language is parsed by the PEG in Specification \S4--\S7.
Section~\ref{sec:front-end} gives that PEG a recognition semantics, maps its
concrete trees to stable AST identities, and defines the static elaboration
judgment that rejects ill-formed uses of \kw{self}, \kw{outer}, \kw{super},
returns, and class names.

The semantics is a small-step transition system.  This choice makes evaluation
order, abrupt control transfer, actors, and transactions visible without
encoding them as host-language accidents.  Metalevel records below are not a
claim that runtime entities are inaccessible: Section~\ref{sec:reflection}
reifies them through mirror capabilities and preserves that causal connection.

\paragraph{Base/reflective separation.}
The base relation holds the program definition $P$ fixed.  In a machine
configuration, $P$ is the static program, $H$ is the heap, and $\mathcal A$ is
the activation stack.  The latter includes the active activations and their
continuation links.  At the sequential boundary a reflective step has the
schematic form
\begin{equation}\label{eq:reflective-step}
  \langle P,H,\mathcal A\rangle
    \xrightarrow{\;\mathsf{reflectiveStep}(\Delta)\;}
  \langle P',H',\mathcal A'\rangle,
\end{equation}
where validation and commit are atomic; Section~\ref{sec:reflection} refines
this schema for the actor/VM state.  This separation is methodological,
not semantic simplification: the base domains already carry stable declaration,
mixin, class, object, and activation identities needed by reflection.

\section{Surface and core phrases}
\label{sec:surface}

Let $x$ range over identifiers, $s$ over selector symbols, and $\bar e$ over
ordered expression lists.  The dispatch-relevant core expressions are
\begin{equation}\label{eq:core-expressions}
\begin{array}{rcl}
 e &::=& v \mid \kw{self} \mid \kw{atom}(\omega)
          \mid \kw{tuple}(\bar e) \mid \kw{pattern}(p,d) \\
   &\mid& \kw{objectLiteral}(\chi_o,e_S)
          \mid \kw{closureLiteral}(q,\bar x,\ell,b) \\
   &\mid& \kw{send}(e,s,\bar e)
          \mid \kw{implicit}(s,\bar e,d) \\
   &\mid& \kw{selfSend}(s,\bar e,D) \mid
          \kw{outerSend}(s,\bar e,D_t,D_0) \\
   &\mid& \kw{superSend}(s,\bar e) \mid
          \kw{asyncSend}(e,s,\bar e) \mid \cdots
\end{array}
\end{equation}
There is deliberately no variable-reference case $x$.  A standalone identifier
$x$ is a unary message clause with an implicit receiver and elaborates to
$\kw{implicit}(\#x,\epsilon,d)$, where $d$ is the lexical declaration selected
for that send.  This remains true when $d$ declares a parameter or local slot;
the resulting send is directed to the appropriate activation.  The
metavariable $x$ is retained only for binding occurrences such as formal
parameters and slot declarations.

Here $d$ records the innermost lexical declaration selected for an implicit
send, if any; $D_0$ is the immediately enclosing class declaration and $D_t$
is the named lexical target of an outer send.  These annotations are computed
from the AST, not guessed dynamically.  Retaining distinct send nodes is
essential: ordinary, self, outer, and super sends grant different authority.

A super send deliberately carries no declaration parameter.  Its receiver is
the current instance, and its lookup start is the superclass of the current
activation's runtime current class.  That current class is installed by method
dispatch and captured by closures.  A lexical declaration identity would not
identify the particular runtime application of its mixin from which super
lookup must continue; using it for that purpose would be incorrect when one
declaration has multiple applications.

Elaboration may remove syntax whose specified expansion preserves all
observable behavior.  For example, a cascade evaluates by the specified
closure expansion, which evaluates its receiver exactly once.  Elaboration
must not turn a self send into an ordinary send, including when \kw{self} is
parenthesized.

Slot declarations and nested-class declarations induce methods before runtime
evaluation:
\begin{itemize}
\item an immutable slot induces a getter with the slot's access;
\item a mutable slot induces a getter and a one-argument setter, both with the
      slot's access; the setter returns its receiver;
\item a lazy slot induces a getter with the slot's access and is physically
      initialized to $\nil$.  On a read, the getter returns the stored value if
      it is non-$\nil$; otherwise it evaluates the slot initializer, writes the
      result directly into the slot, and returns it.  The $\kw{::=}$ form also
      induces the same-access setter described above, whereas the $\kw{=}$ form
      induces no setter;
\item a nested class induces a memoizing getter with the declaration's access.
\end{itemize}
Initialization writes slots directly and therefore does not invoke or dispatch
through these synthesized setters.  Because $\nil$ is the lazy-state sentinel,
an initializer that returns $\nil$ is evaluated again on the next read.
Resetting a lazy slot to $\nil$ for orthogonal synchronization will be modeled
as reflective change rather than ordinary initialization or mutation.

\section{Semantic domains}

We assume pairwise disjoint, countably infinite metalevel identity sets
\begin{equation}\label{eq:runtime-identities}
\begin{aligned}
  o&\in\mathit{ObjId}, & C&\in\mathit{ClassId}, &
  M&\in\mathit{MixinId}, & a&\in\mathit{ActId},\\
  c&\in\mathit{ClosureId}, & g&\in\mathit{MirrorId}, &
  A&\in\mathit{ActorId}. &&
\end{aligned}
\end{equation}
Their tagged union is $\mathit{ObjRef}$.  Every member of $\mathit{ObjRef}$ is
a Newspeak object and, when live, has a class.  The disjoint tags prevent a
semantic rule from confusing (for example) the identity of a class object with
the identity of one of its ordinary instances; they do not introduce
non-object runtime values.
Pause tokens $r_s\in\mathit{PauseToken}$ are drawn from a separate countably
infinite set.  They are unforgeable protocol tokens rather than Newspeak
objects: entering a paused actor state allocates a fresh token, and every
successful paused-stack edit allocates another.  A stale debugger request
therefore cannot mutate a stack that has since been edited or resumed.
Continuation links are not a separate identity domain: throughout,
$k\in\mathit{ActId}\uplus\{\nil\}$.  In accordance with the Specification, a
non-$\nil$ continuation object is the sender activation at which execution
will resume.  Mirror identities are included in $\mathit{ObjRef}$; their
unforgeable heap records are introduced in Section~\ref{sec:reflection}.
Source declarations also have stable AST identities.  We write
$d\in\mathit{DeclId}$ for an arbitrary declaration,
$K\in\mathit{ClassBodyDeclId}\subseteq\mathit{DeclId}$ for a declaration that
induces a mixin, with the disjoint partition
$\mathit{ClassBodyDeclId}=\mathit{ClassDeclId}\uplus\mathit{ObjLitDeclId}$.
Write $D\in\mathit{ClassDeclId}$ for a named class declaration,
$L\in\mathit{ObjLitDeclId}$ for an object-literal declaration, and
$q\in\mathit{ActDeclId}\subseteq\mathit{DeclId}$ for a method or closure
declaration.  The domains $\mathit{ClassBodyDeclId}$ and
$\mathit{ActDeclId}$ are disjoint.  Finally, $z\in\mathit{SiteId}$ is an
expression site and
$i_f\in\mathit{MethodId}\subseteq\mathit{ActDeclId}$ is a method declaration.
A class-declaration identity is not a class name: evaluation of one declaration
may produce many distinct runtime classes.
Class creation provenance is the tagged domain
\begin{equation}\label{eq:class-origin}
 \chi_C\in\mathit{ClassOrigin}::=
   \kw{expression}(K)\mid\kw{actor}(K).
\end{equation}
The tag distinguishes evaluation of a source expression from application of a
top-level mixin while seeding a new actor.  Both retain the inducing source
declaration for reflection.

An access level is $\alpha\in\{\Public,\Protected,\Private\}$.  Absence of an
access modifier elaborates to $\Protected$; a top-level class declaration is
always public.

A message is $\mu=\langle s,\bar o\rangle$.  A method definition is
\begin{equation}\label{eq:method-definition}
 f=\langle i_f,s,\alpha,\bar x,\ell,b,K\rangle,
\end{equation}
where $\ell$ are local-slot declarations, $b$ is the body, and $K$ is the
lexically enclosing class or object-literal declaration.  The identity $i_f$ names the method
declaration independently of its selector and current body.  This lets
$P.\mathit{body}$ and, later, method mirrors refer to the same declaration when
its body is reflectively replaced.  The reflective rules will specify when an
edit preserves $i_f$ and when deletion followed by addition allocates a fresh
identity.  Elaboration assigns each synthesized getter, setter, or nested-class
getter a method identity derived from its inducing declaration and selector.
Each mixin contains at most one direct method per selector after synthesized
members have been installed.

The program definition $P$ contains immutable source structure:
\begin{equation}\label{eq:program-definition}
\begin{array}{rcl}
 P.\mathit{source} &: & \mathit{DeclId}\rightharpoonup\mathit{AST},\\
 P.\mathit{unit} &: & \mathit{AST},\\
 P.\mathit{declMixin} &: & \mathit{ClassBodyDeclId}\to\mathit{MixinId},\\
 P.\mathit{mixin} &: & M \rightharpoonup
   \langle K,\mathit{methodMap},\mathit{slotDecls},
                 \mathit{nestedDecls},\mathit{initializer}\rangle,\\
 P.\mathit{lexParent} &: & D \rightharpoonup D,\\
P.\mathit{body} &: & \mathit{MethodId}\to\mathit{AST},\\
P.\mathit{scopeChain} &: & \mathit{SiteId}\to\mathit{DeclId}^{*},\\
P.\mathit{activationMembers} &: & \mathit{ActDeclId}\to
  \mathcal P_{\mathrm{fin}}(\mathit{Selector}),\\
P.\mathit{owner} &: & \mathit{ClassBodyDeclId}\to
 \left\{\begin{gathered}\kw{activationOwner},\kw{classOwner},\\[-1mm]
   \kw{objectLiteralOwner},\kw{topOwner}\end{gathered}\right\},\\
P.\mathit{actorSeed} &: & \mathit{MixinId}\rightharpoonup
 \langle\mathit{ClassBodyDeclId},\mathit{MixinId},
        \mathit{FactoryDescriptor}\rangle,\\
P.\mathit{atom} &: & \mathit{Atom}\to\mathit{ObjRef},\\
P.\mathit{platform} &: & \mathit{PlatformName}\rightharpoonup\mathit{ObjRef},\\
P.\mathit{patternVariable} &: & \mathit{SiteId}\rightharpoonup\mathit{AST}.
\end{array}
\end{equation}
In the base relation these maps are immutable; reflective installation replaces
the complete $P$ atomically.  The map $P.\mathit{declMixin}$ is injective and
the first component of $P.\mathit{mixin}(P.\mathit{declMixin}(K))$ is $K$.
Every subsequent occurrence of a field such as $P.\mathit{mixin}$,
$P.\mathit{scopeChain}$, $P.\mathit{owner}$, or $P.\mathit{actorSeed}$ refers to
the maps collected in
Equation~\eqref{eq:program-definition}.
More precisely, if
\begin{equation}\label{eq:mixin-record}
  P.\mathit{mixin}(M)
    =\langle K,\phi,\delta,\nu,\iota\rangle,
\end{equation}
then $\methods_P(M)=\phi$, where
$\phi:\mathit{Selector}\rightharpoonup\mathit{MethodDef}$ is the finite map
of methods defined directly by $M$, including methods introduced by
elaboration.  The remaining components are the slot declarations $\delta$,
nested-class declarations $\nu$, and normalized mixin initializer
$\iota$, whose form is
\begin{equation}\label{eq:mixin-initializer}
 \iota=\langle i_I,s_I,\bar x_I,\bar g,b_I\rangle.
\end{equation}
Here $i_I$ is its stable
synthetic method identity, $s_I$ and $\bar x_I$ are its selector and formals,
$\bar g$ is an ordered list of sequential or simultaneous slot groups, and
$b_I$ is its initializer body.  For an $f$ having the form in
Equation~\eqref{eq:method-definition}, the projections are
\begin{equation}\label{eq:method-projections}
\begin{aligned}
 \iota_P(M)&=\iota, & \mathsf{methodId}(f)&=i_f,
   & \selector(f)&=s,\\
 \access(f)&=\alpha, & \mathsf{ownerDecl}(f)&=K.
\end{aligned}
\end{equation}
Write $\mathsf{ownSlots}_P(M)$ for the source-ordered sequence of every
physical slot identity declared by $\delta$, including the backing slots of
lazy declarations.  The groups in $\iota$ refer to the non-lazy members of
that sequence; lazy backing slots have no eager initializer group.

The heap $H$ is a finite map from the identities in
Equation~\eqref{eq:runtime-identities} to tagged records.  A factory descriptor
is
\begin{equation}\label{eq:factory-descriptor}
 q=\langle i_I,s_f,\bar x,\kappa\rangle,
 \qquad
 \kappa::=\kw{bodyInitializer}(m_S)\mid
            \kw{appliedMixinInitializer}(m_S,m_M),
\end{equation}
where $i_I$ identifies the class's synthesized instance initializer, $s_f$ and
$\bar x$ describe its primary factory, $m_S$ is the deferred superclass
initializer message, and $m_M$ is the deferred message for an applied mixin.
The records relevant to the base sequential semantics are:
\begin{equation}\label{eq:heap-records}
\begin{array}{rcl}
H(o)&=&\kw{object}\langle C,\sigma,\eta\rangle,\\
H(C)&=&\kw{class}\langle M,S,eo,C_m,q,\chi_C\rangle,\\
H(a)&=&\kw{activation}\langle p,\rho,\ell,o,C,k,h,q_a,\gamma\rangle,\\
H(c)&=&\kw{closure}\langle q,\bar x,\ell,b,a,o,C,h\rangle.
\end{array}
\end{equation}
$\sigma$ maps physical slot identities to values, and $\eta$ maps nested-class
declaration identities to the class objects already created for this receiver.
A class record contains its mixin $M$, superclass $S$, enclosing object $eo$,
metaclass $C_m$, primary factory description $q$, and stable creation
provenance $\chi_C$ from Equation~\eqref{eq:class-origin};
$q=\kw{none}$ only for a
non-instantiable platform class such as a metaclass.  An activation
records its method/closure provenance $p$, parameter environment $\rho$, local
slots $\ell$, current instance $o$, current class $C$, continuation object $k$,
home method activation $h$, continuability flag $q_a$, and lexical environment
$\gamma$.  The latter maps an enclosing method or closure declaration to its
current activation and an enclosing object-literal declaration to that
literal's runtime value; write $\mathsf{scope}_H(a)=\gamma$.  A closure captures
its source declaration $q$, the instantiating activation, current instance,
current class, and home method activation.
In particular, $\classof_H(C)=C_m$ for a class object.  The standard platform
classes of mixins, activations, closures, and continuations supply the
corresponding cases of $\classof_H$; their identities are left as platform
parameters until their reflective protocols are formalized.

We write $\classof_H(o)=C$, $\mixin_H(C)=M$, $\super_H(C)=S$, and
$\decl_P(M)=D$ for the corresponding projections of
Equations~\eqref{eq:mixin-record} and~\eqref{eq:heap-records}.  The distinguished
class $\Top$ terminates lookup and is not an instantiable Newspeak class.  In
the non-reflective semantics, every live class chain is finite and ends at
$\Top$.
The distinguished names used throughout abbreviate platform entries; in
particular,
\begin{equation}\label{eq:distinguished-platform-classes}
 P.\mathit{platform}(\kw{Object})=\Object,
 \qquad P.\mathit{platform}(\kw{Top})=\Top.
\end{equation}

\section{Class-chain operations}
\label{sec:class-chain}

Direct method selection, using the mixin map in
Equation~\eqref{eq:mixin-record}, is
\begin{equation}\label{eq:direct}
 \mathsf{direct}_{P,H}(C,s)=
 \begin{cases}
   f & \methods_P(\mixin_H(C))(s)=f,\\
   \undef & \text{otherwise.}
 \end{cases}
\end{equation}
Unrestricted definition lookup follows the superclass chain and ignores
accessibility.  Its named rules use direct selection from
Equation~\eqref{eq:direct}:
\[
\begin{array}{c}
\dfrac{C=\Top}{P,H\vdash \lookup_U(s,C)=\undef}
\quad(\rulelabel{U-Top})\\[2ex]
\dfrac{\mathsf{direct}(C,s)=f}{P,H\vdash \lookup_U(s,C)=\langle f,C\rangle}
\quad(\rulelabel{U-Here})\\[2ex]
\dfrac{C\ne\Top\quad \mathsf{direct}(C,s)=\undef\quad
       P,H\vdash\lookup_U(s,\super_H(C))=r}
      {P,H\vdash\lookup_U(s,C)=r}
\quad(\rulelabel{U-Super})
\end{array}
\]
The defining class is independently defined by the following partial
function:
\[
\begin{array}{c}
\dfrac{C=\Top}
      {P,H\vdash\definingClass(s,C)=\undef}
\quad(\rulelabel{Defining-Top})\\[2ex]
\dfrac{\mathsf{direct}(C,s)=f}
      {P,H\vdash\definingClass(s,C)=C}
\quad(\rulelabel{Defining-Here})\\[2ex]
\dfrac{C\ne\Top\quad \mathsf{direct}(C,s)=\undef\quad
       P,H\vdash\definingClass(s,\super_H(C))=C_f}
      {P,H\vdash\definingClass(s,C)=C_f}
\quad(\rulelabel{Defining-Super})
\end{array}
\]
It follows directly from these rules and the unrestricted-lookup rules that
$P,H\vdash\lookup_U(s,C)=\langle f,C_f\rangle$ implies
$P,H\vdash\definingClass(s,C)=C_f$.

\subsection{Public lookup}
\label{sec:public-lookup}

Public lookup implements ordinary sends.  A protected direct declaration is a
\emph{barrier}: it makes lookup fail rather than exposing a public superclass
method.  A private direct declaration is not inherited and is transparent to
this lookup.  Every occurrence of $\mathsf{direct}$ below denotes
Equation~\eqref{eq:direct}.
\[
\begin{array}{c}
\dfrac{C=\Top}{P,H\vdash\lookup_{pub}(s,C)=\undef}
\quad(\rulelabel{Pub-Top})\\[2ex]
\dfrac{\mathsf{direct}(C,s)=f\quad\access(f)=\Public}
      {P,H\vdash\lookup_{pub}(s,C)=\langle f,C\rangle}
\quad(\rulelabel{Pub-Hit})\\[2ex]
\dfrac{\mathsf{direct}(C,s)=f\quad\access(f)=\Protected}
      {P,H\vdash\lookup_{pub}(s,C)=\undef}
\quad(\rulelabel{Pub-Protected-Barrier})\\[2ex]
\dfrac{\substack{C\ne\Top\quad
       (\mathsf{direct}(C,s)=\undef\ \lor\
        \access(\mathsf{direct}(C,s))=\Private)\quad
       P,H\vdash\lookup_{pub}(s,\super_H(C))=r}}
      {P,H\vdash\lookup_{pub}(s,C)=r}
\quad(\rulelabel{Pub-Super})
\end{array}
\]

\subsection{Protected lookup}
\label{sec:protected-lookup}

Protected lookup implements the virtual part of self, outer, and super sends.
Both public and protected methods participate; private declarations remain
lexical and therefore are transparent.  These rules likewise use direct
selection from Equation~\eqref{eq:direct}.
\[
\begin{array}{c}
\dfrac{C=\Top}{P,H\vdash\lookup_{prot}(s,C)=\undef}
\quad(\rulelabel{Prot-Top})\\[2ex]
\dfrac{\mathsf{direct}(C,s)=f\quad
       \access(f)\in\{\Public,\Protected\}}
      {P,H\vdash\lookup_{prot}(s,C)=\langle f,C\rangle}
\quad(\rulelabel{Prot-Hit})\\[2ex]
\dfrac{\substack{C\ne\Top\quad
       (\mathsf{direct}(C,s)=\undef\ \lor\
        \access(\mathsf{direct}(C,s))=\Private)\quad
       P,H\vdash\lookup_{prot}(s,\super_H(C))=r}}
      {P,H\vdash\lookup_{prot}(s,C)=r}
\quad(\rulelabel{Prot-Super})
\end{array}
\]

\section{Lexical classes and enclosing receivers}

For a well-formed heap $H$, define the finite superclass-chain sequence by
primitive recursion:
\begin{equation}\label{eq:class-chain}
\begin{aligned}
 \mathsf{chain}_H(\Top) &= \langle\Top\rangle,\\
 \mathsf{chain}_H(R) &=
   \langle R\rangle\mathbin{\cdot}\mathsf{chain}_H(\super_H(R))
   &&\text{if }R\ne\Top.
\end{aligned}
\end{equation}
Here $\cdot$ is sequence concatenation and indices are zero-based.  The base
invariant that superclass links are acyclic and reach $\Top$ makes this
definition total and gives a unique finite sequence for every live class.
Superclass membership and depth are then, from the chain in
Equation~\eqref{eq:class-chain},
\begin{equation}\label{eq:class-membership-depth}
\begin{aligned}
 C\preceq_H R
 &\iff \exists j\,.
   0\le j<|\mathsf{chain}_H(R)|\ \land\
   \mathsf{chain}_H(R)[j]=C,\\
 \mathsf{depth}_H(C,R)
 &=\min\{j\mid \mathsf{chain}_H(R)[j]=C\},
\end{aligned}
\end{equation}
where depth is undefined when $C\not\preceq_H R$.

The lexical-class chain is the corresponding static construction.  Define the
partial function
\begin{equation}\label{eq:lex-ancestor}
\begin{aligned}
 \mathsf{lexAncestor}_P(D,0)&=D,\\
 \mathsf{lexAncestor}_P(D,n+1)&=
   P.\mathit{lexParent}(\mathsf{lexAncestor}_P(D,n)),
\end{aligned}
\end{equation}
where the second clause is undefined if either operand on its right is
undefined.  Using Equation~\eqref{eq:lex-ancestor}, define for a lexical
ancestor $D_t$ of $D_0$
\begin{equation}\label{eq:lex-depth}
 \mathsf{lexDepth}_P(D_t,D_0)=
   \min\{k\mid \mathsf{lexAncestor}_P(D_0,k)=D_t\}.
\end{equation}

For a declaration $D$, let the indices of its applications in $R$'s chain
from Equation~\eqref{eq:class-chain} be
\begin{equation}\label{eq:application-indices}
 \mathsf{appIndices}_{P,H}(R,D)=
 \left\{j\ \middle|\
 \begin{array}{l}
 0\le j<|\mathsf{chain}_H(R)|,\quad
 \mathsf{chain}_H(R)[j]\ne\Top,\\
 \decl_P(\mixin_H(\mathsf{chain}_H(R)[j]))=D
 \end{array}\right\}.
\end{equation}
Using Equation~\eqref{eq:application-indices}, the nearest application is the
partial function
\begin{equation}\label{eq:nearest-application}
 \mathsf{nearestApp}_{P,H}(R,D)=
 \begin{cases}
  \mathsf{chain}_H(R)[\min I]
    & I=\mathsf{appIndices}_{P,H}(R,D)\ne\varnothing,\\
  \undef & I=\varnothing.
 \end{cases}
\end{equation}
Minimal index means most specific.  This is significant when the same mixin
occurs repeatedly in one hierarchy.

For a runtime receiver $o_t$ and lexical target declaration $D_t$, define the
target application from Equation~\eqref{eq:nearest-application}:
\begin{equation}\label{eq:target-application}
  \mathsf{targetApp}_{P,H}(o_t,D_t)
    =\mathsf{nearestApp}_{P,H}(\classof_H(o_t),D_t).
\end{equation}
Well-formed outer sends preserve the invariant that this application exists.

The partial relation
$P,H\vdash \mathsf{enc}(o,C,n)=o_n$ formalizes the Specification's $n$th
enclosing object of $o$ with respect to $C$.  Its premises use superclass
membership from Equation~\eqref{eq:class-membership-depth}, lexical ancestry
from Equation~\eqref{eq:lex-ancestor}, and nearest application from
Equation~\eqref{eq:nearest-application}:
\[
\begin{array}{c}
\dfrac{C\preceq_H\classof_H(o)}
      {P,H\vdash\mathsf{enc}(o,C,0)=o}
\quad(\rulelabel{Enc-Zero})\\[2ex]
\dfrac{C\preceq_H\classof_H(o)\quad H(C).eo=o_1}
      {P,H\vdash\mathsf{enc}(o,C,1)=o_1}
\quad(\rulelabel{Enc-One})\\[2ex]
\dfrac{\substack{n>1\quad P,H\vdash\mathsf{enc}(o,C,n-1)=o_{n-1}\quad
       D_{n-1}=\mathsf{lexAncestor}_P(\decl_P(\mixin_H(C)),n-1)\\
       S_{n-1}=\mathsf{nearestApp}(\classof_H(o_{n-1}),D_{n-1})\quad
       H(S_{n-1}).eo=o_n}}
      {P,H\vdash\mathsf{enc}(o,C,n)=o_n}
\quad(\rulelabel{Enc-Step})
\end{array}
\]
The separate $n=1$ rule is intentional.  It preserves the Specification's
behavior under repeated mixin application rather than recomputing a possibly
different application from the receiver's dynamic class.

\section{Dispatch}

A dispatch result is $\kw{invoke}(f,o,C_f,\mu)$,
$\kw{missing}(o,C_0,\mu)$, or
$\kw{runtimeError}(\kw{invalidSuper}(C))$.  $C_f$ is the current class to
install in the new method activation; $C_0$ is the class at which DNU lookup
must begin.  The third form is an internal failure, deliberately not a DNU
request; Section~\ref{sec:exception-boundary} reifies it as a platform
exception and signals it through ordinary dispatch.
For compact rules, $\mathsf{outerDispatch}_a$ denotes outer dispatch in activation $a$.

\subsection{Ordinary sends}

After receiver and arguments have evaluated left-to-right to $o$ and $\bar o$,
let $\mu=\langle s,\bar o\rangle$ and $R=\classof_H(o)$.  Dispatch uses the
public lookup of Section~\ref{sec:public-lookup}:
\[
\dfrac{P,H\vdash\lookup_{pub}(s,R)=\langle f,C_f\rangle}
      {P,H\vdash\mathsf{ordinaryDispatch}(o,\mu)
       \Downarrow\kw{invoke}(f,o,C_f,\mu)}
\quad(\rulelabel{Send-Public})
\]
\[
\dfrac{P,H\vdash\lookup_{pub}(s,R)=\undef}
      {P,H\vdash\mathsf{ordinaryDispatch}(o,\mu)
       \Downarrow\kw{missing}(o,R,\mu)}
\quad(\rulelabel{Send-Missing})
\]

\subsection{Outer and self sends}
\label{sec:outer-self}

Private selection is a static operation on the program.  Define
\begin{equation}\label{eq:private-method}
 \mathsf{privateMethod}_P(D,s)=
 \begin{cases}
  f & \begin{array}{l}
        \text{if there exists }M\in\dom(P.\mathit{mixin})\text{ such that}\\[-1mm]
        \decl_P(M)=D,\quad \methods_P(M)(s)=f,\quad
        \access(f)=\Private,
      \end{array}\\[3mm]
  \undef & \text{otherwise.}
 \end{cases}
\end{equation}
Well-formedness gives each class declaration one static mixin identity, and
each mixin at most one direct method per selector, so the first case determines
a unique $f$.  Thus this operation selects only a private method declared
directly by $D$; it does not search a runtime superclass chain.  In particular,
a direct public or protected method of selector $s$ makes
$\mathsf{privateMethod}_P(D,s)=\undef$.

Suppose the current activation has current instance $o_0$ and current class
$C_0$.  If the immediately enclosing class declaration is $D_0$ and target
$D_t$ is its $k$th lexical class ancestor, obtain $o_t$ from
$\mathsf{enc}(o_0,C_0,k)$.  Let $N_t$ be the corresponding class
$N_t=\mathsf{targetApp}(o_t,D_t)$ from
Equation~\eqref{eq:target-application}, and let
$\mu=\langle s,\bar o\rangle$.

If Equation~\eqref{eq:private-method} selects a private method $f$ directly
declared by $D_t$, dispatch is lexically bound:
\[
\dfrac{P,H\vdash\mathsf{enc}(o_0,C_0,k)=o_t\quad
       \mathsf{privateMethod}_P(D_t,s)=f}
      {P,H\vdash\mathsf{outerDispatch}_a(D_t,D_0,\mu)
       \Downarrow\kw{invoke}(f,o_t,N_t,\mu)}
\quad(\rulelabel{Outer-Private})
\]
Otherwise dispatch remains virtual and sees public or protected methods via
Section~\ref{sec:protected-lookup}:
{\small
\[
\dfrac{\substack{P,H\vdash\mathsf{enc}(o_0,C_0,k)=o_t\quad
       \mathsf{privateMethod}_P(D_t,s)=\undef\quad
       R=\classof_H(o_t)\quad
       P,H\vdash\lookup_{prot}(s,R)=\langle f,C_f\rangle}}
      {P,H\vdash\mathsf{outerDispatch}_a(D_t,D_0,\mu)
       \Downarrow\kw{invoke}(f,o_t,C_f,\mu)}
\quad(\rulelabel{Outer-Protected})
\]
}
A failed protected lookup yields $\kw{missing}(o_t,R,\mu)$.

A self send lexically enclosed immediately by $D_0$ is the depth-zero case
\begin{equation}\label{eq:self-send-derived}
  \mathsf{selfDispatch}_a(\mu)
    \equiv \mathsf{outerDispatch}_a(D_0,D_0,\mu).
\end{equation}
It is therefore not an ordinary send: it can select a private method declared
in $D_0$ and can dynamically select protected methods.

\paragraph{Clarification obligation.}
The $N_t$ current class in \rulelabel{Outer-Private} is the actual class
application on the enclosing path, not merely the class found by unrestricted
lookup of $s$ on $o_t$.  Otherwise a more-derived declaration of the same
selector could cause a statically selected private method to execute with the
wrong current class.  The Specification appears to require this result but does
not state it directly.

\subsection{Super sends}

For current instance $o$, current class $C\notin\{\Object,\Top\}$,
$S=\super_H(C)$, and $\mu=\langle s,\bar o\rangle$, dispatch uses the protected
lookup of Section~\ref{sec:protected-lookup}:
\[
\dfrac{P,H\vdash\lookup_{prot}(s,S)=\langle f,C_f\rangle}
      {P,H,a\vdash\mathsf{superDispatch}(\mu)
       \Downarrow\kw{invoke}(f,o,C_f,\mu)}
\quad(\rulelabel{Super-Hit})
\]
\[
\dfrac{P,H\vdash\lookup_{prot}(s,S)=\undef}
      {P,H,a\vdash\mathsf{superDispatch}(\mu)
       \Downarrow\kw{missing}(o,S,\mu)}
\quad(\rulelabel{Super-Missing})
\]
The $\kw{missing}$ result is processed by the DNU fallback of
Section~\ref{sec:dnu}; in particular, $S$ is the class at which DNU lookup
begins.

The two excluded current classes have explicit error rules:
\[
\dfrac{C=\Object}
      {P,H,a\vdash\mathsf{superDispatch}(\mu)
       \Downarrow\kw{runtimeError}(\kw{invalidSuper}(\Object))}
\quad(\rulelabel{Super-Object-Error})
\]
\[
\dfrac{C=\Top}
      {P,H,a\vdash\mathsf{superDispatch}(\mu)
       \Downarrow\kw{runtimeError}(\kw{invalidSuper}(\Top))}
\quad(\rulelabel{Super-Top-Error})
\]
Neither error proceeds to DNU fallback.

\subsection{Implicit receiver selection}
\label{sec:implicit}

Lexical resolution precedes dynamic inheritance.  If
\begin{equation}\label{eq:scope-chain}
 P.\mathit{scopeChain}(z)=\langle d_0,\ldots,d_n\rangle,
\end{equation}
then the declarations are ordered from the construct immediately surrounding
site $z$ outwards.  The static predicate $\mathsf{declares}_P(d,s)$ is
\begin{equation}\label{eq:declares}
\begin{array}{rcl}
 \mathsf{declares}_P(K,s)
 &\Longleftrightarrow&
   \exists M\in\dom(P.\mathit{mixin})\,.
   \decl_P(M)=K\land s\in\dom(\methods_P(M)),\\
 \mathsf{declares}_P(q,s)
 &\Longleftrightarrow&
   s\in P.\mathit{activationMembers}(q).
\end{array}
\end{equation}
The method maps include accessors synthesized for class or object-literal slots
and nested classes; activation members include accessors for parameters and
locals.  Using the scope sequence in Equation~\eqref{eq:scope-chain} and the
predicate in Equation~\eqref{eq:declares}, define
\begin{equation}\label{eq:lexical-binding}
\begin{aligned}
 J_{P,z,s}&=\{j\mid 0\le j\le n\land
                         \mathsf{declares}_P(d_j,s)\},\\
 \mathsf{bind}_P(z,s)&=
 \begin{cases}
   d_{\min J_{P,z,s}} & J_{P,z,s}\ne\varnothing,\\
   \undef & J_{P,z,s}=\varnothing.
 \end{cases}
\end{aligned}
\end{equation}
Thus $\mathsf{bind}$ selects the innermost lexical declaration, independently
of access level.  Elaboration annotates the core send with
$d=\mathsf{bind}_P(z,s)$ as defined by
Equation~\eqref{eq:lexical-binding}.  It also retains the immediately enclosing class
declaration $D_0$ as static context.  This context is written on the judgment
rather than repeated inside the core node below.

Suppose the message clause has evaluated to $\mu=\langle s,\bar o\rangle$ and
the current activation is $a$.  The judgment
\begin{equation}\label{eq:implicit-judgment}
 P,H,a,D_0\vdash\mathsf{implicitDispatch}(d,\mu)\Downarrow r
\end{equation}
is defined by the following four rules, where $r$ ranges over dispatch results.
For a class declaration, implicit dispatch is exactly outer dispatch:
\[
\dfrac{d=D_t\in\mathit{ClassDeclId}\quad
       P,H\vdash\mathsf{outerDispatch}_a(D_t,D_0,\mu)\Downarrow r}
      {P,H,a,D_0\vdash\mathsf{implicitDispatch}(d,\mu)\Downarrow r}
\quad(\rulelabel{Implicit-Class})
\]
For an enclosing object literal, its lexical environment supplies the literal
value and the statically selected direct method is invoked without an access
check:
\[
\dfrac{\substack{d=L\in\mathit{ObjLitDeclId}\\
       \mathsf{scope}_H(a)(L)=o_L\\
       P,H\vdash\lookup_U(s,\classof_H(o_L))=\langle f,C_f\rangle\\
       \decl_P(\mixin_H(C_f))=L}}
      {P,H,a,D_0\vdash\mathsf{implicitDispatch}(d,\mu)
       \Downarrow\kw{invoke}(f,o_L,C_f,\mu)}
\quad(\rulelabel{Implicit-Literal})
\]
For a method or closure declaration, the lexical environment supplies that
declaration's current activation, and the specified send to the activation is
an ordinary send:
\[
\dfrac{\substack{d=q\in\mathit{ActDeclId}\\
       \mathsf{scope}_H(a)(q)=a_q\\
       P,H\vdash\mathsf{ordinaryDispatch}(a_q,\mu)\Downarrow r}}
      {P,H,a,D_0\vdash\mathsf{implicitDispatch}(d,\mu)\Downarrow r}
\quad(\rulelabel{Implicit-Activation})
\]
Finally, failure of lexical binding is a self send:
\[
\dfrac{d=\undef\quad
       P,H\vdash\mathsf{outerDispatch}_a(D_0,D_0,\mu)\Downarrow r}
      {P,H,a,D_0\vdash\mathsf{implicitDispatch}(d,\mu)\Downarrow r}
\quad(\rulelabel{Implicit-Self})
\]
These rules are mutually exclusive because the declaration-identity domains
are disjoint.  Their premises are total for a well-formed annotation.  A
reflective edit that changes a scope chain or member map must rebuild the
affected $\mathsf{bind}$ annotations and lexical environments in the same
atomic install transaction.

\paragraph{ECOOP-style derived presentation.}
After nearer activation and object-literal scopes have failed to declare $s$,
the class-only part of the search can instead be written recursively using
Equation~\eqref{eq:declares}:
\begin{equation}\label{eq:class-scan}
 \mathsf{scan}_P(D,s)=
 \begin{cases}
  D & \mathsf{declares}_P(D,s),\\
  \mathsf{scan}_P(D',s) & \neg\mathsf{declares}_P(D,s)\ \land\
                           P.\mathit{lexParent}(D)=D',\\
  \undef & \neg\mathsf{declares}_P(D,s)\ \land\
            P.\mathit{lexParent}(D)=\undef.
 \end{cases}
\end{equation}
Unlike the 2010 algorithm, the result must retain both the receiver and the
declaration that granted lexical authority.  Write $\mathsf{irecv}$ for this
access-aware form of \emph{implicitReceiver}.  Using
Equations~\eqref{eq:class-scan} and~\eqref{eq:lex-depth}, define
\begin{equation}\label{eq:implicit-receiver}
 \mathsf{irecv}_{P,H}(o_0,C_0,D_0,s)=
 \begin{cases}
  \langle o_t,D_t\rangle &
    \begin{array}{l}
      \mathsf{scan}_P(D_0,s)=D_t,\\[-1mm]
      k=\mathsf{lexDepth}_P(D_t,D_0),\quad
      P,H\vdash\mathsf{enc}(o_0,C_0,k)=o_t,
    \end{array}\\[3mm]
  \langle o_0,\undef\rangle & \mathsf{scan}_P(D_0,s)=\undef.
 \end{cases}
\end{equation}
The first result dispatches by $\mathsf{outerDispatch}_a(D_t,D_0,\mu)$; the second by
$\mathsf{outerDispatch}_a(D_0,D_0,\mu)$.  Consequently a private method is considered
only as the direct declaration at the selected lexical target.  If that check
fails, $\lookup_{prot}$ performs virtual lookup and skips every private method
encountered on the runtime class chain.

\paragraph{Equivalence obligation.}
For a well-formed fixed $P,H$, and after ruling out a nearer activation or
object-literal binding, the elaborated rules and the ECOOP-style presentation
must produce the same dispatch result.  In terms of
Equations~\eqref{eq:lexical-binding} and~\eqref{eq:class-scan}, the static part
reduces to
\begin{equation}\label{eq:binding-equivalence}
 \mathsf{bind}_P(z,s)=D_t
 \quad\Longleftrightarrow\quad
 \mathsf{scan}_P(D_0,s)=D_t,
\end{equation}
with $D_0$ the first class scope at $z$; the receiver equality follows from
$\mathsf{lexDepth}$ in Equation~\eqref{eq:lex-depth} and the rules for
$\mathsf{enc}$.  A full proof is deferred
until the static well-formedness judgment fixes the precise correspondence
between $P.\mathit{scopeChain}$ and $P.\mathit{lexParent}$.

\subsection{DNU fallback}
\label{sec:dnu}

Missing dispatch does not perform an ordinary send of \dnu.  It uses
unrestricted lookup from Section~\ref{sec:class-chain}, starting at the supplied
class $C_0$, and invokes the result
on the original receiver with a freshly reified mirror of the original message:
\[
\dfrac{P,H\vdash\lookup_U(\dnu,C_0)=\langle f,C_f\rangle\quad
       H' = H[m\mapsto\mathsf{messageMirror}(\mu)]}
      {P,H\vdash\kw{missing}(o,C_0,\mu)
       \longrightarrow
       P,H'\vdash\kw{invoke}(f,o,C_f,\langle\dnu,m\rangle)}
\quad(\rulelabel{DNU})
\]
Consequently a class's DNU implementation is honored regardless of its access
modifier.  The required protected implementation on \Object{} ensures that
unrestricted lookup succeeds in every well-formed base program and ultimately
throws \kw{MessageNotUnderstood} unless overridden.

\section{Sequential abstract machine}
\label{sec:machine}

This section closes the dispatch loop with a deterministic small-step machine.
It covers left-to-right message construction, sequencing, closure creation,
method and closure invocation, tail-call-observable continuation links, and
local and non-local return.  Simultaneous slot clauses remain with the class
initialization layer; their absence is stated explicitly below.

\subsection{Configurations and frames}

A request records the authority of a send after its receiver, if any, has been
evaluated:
\begin{equation}\label{eq:request-grammar}
 \begin{aligned}
 \theta ::={}&\kw{ordinaryRequest}(o)\mid\kw{implicitRequest}(d,D_0)\\
 &\mid\kw{outerRequest}(D_t,D_0)\mid\kw{selfRequest}(D_0)
 \mid\kw{superRequest}.
 \end{aligned}
\end{equation}
An intra-activation frame stack $\pi$ is built from
\begin{equation}\label{eq:frame-grammar}
\begin{array}{rcl}
 F ::=&& \kw{recv}(s,\bar e)
 \mid \kw{args}(\theta,s,\bar o,\bar e)
 \mid \kw{seq}(\tau,\bar s)\\
 &&\mid \kw{returnFrame}
 \mid \kw{init}(j,\bar\ell,b,\tau),\\
 \pi ::=&& \epsilon\mid\pi\mathbin{\cdot}F,
 \qquad
 \tau ::= \kw{methodBody}(o)\mid\kw{closureBody}.
\end{array}
\end{equation}
Here $\bar s$ is a statement list, $\bar\ell$ is a normalized local-slot list,
and $j$ is a physical local-slot identity.  An activation stack and machine
configuration are
\begin{equation}\label{eq:configuration}
 \mathcal A::=\epsilon\mid\mathcal A\mathbin{\cdot}\langle a,\pi\rangle,
 \qquad
 \Sigma=\langle P,H,\mathcal A,t\rangle.
\end{equation}
We abbreviate the configuration of Equation~\eqref{eq:configuration} as
$\langle H;\mathcal A;t\rangle_P$ when this improves the layout.
The rightmost activation is current; $t$ is its control term.  Suspended
evaluation frames live in $\pi$, while the corresponding activation and its
observable continuation link live in $H$.  We write $\longmapsto$ for one
machine step and omit unchanged $P$ below.

\subsection{Expression and message order}

Evaluation of an ordinary send begins with its receiver:
\[
 \langle P,H,\mathcal A\cdot\langle a,\pi\rangle,
        \kw{send}(e,s,\bar e)\rangle
 \longmapsto
 \langle P,H,\mathcal A\cdot\langle a,\pi\cdot\kw{recv}(s,\bar e)\rangle,
        e\rangle.
 \tag{Send-Receiver}
\]
Once the receiver is $o$, the request is $\kw{ordinaryRequest}(o)$.  All other send
forms in Equation~\eqref{eq:core-expressions} determine their request without
evaluating a receiver:
\begin{equation}\label{eq:request-construction}
\begin{array}{rcl}
 \mathsf{request}_{D_0}(\kw{implicit}(s,\bar e,d))
   &=&\langle\kw{implicitRequest}(d,D_0),s,\bar e\rangle,\\
 \mathsf{request}_{D_0}(\kw{selfSend}(s,\bar e,D_0))
   &=&\langle\kw{selfRequest}(D_0),s,\bar e\rangle,\\
 \mathsf{request}_{D_0}(\kw{outerSend}(s,\bar e,D_t,D_0))
   &=&\langle\kw{outerRequest}(D_t,D_0),s,\bar e\rangle,\\
 \mathsf{request}_{D_0}(\kw{superSend}(s,\bar e))
   &=&\langle\kw{superRequest},s,\bar e\rangle.
\end{array}
\end{equation}
The following rules use Equation~\eqref{eq:request-construction} to start an
argument list; $e_s$ ranges over its four non-ordinary send forms:
\[
\dfrac{\mathsf{request}_{D_0}(e_s)=\langle\theta,s,\epsilon\rangle}
 {\langle P,H,\mathcal A\cdot\langle a,\pi\rangle,e_s\rangle
  \longmapsto
  \langle P,H,\mathcal A\cdot\langle a,\pi\rangle,
          \kw{dispatch}(\theta,\langle s,\epsilon\rangle)\rangle}
\quad(\rulelabel{Send-No-Args})
\]
\[
\dfrac{\mathsf{request}_{D_0}(e_s)=
       \langle\theta,s,\langle e_1,\bar e\rangle\rangle}
 {\langle P,H,\mathcal A\cdot\langle a,\pi\rangle,e_s\rangle
  \longmapsto
  \langle P,H,\mathcal A\cdot
    \langle a,\pi\cdot\kw{args}(\theta,s,\epsilon,\bar e)\rangle,e_1\rangle}
\quad(\rulelabel{Send-First-Arg})
\]
The receiver frame from Equation~\eqref{eq:frame-grammar} uses the same two
continuations:
\begin{equation}\label{eq:receiver-continuations}
\begin{array}{c}
\langle P,H,\mathcal A\cdot\langle a,\pi\cdot\kw{recv}(s,\epsilon)\rangle,o\rangle
\longmapsto
\langle P,H,\mathcal A\cdot\langle a,\pi\rangle,
        \kw{dispatch}(\kw{ordinaryRequest}(o),\langle s,\epsilon\rangle)\rangle,
\\[2ex]
\langle P,H,\mathcal A\cdot
 \langle a,\pi\cdot\kw{recv}(s,\langle e_1,\bar e\rangle)\rangle,o\rangle
\longmapsto
\langle P,H,\mathcal A\cdot
 \langle a,\pi\cdot\kw{args}(\kw{ordinaryRequest}(o),s,\epsilon,\bar e)\rangle,e_1\rangle.
\end{array}
\end{equation}
Argument values are accumulated on the left, so these rules enforce source
order:
\begin{equation}\label{eq:argument-continuations}
\begin{array}{c}
\langle P,H,\mathcal A\cdot
 \langle a,\pi\cdot\kw{args}(\theta,s,\bar o,
                              \langle e,\bar e\rangle)\rangle,o'\rangle
\longmapsto
\langle P,H,\mathcal A\cdot
 \langle a,\pi\cdot\kw{args}(\theta,s,
                 \bar o\cdot\langle o'\rangle,\bar e)\rangle,e\rangle,
\\[2ex]
\langle P,H,\mathcal A\cdot
 \langle a,\pi\cdot\kw{args}(\theta,s,\bar o,\epsilon)\rangle,o'\rangle
\longmapsto
\langle P,H,\mathcal A\cdot\langle a,\pi\rangle,
 \kw{dispatch}(\theta,\langle s,\bar o\cdot\langle o'\rangle\rangle)\rangle.
\end{array}
\end{equation}
Standalone $\kw{self}$ reduces to the current instance recorded in $H(a)$;
parentheses have already been eliminated without changing a parenthesized self
send's syntactic classification.
Formally,
\[
 \dfrac{H(a).o=o}
 {\langle P,H,\mathcal A\cdot\langle a,\pi\rangle,\kw{self}\rangle
  \longmapsto
  \langle P,H,\mathcal A\cdot\langle a,\pi\rangle,o\rangle}
 \quad(\rulelabel{Self-Value}).
\]

\subsection{From requests to invocation}

For $n\geq0$, let $\mathsf{valueSelector}(n)$ be $\#\kw{value}$ for $n=0$
and the symbol consisting of $n$ repetitions of the keyword \kws{value}
otherwise.  Let $\mathsf{closureCall}_H(o,\mu)$ hold exactly when $H(o)$ is a
closure with $n$ parameters and
$\selector(\mu)=\mathsf{valueSelector}(n)$.  Define the partial
request-resolution function, over the requests of
Equation~\eqref{eq:request-grammar}, by the existing dispatch judgments:
\begin{equation}\label{eq:request-resolution}
\begin{aligned}
 &\mathsf{resolve}_{P,H,a}(\kw{ordinaryRequest}(o),\mu)=r\\[-1mm]
 &\quad\Longleftrightarrow
   \neg\mathsf{closureCall}_H(o,\mu)\ \land\
   P,H\vdash\mathsf{ordinaryDispatch}(o,\mu)\Downarrow r,\\
 &\mathsf{resolve}_{P,H,a}(\kw{implicitRequest}(d,D_0),\mu)=r\\[-1mm]
 &\quad\Longleftrightarrow
   P,H,a,D_0\vdash\mathsf{implicitDispatch}(d,\mu)\Downarrow r,\\
 &\mathsf{resolve}_{P,H,a}(\kw{outerRequest}(D_t,D_0),\mu)=r\\[-1mm]
 &\quad\Longleftrightarrow
   P,H\vdash\mathsf{outerDispatch}_a(D_t,D_0,\mu)\Downarrow r,\\
 &\mathsf{resolve}_{P,H,a}(\kw{selfRequest}(D_0),\mu)=r\\[-1mm]
 &\quad\Longleftrightarrow
   P,H\vdash\mathsf{outerDispatch}_a(D_0,D_0,\mu)\Downarrow r,\\
 &\mathsf{resolve}_{P,H,a}(\kw{superRequest},\mu)=r\\[-1mm]
 &\quad\Longleftrightarrow
   P,H,a\vdash\mathsf{superDispatch}(\mu)\Downarrow r.
\end{aligned}
\end{equation}
The request-resolution step applies Equation~\eqref{eq:request-resolution}:
\[
\dfrac{\mathsf{resolve}_{P,H,a}(\theta,\mu)=r}
 {\langle P,H,\mathcal A\cdot\langle a,\pi\rangle,
             \kw{dispatch}(\theta,\mu)\rangle
  \longmapsto
  \langle P,H,\mathcal A\cdot\langle a,\pi\rangle,r\rangle}
\quad(\rulelabel{Resolve})
\]
and the DNU rule of Section~\ref{sec:dnu} rewrites a $\kw{missing}$ control
term to $\kw{invoke}$ while allocating its message mirror.  A
$\kw{runtimeError}$ control term is reified and signaled at the standard-library
boundary in Section~\ref{sec:exception-boundary}.

Before ordinary request resolution, a closure receiver has the specified
invocation case
\[
\dfrac{H(c)=\kw{closure}\langle q,\bar x,\ell,b,a_c,o,C,h\rangle\quad
       \mu=\langle\mathsf{valueSelector}(|\bar x|),\bar o\rangle}
 {\begin{aligned}
  &\langle H;\mathcal A\cdot\langle a,\pi\rangle;
       \kw{dispatch}(\kw{ordinaryRequest}(c),\mu)\rangle_P\\[-1mm]
  &\qquad\longmapsto
   \langle H;\mathcal A\cdot\langle a,\pi\rangle;
       \kw{invokeClosure}(c,\mu)\rangle_P
  \end{aligned}}
\quad(\rulelabel{Closure-Value})
\]
which is disjoint from \rulelabel{Resolve} for that request.

\subsection{Closure creation and activation allocation}

Evaluating a closure literal allocates a fresh closure object and captures the
current activation state.  Let
\begin{equation}\label{eq:home-activation}
 \mathsf{home}_H(a)=
 \begin{cases}
  a & H(a).p\in\mathit{MethodId},\\
  H(a).h & \text{otherwise.}
\end{cases}
\end{equation}
For fresh $c$, the allocation rule uses $\mathsf{home}_H$ from
Equation~\eqref{eq:home-activation}:
\[
\dfrac{H(a).o=o\quad H(a).C=C\quad h=\mathsf{home}_H(a)}
 {\begin{aligned}
  &\langle H;\mathcal A\cdot\langle a,\pi\rangle;
       \kw{closureLiteral}(q,\bar x,\ell,b)\rangle_P\\[-1mm]
  &\quad\longmapsto
   \langle H[c\mapsto\kw{closure}\langle q,\bar x,\ell,b,a,o,C,h\rangle];
       \mathcal A\cdot\langle a,\pi\rangle;c\rangle_P
  \end{aligned}}
\quad(\rulelabel{Closure-New})
\]

The continuation installed in a callee is
\begin{equation}\label{eq:callee-continuation}
 \mathsf{calleeKont}_H(a,\pi)=
 \begin{cases}a&\pi\ne\epsilon,\\H(a).k&\pi=\epsilon.\end{cases}
\end{equation}
Thus an observable continuation points to the sender exactly when the sender
has further computation.  Define method-activation allocation for
$f=\langle i_f,s,\alpha,\bar x,\ell,b,K\rangle$ and
$\mu=\langle s,\bar o\rangle$, with $|\bar x|=|\bar o|$, by
\begin{equation}\label{eq:method-activation}
\begin{aligned}
 \rho'&=[\bar x\mapsto\bar o],&
 \lambda'&=[\bar\ell\mapsto\kw{uninit}],\\
 \gamma'&=\{i_f\mapsto a'\}\cup
   \begin{cases}\{K\mapsto o\}&K\in\mathit{ObjLitDeclId},\\
                 \varnothing&\text{otherwise,}\end{cases}\\
 \mathsf{newM}(P,H,a',f,o,C_f,\mu,k)
   &=H[a'\mapsto\kw{activation}\langle
       i_f,\rho',\lambda',o,C_f,k,a',\kw{true},\gamma'\rangle].
\end{aligned}
\end{equation}
where $a'$ is fresh and $\bar\ell$ lists the physical local-slot identities.
Allocation itself does not evaluate local initializers.  The non-tail and tail
invocation rules implement Equations~\eqref{eq:callee-continuation}
and~\eqref{eq:method-activation}:
\[
\dfrac{\pi\ne\epsilon\quad k=a\quad
       H'=\mathsf{newM}(P,H,a',f,o,C_f,\mu,k)}
 {\begin{aligned}
  &\langle H;\mathcal A\cdot\langle a,\pi\rangle;
      \kw{invoke}(f,o,C_f,\mu)\rangle_P\\[-1mm]
  &\quad\longmapsto
   \langle H';\mathcal A\cdot\langle a,\pi\rangle
                    \cdot\langle a',\epsilon\rangle;
      \kw{initialize}(a',\ell,b,\kw{methodBody}(o))\rangle_P
  \end{aligned}}
\quad(\rulelabel{Invoke-NonTail})
\]
\[
\dfrac{\pi=\epsilon\quad k=H(a).k\quad
       H'=\mathsf{newM}(P,H,a',f,o,C_f,\mu,k)}
 {\begin{aligned}
  &\langle H;\mathcal A\cdot\langle a,\epsilon\rangle;
      \kw{invoke}(f,o,C_f,\mu)\rangle_P\\[-1mm]
  &\quad\longmapsto
   \langle H';\mathcal A\cdot\langle a',\epsilon\rangle;
      \kw{initialize}(a',\ell,b,\kw{methodBody}(o))\rangle_P
  \end{aligned}}
\quad(\rulelabel{Invoke-Tail})
\]
The old activation remains in $H$ for mirrors and debugging even when its
dynamic stack entry is replaced.

For closure activation, suppose the captured data satisfy
\begin{equation}\label{eq:closure-invocation-data}
 H(c)=\kw{closure}\langle q,\bar x,\ell,b,a_c,o,C,h\rangle,
 \qquad
 \mu=\langle\mathsf{valueSelector}(|\bar x|),\bar o\rangle.
\end{equation}
Under Equation~\eqref{eq:closure-invocation-data}, with $a'$ fresh, define
\begin{equation}\label{eq:closure-activation}
\begin{aligned}
 \rho_c&=[\bar x\mapsto\bar o],&
 \lambda_c&=[\bar\ell\mapsto\kw{uninit}],\\
 \gamma_c&=\mathsf{scope}_H(a_c)[q\mapsto a'],\\
 \mathsf{newC}(H,a',c,\bar o,k)
   &=H[a'\mapsto\kw{activation}\langle
       c,\rho_c,\lambda_c,o,C,k,h,\kw{true},\gamma_c\rangle].
\end{aligned}
\end{equation}
The two closure invocation cases use the activation allocated by
Equation~\eqref{eq:closure-activation}:
\[
\dfrac{\pi\ne\epsilon\quad k=a\quad
       H'=\mathsf{newC}(H,a',c,\bar o,k)}
 {\begin{aligned}
  &\langle H;\mathcal A\cdot\langle a,\pi\rangle;
      \kw{invokeClosure}(c,\mu)\rangle_P\\[-1mm]
  &\quad\longmapsto
   \langle H';\mathcal A\cdot\langle a,\pi\rangle
                    \cdot\langle a',\epsilon\rangle;
      \kw{initialize}(a',\ell,b,\kw{closureBody})\rangle_P
  \end{aligned}}
\quad(\rulelabel{Closure-Invoke-NonTail})
\]
\[
\dfrac{\pi=\epsilon\quad k=H(a).k\quad
       H'=\mathsf{newC}(H,a',c,\bar o,k)}
 {\begin{aligned}
  &\langle H;\mathcal A\cdot\langle a,\epsilon\rangle;
      \kw{invokeClosure}(c,\mu)\rangle_P\\[-1mm]
  &\quad\longmapsto
   \langle H';\mathcal A\cdot\langle a',\epsilon\rangle;
      \kw{initialize}(a',\ell,b,\kw{closureBody})\rangle_P
  \end{aligned}}
\quad(\rulelabel{Closure-Invoke-Tail})
\]

\subsection{Local initialization and bodies}

Normalize a sequential local declaration as
\begin{equation}\label{eq:local-declaration}
 \lambda::=\kw{immutable}(j,e)\mid\kw{mutable}(j,e)\mid
 \kw{mutable}(j,\kw{none})\mid
 \kw{lazyImmutable}(j,e)\mid\kw{lazyMutable}(j,e).
\end{equation}
The normalized sequence may also contain a simultaneous group:
\begin{equation}\label{eq:local-declaration-sequence}
 \bar\ell::=\epsilon\mid\lambda\cdot\bar\ell
       \mid\kw{simultaneous}(\bar\lambda)\cdot\bar\ell.
\end{equation}
Empty initialization of the declaration sequence in
Equation~\eqref{eq:local-declaration-sequence} enters the body:
\begin{equation}\label{eq:local-init-empty}
 \langle P,H,\mathcal A\cdot\langle a,\pi\rangle,
   \kw{initialize}(a,\epsilon,b,\tau)\rangle
 \longmapsto
 \langle P,H,\mathcal A\cdot\langle a,\pi\rangle,
   \kw{body}(\tau,b)\rangle.
\end{equation}
An immutable or initialized mutable local evaluates its initializer in the new
activation and then writes directly:
\begin{equation}\label{eq:local-init-start}
\dfrac{\lambda\in\{\kw{immutable}(j,e),\kw{mutable}(j,e)\}}
 {\langle P,H,\mathcal A\cdot\langle a,\pi\rangle,
       \kw{initialize}(a,\lambda\cdot\bar\ell,b,\tau)\rangle
  \longmapsto
  \langle P,H,\mathcal A\cdot
       \langle a,\pi\cdot\kw{init}(j,\bar\ell,b,\tau)\rangle,e\rangle}
\end{equation}
\begin{equation}\label{eq:local-init-commit}
 \langle P,H,\mathcal A\cdot
       \langle a,\pi\cdot\kw{init}(j,\bar\ell,b,\tau)\rangle,o\rangle
 \longmapsto
 \langle P,H[a.\ell(j)\mapsto o],
       \mathcal A\cdot\langle a,\pi\rangle,
       \kw{initialize}(a,\bar\ell,b,\tau)\rangle.
\end{equation}
An uninitialized mutable local and either lazy form write $\nil$ without
evaluating the lazy initializer:
\begin{equation}\label{eq:local-init-nil}
\dfrac{\lambda\in\{\kw{mutable}(j,\kw{none}),
        \kw{lazyImmutable}(j,e),\kw{lazyMutable}(j,e)\}}
 {\langle P,H,\mathcal A\cdot\langle a,\pi\rangle,
       \kw{initialize}(a,\lambda\cdot\bar\ell,b,\tau)\rangle
  \longmapsto
  \langle P,H[a.\ell(j)\mapsto\nil],
       \mathcal A\cdot\langle a,\pi\rangle,
       \kw{initialize}(a,\bar\ell,b,\tau)\rangle.}
\end{equation}

Simultaneous locals use direct activation-slot operations, just as instance
initialization uses the direct object-slot operations of
Equation~\eqref{eq:slot-operations}.  Extend the frame grammar in
Equation~\eqref{eq:frame-grammar} with $\kw{localWrite}(a,j)$ and define
\begin{equation}\label{eq:local-slot-operations}
\begin{array}{c}
\dfrac{H(a_0).\ell(j)=v}
 {\langle P,H,\mathcal A\cdot\langle a,\pi\rangle,
       \kw{readLocal}(a_0,j)\rangle
  \longmapsto
  \langle P,H,\mathcal A\cdot\langle a,\pi\rangle,v\rangle},\\[3mm]
 \langle P,H,\mathcal A\cdot\langle a,\pi\rangle,
       \kw{writeLocal}(a_0,j,e)\rangle
 \longmapsto
 \langle P,H,\mathcal A\cdot
       \langle a,\pi\cdot\kw{localWrite}(a_0,j)\rangle,e\rangle,\\[3mm]
 \langle P,H,\mathcal A\cdot
       \langle a,\pi\cdot\kw{localWrite}(a_0,j)\rangle,v\rangle
 \longmapsto
 \langle P,H[a_0.\ell(j)\mapsto v],
       \mathcal A\cdot\langle a,\pi\rangle,v\rangle.
\end{array}
\end{equation}
Let $e_F$ be the annotated core send denoting \PastFuture{} in activation $a$.
For a declaration in a simultaneous group, define its installation term by
\begin{equation}\label{eq:local-simultaneous-install}
\begin{array}{rcl}
 \mathsf{localInstall}_a(\kw{immutable}(j,e))
  &=&\kw{writeLocal}(a,j,\kw{send}(e_F,\kws{computing},[e])),\\
 \mathsf{localInstall}_a(\kw{mutable}(j,e))
  &=&\kw{writeLocal}(a,j,\kw{send}(e_F,\kws{computing},[e])),\\
 \mathsf{localInstall}_a(\kw{mutable}(j,\kw{none}))
  &=&\kw{writeLocal}(a,j,\nil),\\
 \mathsf{localInstall}_a(\kw{lazyImmutable}(j,e))
  &=&\kw{writeLocal}(a,j,\nil),\\
 \mathsf{localInstall}_a(\kw{lazyMutable}(j,e))
  &=&\kw{writeLocal}(a,j,\nil).
\end{array}
\end{equation}
Suppose the declarations with eager right-hand sides in $\bar\lambda$, in
source order, are $j_1=e_1,\ldots,j_n=e_n$.  Using the installation terms of
Equation~\eqref{eq:local-simultaneous-install}, define
\begin{equation}\label{eq:local-simultaneous-expansion}
\begin{split}
 \mathsf{localSimCode}_a(\bar\lambda)=\langle{}
 &\mathsf{localInstall}_a(\lambda_1),\ldots,
   \mathsf{localInstall}_a(\lambda_m),\\
 &\kw{send}(\kw{readLocal}(a,j_1),\#\kw{resolve},\epsilon),\ldots,
  \kw{send}(\kw{readLocal}(a,j_n),\#\kw{resolve},\epsilon)\rangle.
\end{split}
\end{equation}
Using the internal sequence defined in
Equation~\eqref{eq:internal-sequence}, the group rule is
\begin{equation}\label{eq:local-simultaneous-start}
\begin{split}
 &\kw{initialize}(a,
     \kw{simultaneous}(\bar\lambda)\cdot\bar\ell,b,\tau)\\
 &\quad\longmapsto
 \kw{internalSeq}(\mathsf{localSimCode}_a(\bar\lambda),
                  \kw{initialize}(a,\bar\ell,b,\tau)).
\end{split}
\end{equation}
Equations~\eqref{eq:local-slot-operations}--\eqref{eq:local-simultaneous-start}
therefore install every future before the first resolution and choose source
order for resolution, exactly as the instance-slot expansion in
Equation~\eqref{eq:simultaneous-expansion}.

Statements are $\kw{expr}(e)$ or $\kw{return}(e)$.  Define
\begin{equation}\label{eq:body-results}
\begin{array}{rclcrcl}
 \mathsf{empty}(\kw{methodBody}(o))&=&o,&\quad&
 \mathsf{last}(\kw{methodBody}(o),v)&=&o,\\
 \mathsf{empty}(\kw{closureBody})&=&\nil,&&
 \mathsf{last}(\kw{closureBody},v)&=&v.
\end{array}
\end{equation}
For a frame stack $\pi$, define
\begin{equation}\label{eq:statement-start}
 \mathsf{start}(\pi,\kw{expr}(e))=\langle\pi,e\rangle,
 \qquad
 \mathsf{start}(\pi,\kw{return}(e))
    =\langle\pi\cdot\kw{returnFrame},e\rangle.
\end{equation}
Thus Equation~\eqref{eq:statement-start} installs the return frame before its
expression is evaluated.  Using the result functions in
Equation~\eqref{eq:body-results}, the empty-body rule is
\begin{equation}\label{eq:empty-body}
 \kw{body}(\tau,\epsilon)\longmapsto\kw{finish}(\mathsf{empty}(\tau)).
\end{equation}
A nonempty body uses Equation~\eqref{eq:statement-start} to start its first
statement and save the remainder:
\begin{equation}\label{eq:nonempty-body}
\dfrac{\mathsf{start}(\pi\cdot\kw{seq}(\tau,\bar s),s_1)
       =\langle\pi',e_1\rangle}
 {
 \langle P,H,\mathcal A\cdot\langle a,\pi\rangle,
       \kw{body}(\tau,\langle s_1,\bar s\rangle)\rangle
 \longmapsto
 \langle P,H,\mathcal A\cdot\langle a,\pi'\rangle,e_1\rangle}.
\end{equation}
On an ordinary statement value,
\begin{equation}\label{eq:statement-continuations}
\begin{array}{rcl}
 \langle\pi\cdot\kw{seq}(\tau,\epsilon),v\rangle
 &\longmapsto&\langle\pi,\kw{finish}(\mathsf{last}(\tau,v))\rangle,\\
 \mathsf{start}(\pi\cdot\kw{seq}(\tau,\bar s),s)
   =\langle\pi',e\rangle
 &\Longrightarrow&
 \langle\pi\cdot\kw{seq}(\tau,\langle s,\bar s\rangle),v\rangle
   \longmapsto\langle\pi',e\rangle.
\end{array}
\end{equation}
The abbreviated pairs above leave $P,H,\mathcal A$ and current $a$ unchanged.
The installed return frame transfers control non-locally:
\[
 \langle P,H,\mathcal A\cdot\langle a,\pi\cdot\kw{returnFrame}\rangle,v\rangle
 \longmapsto
 \langle P,H,\mathcal A\cdot\langle a,\pi\rangle,
    \kw{transfer}(\mathsf{returnTarget}_H(a),v)\rangle.
 \tag{Return-Value}
\]

\subsection{Completion and return}

For a stack containing $c$, define
\begin{equation}\label{eq:resume}
 \mathsf{resume}(\mathcal A_0\cdot\langle c,\pi_c\rangle
                    \cdot\mathcal A_1,c,v)
   =\langle\mathcal A_0\cdot\langle c,\pi_c\rangle,v\rangle.
\end{equation}
It discards dynamic stack entries above the continuation but retains their heap
records.  Normal completion uses Equation~\eqref{eq:resume}; it first saves
$c=H(a).k$, then severs the link:
\[
\dfrac{c=H(a).k\ne\nil\quad
       \mathsf{resume}(\mathcal A\cdot\langle a,\pi\rangle,c,v)
          =\langle\mathcal A',v\rangle}
 {\langle P,H,\mathcal A\cdot\langle a,\pi\rangle,\kw{finish}(v)\rangle
  \longmapsto
  \langle P,H[a.k\mapsto\nil],\mathcal A',v\rangle}
\quad(\rulelabel{Finish-Resume})
\]
If $H(a).k=\nil$, normal completion yields a terminal result rather than
\kw{cannotReturn}:
\[
 \dfrac{H(a).k=\nil}
 {\langle P,H,\mathcal A\cdot\langle a,\pi\rangle,\kw{finish}(v)\rangle
  \longmapsto
  \langle P,H[a.k\mapsto\nil],\epsilon,\kw{halt}(v)\rangle}
 \quad(\rulelabel{Finish-Halt}).
\]

Let the target of an explicit return from $a$ be
\begin{equation}\label{eq:return-target}
 \mathsf{returnTarget}_H(a)=
 \begin{cases}a&H(a).p\in\mathit{MethodId},\\H(a).h&\text{otherwise.}\end{cases}
\end{equation}
After the return expression evaluates to $v$, its frame produces
$\kw{transfer}(\mathsf{returnTarget}_H(a),v)$ using
Equation~\eqref{eq:return-target}.  Define
$\mathsf{dependents}_H(t)$ as the least set $X$ such that
\begin{equation}\label{eq:dependents}
 t\in X\quad\land\quad
 \forall x\in\mathit{ActId}.\;(H(x).k\in X\Rightarrow x\in X).
\end{equation}
Let $H'=\mathsf{markUncontinuable}(H,t)$ be $H$ updated, for every
$x\in\mathsf{dependents}_H(t)$, by
\begin{equation}\label{eq:mark-uncontinuable}
 H'(x).k=\nil,\qquad H'(x).q_a=\kw{false}.
\end{equation}
Using Equations~\eqref{eq:dependents} and~\eqref{eq:mark-uncontinuable}, for
$c=H(t).k\ne\nil$, explicit return resumes $c$ with $v$ after this marking:
\[
\dfrac{\substack{t=\mathsf{returnTarget}_H(a)\quad c=H(t).k\ne\nil\\
       H'=\mathsf{markUncontinuable}(H,t)\\
       \mathsf{resume}(\mathcal A\cdot\langle a,\pi\rangle,c,v)
          =\langle\mathcal A',v\rangle}}
 {\begin{aligned}
  &\langle H;\mathcal A\cdot\langle a,\pi\rangle;
       \kw{transfer}(t,v)\rangle_P\\[-1mm]
  &\qquad\longmapsto\langle H';\mathcal A';v\rangle_P
  \end{aligned}}
\quad(\rulelabel{Return-Transfer})
\]
If the saved $c$ is $\nil$, the machine instead produces an exception after
marking $t$ uncontinuable:
\[
\dfrac{t=\mathsf{returnTarget}_H(a)\quad H(t).k=\nil\quad
       H'=\mathsf{markUncontinuable}(H,t)}
 {\begin{aligned}
  &\langle H;\mathcal A\cdot\langle a,\pi\rangle;
       \kw{transfer}(t,v)\rangle_P\\[-1mm]
  &\qquad\longmapsto
   \langle H';\mathcal A\cdot\langle a,\pi\rangle;
       \kw{throw}(\kw{cannotReturn})\rangle_P
  \end{aligned}}
\quad(\rulelabel{Return-Cannot})
\]
This applies
both to a repeated method return and to a non-local return whose home method has
already completed.  Section~\ref{sec:exception-boundary} reifies this abrupt
term without discarding its frames; the standard exception library determines
whether signaling unwinds them or resumes at the failure point.

\section{Class construction and initialization}
\label{sec:classes}

This section separates three times that the surface notation intentionally
interleaves: elaboration creates and shares mixins, evaluation of a class
expression determines a superclass and allocates a fresh class, and a later
primary-factory invocation initializes a fresh instance.  In particular, the
superclass expression is not re-evaluated when an instance is made, while its
initializer message is not evaluated when the class is made.

\subsection{Normalized class expressions}

Elaboration gives every class body a descriptor
\begin{equation}\label{eq:body-descriptor}
 \chi_b=\langle D,n,u,s_f,\bar x,M,M_c,m_S\rangle.
\end{equation}
$D$ is the class declaration identity, $n$ its name, $u$ records whether the
superclass was implicit, $s_f,\bar x$ are the primary-factory signature, $M$
and $M_c$ are the instance- and class-side mixins, and $m_S$ is the deferred
superclass-initializer message template.  A mixin application has descriptor
\begin{equation}\label{eq:application-descriptor}
 \chi_a=\langle D,n,s_f,\bar x,M_c,i_A,m_S,m_M\rangle,
\end{equation}
where $i_A$ is the stable identity of its synthesized instance initializer and
$m_M$ is the deferred initializer message for the applied mixin.  From
Equations~\eqref{eq:body-descriptor} and~\eqref{eq:application-descriptor}, the
core forms are
\begin{equation}\label{eq:class-forms}
 \kw{classBody}(\chi_b,e_S),
 \qquad
 \kw{mixinApply}(\chi_a,e_S,e_M).
\end{equation}
Chained surface applications elaborate to left-associated
$\kw{mixinApply}$ nodes.

The surrounding-object annotation on $D$ determines
\begin{equation}\label{eq:class-enclosing-object}
 \mathsf{classEO}_{P,H}(a,D)=
 \begin{cases}
  a & P.\mathit{owner}(D)=\kw{activationOwner},\\
  H(a).o & P.\mathit{owner}(D)\in
      \{\kw{classOwner},\kw{objectLiteralOwner}\},\\
  \nil & P.\mathit{owner}(D)=\kw{topOwner}.
 \end{cases}
\end{equation}
Thus a class expression directly in a method or closure is enclosed by its
activation, while a member class is enclosed by the receiver of its synthesized
getter.  Top-level evaluation supplies $\nil$ without an activation.

Write $\mathsf{isClass}_H(C)$ when the corresponding record in
Equation~\eqref{eq:heap-records} has tag \kw{class}, and define
\begin{equation}\label{eq:factory-selector}
 \mathsf{factorySel}_H(C)=s_f
 \quad\text{when}\quad
 H(C).q=\langle i_I,s_f,\bar x,\kappa\rangle.
\end{equation}
Write $\mathsf{hasFactory}_H(C)$ when $H(C).q$ is a factory descriptor rather
than \kw{none}; ``factory descriptor'' means the form in
Equation~\eqref{eq:factory-descriptor}.  Using the selector projection in
Equation~\eqref{eq:factory-selector}, validation of an inheritance clause is
the total function
\begin{equation}\label{eq:inheritance-error}
 \mathsf{inhErr}_H(v,m)=
 \begin{cases}
  \kw{classExpected} & \neg\mathsf{isClass}_H(v),\\
  \kw{classNotInstantiable} & \mathsf{isClass}_H(v)\ \land\
       \neg\mathsf{hasFactory}_H(v),\\
  \kw{wrongFactory} & \mathsf{hasFactory}_H(v)\ \land\
       \selector(m)\ne\mathsf{factorySel}_H(v),\\
  \undef & \text{otherwise.}
 \end{cases}
\end{equation}
The following allocation is defined only for fresh $C,C_m$, a class $S$, and a
descriptor whose class-side mixin contains the synthesized public primary
factory:
\begin{equation}\label{eq:allocate-class}
\begin{aligned}
 \mathsf{allocClass}(H,C,C_m,\chi_C,M,S,eo,M_c,q)
   =H[&C_m\mapsto\kw{class}\langle
        M_c,\kw{Class},eo,\kw{Metaclass},\kw{none},\chi_C\rangle,\\[-1mm]
      &C\mapsto\kw{class}\langle M,S,eo,C_m,q,\chi_C\rangle].
\end{aligned}
\end{equation}
This makes every created class's metaclass fresh, a direct subclass of the
platform class \kw{Class}, an instance of \kw{Metaclass}, and enclosed by the
same object as its class.

Extend the base frame grammar of Equation~\eqref{eq:frame-grammar} by
\begin{equation}\label{eq:class-construction-frames}
 F ::= \cdots\mid\kw{makeBody}(\chi_b,eo)
       \mid\kw{mixLeft}(\chi_a,e_M)
       \mid\kw{mixRight}(\chi_a,S).
\end{equation}
Class-body evaluation uses Equation~\eqref{eq:class-enclosing-object} to fix the
enclosing object, then evaluates the superclass expression:
\[
\dfrac{eo=\mathsf{classEO}_{P,H}(a,D)}
 {\begin{aligned}
  &\langle H;\mathcal A\cdot\langle a,\pi\rangle;
       \kw{classBody}(\chi_b,e_S)\rangle_P\\[-1mm]
  &\quad\longmapsto
   \langle H;\mathcal A\cdot\langle a,
       \pi\cdot\kw{makeBody}(\chi_b,eo)\rangle;e_S\rangle_P
 \end{aligned}}
\quad(\rulelabel{Class-Super})
\]
Let $\iota_P(M)=\langle i_I,s_f,\bar x,\bar g,b_I\rangle$ as in
Equations~\eqref{eq:mixin-initializer} and~\eqref{eq:method-projections}.  Once $e_S$
has produced $S$, class creation uses Equation~\eqref{eq:allocate-class}:
\[
\dfrac{\substack{\mathsf{inhErr}_H(S,m_S)=\undef\\
       eo\ne\nil\ \lor\ (u=\kw{implicitSuperclass}\land S=\Object)\\
       q=\langle i_I,s_f,\bar x,\kw{bodyInitializer}(m_S)\rangle\\
       H'=\mathsf{allocClass}(H,C,C_m,\kw{expression}(D),M,S,eo,M_c,q)}}
 {\begin{aligned}
  &\langle H;\mathcal A\cdot\langle a,
       \pi\cdot\kw{makeBody}(\chi_b,eo)\rangle;S\rangle_P\\[-1mm]
  &\qquad\longmapsto
   \langle H';\mathcal A\cdot\langle a,\pi\rangle;C\rangle_P
 \end{aligned}}
\quad(\rulelabel{Class-Commit})
\]
where $C,C_m$ are fresh.  A non-class or invalid factory produces
$\kw{runtimeError}$ by the error function in
Equation~\eqref{eq:inheritance-error} and the schema below.  A top-level class with an
explicit superclass produces $\kw{runtimeError}(\kw{topSuperclass})$.
An omitted superclass elaborates to the class \Object{} and the message
$\langle\#\kw{new},\epsilon\rangle$.
For any construction frame $F$ whose associated inheritance template is $m$,
\begin{equation}\label{eq:inheritance-error-step}
\dfrac{\mathsf{inhErr}_H(v,m)=\xi\ne\undef}
 {\langle P,H,\mathcal A\cdot\langle a,\pi\cdot F\rangle,v\rangle
  \longmapsto
  \langle P,H,\mathcal A\cdot\langle a,\pi\rangle,
       \kw{runtimeError}(\xi)\rangle}.
\end{equation}
This schema applies to \kw{makeBody}, \kw{mixLeft}, and \kw{mixRight}.  The
remaining top-level error is
\begin{equation}\label{eq:top-superclass-error}
\dfrac{\mathsf{inhErr}_H(S,m_S)=\undef\quad eo=\nil\quad
       (u\ne\kw{implicitSuperclass}\ \lor\ S\ne\Object)}
 {\langle P,H,\mathcal A\cdot\langle a,
    \pi\cdot\kw{makeBody}(\chi_b,eo)\rangle,S\rangle
  \longmapsto
  \langle P,H,\mathcal A\cdot\langle a,\pi\rangle,
       \kw{runtimeError}(\kw{topSuperclass})\rangle}.
\end{equation}

\subsection{Mixin application}

Using the class forms in Equation~\eqref{eq:class-forms} and the construction
frames in Equation~\eqref{eq:class-construction-frames}, the superclass side is
evaluated before the class supplying the mixin:
\begin{equation}\label{eq:mixin-start}
 \langle P,H,\mathcal A\cdot\langle a,\pi\rangle,
   \kw{mixinApply}(\chi_a,e_S,e_M)\rangle
 \longmapsto
 \langle P,H,\mathcal A\cdot
   \langle a,\pi\cdot\kw{mixLeft}(\chi_a,e_M)\rangle,e_S\rangle,
\end{equation}
\begin{equation}\label{eq:mixin-right}
\dfrac{\mathsf{inhErr}_H(S,m_S)=\undef}
 {\langle P,H,\mathcal A\cdot
    \langle a,\pi\cdot\kw{mixLeft}(\chi_a,e_M)\rangle,S\rangle
  \longmapsto
  \langle P,H,\mathcal A\cdot
    \langle a,\pi\cdot\kw{mixRight}(\chi_a,S)\rangle,e_M\rangle}.
\end{equation}
For a class record $H(R)=\kw{class}\langle M_R,S_R,eo_R,C_{mR},q_R,K_R\rangle$ of
the form in Equation~\eqref{eq:heap-records}, commit uses
Equation~\eqref{eq:allocate-class}:
\[
\dfrac{\substack{\mathsf{inhErr}_H(R,m_M)=\undef\quad
       \mathsf{factorySel}_H(R)=s_f\\
       \iota_P(M_R)=\langle i_M,s_f,\bar y,\bar g,b_M\rangle\\
       q=\langle i_A,s_f,\bar x,\kw{appliedMixinInitializer}(m_S,m_M)\rangle\\
       H'=\mathsf{allocClass}(H,I,I_m,\kw{expression}(D),M_R,S,eo_R,M_c,q)}}
 {\begin{aligned}
  &\langle H;\mathcal A\cdot\langle a,
       \pi\cdot\kw{mixRight}(\chi_a,S)\rangle;R\rangle_P\\[-1mm]
  &\qquad\longmapsto
   \langle H';\mathcal A\cdot\langle a,\pi\rangle;I\rangle_P
 \end{aligned}}
\quad(\rulelabel{Mixin-Commit})
\]
for fresh $I,I_m$.  Hence every application shares $M_R$ and any later
reflective edit to that mixin, but has a fresh class, metaclass, superclass
link, and primary factory.  Its enclosing object is copied from $R$, not from
the activation evaluating the application.  Non-class operands and either
selector mismatch produce the corresponding runtime errors above.

\subsection{Message templates and primary factories}

A deferred message template is
\begin{equation}\label{eq:message-template-domain}
 m=\langle s,\bar e\rangle.
\end{equation}
The control term
$\kw{message}(m)$ evaluates only its arguments, left to right, using
$\kw{msgArgs}(s,\bar o,\bar e)$ frames analogous to the receiver and argument
continuations in Equations~\eqref{eq:receiver-continuations}
and~\eqref{eq:argument-continuations}:
\begin{equation}\label{eq:message-evaluation}
\begin{array}{c}
 \kw{message}(\langle s,\epsilon\rangle)
      \longmapsto\langle s,\epsilon\rangle,\\[1mm]
 \kw{message}(\langle s,\langle e_1,\bar e\rangle\rangle)
      \longmapsto
      \langle\kw{msgArgs}(s,\epsilon,\bar e),e_1\rangle,\\[1mm]
 \langle\kw{msgArgs}(s,\bar o,\langle e,\bar e\rangle),o'\rangle
      \longmapsto
      \langle\kw{msgArgs}(s,\bar o\cdot\langle o'\rangle,\bar e),e\rangle,\\[1mm]
 \langle\kw{msgArgs}(s,\bar o,\epsilon),o'\rangle
      \longmapsto\langle s,\bar o\cdot\langle o'\rangle\rangle.
\end{array}
\end{equation}
These abbreviated transitions preserve the surrounding machine components.
They explain precisely why factory parameters are unavailable while $e_S$ is
evaluated but are available when $m_S$ is evaluated.

Elaboration installs in $M_c$ a public method with selector $s_f$, formals
$\bar x$, and body equivalent to
\begin{equation}\label{eq:primary-factory-body}
 \kw{return}\bigl(\kw{newInstanceCurrent}
   (\langle s_f,\bar x\rangle)\bigr),
\end{equation}
where each binding occurrence in $\bar x$ elaborates to its activation send.
Thus primary factories participate in ordinary public method lookup;
$\kw{newInstanceCurrent}$ is an internal primitive occurring only in this
synthesized body; it obtains the class from the factory activation's current
instance.  Rule \rulelabel{Factory-New} below is reached from the synthesized
body in Equation~\eqref{eq:primary-factory-body}.

Define the root-to-leaf physical layout by
\begin{equation}\label{eq:slot-layout}
 \mathsf{layout}_{P,H}(\Top)=\epsilon,
 \qquad
 \mathsf{layout}_{P,H}(C)=
   \mathsf{layout}_{P,H}(\super_H(C))\cdot
   \mathsf{ownSlots}_P(\mixin_H(C)).
\end{equation}
Physical slot identities make the concatenation disjoint even when source
names override one another.  For
$H(C).q=\langle i_I,s_f,\bar x,\kappa\rangle$ and
$\mu=\langle s_f,\bar o\rangle$, define, with fresh $a_I$,
\begin{equation}\label{eq:instance-initializer-activation}
 \mathsf{newI}(H,a_I,C,o,\mu,k)=
 H[a_I\mapsto\kw{activation}\langle
 i_I,[\bar x\mapsto\bar o],\varnothing,o,C,k,a_I,\kw{true},
 \{i_I\mapsto a_I\}\rangle].
\end{equation}
If the current factory activation is $a$ and its return frame is already
installed, instance allocation uses Equation~\eqref{eq:slot-layout} and
initializer activation allocation uses
Equation~\eqref{eq:instance-initializer-activation}:
\[
\dfrac{\substack{H(a).o=C\quad
       H(C).q=\langle i_I,s_f,\bar x,\kappa\rangle\\
       \mu=\langle s_f,\bar o\rangle\quad |\bar x|=|\bar o|\\
       \sigma_0=[\mathsf{layout}_{P,H}(C)\mapsto\nil]\\
       H_0=H[o\mapsto\kw{object}\langle C,\sigma_0,\varnothing\rangle]\quad
       H'=\mathsf{newI}(H_0,a_I,C,o,\mu,a)}}
 {\begin{aligned}
  &\langle H;\mathcal A\cdot\langle a,
       \pi\cdot\kw{returnFrame}\rangle;\kw{newInstanceCurrent}(\mu)\rangle_P\\[-1mm]
  &\quad\longmapsto
   \langle H';\mathcal A\cdot\langle a,\pi\cdot\kw{returnFrame}\rangle
                    \cdot\langle a_I,\epsilon\rangle;
       \kw{initClass}(C,o,\mu)\rangle_P
 \end{aligned}}
\quad(\rulelabel{Factory-New})
\]
where $o,a_I$ are fresh.  All inherited and direct physical slots therefore
contain $\nil$ before any superclass initializer runs.  Normal completion of
$a_I$ resumes the factory's return frame with $o$; the existing return rule then
returns $o$ from the primary factory.

\subsection{Superclass and mixin initialization}

Extend the frame grammar with the initialization frames
\begin{equation}\label{eq:init-frames}
\begin{split}
 F::=\cdots&\mid\kw{superMsg}(C,o,\mu)
 \mid\kw{afterSuper}(C,o,\mu)\\
 &\mid\kw{mixinMsg}(C,o)
 \mid\kw{afterMixin}(C,o).
\end{split}
\end{equation}
The transitions below use the frames of Equation~\eqref{eq:init-frames}.  If
$\super_H(C)=\Top$, necessarily $C=\Object$ and initialization proceeds
directly to this mixin.  Otherwise the deferred superclass message is evaluated
by Equation~\eqref{eq:message-evaluation} in the new initializer activation:
\begin{equation}\label{eq:init-class}
\begin{array}{c}
\dfrac{\super_H(C)=\Top}
 {\langle P,H,\mathcal A\cdot\langle a,\pi\rangle,
    \kw{initClass}(C,o,\mu)\rangle
  \longmapsto
  \langle P,H,\mathcal A\cdot\langle a,\pi\rangle,
    \kw{ownInit}(C,o,\mu)\rangle},\\[4mm]
\dfrac{\super_H(C)=S\ne\Top\quad H(C).q=\langle i,s,\bar x,\kappa\rangle
       \quad \mathsf{superMsg}(\kappa)=m_S}
 {\langle P,H,\mathcal A\cdot\langle a,\pi\rangle,
    \kw{initClass}(C,o,\mu)\rangle
  \longmapsto
  \langle P,H,\mathcal A\cdot
    \langle a,\pi\cdot\kw{superMsg}(C,o,\mu)\rangle,
    \kw{message}(m_S)\rangle}.
\end{array}
\end{equation}
After that message becomes $\mu_S$, the superclass initializer is invoked on
the same object, but with current class $S$, using the activation constructor
of Equation~\eqref{eq:instance-initializer-activation}:
\[
\dfrac{\substack{S=\super_H(C)\quad
       \selector(\mu_S)=\mathsf{factorySel}_H(S)\\
       H'=\mathsf{newI}(H,a_S,S,o,\mu_S,a)}}
 {\begin{aligned}
  &\langle H;\mathcal A\cdot\langle a,
       \pi\cdot\kw{superMsg}(C,o,\mu)\rangle;\mu_S\rangle_P\\[-1mm]
  &\quad\longmapsto
   \langle H';\mathcal A\cdot
       \langle a,\pi\cdot\kw{afterSuper}(C,o,\mu)\rangle
       \cdot\langle a_S,\epsilon\rangle;
       \kw{initClass}(S,o,\mu_S)\rangle_P
 \end{aligned}}
\quad(\rulelabel{Init-Super})
\]
with fresh $a_S$.  Arity mismatch produces $\kw{wrongFactory}$.  When $a_S$
finishes, the ordinary
completion rule resumes the saved frame, and
\begin{equation}\label{eq:after-super}
 \langle\pi\cdot\kw{afterSuper}(C,o,\mu),o\rangle
 \longmapsto\langle\pi,\kw{ownInit}(C,o,\mu)\rangle.
\end{equation}

For a body class, its own initializer runs in that same activation, using the
slot-group form of Equation~\eqref{eq:slot-group}:
\begin{equation}\label{eq:body-own-init}
\dfrac{H(C).q=\langle i,s,\bar x,\kw{bodyInitializer}(m_S)\rangle\quad
       \iota_P(\mixin_H(C))=\langle i,s,\bar x,\bar g,b_I\rangle}
 {\kw{ownInit}(C,o,\mu)
  \longmapsto\kw{initGroups}(o,\bar g,b_I)}.
\end{equation}
For an applied mixin, $m_M$ is first evaluated in the coordinator activation:
\begin{equation}\label{eq:applied-own-init}
\dfrac{H(C).q=\langle i_A,s,\bar x,\kw{appliedMixinInitializer}(m_S,m_M)\rangle}
 {\langle\pi,\kw{ownInit}(C,o,\mu)\rangle
  \longmapsto
  \langle\pi\cdot\kw{mixinMsg}(C,o),\kw{message}(m_M)\rangle}.
\end{equation}
Suppose $M=\mixin_H(C)$.  Using the initializer form and mixin projection in
Equations~\eqref{eq:mixin-initializer} and~\eqref{eq:method-projections}, write
its initializer as
$\iota_P(M)=\langle i_M,s_M,\bar x_M,\bar g,b_M\rangle$.  Define
\begin{equation}\label{eq:mixin-initializer-activation}
\begin{aligned}
 &\mathsf{newMixin}(H,a_M,M,C,o,
                    \langle s_M,\bar o\rangle,k)\\
 &\qquad=H[a_M\mapsto\kw{activation}\langle
 i_M,[\bar x_M\mapsto\bar o],\varnothing,o,C,k,a_M,\kw{true},
 \{i_M\mapsto a_M\}\rangle].
\end{aligned}
\end{equation}
Then, using the activation allocated by
Equation~\eqref{eq:mixin-initializer-activation},
\[
\dfrac{\substack{\mu_M=\langle s_M,\bar o\rangle\quad
       |\bar x_M|=|\bar o|\\
       H'=\mathsf{newMixin}(H,a_M,M,C,o,\mu_M,a)}}
 {\begin{aligned}
  &\langle H;\mathcal A\cdot\langle a,
       \pi\cdot\kw{mixinMsg}(C,o)\rangle;\mu_M\rangle_P\\[-1mm]
  &\quad\longmapsto
   \langle H';\mathcal A\cdot
       \langle a,\pi\cdot\kw{afterMixin}(C,o)\rangle
       \cdot\langle a_M,\epsilon\rangle;
       \kw{initGroups}(o,\bar g,b_M)\rangle_P
 \end{aligned}}
\quad(\rulelabel{Init-Applied-Mixin})
\]
for fresh $a_M$.  Its current class is the new application $C$, not the class
from which $M$ was obtained.  Finally,
\begin{equation}\label{eq:after-mixin}
 \langle\pi\cdot\kw{afterMixin}(C,o),o\rangle
 \longmapsto\langle\pi,\kw{finish}(o)\rangle.
\end{equation}
This yields the specified order: superclass message, superclass initializer,
applied-mixin message if any, mixin slots and body, and factory return.

\subsection{Slot groups}
\label{sec:slot-groups}

Normalize an initializer group as
\begin{equation}\label{eq:slot-group}
\begin{aligned}
 g&::=\kw{sequential}(\bar d)\mid\kw{simultaneous}(\bar d),\\
 d&::=\kw{immutable}(j,e,i_t)\mid\kw{mutable}(j,e,i_t)
       \mid\kw{mutable}(j,\kw{none}).
\end{aligned}
\end{equation}
For an eager declaration, $i_t$ is the stable activation-declaration identity
assigned by elaboration to the zero-argument closure used when that declaration
occurs in a simultaneous group.  Sequential initialization ignores $i_t$.
Extend the frame grammar of Equation~\eqref{eq:frame-grammar} by
\begin{equation}\label{eq:slot-frames}
 F::=\cdots\mid\kw{slotWrite}(o,j)
   \mid\kw{slotStore}(o,j,\bar d,\bar g,b)
   \mid\kw{internal}(\bar e,t)
   \mid\kw{lazyStore}(o,j)
   \mid\kw{nestedStore}(o,N).
\end{equation}
The internal $\kw{readSlot}(o,j)$ and $\kw{writeSlot}(o,j,e)$ operations use
physical identity $j$ and do not send accessors.  They are emitted only by
elaboration of the declaration inducing $j$.  Their rules are
\begin{equation}\label{eq:slot-operations}
\begin{array}{c}
\dfrac{H(o).\sigma(j)=v}
 {\langle P,H,\mathcal A\cdot\langle a,\pi\rangle,
    \kw{readSlot}(o,j)\rangle
  \longmapsto
  \langle P,H,\mathcal A\cdot\langle a,\pi\rangle,v\rangle},\\[4mm]
 \langle P,H,\mathcal A\cdot\langle a,\pi\rangle,
    \kw{writeSlot}(o,j,e)\rangle
 \longmapsto
 \langle P,H,\mathcal A\cdot
    \langle a,\pi\cdot\kw{slotWrite}(o,j)\rangle,e\rangle,\\[3mm]
 \langle P,H,\mathcal A\cdot
    \langle a,\pi\cdot\kw{slotWrite}(o,j)\rangle,v\rangle
 \longmapsto
 \langle P,H[o.\sigma(j)\mapsto v],
    \mathcal A\cdot\langle a,\pi\rangle,v\rangle.
\end{array}
\end{equation}
Empty initialization enters the initializer body:
\begin{equation}\label{eq:init-groups-empty}
 \kw{initGroups}(o,\epsilon,b)
   \longmapsto\kw{body}(\kw{methodBody}(o),b).
\end{equation}
Groups are taken in source order.  A sequential group is replaced by its
declaration state machine:
\begin{equation}\label{eq:init-group-sequential}
 \kw{initGroups}(o,\kw{sequential}(\bar d)\cdot\bar g,b)
 \longmapsto\kw{initSeq}(o,\bar d,\bar g,b).
\end{equation}
An uninitialized mutable declaration advances without a write, because every
slot was already set to $\nil$ by \rulelabel{Factory-New}.  A declaration with
an initializer evaluates it before writing:
\begin{equation}\label{eq:init-sequential}
\begin{array}{c}
 \kw{initSeq}(o,\kw{mutable}(j,\kw{none})\cdot\bar d,\bar g,b)
 \longmapsto\kw{initSeq}(o,\bar d,\bar g,b),\\[1mm]
\dfrac{d\in\{\kw{immutable}(j,e,i_t),\kw{mutable}(j,e,i_t)\}}
 {\langle\pi,\kw{initSeq}(o,d\cdot\bar d,\bar g,b)\rangle
  \longmapsto
  \langle\pi\cdot\kw{slotStore}(o,j,\bar d,\bar g,b),e\rangle},\\[1mm]
 \langle H;\mathcal A\cdot\langle a,
    \pi\cdot\kw{slotStore}(o,j,\bar d,\bar g,b)\rangle;v\rangle_P
 \longmapsto
 \langle H[o.\sigma(j)\mapsto v];\mathcal A\cdot\langle a,\pi\rangle;
    \kw{initSeq}(o,\bar d,\bar g,b)\rangle_P,\\[1mm]
 \kw{initSeq}(o,\epsilon,\bar g,b)
 \longmapsto\kw{initGroups}(o,\bar g,b).
\end{array}
\end{equation}

Simultaneous groups are specified by an elaboration into the ordinary
\NS{} library protocol.  If the declarations with right-hand sides, in source
order, are
$j_1=e_1\ @\ i_{t_1},\ldots,j_n=e_n\ @\ i_{t_n}$, let $e_F$ be the annotated
core send that
denotes \PastFuture{} in this initializer's lexical context.  Using the direct
slot operations of Equation~\eqref{eq:slot-operations}, define
\begin{equation}\label{eq:simultaneous-expansion}
\begin{split}
 \mathsf{simCode}_a(o,\bar d)=\langle{}
 &\kw{writeSlot}(o,j_1,
      \kw{send}(e_F,\kws{computing},[e_1]_{i_{t_1}})),\ldots,\\
 &\kw{writeSlot}(o,j_n,
      \kw{send}(e_F,\kws{computing},[e_n]_{i_{t_n}})),\\
 &\kw{send}(\kw{readSlot}(o,j_1),\#\kw{resolve},\epsilon),\ldots,
  \kw{send}(\kw{readSlot}(o,j_n),\#\kw{resolve},\epsilon)\rangle,
\end{split}
\end{equation}
where $[e_i]_{i_{t_i}}$ denotes the zero-argument closure with declaration
identity $i_{t_i}$ created in activation $a$, and
uninitialized mutable declarations contribute no term.  Then, using the
expansion of Equation~\eqref{eq:simultaneous-expansion},
\begin{equation}\label{eq:init-group-simultaneous}
 \kw{initGroups}(o,\kw{simultaneous}(\bar d)\cdot\bar g,b)
 \longmapsto
 \kw{internalSeq}(\mathsf{simCode}_a(o,\bar d),
                  \kw{initGroups}(o,\bar g,b)).
\end{equation}
$\kw{internalSeq}$ in Equation~\eqref{eq:init-group-simultaneous} is not a new
evaluation mechanism.  It is defined by
\begin{equation}\label{eq:internal-sequence}
\begin{array}{rcl}
 \kw{internalSeq}(\epsilon,t)&\longmapsto&t,\\
 \langle\pi,\kw{internalSeq}(e\cdot\bar e,t)\rangle
   &\longmapsto&
   \langle\pi\cdot\kw{internal}(\bar e,t),e\rangle,\\
 \langle\pi\cdot\kw{internal}(\bar e,t),v\rangle
   &\longmapsto&\langle\pi,\kw{internalSeq}(\bar e,t)\rangle.
\end{array}
\end{equation}
As elsewhere, the pair notation leaves $P,H,\mathcal A$, and the current
activation unchanged.  By Equation~\eqref{eq:internal-sequence}, consequently
every future is installed before the first is resolved, and all
evaluation, forwarding, recursive resolution, and failure behavior is that of
the specified \PastFuture{} class under the ordinary send rules.  The draft
chooses source order for the final \kw{resolve} sends; whether the Specification
intends this order to be observable is recorded as an explicit question in the
coverage ledger.

\subsection{Synthesized slot and nested-class methods}
\label{sec:synthesized-methods}

Ordinary immutable and mutable getters elaborate to
$\kw{return}(\kw{currentRead}(j))$.  The activation-relative primitives reduce
to the direct slot operations of Equation~\eqref{eq:slot-operations}:
\begin{equation}\label{eq:current-slot-operations}
\begin{gathered}
\dfrac{H(a).o=o}
 {\kw{currentRead}(j)\longmapsto\kw{readSlot}(o,j)},
\\[1.1ex]
\dfrac{H(a).o=o}
 {\kw{currentWrite}(j,e)\longmapsto\kw{writeSlot}(o,j,e)}.
\end{gathered}
\end{equation}
The abbreviated rules preserve the surrounding configuration.  A mutable
setter invokes \kw{currentWrite} with slot identifier $j$ and argument $x$, then
returns \kw{self}.  Getter and setter both retain the declared access level.
Thus initialization never invokes an overridable accessor.

A lazy declaration also allocates physical slot $j$, initially $\nil$, and its
getter body contains $\kw{lazyRead}(j,e)$.  In a getter activation whose current
instance is $o$, the $\kw{lazyStore}$ frame declared in
Equation~\eqref{eq:slot-frames} completes the direct write:
\begin{equation}\label{eq:lazy-read}
\begin{array}{c}
\dfrac{H(a).o=o\quad H(o).\sigma(j)=v\ne\nil}
 {\langle P,H,\mathcal A\cdot\langle a,\pi\rangle,
    \kw{lazyRead}(j,e)\rangle
  \longmapsto
  \langle P,H,\mathcal A\cdot\langle a,\pi\rangle,v\rangle},\\[4mm]
\dfrac{H(a).o=o\quad H(o).\sigma(j)=\nil}
 {\langle P,H,\mathcal A\cdot\langle a,\pi\rangle,
    \kw{lazyRead}(j,e)\rangle
  \longmapsto
  \langle P,H,\mathcal A\cdot
    \langle a,\pi\cdot\kw{lazyStore}(o,j)\rangle,e\rangle},\\[4mm]
 \langle P,H,\mathcal A\cdot
    \langle a,\pi\cdot\kw{lazyStore}(o,j)\rangle,v\rangle
 \longmapsto
 \langle P,H[o.\sigma(j)\mapsto v],
    \mathcal A\cdot\langle a,\pi\rangle,v\rangle.
\end{array}
\end{equation}
There is deliberately no in-progress sentinel: a recursive read while the slot
is still $\nil$ starts another evaluation, and an initializer returning $\nil$
causes the next read to evaluate it again.  A mutable lazy setter is the same
direct setter as Equation~\eqref{eq:current-slot-operations} and may reset the
slot to $\nil$.

Finally, a nested declaration $N$ induces a getter whose body contains
$\kw{nestedClass}(N,E_C)$, where $E_C$ is one of the normalized class forms in
Equation~\eqref{eq:class-forms}.  For
current receiver $o$, the $\kw{nestedStore}$ frame from
Equation~\eqref{eq:slot-frames} records a newly evaluated class:
\begin{equation}\label{eq:nested-class-cache}
\begin{array}{c}
\dfrac{H(a).o=o\quad H(o).\eta(N)=C}
 {\langle H;\mathcal A\cdot\langle a,\pi\rangle;
    \kw{nestedClass}(N,E_C)\rangle_P
  \longmapsto
  \langle H;\mathcal A\cdot\langle a,\pi\rangle;C\rangle_P},\\[4mm]
\dfrac{H(a).o=o\quad N\notin\dom(H(o).\eta)}
 {\begin{aligned}
  &\langle H;\mathcal A\cdot\langle a,\pi\rangle;
    \kw{nestedClass}(N,E_C)\rangle_P\\[-1mm]
  \longmapsto
  &\quad\langle H;\mathcal A\cdot
    \langle a,\pi\cdot\kw{nestedStore}(o,N)\rangle;
    \kw{classWithEO}(E_C,o)\rangle_P
  \end{aligned}},\\[4mm]
 \langle H;\mathcal A\cdot
    \langle a,\pi\cdot\kw{nestedStore}(o,N)\rangle;C\rangle_P
 \longmapsto
 \langle H[o.\eta(N)\mapsto C];
    \mathcal A\cdot\langle a,\pi\rangle;C\rangle_P.
\end{array}
\end{equation}
$\kw{classWithEO}$ used by Equation~\eqref{eq:nested-class-cache} is precisely
\begin{equation}\label{eq:class-with-enclosing-object}
\begin{aligned}
 &\langle\pi,\kw{classWithEO}
       (\kw{classBody}(\chi_b,e_S),o)\rangle
   \longmapsto\langle\pi\cdot\kw{makeBody}(\chi_b,o),e_S\rangle,\\
 &\kw{classWithEO}(\kw{mixinApply}(\chi_a,e_S,e_M),o)
   \longmapsto\kw{mixinApply}(\chi_a,e_S,e_M).
\end{aligned}
\end{equation}
The second clause of Equation~\eqref{eq:class-with-enclosing-object} ignores the
forced object because the enclosing object of a
mixin application is, by definition, copied from its mixin-source class.  The
cache belongs to the receiver, so two
instances obtain distinct nested classes, while repeated sends to one receiver
return the identical class even if inputs used by its superclass expression
subsequently change.

\section{Remaining sequential language}
\label{sec:remaining-sequential}

This section closes the sequential surface-language layer.  Its rules do not
add special cases to dispatch: compound literals elaborate to sends, an object
literal reuses class and factory construction, and top-level execution merely
supplies a distinguished activation.  The fixed platform entries introduced
below are the current specified meaning of literals; they are not implicit
receivers and consequently cannot be shadowed by user declarations.

\subsection{Atomic and tuple literals}
\label{sec:literals}

An atomic literal payload is
\begin{equation}\label{eq:atom-payload}
\begin{split}
 \omega::={}&\kw{integer}(n)\mid\kw{fixed}(q)\mid
 \kw{boolean}(b)\mid\kw{nilAtom}\mid
 \kw{string}(u)\mid\kw{symbol}(u),
\end{split}
\end{equation}
where $n\in\mathbb Z$, $q\in\mathbb Q$, $b$ is a mathematical Boolean, and
$u$ is a Unicode-normal-form-C string.  The decoder
for a digit character $c$ is $\mathsf{dig}(c)\in\{0,\ldots,35\}$, with letters
denoting 10 through 35.  If the integer digits are
$u_0\ldots u_{n-1}$, the fractional digits are $f_1\ldots f_m$, the radix is
$r$ (10 when omitted), the sign is $\varepsilon\in\{-1,1\}$, and the decimal
exponent is $z$ (zero when omitted), define
\begin{equation}\label{eq:number-decoding}
 \mathsf{num}(r,\varepsilon,u,f,z)=
 \varepsilon\left(
   \sum_{i=0}^{n-1}\mathsf{dig}(u_i)r^{n-1-i}+
   \sum_{i=1}^{m}\mathsf{dig}(f_i)r^{-i}
 \right)10^z.
\end{equation}
Every digit must be smaller than $r$.  A numeral with neither a fraction nor
an exponent elaborates to
$\kw{atom}(\kw{integer}(\mathsf{num}(r,\varepsilon,u,\epsilon,0)))$;
every other numeral elaborates to the corresponding $\kw{fixed}$ payload.
Thus the exponent is decimal even for a radix numeral, matching the specified
grammar and parser behavior, and integer magnitude is unbounded.

The static map $P.\mathit{atom}$ of
Equation~\eqref{eq:program-definition} is canonical:
\begin{equation}\label{eq:atom-canonicality}
 P.\mathit{atom}(\omega_1)=P.\mathit{atom}(\omega_2)
 \quad\Longleftrightarrow\quad \omega_1=\omega_2.
\end{equation}
The initial heap contains all referenced atomic objects with the platform
classes required by their payloads.  String decoding applies Unicode NFC
before Equation~\eqref{eq:atom-canonicality}; symbol spelling is decoded by
the symbol grammar and canonicalized in the same way.  Boolean and nil syntax
maps to the three distinguished payloads.  Specification 0.109 removes both
character-literal syntax and the \kw{Character} platform class; neither occurs
in the semantic domains.  The atomic evaluation rule is
\[
 \langle P,H,\mathcal A\cdot\langle a,\pi\rangle,
          \kw{atom}(\omega)\rangle
 \longmapsto
 \langle P,H,\mathcal A\cdot\langle a,\pi\rangle,
          P.\mathit{atom}(\omega)\rangle.
 \tag{Atom}
\]
Numeric, Boolean, nil, string, and symbol atoms are value objects.
The mapping is an implementation freedom only as to representation: it must
preserve the exact payload, the specified class protocol, and the identity
condition in Equation~\eqref{eq:atom-canonicality}.

Tuple literals use the fixed platform bindings
\begin{equation}\label{eq:tuple-platform-bindings}
 R=P.\mathit{platform}(\kw{ReadOnlyTuple}),\qquad
 B=P.\mathit{platform}(\kw{Array}).
\end{equation}
Let $\underline n$ abbreviate the integer atom for $n$.  Using the bindings in
Equation~\eqref{eq:tuple-platform-bindings}, their specified expansion is
\begin{equation}\label{eq:tuple-expansion}
\begin{array}{rcl}
 \mathsf{tupleElab}_P(\epsilon)
  &=&\kw{send}(R,\kws{new},\langle\underline0\rangle),\\[1mm]
 \mathsf{tupleElab}_P(\langle e_1,\ldots,e_n\rangle)
  &=&\kw{send}\bigl(R,\kws{fromArray},
       \langle\mathsf{cascade}(A_n,\bar m_n)\rangle\bigr),\\
 A_n&=&\kw{send}(B,\kws{new},\langle\underline n\rangle),\\
 \bar m_n&=&\langle
   \langle\selsym{at:put:},\langle\underline1,e_1\rangle\rangle,\ldots,
   \langle\selsym{at:put:},\langle\underline n,e_n\rangle\rangle,
   \langle\selsym{yourself},\epsilon\rangle\rangle.
\end{array}
\end{equation}
Here $\mathsf{cascade}$ is exactly the single-receiver surface cascade whose
closure elaboration is fixed in Section~\ref{sec:surface}.  Hence $A_n$ is
evaluated once and the $e_i$ are evaluated in source order by the ordinary
send rules of Section~\ref{sec:machine}.  Formally,
\begin{equation}\label{eq:tuple-step}
 \kw{tuple}(\bar e)\longmapsto\mathsf{tupleElab}_P(\bar e).
\end{equation}
The surrounding configuration is unchanged.  Tuple objects are shallowly
immutable; unlike the atomic literals above, the Specification does not make
them deeply immutable merely by virtue of being literal tuples.

\subsection{Pattern literals}
\label{sec:patterns}

Pattern literals are syntax-directed library expressions.  Let
\begin{equation}\label{eq:pattern-binding}
 \mathsf{Pattern}_d=
   \kw{implicit}(\#\kw{Pattern},\epsilon,d),
\end{equation}
where $d$ is the lexical annotation selected during elaboration as in
Equation~\eqref{eq:core-expressions}.  For a wildcard, a literal, and a nested
pattern, respectively, define
\begin{equation}\label{eq:pattern-components}
\begin{array}{rcl}
 \mathsf{pat}_P(\_,d)&=&
   \kw{send}(\mathsf{Pattern}_d,\#\kw{wildcard},\epsilon),\\
 \mathsf{pat}_P(l,d)&=&
   \kw{send}(\mathsf{Pattern}_d,\kws{literal},\langle l\rangle),\\
 \mathsf{pat}_P(\langle p\rangle,d)&=&\kw{pattern}(p,d).
\end{array}
\end{equation}
For keyword pairs $k_1p_1\ldots k_np_n$, let $\bar k$ be the tuple of symbol
atoms $\#k_1,\ldots,\#k_n$ and let $\bar p$ be the tuple of the corresponding
component expansions from Equation~\eqref{eq:pattern-components}.  Then
\begin{equation}\label{eq:keyword-pattern-expansion}
\begin{split}
 &\mathsf{pat}_P(k_1p_1\ldots k_np_n,d)=\\[-1mm]
 &\quad\kw{send}\bigl(\mathsf{Pattern}_d,
      \selsym{keywords:patterns:},
      \langle\kw{tuple}(\bar k),\kw{tuple}(\bar p)\rangle\bigr).
\end{split}
\end{equation}
The tuples in Equation~\eqref{eq:keyword-pattern-expansion} are ordered
collections; a conforming Pattern library may copy them into its preferred
list representation.  Pattern evaluation is now only the expansion
\begin{equation}\label{eq:pattern-step}
 \kw{pattern}(p,d)\longmapsto\mathsf{pat}_P(p,d),
\end{equation}
after which ordinary lookup, access control, and argument order apply.  In
particular, rebinding \kw{Pattern} changes the meaning of the literal, and a
pattern literal is therefore not a top-level-independent literal.

The Specification gives syntax for a variable component $?x$ but gives no
expression that constructs it and no required representation beyond its role
in a later binding object.  For a variable occurrence at site $z$, this draft
therefore uses the explicit elaboration parameter
\begin{equation}\label{eq:variable-pattern-parameter}
 \mathsf{pat}_P(?x_z,d)=P.\mathit{patternVariable}(z).
\end{equation}
Equation~\eqref{eq:variable-pattern-parameter} is not implementation latitude
for the other pattern forms: it isolates the one missing normative clause.
Once that constructor protocol is specified, the map is replaced by its
ordinary-send expansion.  The current library's \kw{NamedPattern} suggests an
implementation, but does not determine language semantics.

\subsection{Object literals}
\label{sec:object-literals}

Elaboration assigns an object literal the descriptor
\begin{equation}\label{eq:object-literal-descriptor}
 \chi_o=\langle L,u,M,M_c,m_S\rangle.
\end{equation}
$L$ is its stable object-literal declaration, $u$ records whether its
superclass clause was implicit, $M$ is the shared body mixin, $M_c$ is a
synthetic class-side mixin containing a public nullary \kw{new} factory, and
$m_S$ is the deferred superclass initializer message.  Its own initializer
has the normalized form
$\iota_P(M)=\langle i_I,\#\kw{new},\epsilon,\bar g,b_I\rangle$ from
Equations~\eqref{eq:mixin-initializer} and~\eqref{eq:method-projections}.
Thus each evaluation will allocate a class,
but every evaluation of the same $L$ applies the identical $M$.

The current activation itself encloses a non-top-level object literal.  A
top-level activation is deliberately excluded because top-level expressions
have no enclosing object:
\begin{equation}\label{eq:object-literal-enclosing-object}
 \mathsf{literalEO}_H(a)=
 \begin{cases}
   \nil & H(a).p=\kw{top},\\
   a & \text{otherwise.}
 \end{cases}
\end{equation}
Extend the frame grammar in Equation~\eqref{eq:frame-grammar} by
\begin{equation}\label{eq:object-literal-frame}
 F::=\cdots\mid\kw{makeObject}(\chi_o,eo).
\end{equation}
Evaluation first fixes the enclosing object using
Equation~\eqref{eq:object-literal-enclosing-object}, then evaluates the
normalized superclass expression:
\[
\dfrac{eo=\mathsf{literalEO}_H(a)}
 {\begin{aligned}
  &\langle H;\mathcal A\cdot\langle a,\pi\rangle;
       \kw{objectLiteral}(\chi_o,e_S)\rangle_P\\[-1mm]
  &\quad\longmapsto
   \langle H;\mathcal A\cdot\langle a,
       \pi\cdot\kw{makeObject}(\chi_o,eo)\rangle;e_S\rangle_P
 \end{aligned}}
\quad(\rulelabel{Object-Super})
\]
If $eo=\nil$, omission elaborates $e_S$ to the fixed
$P.\mathit{platform}(\kw{Object})$.  Otherwise omission elaborates $e_S$ to
the annotated implicit send \kw{Object}; hence a surrounding module may
rebind the default.  In both cases omission supplies the deferred message
$\langle\#\kw{new},\epsilon\rangle$.

After a valid superclass result, class allocation uses
Equation~\eqref{eq:allocate-class}; evaluation immediately continues as an
ordinary send of \kw{new} to that fresh class:
\[
\dfrac{\substack{\mathsf{inhErr}_H(S,m_S)=\undef\\
       eo\ne\nil\ \lor\ (u=\kw{implicitSuperclass}\land
          S=P.\mathit{platform}(\kw{Object}))\\
       \iota_P(M)=\langle i_I,\#\kw{new},\epsilon,\bar g,b_I\rangle\\
       q=\langle i_I,\#\kw{new},\epsilon,\kw{bodyInitializer}(m_S)\rangle\\
       H'=\mathsf{allocClass}(H,C,C_m,\kw{expression}(L),M,S,eo,M_c,q)}}
 {\begin{aligned}
  &\langle H;\mathcal A\cdot\langle a,
       \pi\cdot\kw{makeObject}(\chi_o,eo)\rangle;S\rangle_P\\[-1mm]
  &\qquad\longmapsto
   \langle H';\mathcal A\cdot\langle a,\pi\rangle;
       \kw{dispatch}(\kw{ordinaryRequest}(C),
          \langle\#\kw{new},\epsilon\rangle)\rangle_P
 \end{aligned}}
\quad(\rulelabel{Object-Commit})
\]
where $C,C_m$ are fresh.  The ordinary dispatch reaches the synthesized
factory body in Equation~\eqref{eq:primary-factory-body}, so
\rulelabel{Factory-New} allocates the fresh result object and all superclass
and mixin initialization rules apply unchanged.  Consequently the result
object, its implicit class, and that class's metaclass are fresh on every
evaluation, while the mixin named by $L$ is shared.

The invalid-inheritance schema for object literals is
\begin{equation}\label{eq:object-inheritance-error}
\dfrac{\mathsf{inhErr}_H(v,m_S)=\xi\ne\undef}
 {\langle P,H,\mathcal A\cdot\langle a,
       \pi\cdot\kw{makeObject}(\chi_o,eo)\rangle,v\rangle
  \longmapsto
  \langle P,H,\mathcal A\cdot\langle a,\pi\rangle,
       \kw{runtimeError}(\xi)\rangle}.
\end{equation}
When the inheritance itself is valid but a top-level object literal has an
explicit superclass, the corresponding rule is
\begin{equation}\label{eq:object-top-superclass-error}
\dfrac{\mathsf{inhErr}_H(S,m_S)=\undef\quad eo=\nil\quad
       (u\ne\kw{implicitSuperclass}\ \lor\
        S\ne P.\mathit{platform}(\kw{Object}))}
 {\langle P,H,\mathcal A\cdot\langle a,
       \pi\cdot\kw{makeObject}(\chi_o,eo)\rangle,S\rangle
  \longmapsto
  \langle P,H,\mathcal A\cdot\langle a,\pi\rangle,
       \kw{runtimeError}(\kw{topSuperclass})\rangle}.
\end{equation}

\subsection{Modules and value status}
\label{sec:modules}

The creation provenance retained by every class record in
Equation~\eqref{eq:heap-records} makes the module distinction stable and
heap-local.  Define
\begin{equation}\label{eq:module-definition}
\begin{split}
 \mathsf{moduleDefinition}_{P,H}(C)\quad\Longleftrightarrow\quad{}
 &\mathsf{isClass}_H(C)\ \land\ \mathsf{hasFactory}_H(C)\ \land\\[-1mm]
 &H(C).\chi_C=\kw{expression}(D)\ \land\
   P.\mathit{owner}(D)=\kw{topOwner}.
\end{split}
\end{equation}
The factory test excludes the metaclass allocated at the same declaration.
Equation~\eqref{eq:module-definition} includes a top-level mixin application:
its class retains the application declaration $D$ even though it shares the
source class's mixin.  A module instance is exactly
\begin{equation}\label{eq:module-instance}
 \mathsf{module}_{P,H}(o)\quad\Longleftrightarrow\quad
 \mathsf{moduleDefinition}_{P,H}(\classof_H(o)).
\end{equation}
Every module definition satisfies the Specification's primitive value-object
clause, irrespective of reflective mutability.  This does not make its module
instances values.  More generally, base-level value status is the greatest
fixed point containing the specified atomic objects and module definitions and
closed under the structural value conditions (initialized immutable slots,
value contents and enclosing objects, inheritance from \kw{Value}, and no
direct identity-method override).  Reflection may alter behavior without
altering this base-level classification; its causal effect is deferred to the
reflective layer.

\subsection{Top-level programs, compilation units, and entry}
\label{sec:top-level}

The least top-level-expression predicate is
\begin{equation}\label{eq:top-level-expression}
\begin{array}{@{}rcl@{}}
 \mathsf{TL}_P(\kw{atom}(\omega))&\Leftarrow&\kw{true},\\
 \mathsf{TL}_P(\kw{closureLiteral}(q,\bar x,\ell,b))&\Leftarrow&\kw{true},\\
 \mathsf{TL}_P(\kw{tuple}(\bar e))&\Leftarrow&
   \forall e_i\in\bar e.\ \mathsf{TL}_P(e_i),\\
 \mathsf{TL}_P(\kw{objectLiteral}(\chi_o,e_S))
   &\Leftarrow&\chi_o=\langle L,u,M,M_c,m_S\rangle\ \land\
     P.\mathit{owner}(L)=\kw{topOwner},\\
 \mathsf{TL}_P(E_C)&\Leftarrow&
   E_C\text{ is a class form induced by }D\ \land\
   P.\mathit{owner}(D)=\kw{topOwner},\\
 \mathsf{TL}_P(\kw{send}(e,s,\bar e))&\Leftarrow&
   \mathsf{TL}_P(e)\ \land\
   \forall e_i\in\bar e.\ \mathsf{TL}_P(e_i),\\
 \mathsf{TL}_P((e))&\Leftarrow&\mathsf{TL}_P(e).
\end{array}
\end{equation}
There are no other clauses.  Thus implicit, self, outer, super, and eventual
sends are rejected at top level.  Pattern literals are also excluded because
their expansion in Equation~\eqref{eq:pattern-binding} depends on an implicit
send.  The ownership premises exclude a class
or object literal lexically nested in another body, even if its AST is
considered in isolation.  Parentheses disappear after this judgment.

A top-level activation has no parameters or slots and has nil current
instance, current class, continuation, and home.  With fresh $a_T$, define
\begin{equation}\label{eq:new-top-activation}
 \mathsf{newTop}(H,a_T)=H[a_T\mapsto\kw{activation}\langle
 \kw{top},\varnothing,\varnothing,\nil,\nil,\nil,\nil,
 \kw{true},\varnothing\rangle].
\end{equation}
For a well-formed source program, execution begins at
\begin{equation}\label{eq:run-program}
 \mathsf{run}(P,H,e)=
 \langle P,\mathsf{newTop}(H,a_T),
   \epsilon\cdot\langle a_T,\epsilon\rangle,e\rangle
 \quad\text{if }\mathsf{TL}_P(e),
\end{equation}
where $a_T$ is fresh.  A value with no pending top-level frame halts:
\[
\dfrac{H(a).p=\kw{top}}
 {\langle P,H,\epsilon\cdot\langle a,\epsilon\rangle,v\rangle
  \longmapsto
  \langle P,H[a.q_a\mapsto\kw{false}],\epsilon,
       \kw{halt}(v)\rangle}.
\quad(\rulelabel{Top-Halt})
\]
Closures created during Equation~\eqref{eq:run-program} capture nil current
instance and class through the ordinary closure-creation rule; the top
activation remains their lexical activation record.

A compilation unit is source text paired with a language id:
\begin{equation}\label{eq:compilation-unit}
 U=\langle\lambda,u\rangle
 \in\mathit{LanguageId}\times\mathit{SourceText}.
\end{equation}
Section~\ref{sec:front-end} defines PEG recognition, language-specific
elaboration, compilation by reified Newspeak objects, and the resulting
$\langle P,e\rangle$.  The current implementation's restriction of
Equation~\eqref{eq:compilation-unit} to one top-level class is a temporary
accepted-syntax restriction, not a different dynamic semantics.  Serialization
must preserve the stable identities in $P$, the referenced atomic values, and
all reachable heap records.

Finally, if $p=P.\mathit{platform}(\kw{platform})$ is the supplied platform
object and $r$ is the host's argument-list object, application entry is the
top-level expression
\begin{equation}\label{eq:application-entry}
 \mathsf{entry}(o,p,r)=
   \kw{send}(o,\selsym{main:args:},\langle p,r\rangle).
\end{equation}
It is admitted exactly when
\begin{equation}\label{eq:application-object}
 \mathsf{application}_{P,H}(o)\quad\Longleftrightarrow\quad
 \exists f,C_f.\ P,H\vdash
 \lookup_{pub}(\selsym{main:args:},\classof_H(o))=\langle f,C_f\rangle
 \ \land\ |f.\bar x|=2.
\end{equation}
Execution of Equation~\eqref{eq:application-entry} is then
Equation~\eqref{eq:run-program}.  Linking and module instantiation performed
by that method are ordinary sends.

\subsection{The exception-library boundary}
\label{sec:exception-boundary}

Newspeak has no exception declaration, throw expression, catch expression, or
handler stack in its core syntax.  Signaling, catching, passing, returning,
and resuming are library operations implemented with ordinary objects and
reflective access to activations and continuations.  Adding built-in handler
rules here would therefore duplicate, and potentially contradict, that
library semantics.

The VM can nevertheless detect failures such as those produced by
Equations~\eqref{eq:inheritance-error-step},
\eqref{eq:top-superclass-error}, \eqref{eq:object-inheritance-error}, and the
failed return in \rulelabel{Return-Cannot}.  The platform supplies a total
reification operation
\begin{equation}\label{eq:exception-reification}
 \mathsf{reifyError}_P(H,\xi)=\langle H',x\rangle,
\end{equation}
where $H'$ extends $H$ without changing an existing record, $x$ is a near
exception object representing $\xi$, and ordinary public dispatch of
$\langle\#\kw{signal},\epsilon\rangle$ to $x$ is defined.  Reification may
allocate a fresh exception and its immutable diagnostic data; it may not
invoke user code.

Both internal control terms cross the boundary in the same way:
\[
\dfrac{\mathsf{reifyError}_P(H,\xi)=\langle H',x\rangle}
 {\begin{aligned}
  &\langle H;\mathcal A\cdot\langle a,\pi\rangle;
       \kw{runtimeError}(\xi)\rangle_P\\[-1mm]
  &\quad\longmapsto
   \langle H';\mathcal A\cdot\langle a,\pi\rangle;
       \kw{dispatch}(\kw{ordinaryRequest}(x),
          \langle\#\kw{signal},\epsilon\rangle)\rangle_P
 \end{aligned}}
\quad(\rulelabel{Runtime-Error-Signal})
\]
\[
\dfrac{\mathsf{reifyError}_P(H,\xi)=\langle H',x\rangle}
 {\begin{aligned}
  &\langle H;\mathcal A\cdot\langle a,\pi\rangle;
       \kw{throw}(\xi)\rangle_P\\[-1mm]
  &\quad\longmapsto
   \langle H';\mathcal A\cdot\langle a,\pi\rangle;
       \kw{dispatch}(\kw{ordinaryRequest}(x),
          \langle\#\kw{signal},\epsilon\rangle)\rangle_P
 \end{aligned}}
\quad(\rulelabel{Throw-Signal})
\]
The original frames $\pi$ remain in place.  If the library resumes the
exception, the \kw{signal} send returns a value at the failure point; if it
handles, passes, or returns from it, its ordinary sends and reflective
continuation operations determine the transfer.  In particular, no base rule
silently discards sequence frames.  The activation and stack-mirror operations
in Equations~\eqref{eq:actor-stack-snapshot},
\eqref{eq:replace-current-actor-stack}, and
\eqref{eq:replace-resume-paused-actor-stack} provide the control primitives
used by the standard library; they do not add a second notion of exception
handling.

\section{Actors and eventual references}
\label{sec:actors}

This section lifts the sequential relation of Section~\ref{sec:machine} to a
system of isolated, single-threaded actors.  The lift does not make an
asynchronous send an ordinary send in disguise: evaluation returns a promise
immediately, the request becomes a mailbox item, and ordinary dispatch begins
only when the target actor starts a later turn.

\subsection{Global domains and isolation}

Let the following identity sets be pairwise disjoint and disjoint from those in
Equation~\eqref{eq:runtime-identities}:
\begin{equation}\label{eq:eventual-identities}
 p\in\mathit{PromiseId},\qquad r\in\mathit{FarRefId},\qquad
 u\in\mathit{EventId},\qquad
 \mathit{Ref}=\mathit{ObjRef}\uplus\mathit{PromiseId}\uplus
                 \mathit{FarRefId}.
\end{equation}
Promises and far references are Newspeak objects with platform-supplied
classes, but their routing state is recorded globally below.  An ordinary send
to such a handle invokes its local platform protocol; only
$\kw{asyncSend}$ dereferences it for remote delivery.

An actor has a private heap $H_A$.  A platform may additionally maintain a
value heap $V$ shared by actors managed by one virtual machine.  The heap visible to the
sequential machine running actor $A$ is
\begin{equation}\label{eq:actor-heap-view}
 \widehat H_A=V\uplus H_A,\qquad
 A\ne B\Longrightarrow
 \dom(H_A)\cap\dom(H_B)=\varnothing,\qquad
 \dom(V)\cap\dom(H_A)=\varnothing.
\end{equation}
Only objects satisfying the value predicate of Section~\ref{sec:modules} may
enter $V$.  A distributed implementation may instead copy their complete
immutable reachable graph.  We write
\begin{equation}\label{eq:value-transfer-contract}
 \Omega\vdash\mathsf{valTransfer}_{A\Rightarrow B}(x)
   \Downarrow\langle\Omega',y\rangle
\end{equation}
for either permitted choice; it is defined only for a fully initialized value
$x$, preserves its class and immutable reachable graph up to identity
isomorphism, and neither aliases nor modifies mutable actor-local state.
Sharing is the identity case.  The reflective layer will distinguish actors in
one VM, which may share values and observe common code change, from actors in
different VMs, which must use the copying case.

A far-reference table and promise table have the forms
\begin{equation}\label{eq:eventual-tables}
\begin{aligned}
 \Phi(r)&=\langle A,o\rangle,\\
 \Pi(p)&::=\kw{pending}(A,W)
       \mid\kw{fulfilled}(A,x)\mid\kw{broken}(A,x),\\
 W&::=\epsilon\mid W\mathbin{\cdot}
       \langle B,\mu,p_r,\kappa\rangle.
\end{aligned}
\end{equation}
$\Phi(r)=\langle A,o\rangle$ means that $r$ designates $o$ in $H_A$.
The actor in a promise state is the actor relative to whose heap its result is
represented.  Pending sends are retained in registration order; a waiter
$\langle B,\mu,p_r,\kappa\rangle$ requests that, after settlement, actor $B$
route $\mu$ to the result and settle $p_r$.  The causal history $\kappa$ is
defined below.

Remote representation is an actor-relative relation
$\Omega\vdash\mathsf{rr}_{A\Rightarrow B}(x)
\Downarrow\langle\Omega',y\rangle$.  It is the least relation given by the
following rules:
\[
\dfrac{A=B}
 {\Omega\vdash\mathsf{rr}_{A\Rightarrow B}(x)
   \Downarrow\langle\Omega,x\rangle}
\quad(\rulelabel{RR-Same})
\]
\[
\dfrac{A\ne B\quad \mathsf{value}_{P,\widehat H_A}(x)\quad
       \Omega\vdash\mathsf{valTransfer}_{A\Rightarrow B}(x)
         \Downarrow\langle\Omega',y\rangle}
 {\Omega\vdash\mathsf{rr}_{A\Rightarrow B}(x)
   \Downarrow\langle\Omega',y\rangle}
\quad(\rulelabel{RR-Value})
\]
\[
\dfrac{A\ne B\quad \Phi(r)=\langle B,o\rangle}
 {\Omega\vdash\mathsf{rr}_{A\Rightarrow B}(r)
   \Downarrow\langle\Omega,o\rangle}
\quad(\rulelabel{RR-Far-Home})
\qquad
\dfrac{A\ne B\quad \Phi(r)=\langle C,o\rangle\quad C\ne B}
 {\Omega\vdash\mathsf{rr}_{A\Rightarrow B}(r)
   \Downarrow\langle\Omega,r\rangle}
\quad(\rulelabel{RR-Far-Pass})
\]
\[
\dfrac{A\ne B\quad p\in\dom(\Pi)}
 {\Omega\vdash\mathsf{rr}_{A\Rightarrow B}(p)
   \Downarrow\langle\Omega,p\rangle}
\quad(\rulelabel{RR-Promise})
\]
\[
\dfrac{A\ne B\quad o\in\dom(H_A)\quad
       \neg\mathsf{value}_{P,\widehat H_A}(o)\quad r\text{ fresh}\quad
       \Omega'=\Omega[\Phi(r)\mapsto\langle A,o\rangle]}
 {\Omega\vdash\mathsf{rr}_{A\Rightarrow B}(o)
   \Downarrow\langle\Omega',r\rangle}
\quad(\rulelabel{RR-Near-Far})
\]
Thus a far reference passed back to its target actor becomes the near referent,
and a far reference passed to a third actor does not acquire a proxy chain.
Promise identities are monotonic eventual capabilities and likewise pass
without a proxy layer.  Applying remote representation componentwise defines
message transport:
\begin{equation}\label{eq:message-transfer}
 \dfrac{\Omega_0=\Omega\quad
   \Omega_{i-1}\vdash\mathsf{rr}_{A\Rightarrow B}(o_i)
      \Downarrow\langle\Omega_i,o_i'\rangle\ (1\le i\le n)}
 {\Omega\vdash\mathsf{msgRR}_{A\Rightarrow B}
      (\langle s,\langle o_1,\ldots,o_n\rangle\rangle)
   \Downarrow
      \langle\Omega_n,\langle s,\langle o_1',\ldots,o_n'\rangle\rangle\rangle}.
\end{equation}
The selector is a symbol value and is therefore unchanged.

\subsection{Actor configurations and causal delivery}

A mailbox or network packet is
\begin{equation}\label{eq:packet-grammar}
\begin{split}
 b::={}&\kw{app}\langle u,A,B,o,\mu,p,\kappa\rangle\\
  \mid{}&\kw{settlementPacket}\langle u,A,B,p,\chi,x,\kappa\rangle\\
  \mid{}&\kw{wake}\langle u,A,B,p,p_r,\chi,x,\mu,\kappa_w,\kappa\rangle,
 \qquad \chi\in\{\kw{ok},\kw{error}\}.
\end{split}
\end{equation}
The first two actor fields are source and destination.  Every packet has a
fresh event identity $u$ and carries the source actor's causally prior event
set $\kappa$.  A wake names both the settled promise $p$ that caused the wake
and the dependent result promise $p_r$; it also retains $\kappa_w$, the
history at which its dependent send was originally issued.

An actor is idle, is executing one sequential configuration on behalf of an
application packet, or is suspended at a debugger safepoint:
\begin{equation}\label{eq:actor-run-state}
 \begin{aligned}
 \mathcal X(A)::={}&\kw{idle}\mid
   \kw{runningTurn}\langle\mathcal A,t,p,u\rangle\mid\\
 &\kw{pausedTurn}\langle\mathcal A,t,p,u,r_s,d_s\rangle.
 \end{aligned}
\end{equation}
Here $d_s\in\mathit{StopReason}$.  The ordinary actor rules have no premise
for \kw{pausedTurn}; deliveries may still append to its mailbox, but the actor
executes no turn step until an authorized debugger operation resumes it.
Section~\ref{sec:debugging} defines suspension, stack replacement, and
resumption.
The global state is
\begin{equation}\label{eq:actor-configuration}
 \Omega=\langle P,V,\mathcal H,\mathcal X,\mathcal Q,\mathcal N,
              \Pi,\Phi,\Gamma,\mathcal E,\mathcal D,w,n\rangle,
 \qquad w\in\mathit{VMId},\quad n\in\mathbb N.
\end{equation}
$w$ identifies the VM cohort and $n$ is its committed program version.  Every
actor in $\dom(\mathcal H)$ is managed by $w$ and observes the same $P$ and
$V$.  Section~\ref{sec:reflection} changes $P$ and $n$ atomically.  A system of
multiple VMs is a map from distinct $w$ to configurations of this form; its
program components evolve independently.
$\Gamma_\Omega$, $\Pi_\Omega$, and analogous subscripts name a component of a
particular configuration when more than one configuration occurs in a rule;
the subscript is omitted otherwise.
$\mathcal Q(A)$ is a FIFO mailbox, $\mathcal N$ is the multiset of undelivered
cross-actor packets, and $\Gamma(A)$ is a finite downward-closed causal
history.  The permanent event ledger records
\begin{equation}\label{eq:event-ledger}
 \mathcal E(u)=\langle A,B,\kappa\rangle,\qquad
 \mathsf{dst}_{\Omega}(u)=B,\qquad
 \mathcal D(B)\in\mathit{EventId}^{*}.
\end{equation}
$\mathcal D(B)$ lists, in arrival order, events delivered to $B$; a failed
network event is marked retired at the same position.  Processed event records
are retained because later causal histories may still mention them.

All application and system packets are emitted by the helper
\begin{equation}\label{eq:emit-packet}
\begin{split}
 &\mathsf{emit}(\Omega,A,B,\beta,\kappa)=\Omega_2,\\
 &u\text{ fresh},\quad b=\beta(u,A,B,\kappa),\quad
 \Omega_1=\Omega[\mathcal E(u)\mapsto\langle A,B,\kappa\rangle,
                  \Gamma(A)\mapsto\Gamma(A)\cup\{u\}],\\
 &\Omega_2=
 \begin{cases}
  \Omega_1[\mathcal Q(B)\mapsto\mathcal Q(B)\cdot b,
           \mathcal D(B)\mapsto\mathcal D(B)\cdot u] & A=B,\\
  \Omega_1[\mathcal N\mapsto\mathcal N\uplus\{b\}] & A\ne B.
 \end{cases}
\end{split}
\end{equation}
Here $\beta$ is a packet template supplying every field except
$u,A,B,\kappa$.  A same-actor asynchronous send arrives immediately at the
tail of its own mailbox, but cannot start until the current turn finishes.

A network packet $b$ with identity $u$, destination $B$, and history $\kappa$
is eligible precisely when
\begin{equation}\label{eq:e-order-eligibility}
 \mathsf{eligible}_{\Omega}(b)\quad\Longleftrightarrow\quad
 \{u'\in\kappa\mid\mathsf{dst}_{\Omega}(u')=B\}
       \subseteq\mathsf{set}(\mathcal D(B)).
\end{equation}
An eligible packet may arrive:
\[
\dfrac{b\in\mathcal N\quad \mathsf{id}(b)=u\quad
       \mathsf{destination}(b)=B\quad
       \mathsf{eligible}_{\Omega}(b)}
 {\Omega\longrightarrow_{\mathsf{act}}
  \Omega[\mathcal N\mapsto\mathcal N\setminus\{b\},
         \mathcal Q(B)\mapsto\mathcal Q(B)\cdot b,
         \mathcal D(B)\mapsto\mathcal D(B)\cdot u]}
\quad(\rulelabel{Network-Deliver})
\]
Because mailboxes are FIFO and an actor runs at most one turn, delivery
in this order implies processing in this order.

Equation~\eqref{eq:e-order-eligibility} is a constructive form of E-order.
After $A_1$ emits $u_1$ to $A_2$, $u_1\in\Gamma(A_1)$, so a later $u_2$ from
$A_1$ to $A_2$ carries $u_1$ and cannot arrive first.  If the later packet
instead passes a far reference for $A_2$ to $A_3$, delivery merges its history
into $\Gamma(A_3)$; a subsequent send by $A_3$ to that referent therefore also
carries $u_1$ and cannot arrive at $A_2$ first.  These are exactly the two
ordering clauses in Specification \S5.6, generalized by transitive causal
history.

Maximal executions are required to be weakly fair:
\begin{equation}\label{eq:actor-fairness}
 \mathsf{continuouslyEnabled}(\alpha,i)
   \Longrightarrow\exists j\ge i.\ \mathsf{taken}(\alpha,j),
\end{equation}
for packet delivery, dequeue of the head of a nonempty idle mailbox, and an
enabled step of a running actor.  Fairness does not rescue an actor whose
current turn diverges; its later mailbox entries may consequently remain
unprocessed.

\subsection{Asynchronous-send evaluation and routing}

Extend the request and frame grammars in
Equations~\eqref{eq:request-grammar} and~\eqref{eq:frame-grammar} by
\begin{equation}\label{eq:async-frames}
 \theta::=\cdots\mid\kw{eventualRequest}(x),\qquad
 F::=\cdots\mid\kw{asyncRecv}(s,\bar e).
\end{equation}
The receiver is evaluated first:
\[
 \langle P,H,\mathcal A\cdot\langle a,\pi\rangle,
        \kw{asyncSend}(e,s,\bar e)\rangle
 \longmapsto
 \langle P,H,\mathcal A\cdot
        \langle a,\pi\cdot\kw{asyncRecv}(s,\bar e)\rangle,e\rangle
 \quad(\rulelabel{Async-Receiver}).
\]
After it yields $x$, the existing argument frames of
Equation~\eqref{eq:argument-continuations} evaluate the arguments from left to
right:
\begin{equation}\label{eq:async-receiver-continuations}
\begin{array}{c}
 \langle\pi\cdot\kw{asyncRecv}(s,\epsilon),x\rangle
  \longmapsto
 \langle\pi,\kw{dispatch}(\kw{eventualRequest}(x),
                          \langle s,\epsilon\rangle)\rangle,\\[2mm]
 \langle\pi\cdot\kw{asyncRecv}(s,\langle e_1,\bar e\rangle),x\rangle
  \longmapsto
 \langle\pi\cdot\kw{args}(\kw{eventualRequest}(x),s,\epsilon,\bar e),e_1\rangle.
\end{array}
\end{equation}

For a non-promise receiver, define its target by
\begin{equation}\label{eq:eventual-target}
 \mathsf{target}_{\Omega,A}(x)=
 \begin{cases}
  \langle A,x\rangle & x\in\dom(\widehat H_A),\\
  \langle B,o\rangle & \Phi(x)=\langle B,o\rangle.
 \end{cases}
\end{equation}
The routing judgment
$\Omega,A\vdash\mathsf{route}(x,\mu,p_r,\kappa)\Downarrow\Omega'$ is
defined by the following four cases.  A known target transports the arguments
and emits an application packet:
{\small
\[
\dfrac{\begin{gathered}
 \mathsf{target}_{\Omega,A}(x)=\langle B,o\rangle\\
 \Omega\vdash\mathsf{msgRR}_{A\Rightarrow B}(\mu)
    \Downarrow\langle\Omega_1,\mu'\rangle\\
 \Omega'=\mathsf{emit}(\Omega_1,A,B,
    \lambda u,A,B,\kappa.\kw{app}\langle
       u,A,B,o,\mu',p_r,\kappa\rangle,\kappa)
\end{gathered}}
 {\Omega,A\vdash\mathsf{route}(x,\mu,p_r,\kappa)\Downarrow\Omega'}
\quad(\rulelabel{Route-Target})
\]}
An unresolved promise records the send without blocking its actor:
\[
\dfrac{\Pi(p)=\kw{pending}(C,W)\quad
       \Omega'=\Omega[\Pi(p)\mapsto
          \kw{pending}(C,W\cdot\langle A,\mu,p_r,\kappa\rangle)]}
 {\Omega,A\vdash\mathsf{route}(p,\mu,p_r,\kappa)\Downarrow\Omega'}
\quad(\rulelabel{Route-Pending})
\]
For a fulfilled promise, its result is first represented in the requesting
actor and routing repeats:
{\small\[
\dfrac{\Pi(p)=\kw{fulfilled}(C,x)\quad
       \Omega\vdash\mathsf{rr}_{C\Rightarrow A}(x)
          \Downarrow\langle\Omega_1,y\rangle\quad
       \Omega_1,A\vdash\mathsf{route}(y,\mu,p_r,\kappa)
          \Downarrow\Omega'}
 {\Omega,A\vdash\mathsf{route}(p,\mu,p_r,\kappa)\Downarrow\Omega'}
\quad(\rulelabel{Route-Fulfilled})
\]}
For a broken promise, the result promise is broken with the same represented
exception:
{\small\[
\dfrac{\begin{gathered}\Pi(p)=\kw{broken}(C,x)\\
       \Omega\vdash\mathsf{rr}_{C\Rightarrow A}(x)
          \Downarrow\langle\Omega_1,y\rangle\\
       \Omega_1,A\vdash\mathsf{settlePromise}(p_r,\kw{error},A,y,\kappa)
          \Downarrow\Omega'\end{gathered}}
 {\Omega,A\vdash\mathsf{route}(p,\mu,p_r,\kappa)\Downarrow\Omega'}
\quad(\rulelabel{Route-Broken})
\]}
The last rule makes explicit the otherwise unspecified meaning of sends to a
broken promise: break propagation is the conventional eventual-send behavior,
but is recorded as a language-design choice in the coverage ledger.

Routing is an actor-level action.  The asynchronous expression itself
allocates its result promise and immediately leaves that promise as the value
of the expression:
{\small
\[
\dfrac{\begin{gathered}
 \mathcal X(A)=\kw{runningTurn}\langle\mathcal A\cdot\langle a,\pi\rangle,
   \kw{dispatch}(\kw{eventualRequest}(x),\mu),p_0,u_0\rangle\\
 p_r\text{ fresh}\qquad
 \Omega_0=\Omega[\Pi(p_r)\mapsto\kw{pending}(A,\epsilon)]\\
 \Omega_0,A\vdash\mathsf{route}(x,\mu,p_r,\Gamma(A))\Downarrow\Omega_1
\end{gathered}}
 {\begin{aligned}
  \Omega\longrightarrow_{\mathsf{act}}\Omega_1[\mathcal X(A)\mapsto{}&
   \kw{runningTurn}\langle\mathcal A\cdot\langle a,\pi\rangle,\\[-1mm]
   &p_r,p_0,u_0\rangle]
 \end{aligned}}
\quad(\rulelabel{Async-Return})
\]}
The displayed run record in the conclusion abbreviates replacement of its
control component by $p_r$; its reply promise and triggering event remain
$p_0,u_0$.

\subsection{Turns and promise settlement}

When an application packet reaches the head of an idle actor's mailbox, it
starts a turn.  With fresh top activation $a_T$ and
$H' =\mathsf{newTop}(\widehat H_B,a_T)$ from
Equation~\eqref{eq:new-top-activation}:
{\small
\[
\dfrac{\mathcal X(B)=\kw{idle}\quad
 \mathcal Q(B)=\kw{app}\langle u,A,B,o,\mu,p,\kappa\rangle\cdot Q}
 {\begin{aligned}
  \Omega\longrightarrow_{\mathsf{act}}\Omega[&\mathcal X(B)\mapsto
    \kw{runningTurn}\langle\epsilon\cdot\langle a_T,\epsilon\rangle,
       \kw{dispatch}(\kw{ordinaryRequest}(o),\mu),p,u\rangle,\\[-1mm]
   &\mathcal Q(B)\mapsto Q,
    \Gamma(B)\mapsto\Gamma(B)\cup\kappa\cup\{u\},
    \mathcal H(B)\mapsto H'\setminus V]
 \end{aligned}}
\quad(\rulelabel{Turn-Start})
\]}
While $B$ is running, each applicable sequential step is lifted pointwise:
\[
\dfrac{\langle P,\widehat H_B,\mathcal A,t\rangle
          \longmapsto\langle P,H',\mathcal A',t'\rangle}
 {\Omega\longrightarrow_{\mathsf{act}}
  \Omega[\mathcal H(B)\mapsto H'\setminus V,
         \mathcal X(B)\mapsto
          \kw{runningTurn}\langle\mathcal A',t',p,u\rangle]}
\quad(\rulelabel{Actor-Sequential})
\]
It is inapplicable to the eventual dispatch intercepted by
\rulelabel{Async-Return}.  No second packet can start in $B$ until this turn
ends.

Settlement wakes every recorded dependent send.  Define the ordered helper
{\small
\begin{equation}\label{eq:wake-all}
\begin{aligned}
 \mathsf{wakeAll}(\Omega,B,p,\chi,y,\epsilon)&=\Omega,\\
 \mathsf{wakeAll}(\Omega,B,p,\chi,y,
    \langle C,\mu,p_r,\kappa_w\rangle\cdot W)
   &=\mathsf{wakeAll}(\Omega_1,B,p,\chi,y,W),\\
 \Omega_1&=\mathsf{emit}(\Omega,B,C,
   \lambda u,B,C,\kappa.\kw{wake}\langle
      u,B,C,p,p_r,\chi,y,\mu,\kappa_w,\kappa\rangle,\Gamma(B)).
\end{aligned}
\end{equation}}
Earlier waiters for the same actor are emitted first and therefore ordered by
Equation~\eqref{eq:e-order-eligibility}.  Promise settlement is the partial
judgment
{\small
\begin{equation}\label{eq:promise-settlement}
\dfrac{\begin{gathered}
 \Pi(p)=\kw{pending}(B,W)\qquad \kappa\subseteq\Gamma(B)\\
 \Omega\vdash\mathsf{rr}_{A\Rightarrow B}(x)
    \Downarrow\langle\Omega_0,y\rangle\\
 \Omega_1=
 \begin{cases}
  \Omega_0[\Pi(p)\mapsto\kw{fulfilled}(B,y)]&\chi=\kw{ok},\\
  \Omega_0[\Pi(p)\mapsto\kw{broken}(B,y)]&\chi=\kw{error}
 \end{cases}
 \qquad \Omega'=\mathsf{wakeAll}(\Omega_1,B,p,\chi,y,W)
\end{gathered}}
 {\Omega,B\vdash\mathsf{settlePromise}(p,\chi,A,x,\kappa)
    \Downarrow\Omega'}.
\end{equation}}
There is no rule for settling an already settled promise, so fulfillment and
breaking are mutually exclusive and occur at most once.

When a turn returns normally, a promise owned by the same actor settles
immediately; a remote owner's settlement is sent subsequently:
{\small\[
\dfrac{\begin{gathered}
 \mathcal X(B)=\kw{runningTurn}\langle\epsilon,\kw{halt}(v),p,u\rangle
 \qquad \Pi(p)=\kw{pending}(B,W)\\
 \Omega,B\vdash\mathsf{settlePromise}(p,\kw{ok},B,v,\Gamma(B))
    \Downarrow\Omega'
\end{gathered}}
 {\Omega\longrightarrow_{\mathsf{act}}
  \Omega'[\mathcal X(B)\mapsto\kw{idle}]}
\quad(\rulelabel{Turn-Fulfill-Local})
\]}
{\small\[
\dfrac{\begin{gathered}
 \mathcal X(B)=\kw{runningTurn}\langle\epsilon,\kw{halt}(v),p,u\rangle\\
 \Pi(p)=\kw{pending}(A,W)\qquad A\ne B\\
 \Omega'=\mathsf{emit}(\Omega,B,A,
  \lambda u,B,A,\kappa.\kw{settlementPacket}\langle
     u,B,A,p,\kw{ok},v,\kappa\rangle,\Gamma(B))
\end{gathered}}
 {\Omega\longrightarrow_{\mathsf{act}}
  \Omega'[\mathcal X(B)\mapsto\kw{idle}]}
\quad(\rulelabel{Turn-Fulfill-Remote})
\]}
If exception signaling escapes the turn's top activation, the
exception-library boundary produces $\kw{abort}(x)$ instead of
$\kw{halt}(v)$.  A local owner's promise is broken immediately:
{\small\[
\dfrac{\begin{gathered}
 \mathcal X(B)=\kw{runningTurn}\langle\epsilon,\kw{abort}(x),p,u\rangle
 \qquad \Pi(p)=\kw{pending}(B,W)\\
 \Omega,B\vdash\mathsf{settlePromise}(p,\kw{error},B,x,\Gamma(B))
    \Downarrow\Omega'
\end{gathered}}
 {\Omega\longrightarrow_{\mathsf{act}}
  \Omega'[\mathcal X(B)\mapsto\kw{idle}]}
\quad(\rulelabel{Turn-Break-Local})
\]}
A remote owner receives an error settlement subsequently:
{\small\[
\dfrac{\begin{gathered}
 \mathcal X(B)=\kw{runningTurn}\langle\epsilon,\kw{abort}(x),p,u\rangle\\
 \Pi(p)=\kw{pending}(A,W)\qquad A\ne B\\
 \Omega'=\mathsf{emit}(\Omega,B,A,
  \lambda u,B,A,\kappa.\kw{settlementPacket}\langle
     u,B,A,p,\kw{error},x,\kappa\rangle,\Gamma(B))
\end{gathered}}
 {\Omega\longrightarrow_{\mathsf{act}}
  \Omega'[\mathcal X(B)\mapsto\kw{idle}]}
\quad(\rulelabel{Turn-Break-Remote})
\]}
These rules are the only additional obligation placed on the resumable
exception library by actors.

An idle promise owner processes a settlement packet at the head of its FIFO
mailbox by merging the packet history and applying
Equation~\eqref{eq:promise-settlement}:
{\small\[
\dfrac{\begin{gathered}
 \mathcal X(B)=\kw{idle}\qquad
 \mathcal Q(B)=\kw{settlementPacket}\langle
    u,A,B,p,\chi,x,\kappa\rangle\cdot Q\\
 \Omega_0=\Omega[\mathcal Q(B)\mapsto Q,
                  \Gamma(B)\mapsto\Gamma(B)\cup\kappa\cup\{u\}]\\
 \Omega_0,B\vdash\mathsf{settlePromise}(p,\chi,A,x,\Gamma_0(B))
    \Downarrow\Omega'
\end{gathered}}
 {\Omega\longrightarrow_{\mathsf{act}}\Omega'}
\quad(\rulelabel{Settlement-Packet-Dequeue})
\]}
A fulfilled wake represents the result in the waiting actor and resumes the
stored routing operation; a broken wake settles the dependent promise as
broken:
{\small\[
\dfrac{\begin{gathered}
 \mathcal X(C)=\kw{idle}\qquad
 \mathcal Q(C)=\kw{wake}\langle
  u,B,C,p,p_r,\kw{ok},x,\mu,\kappa_w,\kappa\rangle\cdot Q\\
 \Omega_0=\Omega[\mathcal Q(C)\mapsto Q,
   \Gamma(C)\mapsto\Gamma(C)\cup\kappa\cup\kappa_w\cup\{u\}]\\
 \Omega_0\vdash\mathsf{rr}_{B\Rightarrow C}(x)
    \Downarrow\langle\Omega_1,y\rangle\\
 \Omega_1,C\vdash\mathsf{route}(y,\mu,p_r,\Gamma_1(C))
    \Downarrow\Omega'
\end{gathered}}
 {\Omega\longrightarrow_{\mathsf{act}}\Omega'}
\quad(\rulelabel{Wake-Fulfilled})
\]}
{\small\[
\dfrac{\begin{gathered}
 \mathcal X(C)=\kw{idle}\qquad
 \mathcal Q(C)=\kw{wake}\langle
  u,B,C,p,p_r,\kw{error},x,\mu,\kappa_w,\kappa\rangle\cdot Q\\
 \Omega_0=\Omega[\mathcal Q(C)\mapsto Q,
   \Gamma(C)\mapsto\Gamma(C)\cup\kappa\cup\kappa_w\cup\{u\}]\\
 \Omega_0,C\vdash\mathsf{settlePromise}(p_r,\kw{error},B,x,\Gamma_0(C))
    \Downarrow\Omega'
\end{gathered}}
 {\Omega\longrightarrow_{\mathsf{act}}\Omega'}
\quad(\rulelabel{Wake-Broken})
\]}
These system packets are handled only between user turns, so remote settlement
cannot mutate an actor's heap or promise observations during an executing turn.
A communication subsystem may retire an undelivered application event and
emit an error settlement instead of delivering it.  Its failure object must
satisfy the same remote-representation boundary, the event is entered in
$\mathcal D$ as retired so it cannot block causal successors, and delivery and
failure of the same event are mutually exclusive.

\subsection{Actor creation}

Actor creation is a platform primitive rather than a new surface expression.
Elaboration supplies, for a top-level mixin, the class-side mixin and factory
descriptor needed to apply it to \Object:
\begin{equation}\label{eq:actor-seed}
 P.\mathit{actorSeed}(M)=\langle D,M_c,q\rangle
 \Longrightarrow
 \decl_P(M)=D\ \land\ P.\mathit{owner}(D)=\kw{topOwner}.
\end{equation}
Let $\mathsf{base}_P(B)$ be a fresh actor's platform heap, containing the
distinguished classes and canonical values but no user-created mutable object.
For fresh $B,C,C_m,r$, actor creation is
{\small
\begin{equation}\label{eq:create-actor}
\dfrac{\begin{gathered}
 P.\mathit{actorSeed}(M)=\langle D,M_c,q\rangle\qquad
 H_0=\mathsf{base}_P(B)\\
 H_1=\mathsf{allocClass}(H_0,C,C_m,\kw{actor}(D),M,
          P.\mathit{platform}(\kw{Object}),\nil,M_c,q)\\
 \Omega'=\Omega[\mathcal H(B)\mapsto H_1,
   \mathcal X(B)\mapsto\kw{idle},\mathcal Q(B)\mapsto\epsilon,
   \Gamma(B)\mapsto\varnothing,\\[-1mm]
 \hspace{27mm}\mathcal D(B)\mapsto\epsilon,
   \Phi(r)\mapsto\langle B,C\rangle]
\end{gathered}}
 {\Omega,A\vdash\mathsf{spawn}(M)\Downarrow\langle\Omega',r\rangle}.
\end{equation}}
The result is explicitly a far reference to the new class, as required by the
Specification.  The $\kw{actor}(D)$ provenance distinguishes that class from
evaluation of the top-level class expression itself, so it is not accidentally
classified as a module definition by Equation~\eqref{eq:module-definition}.
Its primary factory can then be invoked asynchronously in the ordinary way.


\section{Mirrors and reflective program change}
\label{sec:reflection}

This section refines the schematic transition in
Equation~\eqref{eq:reflective-step}.  It deliberately specifies a semantic
mirror protocol rather than a concrete library API: the Specification requires
mirror-based capability control, transactional change, object-class change,
and activation access, but explicitly leaves the VM-mirror selector set open.
An implementation may choose selectors freely provided that each operation
below is reached only through the corresponding mirror capability.

\subsection{VM cohorts, mirror objects, and authority}

One actor configuration in Equation~\eqref{eq:actor-configuration} is one VM
cohort.  Define its complete object store by
\begin{equation}\label{eq:vm-object-store}
 \mathsf{vmStore}(\Omega)=
 V\uplus\biguplus_{A\in\dom(\mathcal H)}\mathcal H(A).
\end{equation}
The disjointness condition in Equation~\eqref{eq:actor-heap-view} makes this a
map.  Write $\mathsf{ownerActor}_{\Omega}(x)=A$ exactly when
$x\in\dom(\mathcal H(A))$, and write
$\mathsf{ownerActor}_{\Omega}(x)=\kw{shared}$ exactly when $x\in\dom(V)$.
These two cases are disjoint.

A mirror target and a reflective right are drawn from
\begin{equation}\label{eq:mirror-targets-rights}
\begin{split}
 \zeta::={}&\kw{vmTarget}(w)\mid\kw{programTarget}
  \mid\kw{mixinTarget}(M)\mid\kw{methodTarget}(i_f)\\
 &\mid\kw{classTarget}(C)\mid\kw{objectTarget}(o)
  \mid\kw{activationTarget}(a)\mid\kw{actorTarget}(A),\\
 \varrho\in\mathit{ReflectRight}={}&
 \{\kw{inspect},\kw{editCode},\kw{editClass},\kw{reclassifyObject},\\
 &\quad\kw{editObjectSlots},\kw{editActivation},\kw{controlContinuation},\\
 &\quad\kw{controlActorStack},\\
 &\quad
       \kw{primitive},\kw{delegate}\}.
\end{split}
\end{equation}
Extend the heap-record grammar of Equation~\eqref{eq:heap-records} by
\begin{equation}\label{eq:mirror-record}
 H(g)=\kw{mirror}\langle C_g,w,\zeta,R\rangle,
 \qquad R\in\mathcal P_{\mathrm{fin}}(\mathit{ReflectRight}),
\end{equation}
where $C_g$ is a platform mirror class.  Thus $\classof_H(g)=C_g$.
Mirror identities are fresh and unforgeable.  A mirror is actor-local mutable
state, never a value; ordinary cross-actor transfer therefore yields a far
reference under Rules \rulelabel{RR-Near-Far} and \rulelabel{RR-Far-Pass}.
Possession of such a reference can still confer authority, because the
eventual request executes in the mirror's owning actor.

Target containment is the least reflexive relation $\sqsubseteq_{P,H,w}$
satisfying
\begin{equation}\label{eq:mirror-target-containment}
\begin{aligned}
 \zeta&\sqsubseteq_{P,H,w}\kw{vmTarget}(w),\\
 \kw{mixinTarget}(M)&\sqsubseteq_{P,H,w}\kw{programTarget},\\
 \kw{methodTarget}(i_f)&\sqsubseteq_{P,H,w}\kw{mixinTarget}(M)
   &&\text{if }\mathsf{methodOwner}_P(i_f)=M,\\
 \kw{classTarget}(C)&\sqsubseteq_{P,H,w}
       \kw{mixinTarget}(\mixin_H(C)),\\
 \kw{objectTarget}(o)&\sqsubseteq_{P,H,w}
       \kw{classTarget}(\classof_H(o)),\\
 \kw{activationTarget}(a)&\sqsubseteq_{P,H,w}
       \kw{objectTarget}(H(a).o).
\end{aligned}
\end{equation}
The first, generic clause places an actor target occurring in this cohort below
its VM target.  An actor target does not contain the actor's whole heap:
activation and object mirrors remain separate capabilities.  It authorizes
only the stack operations explicitly requiring \kw{controlActorStack} below.
The method owner used here is the partial function
\begin{equation}\label{eq:method-owner}
 \mathsf{methodOwner}_P(i_f)=M
 \quad\Longleftrightarrow\quad
 \exists s,f.\;\methods_P(M)(s)=f\land\mathsf{methodId}(f)=i_f.
\end{equation}
Uniqueness follows from stable method identities.  VM containment includes all
targets in Equation~\eqref{eq:mirror-targets-rights} that occur in the store of
that VM; it does not include any target belonging to another VM.

The platform may create an initial VM mirror only as part of the authority
explicitly supplied to an application.  There is no operation that obtains it
from an ambient global.  A holder may attenuate, but not amplify, authority:
\[
\dfrac{\begin{gathered}
 H_A(g)=\kw{mirror}\langle C_g,w,\zeta,R\rangle\qquad
 \kw{delegate}\in R\\
 \zeta'\sqsubseteq_{P,\widehat H_A,w}\zeta\qquad R'\subseteq R
 \qquad g'\text{ fresh}\\
 H_A'=H_A[g'\mapsto\kw{mirror}\langle C_g,w,\zeta',R'\rangle]
\end{gathered}}
 {\Omega,A,g\vdash\mathsf{deriveMirror}(\zeta',R')
       \Downarrow\langle\Omega[\mathcal H(A)\mapsto H_A'],g'\rangle}
\quad(\rulelabel{Mirror-Derive})
\]
Access modifiers play no role in this rule.  A private method is visible to a
method or enclosing-mixin mirror precisely because the mirror, rather than a
source-level send, is the authority check.

\subsection{Reified control and stack images}
\label{sec:debugging}

The sequential notation of Equation~\eqref{eq:configuration} stores the active
term once, beside the activation stack.  Debugging requires a uniform
per-frame view.  Let $b_c\in\mathit{CodeImageId}$ name an immutable source,
bytecode, or native-code image and define a control location by
\begin{equation}\label{eq:control-location-domain}
 \lambda_c::=\kw{sourcePoint}(z)\mid
   \kw{machinePoint}(b_c,n_c),\qquad n_c\in\mathbb N.
\end{equation}
Introduce the internal waiting term \kw{awaitResult} and the control-image
domain
\begin{equation}\label{eq:frame-control-domain}
 \kappa::=\kw{activeControl}\langle\pi,t\rangle
       \mid\kw{suspendedControl}\langle\pi\rangle.
\end{equation}
The current term and source-level control location of a frame are the total
functions
\begin{equation}\label{eq:frame-current-term-and-location}
\begin{aligned}
 \mathsf{currentTerm}(\kw{activeControl}\langle\pi,t\rangle)&=t,\\
 \mathsf{controlLocation}(\kw{activeControl}\langle\pi,t\rangle)
   &=\mathsf{site}(\pi,t),\\
 \mathsf{currentTerm}(\kw{suspendedControl}\langle\pi\rangle)
   &=\kw{awaitResult},\\
 \mathsf{controlLocation}(\kw{suspendedControl}\langle\pi\rangle)
   &=\mathsf{resumeSite}(\pi).
\end{aligned}
\end{equation}
$\mathsf{site}$ and $\mathsf{resumeSite}$ are the control-map functions emitted
by elaboration or compilation and return a $\lambda_c$.  A source evaluator
normally returns \kw{sourcePoint}; a bytecode or native implementation returns
\kw{machinePoint}, whose second component is its actual program counter.
Materialized terms and frames are not changed by a later code
install, so the map for a live or debugger-retained frame must be retained as
long as that frame is resumable.

For
$\mathcal A=\langle a_1,\pi_1\rangle\cdots
 \langle a_m,\pi_m\rangle$ with $m>0$, define its canonical control image by
\begin{equation}\label{eq:stack-control-image}
\begin{split}
 \mathsf{controlImage}(\mathcal A,t)=
 &\langle a_1,\kw{suspendedControl}\langle\pi_1\rangle\rangle\cdots\\[-1mm]
 &\langle a_{m-1},\kw{suspendedControl}\langle\pi_{m-1}\rangle\rangle
 \cdot\langle a_m,\kw{activeControl}\langle\pi_m,t\rangle\rangle.
\end{split}
\end{equation}
Thus every live frame has a current term: a lower frame is explicitly waiting
for its callee's result at the location determined by its continuation stack.
The inverse partial function is
\begin{equation}\label{eq:decode-control-image}
 \mathsf{decodeControlImage}
   (\mathsf{controlImage}(\mathcal A,t))=\langle\mathcal A,t\rangle.
\end{equation}
It is undefined for an empty image, for an active frame below the top, or for a
suspended top frame.

A reified frame contains both the activation record and its control image:
\begin{equation}\label{eq:debug-frame-image}
 \varphi=\kw{frameImage}\langle a,H(a),\kappa\rangle.
\end{equation}
The actor-stack snapshot is the immutable value
\begin{equation}\label{eq:actor-stack-snapshot}
\begin{split}
 \mathsf{actorStackSnapshot}_{\Omega}(B)
   &\simeq\langle\kw{running},
      \mathsf{reifyFrames}_{U}(
        \mathsf{controlImage}(\mathcal A,t))\rangle\\[-1mm]
   &\hspace{30mm}\text{if }\mathcal X(B)=
      \kw{runningTurn}\langle\mathcal A,t,p,u\rangle,\\
 \mathsf{actorStackSnapshot}_{\Omega}(B)
   &\simeq\langle\kw{paused}(r_s,d_s),
      \mathsf{reifyFrames}_{U}(
        \mathsf{controlImage}(\mathcal A,t))\rangle\\[-1mm]
   &\hspace{30mm}\text{if }\mathcal X(B)=
      \kw{pausedTurn}\langle\mathcal A,t,p,u,r_s,d_s\rangle,
\end{split}
\end{equation}
where $U=\mathsf{vmStore}(\Omega)$ and $\mathsf{reifyFrames}$ replaces every
$\langle a,\kappa\rangle$ by Equation~\eqref{eq:debug-frame-image}.  It is
undefined for an idle actor.  The paused form exposes the unforgeable token
that must accompany any later paused-stack mutation.

The unique live control of an activation is
\begin{equation}\label{eq:live-activation-control}
\begin{split}
 \mathsf{liveControl}_{\Omega}(a)=\kappa
 \quad\Longleftrightarrow\quad&
 \exists B,\bar\varphi_0,\bar\varphi_1,R.\\[-1mm]
 &\mathsf{actorStackSnapshot}_{\Omega}(B)=
 \langle m,\bar\varphi_0\cdot
   \langle\kw{frameImage}\langle a,R,\kappa\rangle\rangle
   \cdot\bar\varphi_1\rangle.
\end{split}
\end{equation}
Uniqueness follows from actor ownership and the prohibition on duplicate live
activation identities.  The function is undefined for a retained activation
that is no longer on a stack.

\paragraph{Editable stack templates.}
A debugger constructs an immutable template before requesting an atomic
install.  Template-local keys and frames are
\begin{equation}\label{eq:stack-template-domain}
 \eta::=\kw{retain}(a)\mid\kw{freshFrame}(u_f),
 \qquad
 \psi=\kw{frameTemplate}\langle\eta,R,\kappa\rangle,
 \qquad
 \mathcal S=\langle\psi_1,\ldots,\psi_m\rangle.
\end{equation}
$R$ is a complete activation-record template.  Its continuation, home, and
lexical-environment activation references may name other keys in the same
$\mathcal S$.  A retained key preserves an existing activation identity; a
fresh-frame key allocates a new identity at installation.  Consequently a
copied frame is never made live twice under one activation identity.

The three basic stack-editing operations have distinct names.  First, promote
and demote control without overloading either operation:
\begin{equation}\label{eq:promote-demote-frame-control}
\begin{aligned}
 &\mathsf{promoteFrame}(\kw{frameTemplate}\langle
      \eta,R,\kw{suspendedControl}\langle\pi\rangle\rangle,v)
   \\[-1mm]&\qquad={}\kw{frameTemplate}\langle
      \eta,R,\kw{activeControl}\langle\pi,v\rangle\rangle,\\
 &\mathsf{demoteFrame}(\kw{frameTemplate}\langle
      \eta,R,\kw{activeControl}\langle\pi,t\rangle\rangle)
   \\[-1mm]&\qquad={}\kw{frameTemplate}\langle
      \eta,R,\kw{suspendedControl}\langle\pi\rangle\rangle.
\end{aligned}
\end{equation}
Popping $j$ frames supplies the chosen continuation value to the new top:
\begin{equation}\label{eq:pop-stack-frames}
 \mathsf{popStackFrames}
   (\mathcal S_0\cdot\langle\psi\rangle\cdot\mathcal S_1,
      |\mathcal S_1|,v)
 =\mathcal S_0\cdot\langle\mathsf{promoteFrame}(\psi,v)\rangle,
 \qquad |\mathcal S_1|<|\mathcal S|.
\end{equation}
Pushing a prepared active frame makes the old top wait for the pushed frame's
result:
\begin{equation}\label{eq:push-stack-frame}
 \mathsf{pushStackFrame}
  (\mathcal S_0\cdot\langle\psi\rangle,\psi')
 =\mathcal S_0\cdot\langle\mathsf{demoteFrame}(\psi),\psi'\rangle,
 \quad
 \mathsf{currentTerm}(\psi'.\kappa)\ne\kw{awaitResult}.
\end{equation}
Finally, copying a contiguous segment allocates fresh template keys and rewires
every stack-structural activation link internal to the copied segment:
\begin{equation}\label{eq:copy-stack-frame-segment}
\begin{aligned}
 \mathsf{copyStackFrameSegment}(\mathcal S,i,j,\bar u)
   &=\mathsf{rename}_{\nu}
       (\langle\psi_i,\ldots,\psi_j\rangle),\\
 \nu(\psi_k.\eta)&=\kw{freshFrame}(u_k)\qquad(i\le k\le j),
 \qquad |\bar u|=j-i+1.
\end{aligned}
\end{equation}
Continuation, home, and lexical-environment links outside the copied segment
remain references to their original retained activations.  Parameter and local
values are shallow-copied; an activation object stored there is not implicitly
retargeted.  A debugger may concatenate the returned segment at any position,
edit its records or controls, and submit the complete result for validation.
These pure operations therefore express push, pop, single-frame copy, and
multi-frame copy without making a partially edited live stack observable.

\subsection{Inspection}

Mirror queries are
\begin{equation}\label{eq:mirror-query-grammar}
\begin{split}
 q_r::={}&\kw{programVersionQuery}\mid\kw{mixinDescriptorQuery}(M)\\
 &\mid\kw{methodDescriptorQuery}(i_f)\mid\kw{classDescriptorQuery}(C)\\
 &\mid\kw{objectDescriptorQuery}(o)\mid\kw{activationDescriptorQuery}(a)\\
 &\mid\kw{activationControlQuery}(a)\mid\kw{currentActivationQuery}(A)\\
 &\mid\kw{actorStackQuery}(A)\mid\kw{allInstancesOfQuery}(C)\\
 &\mid\kw{mixinApplicationsQuery}(M).
\end{split}
\end{equation}
For a running actor, the current activation is
\begin{equation}\label{eq:current-activation}
 \mathsf{currentActivation}_{\Omega}(A)=a
 \quad\Longleftrightarrow\quad
 \begin{aligned}[t]
 \mathcal X(A)&=\kw{runningTurn}\langle
   \mathcal A_0\mathbin{\cdot}\langle a,\pi\rangle,t,p,u\rangle
 \quad\lor\\
 \mathcal X(A)&=\kw{pausedTurn}\langle
   \mathcal A_0\mathbin{\cdot}\langle a,\pi\rangle,t,p,u,r_s,d_s\rangle.
 \end{aligned}
\end{equation}
It is undefined for an idle actor.  Descriptor construction is given by the
following separate functions, each of which returns an immutable value object:
\begin{equation}\label{eq:reflective-descriptors}
\begin{aligned}
 \mathsf{describeMixin}_{P}(M)&\simeq
   \left\langle\begin{gathered}\decl_P(M),\methods_P(M),
       \mathsf{ownSlots}_P(M),\\
       P.\mathit{mixin}(M).\mathit{nestedDecls},\iota_P(M)
   \end{gathered}\right\rangle,\\
 \mathsf{describeMethod}_{P}(i_f)
   &\simeq\langle M,f,P.\mathit{body}(i_f)\rangle,\\[-1mm]
 &\quad\text{where }M=\mathsf{methodOwner}_P(i_f),\quad
       \methods_P(M)(\selector(f))=f,\\
 \mathsf{describeClass}_{H}(C)&\simeq
   \left\langle\begin{gathered}\mixin_H(C),\super_H(C),H(C).eo,\\
       H(C).C_m,H(C).q,H(C).\chi_C
   \end{gathered}\right\rangle,\\
 \mathsf{describeObject}_{H}(o)
   &\simeq\langle\classof_H(o),H(o).\sigma,H(o).\eta\rangle,\\
 \mathsf{describeActivation}_{H}(a)&\simeq
   \left\langle\begin{gathered}H(a).p,H(a).\rho,H(a).\ell,H(a).o,H(a).C,\\
       H(a).k,H(a).h,H(a).q_a,H(a).\gamma
   \end{gathered}\right\rangle.
\end{aligned}
\end{equation}
The complete debugger-facing activation descriptor adds its current control
when the activation is live:
\begin{equation}\label{eq:debug-activation-descriptor}
 \mathsf{describeDebugActivation}_{\Omega}(a)
  \simeq\langle
    \mathsf{describeActivation}_{\mathsf{vmStore}(\Omega)}(a),
    \mathsf{option}(\mathsf{liveControl}_{\Omega}(a))\rangle.
\end{equation}
$\mathsf{option}$ is \kw{none} when
Equation~\eqref{eq:live-activation-control} is undefined.  Thus retained dead
activations remain inspectable, while every live activation exposes both its
current term and the source/PC location of
Equation~\eqref{eq:frame-current-term-and-location}.
The symbol $\simeq$ means an immutable, recursively reified snapshot, not a
mutable alias to a metalevel record.  In particular, a continuation component
is the activation identity specified in Section~3, reified as an activation
mirror only if the caller has authority for that target.

Fix an implementation-chosen total order $\prec$ on the object identities of
one VM cohort.  The exact-instance query returns a snapshot sequence, not a
live view:
\begin{equation}\label{eq:all-instances-of}
 \mathsf{allInstancesOf}_{\Omega}(C)=
 \mathsf{sort}_{\prec}
 \{x\in\dom(\mathsf{vmStore}(\Omega))
       \mid \classof_{\mathsf{vmStore}(\Omega)}(x)=C\}.
\end{equation}
Thus subclasses are excluded.  Presence in the VM store is the abstract
notion of liveness.  The optional garbage-collection layer makes this precise
by removing eligible identities from that store; without that layer, the query
observes every exact instance allocated so far.  The order is deliberately a
platform choice because the mirror protocol promises no allocation order.

The applications of a mixin are the live class objects whose class records
name that exact static mixin identity:
\begin{equation}\label{eq:mixin-applications}
 \mathsf{mixinApplications}_{\Omega}(M)=
 \mathsf{sort}_{\prec}
 \{C\in\dom(\mathsf{vmStore}(\Omega))
       \mid \mathsf{isClass}_{\mathsf{vmStore}(\Omega)}(C)
       \land \mixin_{\mathsf{vmStore}(\Omega)}(C)=M\}.
\end{equation}
This is an exact-mixin query: applications of other mixins and ordinary
instances of the resulting classes are excluded.  A metaclass is included
when its class record applies $M$, so a class-side mixin is treated uniformly
with an instance-side mixin.  As with Equation~\eqref{eq:all-instances-of}, the
result is a snapshot in the fixed order $\prec$.  Working Specification 0.109
requires both population queries.

For $\Omega=\langle P,V,\mathcal H,\mathcal X,\mathcal Q,\mathcal N,
\Pi,\Phi,\Gamma,\mathcal E,\mathcal D,w,n\rangle$, put
$U=\mathsf{vmStore}(\Omega)$.  The authority target of every query is
\begin{equation}\label{eq:mirror-query-target}
\begin{aligned}
 \mathsf{queryTarget}_{P,U,w}(\kw{programVersionQuery})
   &=\kw{vmTarget}(w),\\
 \mathsf{queryTarget}_{P,U,w}(\kw{mixinDescriptorQuery}(M))
   &=\kw{mixinTarget}(M),\\
 \mathsf{queryTarget}_{P,U,w}(\kw{methodDescriptorQuery}(i_f))
   &=\kw{methodTarget}(i_f),\\
 \mathsf{queryTarget}_{P,U,w}(\kw{classDescriptorQuery}(C))
   &=\kw{classTarget}(C),\\
 \mathsf{queryTarget}_{P,U,w}(\kw{objectDescriptorQuery}(o))
   &=\kw{objectTarget}(o),\\
 \mathsf{queryTarget}_{P,U,w}(\kw{activationDescriptorQuery}(a))
   &=\kw{activationTarget}(a),\\
 \mathsf{queryTarget}_{P,U,w}(\kw{activationControlQuery}(a))
   &=\kw{activationTarget}(a),\\
 \mathsf{queryTarget}_{P,U,w}(\kw{currentActivationQuery}(A))
   &=\kw{actorTarget}(A),\\
 \mathsf{queryTarget}_{P,U,w}(\kw{actorStackQuery}(A))
   &=\kw{actorTarget}(A),\\
 \mathsf{queryTarget}_{P,U,w}(\kw{allInstancesOfQuery}(C))
   &=\kw{classTarget}(C),\\
 \mathsf{queryTarget}_{P,U,w}(\kw{mixinApplicationsQuery}(M))
   &=\kw{mixinTarget}(M).
\end{aligned}
\end{equation}
Its result is the corresponding immutable snapshot:
\begin{equation}\label{eq:mirror-query-result}
\begin{aligned}
 \mathsf{queryResult}_{\Omega}(\kw{programVersionQuery})&=n,\\
 \mathsf{queryResult}_{\Omega}(\kw{mixinDescriptorQuery}(M))
   &=\mathsf{describeMixin}_{P}(M),\\
 \mathsf{queryResult}_{\Omega}(\kw{methodDescriptorQuery}(i_f))
   &=\mathsf{describeMethod}_{P}(i_f),\\
 \mathsf{queryResult}_{\Omega}(\kw{classDescriptorQuery}(C))
   &=\mathsf{describeClass}_{U}(C),\\
 \mathsf{queryResult}_{\Omega}(\kw{objectDescriptorQuery}(o))
   &=\mathsf{describeObject}_{U}(o),\\
 \mathsf{queryResult}_{\Omega}(\kw{activationDescriptorQuery}(a))
   &=\mathsf{describeDebugActivation}_{\Omega}(a),\\
 \mathsf{queryResult}_{\Omega}(\kw{activationControlQuery}(a))
   &=\mathsf{liveControl}_{\Omega}(a),\\
 \mathsf{queryResult}_{\Omega}(\kw{currentActivationQuery}(A))
   &=\mathsf{currentActivation}_{\Omega}(A),\\
 \mathsf{queryResult}_{\Omega}(\kw{actorStackQuery}(A))
   &=\mathsf{actorStackSnapshot}_{\Omega}(A),\\
 \mathsf{queryResult}_{\Omega}(\kw{allInstancesOfQuery}(C))
   &=\mathsf{allInstancesOf}_{\Omega}(C),\\
 \mathsf{queryResult}_{\Omega}(\kw{mixinApplicationsQuery}(M))
   &=\mathsf{mixinApplications}_{\Omega}(M).
\end{aligned}
\end{equation}
Inspection is the partial judgment
\[
\dfrac{\begin{gathered}
 \widehat H_A(g)=\kw{mirror}\langle C_g,w,\zeta,R\rangle\qquad
 \kw{inspect}\in R\\
 \mathsf{queryTarget}_{P,U,w}(q_r)
       \sqsubseteq_{P,U,w}\zeta\qquad
 d=\mathsf{queryResult}_{\Omega}(q_r)
\end{gathered}}
 {\Omega,A,g\vdash\mathsf{inspectMirror}(q_r)\Downarrow d}
\quad(\rulelabel{Mirror-Inspect})
\]
There is no failed-premise fallback that reveals a target.  The concrete
mirror library may reify denial as an exception, but must not disclose the
descriptor.

\subsection{Reflective transactions}

A transaction $\Delta$ is a finite sequence of commands, each explicitly
carrying the mirror that authorizes it:
\begin{equation}\label{eq:debug-stop-reasons}
 d_s::=\kw{exceptionStop}(x)\mid\kw{breakpointStop}(\lambda_c)
       \mid\kw{requestedStop}(o).
\end{equation}
\begin{equation}\label{eq:reflection-command-grammar}
\begin{split}
 \delta_{\mathrm{code}}::={}&
 \kw{replaceMethodBody}(g,i_f,b)
 \mid\kw{replaceMethodDefinition}(g,M,i_f,f)\\
 &\mid\kw{addMethodDefinition}(g,M,f)
 \mid\kw{removeMethodDefinition}(g,M,i_f)\\
 &\mid\kw{replaceSlotDeclarations}(g,M,\delta_s)
 \mid\kw{replaceNestedDeclarations}(g,M,\delta_n)\\
 &\mid\kw{replaceMixinInitializer}(g,M,\iota),\\
 \delta_{\mathrm{class}}::={}&
 \kw{changeSuperclass}(g,C,S)
 \mid\kw{changeClassEnclosingObject}(g,C,eo),\\
 \delta_{\mathrm{objectClass}}::={}&
 \kw{changeObjectClass}(g,o,C),\\
 \delta_{\mathrm{objectSlot}}::={}&
 \kw{writeObjectSlot}(g,o,j,v),\\
 \delta_{\mathrm{activation}}::={}&
 \kw{writeActivationParameter}(g,a,x,v)
 \mid\kw{writeActivationLocal}(g,a,j,v)\\
 &\mid\kw{changeActivationCurrentClass}(g,a,C)\\
 &\mid\kw{changeActivationContinuation}(g,a,k)
 \mid\kw{markActivationUncontinuable}(g,a),\\
 \delta_{\mathrm{continuation}}::={}&
 \kw{continueAtActivation}(g,A,a,v),\\
 \delta_{\mathrm{debug}}::={}&
 \kw{pauseActor}(g,B,d_s)\\
 &\mid\kw{replaceCurrentActorStack}(g,B,\mathcal S)\\
 &\mid\kw{replacePausedActorStack}(g,B,r_s,\mathcal S)\\
 &\mid\kw{resumePausedActorAtFullSpeed}(g,B,r_s)\\
 &\mid\kw{replaceAndResumePausedActorStack}(g,B,r_s,\mathcal S),\\
 \delta_r::={}&\delta_{\mathrm{code}}
 \mid\delta_{\mathrm{class}}
 \mid\delta_{\mathrm{objectClass}}
 \mid\delta_{\mathrm{objectSlot}}
 \mid\delta_{\mathrm{activation}}
 \mid\delta_{\mathrm{continuation}}
 \mid\delta_{\mathrm{debug}},\\
 \Delta::={}&\langle\delta_1,\ldots,\delta_m\rangle.
\end{split}
\end{equation}
By definition, $\mathit{CodeCommand}$ is generated by
$\delta_{\mathrm{code}}$, and $\mathit{ClassCommand}$ by
$\delta_{\mathrm{class}}$.  The sets $\mathit{ObjectClassCommand}$ and
$\mathit{ObjectSlotCommand}$ are generated by
$\delta_{\mathrm{objectClass}}$ and $\delta_{\mathrm{objectSlot}}$,
respectively.  Finally, $\mathit{ActivationCommand}$ is generated by
$\delta_{\mathrm{activation}}$, and $\mathit{ContinuationCommand}$ by
$\delta_{\mathrm{continuation}}$, and $\mathit{DebuggerCommand}$ by
$\delta_{\mathrm{debug}}$.  These seven sets are disjoint by outer
constructor and form an exhaustive partition of $\delta_r$.  In particular,
a class command edits a runtime class record, an
object-class command reclassifies one object, and an object-slot command
changes state within that object's existing final layout.  These are distinct
operations with distinct rights and distinct preparation phases.  The names
are intentionally not overloaded: changing a whole method definition is also
distinct from replacing only its body.

Define $\mathsf{commandMirror}(\delta_r)$ as the first field of every command.
The required right is the total function
\begin{equation}\label{eq:reflection-command-right}
 \mathsf{requiredRight}(\delta_r)=
 \begin{cases}
  \kw{editCode}&\delta_r\in\mathit{CodeCommand},\\
  \kw{editClass}&\delta_r\in\mathit{ClassCommand},\\
  \kw{reclassifyObject}&\delta_r\in\mathit{ObjectClassCommand},\\
  \kw{editObjectSlots}&\delta_r\in\mathit{ObjectSlotCommand},\\
  \kw{editActivation}&\delta_r\in\mathit{ActivationCommand},\\
  \kw{controlContinuation}&\delta_r\in\mathit{ContinuationCommand},\\
  \kw{controlActorStack}&\delta_r\in\mathit{DebuggerCommand}.
 \end{cases}
\end{equation}
The target function is
\begin{equation}\label{eq:reflection-command-target}
 \mathsf{commandTarget}(\delta_r)=
 \begin{cases}
  \kw{methodTarget}(i_f)&\delta_r\text{ names an existing }i_f,\\
  \kw{mixinTarget}(M)&\delta_r\text{ adds a method or edits mixin structure},\\
  \kw{classTarget}(C)&\delta_r\in\mathit{ClassCommand},\\
  \kw{objectTarget}(o)&\delta_r\in
       \mathit{ObjectClassCommand}\cup\mathit{ObjectSlotCommand},\\
  \kw{activationTarget}(a)&\delta_r\in
       \mathit{ActivationCommand}\cup\mathit{ContinuationCommand},\\
  \kw{actorTarget}(B)&\delta_r\in\mathit{DebuggerCommand}
       \text{ and names actor }B.
 \end{cases}
\end{equation}
For method addition the target is its mixin; for removal, definition
replacement, and body replacement the method identity is the target.  The
authorization predicate is
\begin{equation}\label{eq:reflection-authorized}
\begin{split}
 \mathsf{authorized}_{\Omega,A}(\delta_r)\quad\Longleftrightarrow\quad{}
 &\widehat H_A(g)=\kw{mirror}\langle C_g,w,\zeta,R\rangle\ \land\\
 &g=\mathsf{commandMirror}(\delta_r)\ \land\
 \mathsf{requiredRight}(\delta_r)\in R\ \land\\
 &\mathsf{commandTarget}(\delta_r)
   \sqsubseteq_{P,\mathsf{vmStore}(\Omega),w}\zeta.
\end{split}
\end{equation}
Thus syntactic possession of a mirror variable is not sufficient: the
unforgeable heap record, VM identity, target containment, and operation right
must all agree.

Transactions are simultaneous.  Let $\lambda$ range over the tagged write
locations below, and define the complete write-footprint function
$\mathsf{writes}$ by
{\footnotesize
\begin{equation}\label{eq:stack-template-record-locations}
 \mathsf{templateRecordLocations}(\mathcal S)=
 \{\kw{activationRecordLoc}(a)\mid
   \kw{frameTemplate}\langle\kw{retain}(a),R,\kappa\rangle
      \text{ occurs in }\mathcal S\}.
\end{equation}
}
{\small
\begin{equation}\label{eq:reflection-command-footprint}
\begin{aligned}
 &\mathsf{writes}(\kw{replaceMethodBody}(g,i_f,b))\\
 &\qquad=\{\kw{methodBodyLoc}(i_f)\},\\
 &\mathsf{writes}(\kw{replaceMethodDefinition}(g,M,i_f,f))\\
 &\qquad=\{\kw{methodDefLoc}(i_f),\kw{methodBodyLoc}(i_f),
       \kw{selectorLoc}(M,\selector(f))\},\\
 &\mathsf{writes}(\kw{addMethodDefinition}(g,M,f))\\
 &\qquad=\{\kw{methodDefLoc}(\mathsf{methodId}(f)),
       \kw{selectorLoc}(M,\selector(f))\},\\
 &\mathsf{writes}(\kw{removeMethodDefinition}(g,M,i_f))\\
 &\qquad=\{\kw{methodDefLoc}(i_f),\kw{methodBodyLoc}(i_f)\},\\
 &\mathsf{writes}(\kw{replaceSlotDeclarations}(g,M,\delta_s))
  =\{\kw{slotDeclsLoc}(M)\},\\
 &\mathsf{writes}(\kw{replaceNestedDeclarations}(g,M,\delta_n))
  =\{\kw{nestedDeclsLoc}(M)\},\\
 &\mathsf{writes}(\kw{replaceMixinInitializer}(g,M,\iota))
  =\{\kw{mixinInitializerLoc}(M)\},\\
 &\mathsf{writes}(\kw{changeSuperclass}(g,C,S))
  =\{\kw{superclassLoc}(C)\},\\
 &\mathsf{writes}(\kw{changeClassEnclosingObject}(g,C,eo))
  =\{\kw{classEnclosingLoc}(C)\},\\
 &\mathsf{writes}(\kw{changeObjectClass}(g,o,C))
  =\{\kw{objectClassLoc}(o)\},\\
 &\mathsf{writes}(\kw{writeObjectSlot}(g,o,j,v))
  =\{\kw{objectSlotLoc}(o,j)\},\\
 &\mathsf{writes}(\kw{writeActivationParameter}(g,a,x,v))\\
 &\qquad=\{\kw{activationParameterLoc}(a,x)\},\\
 &\mathsf{writes}(\kw{writeActivationLocal}(g,a,j,v))
  =\{\kw{activationLocalLoc}(a,j)\},\\
 &\mathsf{writes}(\kw{changeActivationCurrentClass}(g,a,C))\\
 &\qquad=\{\kw{activationClassLoc}(a)\},\\
 &\mathsf{writes}(\kw{changeActivationContinuation}(g,a,k))\\
 &\qquad=\{\kw{activationContinuationLoc}(a)\},\\
 &\mathsf{writes}(\kw{markActivationUncontinuable}(g,a))\\
 &\qquad=\{\kw{activationContinuableLoc}(a)\},\\
 &\mathsf{writes}(\kw{continueAtActivation}(g,A,a,v))\\
 &\qquad=\{\kw{actorStackLoc}(A),\kw{activationContinuationLoc}(a)\},\\
 &\mathsf{writes}(\kw{pauseActor}(g,B,d_s))
   =\{\kw{actorStackLoc}(B)\},\\
 &\mathsf{writes}(\kw{replaceCurrentActorStack}(g,B,\mathcal S))\\
 &\qquad=\{\kw{actorStackLoc}(B)\}\cup
      \mathsf{templateRecordLocations}(\mathcal S),\\
 &\mathsf{writes}(\kw{replacePausedActorStack}(g,B,r_s,\mathcal S))\\
 &\qquad=\{\kw{actorStackLoc}(B)\}\cup
      \mathsf{templateRecordLocations}(\mathcal S),\\
 &\mathsf{writes}(\kw{resumePausedActorAtFullSpeed}(g,B,r_s))
   =\{\kw{actorStackLoc}(B)\},\\
 &\mathsf{writes}(\kw{replaceAndResumePausedActorStack}
        (g,B,r_s,\mathcal S))\\
 &\qquad=\{\kw{actorStackLoc}(B)\}\cup
      \mathsf{templateRecordLocations}(\mathcal S).
\end{aligned}
\end{equation}
}
Location overlap is the least reflexive, symmetric relation
\begin{equation}\label{eq:reflection-location-overlap}
\begin{aligned}
 \lambda&\mathrel{\bowtie}\lambda,\\
 \kw{activationRecordLoc}(a)&\mathrel{\bowtie}
   \kw{activationFieldLoc}(a,\xi),\\
 \kw{activationFieldLoc}(a,\xi)&\mathrel{\bowtie}
   \kw{activationRecordLoc}(a).
\end{aligned}
\end{equation}
Each concrete parameter, local, class, continuation, and continuability
location above is the corresponding instance of
$\kw{activationFieldLoc}(a,\xi)$.  All other distinct tagged locations are
nonoverlapping.  This lets two edits
to different fields of one activation coexist while making a whole-record
stack installation conflict with either edit.
For $\Delta=\langle\delta_1,\ldots,\delta_m\rangle$, conflict and
nonconflict are therefore the predicates
\begin{equation}\label{eq:reflection-conflict}
\begin{aligned}
 \mathsf{conflict}_{\Delta}(i,j)
  &\Longleftrightarrow 1\le i<j\le m\ \land\
  &\qquad\exists\lambda\in\mathsf{writes}(\delta_i),
          \lambda'\in\mathsf{writes}(\delta_j).
          \ \lambda\mathrel{\bowtie}\lambda',\\
 \mathsf{nonconflicting}(\Delta)
  &\Longleftrightarrow
    \neg\exists i,j.\;\mathsf{conflict}_{\Delta}(i,j).
\end{aligned}
\end{equation}
The multi-location cases make whole-method replacement overlap body
replacement, and make continuation transfer overlap another transfer for the
same actor or a direct mutation of the resumed continuation.

\subsection{Source change and re-elaboration}

The code-command projection is
\begin{equation}\label{eq:code-command-projection}
 \mathsf{codeCommands}(\Delta)=
 \mathsf{filter}_{\mathit{CodeCommand}}(\Delta).
\end{equation}
Its direct source effect is the partial function $\mathsf{patchSource}$:
\begin{equation}\label{eq:source-patch}
\begin{array}{rcl}
 \mathsf{patchSource}(Q,\kw{replaceMethodBody}(g,i_f,b))
   &=&Q[\mathit{source}(i_f).\mathit{body}\mapsto b],\\
 \mathsf{patchSource}(Q,\kw{replaceMethodDefinition}(g,M,i_f,f))
   &=&Q[\mathit{source}(i_f)\mapsto f],\\
 \mathsf{patchSource}(Q,\kw{addMethodDefinition}(g,M,f))
   &=&\mathsf{insertMethodSource}(Q,\decl_Q(M),f),\\
 \mathsf{patchSource}(Q,\kw{removeMethodDefinition}(g,M,i_f))
   &=&\mathsf{deleteMethodSource}(Q,\decl_Q(M),i_f),\\
 \mathsf{patchSource}(Q,\kw{replaceSlotDeclarations}(g,M,\delta_s))
   &=&Q[\mathit{source}(\decl_Q(M)).\mathit{slots}\mapsto\delta_s],\\
 \mathsf{patchSource}(Q,\kw{replaceNestedDeclarations}(g,M,\delta_n))
   &=&Q[\mathit{source}(\decl_Q(M)).\mathit{nested}\mapsto\delta_n],\\
 \mathsf{patchSource}(Q,\kw{replaceMixinInitializer}(g,M,\iota))
   &=&Q[\mathit{source}(\decl_Q(M)).\mathit{initializer}\mapsto\iota].
\end{array}
\end{equation}
$\mathsf{insertMethodSource}$ requires a fresh method identity and absent
selector; $\mathsf{deleteMethodSource}$ requires that the identity is a direct
method of the named mixin.  Definition and body replacement preserve the
method identity; delete followed by add cannot reuse it.

Because the transaction is nonconflicting, the fold
\begin{equation}\label{eq:apply-source-commands}
\begin{aligned}
 \mathsf{applySourceCommands}(P,\epsilon)&=P,\\
 \mathsf{applySourceCommands}(P,\delta\cdot\bar\delta)
   &=\mathsf{applySourceCommands}(\mathsf{patchSource}(P,\delta),\bar\delta)
\end{aligned}
\end{equation}
is independent of the chosen canonical locus order.  The patched source
candidate is
\begin{equation}\label{eq:patched-source-candidate}
 Q=\mathsf{applySourceCommands}(P,\mathsf{codeCommands}(\Delta)).
\end{equation}
The ordinary elaborator is rerun on the complete affected dependency closure:
\begin{equation}\label{eq:reflective-reelaboration}
 \begin{aligned}
 &\mathsf{reelaborateProgram}(P,Q)=P'
 \quad\Longleftrightarrow\\[-1mm]
 &\quad P'.\mathit{source}=Q.\mathit{source}\ \land\\
 &\quad P'.\mathit{unit}=Q.\mathit{unit}\ \land\\
 &\quad \forall K\in\dom(P.\mathit{declMixin})\cap\dom(Q.\mathit{source}).\;
    P'.\mathit{declMixin}(K)=P.\mathit{declMixin}(K)\ \land\\
 &\quad \forall f,f'\text{ denoting one surviving declaration}.\;
    \mathsf{methodId}(f')=\mathsf{methodId}(f)\ \land\\
 &\quad \mathsf{coreProgramOK}(P').
 \end{aligned}
\end{equation}
The last clause is Equation~\eqref{eq:annotated-program-validity}; it
recomputes all derived fields of
Equation~\eqref{eq:program-definition}, including synthesized accessors,
lexical parents, scope chains, activation members, owner maps, initializer
normal forms, and method-body maps.  It is undefined on a parse, binding,
access, arity, duplicate-selector, or other static error.  This is why a
transaction cannot expose a half-elaborated program.

\subsection{Class graphs, object layouts, and individual objects}

Let the command projections be
\begin{equation}\label{eq:class-command-projections}
\begin{aligned}
 \mathsf{classCommands}(\Delta)&=
  \mathsf{filter}_{\mathit{ClassCommand}}(\Delta),\\
 \mathsf{objectClassCommands}(\Delta)&=
  \mathsf{filter}_{\mathit{ObjectClassCommand}}(\Delta),\\
 \mathsf{objectSlotCommands}(\Delta)&=
  \mathsf{filter}_{\mathit{ObjectSlotCommand}}(\Delta).
\end{aligned}
\end{equation}
Their heap effects are the distinct partial functions
\begin{equation}\label{eq:class-object-heap-effects}
\begin{aligned}
 \mathsf{applyClassCommand}(H,\kw{changeSuperclass}(g,C,S))
   &=H[C.S\mapsto S],\\
 \mathsf{applyClassCommand}(H,\kw{changeClassEnclosingObject}(g,C,eo))
   &=H[C.eo\mapsto eo],\\
 \mathsf{applyObjectClassCommand}(H,\kw{changeObjectClass}(g,o,C))
   &=H[o.C\mapsto C].
\end{aligned}
\end{equation}
The corresponding phases are
\begin{equation}\label{eq:class-object-command-folds}
\begin{aligned}
 \mathsf{applyClassCommands}(H,\epsilon)&=H,\\
 \mathsf{applyClassCommands}(H,\delta\cdot\bar\delta)
   &=\mathsf{applyClassCommands}
       (\mathsf{applyClassCommand}(H,\delta),\bar\delta),\\
 \mathsf{applyObjectClassCommands}(H,\epsilon)&=H,\\
 \mathsf{applyObjectClassCommands}(H,\delta\cdot\bar\delta)
   &=\mathsf{applyObjectClassCommands}\\
   &\quad(\mathsf{applyObjectClassCommand}(H,\delta),\bar\delta).
\end{aligned}
\end{equation}
The command sequences are in canonical locus order.  Neither phase changes
$H(C).\chi_C$: creation provenance records how the class was created, not its
current superclass or enclosing object.

The three projections in Equation~\eqref{eq:class-command-projections} must
not be conflated.  A class command changes the graph or lexical context shared
by all instances of $C$.  An object-class command changes only $o.C$ and then
subjects that object to layout reconciliation.  An object-slot command leaves
both the class graph and $o.C$ unchanged and writes one slot only after the
final layout has been established.

For every ordinary object, layout reconciliation against the final program
and class graph is
\begin{equation}\label{eq:reconcile-object-layout}
\begin{aligned}
 L_o&=\mathsf{set}(\mathsf{layout}_{P',H_c}(\classof_{H_c}(o))),\\
 \mathsf{reconcileObjectLayout}(P',H_c,o)
   &=H_c[o.\sigma\mapsto
       [j\in L_o\mapsto
          \begin{cases}H_c(o).\sigma(j)&j\in\dom(H_c(o).\sigma),\\
                        \nil&j\notin\dom(H_c(o).\sigma),
          \end{cases}]] .
\end{aligned}
\end{equation}
Thus stable physical slot identities preserve retained state, newly introduced
slots contain $\nil$, and removed slots become unreachable in the same atomic
commit.  The operation applies to every object whose old and new layouts
differ, including every instance affected transitively by a superclass or
mixin-slot edit and an individual object named by \kw{changeObjectClass}.
If $P.\mathit{mixin}(M)=\langle K,\phi,\delta,\nu,\iota\rangle$, define
\begin{equation}\label{eq:nested-declaration-keys}
\begin{aligned}
 \mathsf{ownNestedKeys}_{P}(M)&=\dom(\nu),\\
 \mathsf{nestedKeys}_{P,H}(C)
  &=\bigcup_{R\in\mathsf{set}(\mathsf{chain}_H(C))\setminus\{\Top\}}
       \mathsf{ownNestedKeys}_{P}(\mixin_H(R)).
\end{aligned}
\end{equation}
Nested-class caches are reconciled separately using
Equation~\eqref{eq:nested-declaration-keys}:
\begin{equation}\label{eq:reconcile-nested-cache}
 \mathsf{reconcileNestedCache}(P',H,o)=
 H[o.\eta\mapsto H(o).\eta\mathbin{\upharpoonright}
   \mathsf{nestedKeys}_{P',H}(\classof_H(o))].
\end{equation}
Removing a declaration therefore removes its memoized reachability; adding one
does not eagerly create a class.  Existing class objects remain ordinary live
objects if referenced elsewhere.

Using the VM identity order $\prec$ fixed by
Equation~\eqref{eq:all-instances-of}, define
\begin{equation}\label{eq:reconciliation-phases}
\begin{aligned}
 \mathsf{objectIds}(H)
   &=\mathsf{sort}_{\prec}
       \{o\in\dom(H)\mid H(o)=\kw{object}\langle C,\sigma,\eta\rangle\},\\
 \mathsf{reconcileObjectLayouts}(P',H,\epsilon)&=H,\\
 \mathsf{reconcileObjectLayouts}(P',H,o\cdot\bar O)
   &=\mathsf{reconcileObjectLayouts}\\
   &\quad(P',\mathsf{reconcileObjectLayout}(P',H,o),\bar O),\\
 \mathsf{reconcileAllObjectLayouts}(P',H)
   &=\mathsf{reconcileObjectLayouts}(P',H,\mathsf{objectIds}(H)),\\
 \mathsf{reconcileNestedCaches}(P',H,\epsilon)&=H,\\
 \mathsf{reconcileNestedCaches}(P',H,o\cdot\bar O)
   &=\mathsf{reconcileNestedCaches}\\
   &\quad(P',\mathsf{reconcileNestedCache}(P',H,o),\bar O),\\
 \mathsf{reconcileAllNestedCaches}(P',H)
   &=\mathsf{reconcileNestedCaches}(P',H,\mathsf{objectIds}(H)).
\end{aligned}
\end{equation}

After all reconciliations, an explicit mirror slot write is
\begin{equation}\label{eq:reflective-object-slot-write}
 \begin{aligned}
 &\mathsf{applyObjectSlotWrite}(P',H,\kw{writeObjectSlot}(g,o,j,v))
    =H[o.\sigma(j)\mapsto v]\\
 &\hspace{28mm}\text{if }j\in
    \mathsf{set}(\mathsf{layout}_{P',H}(\classof_H(o))).
 \end{aligned}
\end{equation}
It is undefined otherwise.  This operation is different from the base setter
and from initializer direct write; no source-level access check or synthesized
method is involved.  All actor-isolation checks are nevertheless retained.

The object-slot phase is
\begin{equation}\label{eq:object-slot-command-fold}
\begin{aligned}
 \mathsf{applyObjectSlotWrites}(P',H,\epsilon)&=H,\\
 \mathsf{applyObjectSlotWrites}(P',H,\delta\cdot\bar\delta)
   &=\mathsf{applyObjectSlotWrites}
       (P',\mathsf{applyObjectSlotWrite}(P',H,\delta),\bar\delta).
\end{aligned}
\end{equation}

\subsection{Activations and continuations}

The command projections for the remaining two phases are
\begin{equation}\label{eq:control-command-projections}
\begin{aligned}
 \mathsf{activationCommands}(\Delta)&=
  \mathsf{filter}_{\mathit{ActivationCommand}}(\Delta),\\
 \mathsf{continuationCommands}(\Delta)&=
  \mathsf{filter}_{\mathit{ContinuationCommand}}(\Delta).
\end{aligned}
\end{equation}

For an activation command other than continuation transfer, define
\begin{equation}\label{eq:activation-field-effects}
\begin{aligned}
 \mathsf{applyActivationCommand}(H,\kw{writeActivationParameter}(g,a,x,v))
   &=H[a.\rho(x)\mapsto v],\\
 \mathsf{applyActivationCommand}(H,\kw{writeActivationLocal}(g,a,j,v))
   &=H[a.\ell(j)\mapsto v],\\
 \mathsf{applyActivationCommand}(H,\kw{changeActivationCurrentClass}(g,a,C))
   &=H[a.C\mapsto C],\\
 \mathsf{applyActivationCommand}(H,\kw{changeActivationContinuation}(g,a,k))
   &=H[a.k\mapsto k],\\
 \mathsf{applyActivationCommand}(H,\kw{markActivationUncontinuable}(g,a))
   &=\mathsf{markUncontinuable}(H,a).
\end{aligned}
\end{equation}
The first case is reflectively permitted even though parameter slots are
immutable at base level.  The second requires a physical local identity
already present in $H(a).\ell$.  A changed current class must remain on the
current instance's class chain.  A changed continuation is either $\nil$ or a
strictly lower live stack activation, so continuation links remain acyclic.
Marking uses Equations~\eqref{eq:dependents} and
\eqref{eq:mark-uncontinuable}.

The activation phase is
\begin{equation}\label{eq:activation-command-fold}
\begin{aligned}
 \mathsf{applyActivationCommands}(H,\epsilon)&=H,\\
 \mathsf{applyActivationCommands}(H,\delta\cdot\bar\delta)
   &=\mathsf{applyActivationCommands}
       (\mathsf{applyActivationCommand}(H,\delta),\bar\delta).
\end{aligned}
\end{equation}

Equation~\eqref{eq:resume} supplies the stack operation used by reflective
continuation transfer.  The corresponding state transformer is
\begin{equation}\label{eq:reflective-continuation-transfer}
\begin{aligned}
 &\delta_c=\kw{continueAtActivation}(g,A,a,v),\\
 &\mathsf{applyContinuationTransfer}(H,\mathcal X,\delta_c)
      =\langle H',\mathcal X'\rangle,\\
 &\mathcal X(A)=\kw{runningTurn}\langle\mathcal A,t,p,u\rangle,\qquad
 H(a).q_a=\kw{true},\qquad
 \mathsf{resume}(\mathcal A,a,v)=\langle\mathcal A',v\rangle,\\
 &H'=H[a.k\mapsto\nil],\qquad
 \mathcal X'=\mathcal X[A\mapsto
       \kw{runningTurn}\langle\mathcal A',v,p,u\rangle].
\end{aligned}
\end{equation}
It is undefined for an idle actor, a retained but no longer live activation,
or an uncontinuable activation.  At most one continuation transfer per actor
may occur in a transaction.  The discarded stack suffix remains in the heap,
as in ordinary return, so debuggers may retain it; its continuation dependents
are marked uncontinuable before commit.  This gives reflection continuation
power without introducing a second, incompatible continuation representation.

The continuation-transfer phase is
\begin{equation}\label{eq:continuation-command-fold}
\begin{aligned}
 \mathsf{applyContinuationTransfers}(H,\mathcal X,\epsilon)
   &=\langle H,\mathcal X\rangle,\\
 \mathsf{applyContinuationTransfers}
      (H,\mathcal X,\delta\cdot\bar\delta)
   &=\mathsf{applyContinuationTransfers}(H',\mathcal X',\bar\delta)\\
 &\quad\text{where }\langle H',\mathcal X'\rangle=
       \mathsf{applyContinuationTransfer}(H,\mathcal X,\delta).
\end{aligned}
\end{equation}

\subsection{Debugger stack installation and full-speed resumption}

The debugger-command projection is
\begin{equation}\label{eq:debugger-command-projection}
 \mathsf{debuggerCommands}(\Delta)=
   \mathsf{filter}_{\mathit{DebuggerCommand}}(\Delta).
\end{equation}
For a template $\mathcal S=\langle\psi_1,\ldots,\psi_m\rangle$, key
resolution is the partial function
\begin{equation}\label{eq:resolve-stack-template-keys}
\begin{aligned}
 \mathsf{resolveTemplateKeys}(H,\mathcal S)&=\nu,\\
 \nu(\kw{retain}(a))&=a
   &&\text{only if }H(a)\text{ is an activation record},\\
 \nu(\kw{freshFrame}(u_i))&=a_i
   &&a_i\text{ is fresh},\\
 \nu(\eta_i)&\ne\nu(\eta_j)&&i\ne j.
\end{aligned}
\end{equation}
Fresh choices are simultaneous.  Substitution $\nu(R)$ replaces every
template-key occurrence in a continuation, home, or lexical environment by
its resolved activation identity.

The well-formedness predicate for a resolved stack image is
{\footnotesize
\begin{equation}\label{eq:debug-stack-image-validity}
\begin{split}
 \mathsf{debugStackOK}(P,H,
    \langle a_1,\kappa_1\rangle\cdots\langle a_m,\kappa_m\rangle)
 \quad\Longleftrightarrow\quad{}
 &m>0\ \land\ \mathsf{pairwiseDistinct}(a_1,\ldots,a_m)\ \land\\
 &\mathsf{decodeControlImage}
    (\langle a_1,\kappa_1\rangle\cdots\langle a_m,\kappa_m\rangle)
       \downarrow\ \land\\
 &\forall i\le m.\ H(a_i)\text{ is an activation record}\ \land
       H(a_i).q_a=\kw{true}\ \land\\
 &\forall i\le m.\ \mathsf{controlLocation}(\kappa_i)\downarrow
       \ \land\ \mathsf{controlRefs}(\kappa_i)\subseteq\dom(H)\ \land\\
 &\forall i\le m.\ H(a_i).k=\nil\ \lor\
       \exists j<i.\ H(a_i).k=a_j\ \land\\
 &\forall i\le m.\ H(a_i).h=\nil\ \lor\ H(a_i).h\in\dom(H).
\end{split}
\end{equation}
}
$\mathsf{controlRefs}$ is the instance of the structural reference extractor
of Equation~\eqref{eq:structural-reference-extractor} for terms and
intra-activation frames.  The control-map definedness premise validates a
source or machine location; a bytecode platform additionally requires the operand-stack and
temporary layout encoded by $\pi_i$ to match the stack map at the corresponding
PC.  This is the implementation representation of the same predicate, not a
weaker alternate rule.

Removed live activations remain available to retained mirrors but cease to be
continuable:
\begin{equation}\label{eq:retire-removed-stack-activations}
 \mathsf{retireRemoved}(H,R)=
 H[\,a.k\mapsto\nil,\ a.q_a\mapsto\kw{false}\mid a\in R\,].
\end{equation}
Define atomic template materialization by
\begin{equation}\label{eq:materialize-stack-template}
\begin{split}
 &\mathsf{materializeStackTemplate}(P,H,\mathcal A,\mathcal S)
      =\langle H',\mathcal A',t'\rangle
 \quad\Longleftrightarrow\\
 &\quad\nu=\mathsf{resolveTemplateKeys}(H,\mathcal S)\ \land\\
 &\quad\mathcal K=
   \langle\nu(\psi_1.\eta),\psi_1.\kappa\rangle\cdots
   \langle\nu(\psi_m.\eta),\psi_m.\kappa\rangle\ \land\\
 &\quad\mathsf{decodeControlImage}(\mathcal K)
      =\langle\mathcal A',t'\rangle\ \land\\
 &\quad R=\mathsf{set}(\mathsf{activationIds}(\mathcal A))
       \setminus\mathsf{set}(\mathsf{activationIds}(\mathcal A'))\ \land\\
 &\quad H_0=\mathsf{retireRemoved}(H,R)\ \land\\
 &\quad H'=H_0[\,\nu(\psi_i.\eta)\mapsto\nu(\psi_i.R)
          \mid1\le i\le m\,]\ \land\\
 &\quad\mathsf{debugStackOK}(P,H',\mathcal K).
\end{split}
\end{equation}
It is undefined if any key, record, control point, reference, continuation, or
stack-shape condition fails.  Since only $H'$ and the reconstructed
$\langle\mathcal A',t'\rangle$ leave the function, neither allocation nor
retirement is visible on failure.

The effect of one debugger command, issued by actor $A$, is the following
partial function.  Pausing allocates a fresh token:
\begin{equation}\label{eq:pause-actor-effect}
\begin{aligned}
 &\mathsf{applyDebuggerCommand}
   (P,H,\mathcal X,A,\kw{pauseActor}(g,B,d_s))
       =\langle H,\mathcal X'\rangle,\\
 &\mathcal X(B)=\kw{runningTurn}\langle\mathcal A,t,p,u\rangle,
 \quad A\ne B,\quad\mathcal A\ne\epsilon,\quad r_s\text{ fresh},\\
 &\mathcal X'=\mathcal X[B\mapsto
   \kw{pausedTurn}\langle\mathcal A,t,p,u,r_s,d_s\rangle].
\end{aligned}
\end{equation}
A debugger running in the actor that it is editing uses the nonreturning
current-stack replacement:
\begin{equation}\label{eq:replace-current-actor-stack}
\begin{aligned}
 &\mathsf{applyDebuggerCommand}
   (P,H,\mathcal X,A,\kw{replaceCurrentActorStack}(g,B,\mathcal S))
       =\langle H',\mathcal X'\rangle,\\
 &A=B,\quad
 \mathcal X(B)=\kw{runningTurn}\langle\mathcal A,t,p,u\rangle,\\
 &\mathsf{materializeStackTemplate}(P,H,\mathcal A,\mathcal S)
       =\langle H',\mathcal A',t'\rangle,\\
 &\mathcal X'=\mathcal X[B\mapsto
       \kw{runningTurn}\langle\mathcal A',t',p,u\rangle].
\end{aligned}
\end{equation}
The reflective invocation that requested this command is among the discarded
or replaced frames, so it does not return.  This is the operation used by an
in-actor debugger after a top-level resumable-exception handler has constructed
the desired continuation.
For the common resumable-exception case, the derived constructor is
\begin{equation}\label{eq:resume-exception-stack-template}
 \mathsf{resumeExceptionStack}
  (\mathcal S_0\cdot\langle\psi_f\rangle\cdot\mathcal S_d,a_f,v)
 =\mathcal S_0\cdot
   \langle\mathsf{promoteFrame}(\psi_f,v)\rangle
 \quad\text{if }\psi_f.\eta=\kw{retain}(a_f).
\end{equation}
It removes the handler and debugger suffix $\mathcal S_d$ and makes $v$ the
result of the original \kws{signal} send.  Submitting this template through
Equation~\eqref{eq:replace-current-actor-stack} therefore resumes precisely at
the protected failure point under the normal VM relation.

For a paused actor, replacement checks the exact pause token and issues a new
one:
\begin{equation}\label{eq:replace-paused-actor-stack}
\begin{aligned}
 &\mathsf{applyDebuggerCommand}
  (P,H,\mathcal X,A,\kw{replacePausedActorStack}(g,B,r_s,\mathcal S))
      =\langle H',\mathcal X'\rangle,\\
 &\mathcal X(B)=\kw{pausedTurn}\langle
      \mathcal A,t,p,u,r_s,d_s\rangle,\\
 &\mathsf{materializeStackTemplate}(P,H,\mathcal A,\mathcal S)
      =\langle H',\mathcal A',t'\rangle,\quad r_s'\text{ fresh},\\
 &\mathcal X'=\mathcal X[B\mapsto\kw{pausedTurn}\langle
      \mathcal A',t',p,u,r_s',d_s\rangle].
\end{aligned}
\end{equation}
Resuming an unchanged paused stack selects ordinary VM execution:
\begin{equation}\label{eq:resume-paused-actor-full-speed}
\begin{aligned}
 &\mathsf{applyDebuggerCommand}
  (P,H,\mathcal X,A,\kw{resumePausedActorAtFullSpeed}(g,B,r_s))
      =\langle H,\mathcal X'\rangle,\\
 &\mathcal X(B)=\kw{pausedTurn}\langle
      \mathcal A,t,p,u,r_s,d_s\rangle,\\
 &\mathcal X'=\mathcal X[B\mapsto
      \kw{runningTurn}\langle\mathcal A,t,p,u\rangle].
\end{aligned}
\end{equation}
Replacement and full-speed resumption may instead be one atomic operation:
\begin{equation}\label{eq:replace-resume-paused-actor-stack}
\begin{aligned}
 &\mathsf{applyDebuggerCommand}(P,H,\mathcal X,A,
   \kw{replaceAndResumePausedActorStack}(g,B,r_s,\mathcal S))
      =\langle H',\mathcal X'\rangle,\\
 &\mathcal X(B)=\kw{pausedTurn}\langle
      \mathcal A,t,p,u,r_s,d_s\rangle,\\
 &\mathsf{materializeStackTemplate}(P,H,\mathcal A,\mathcal S)
      =\langle H',\mathcal A',t'\rangle,\\
 &\mathcal X'=\mathcal X[B\mapsto
      \kw{runningTurn}\langle\mathcal A',t',p,u\rangle].
\end{aligned}
\end{equation}
``Full speed'' is not a second language semantics: after either resume equation,
Rule \rulelabel{Actor-Sequential} again lifts the ordinary sequential relation.
An interpreter, baseline VM, or optimizing VM conforms by implementing that
same relation from the installed control image.

The debugger phase is the explicit fold
\begin{equation}\label{eq:debugger-command-fold}
\begin{aligned}
 \mathsf{applyDebuggerCommands}(P,H,\mathcal X,A,\epsilon)
   &=\langle H,\mathcal X\rangle,\\
 \mathsf{applyDebuggerCommands}
   (P,H,\mathcal X,A,\delta\cdot\bar\delta)
   &=\mathsf{applyDebuggerCommands}
       (P,H',\mathcal X',A,\bar\delta)\\
 &\quad\text{where }\langle H',\mathcal X'\rangle={}\\[-1mm]
 &\hspace{29mm}\mathsf{applyDebuggerCommand}(P,H,\mathcal X,A,\delta).
\end{aligned}
\end{equation}

\paragraph{Simulator correspondence.}
Let $D$ be a Simulator state represented by ordinary Newspeak objects and let
$D\approx_{\mathsf{sim}}\Sigma$ mean that its bytecode frames and program
counters decode to the stack and control image of
Equation~\eqref{eq:stack-control-image}.  Write
$\longmapsto_{\mathsf{Simulator}}^{+}$ for a nonempty finite Simulator
execution and $\longmapsto^{*}$ for a possibly empty finite semantic
execution.  An admitted compiler/Simulator pair supplies the three-part
certificate
\begin{equation}\label{eq:simulator-step-correspondence}
\begin{aligned}
 D\approx_{\mathsf{sim}}\Sigma\ \land\
 D\longmapsto_{\mathsf{Simulator}}D'
 &\Rightarrow
 \exists\Sigma'.\ \Sigma\longmapsto^{*}\Sigma'
       \land D'\approx_{\mathsf{sim}}\Sigma',\\
 D\approx_{\mathsf{sim}}\Sigma\ \land\ \Sigma\longmapsto\Sigma'
 &\Rightarrow
 \exists D'.\ D\longmapsto_{\mathsf{Simulator}}^{+}D'
       \land D'\approx_{\mathsf{sim}}\Sigma',\\
 D\approx_{\mathsf{sim}}\Sigma
 &\Rightarrow \mathsf{resumeAtFullSpeed}(D)=\Sigma.
\end{aligned}
\end{equation}
The first clause permits bytecode-level stuttering while proving instruction
soundness; the second prevents the Simulator from omitting a semantic step;
the third closes the critical boundary back to native execution without a
change of represented configuration.  The certificate lifts the first clause
from one instruction to every finite Simulator execution.
The Simulator can inspect, push, pop, and copy its ordinary frame-image values;
Equations~\eqref{eq:replace-current-actor-stack} and
\eqref{eq:replace-resume-paused-actor-stack} cross back to normal VM execution.
Proving Equation~\eqref{eq:simulator-step-correspondence} for the actual
Simulator and bytecode set is an implementation-refinement proof; the language
semantics fixes the state and transfer contract it must implement.

\subsection{Preparation, validation, and atomic commit}

All phase ordering is semantic, not user-visible.  Define the partial
preparation function
\begin{equation}\label{eq:prepare-reflection}
\begin{split}
 &\mathsf{prepareReflection}(\Omega,A,\Delta)=\Omega'\\[-1mm]
 &\quad\Longleftrightarrow\quad
 P'=\mathsf{reelaborateProgram}(P,Q)\ \land\\
 &H_1=\mathsf{applyClassCommands}
       (\mathsf{vmStore}(\Omega),\mathsf{classCommands}(\Delta))\ \land\\
 &H_c=\mathsf{applyObjectClassCommands}
       (H_1,\mathsf{objectClassCommands}(\Delta))\ \land\\
 &H_l=\mathsf{reconcileAllObjectLayouts}(P',H_c)\ \land\\
 &H_n=\mathsf{reconcileAllNestedCaches}(P',H_l)\ \land\\
 &H_o=\mathsf{applyObjectSlotWrites}
       (P',H_n,\mathsf{objectSlotCommands}(\Delta))\ \land\\
 &H_a=\mathsf{applyActivationCommands}
       (H_o,\mathsf{activationCommands}(\Delta))\ \land\\
 &\langle H_d,\mathcal X_d\rangle=
     \mathsf{applyDebuggerCommands}
       (P',H_a,\mathcal X,A,\mathsf{debuggerCommands}(\Delta))\ \land\\
 &\langle H_t,\mathcal X'\rangle=
     \mathsf{applyContinuationTransfers}
       (H_d,\mathcal X_d,\mathsf{continuationCommands}(\Delta))\ \land\\
 &\Omega_0=\Omega[P\mapsto P',\mathcal X\mapsto\mathcal X',n\mapsto n+1]\ \land\\
 &\Omega'=\mathsf{repartitionVMStore}(\Omega_0,H_t).
\end{split}
\end{equation}
The fresh activations introduced by stack installation are
\begin{equation}\label{eq:fresh-installed-activations}
 \mathsf{freshInstalled}_{\Omega}(\mathcal X',B)=
 \mathsf{set}(\mathsf{stackActivationIds}_{\mathcal X'}(B))
 \setminus\dom(\mathsf{vmStore}(\Omega)).
\end{equation}
$\mathsf{repartitionVMStore}$ returns every pre-existing identity to the same
$V$ or $\mathcal H(A)$ domain from which it came, and places each member of
Equation~\eqref{eq:fresh-installed-activations} in $\mathcal H(B)$.  The sets
for distinct actors must be disjoint.  Actor-isolation validation below rejects
any installed record containing an illegal near reference; reflection creates
no cross-actor near reference.

The validation predicates below are over tagged record sets, so their
quantifiers have explicit domains:
\begin{equation}\label{eq:reflective-record-id-sets}
\begin{aligned}
 \mathsf{classRecordIds}(H)
  &=\{C\in\dom(H)\mid
       H(C)=\kw{class}\langle M,S,eo,C_m,q,\chi_C\rangle\},\\
 \mathsf{ordinaryObjectIds}(H)
  &=\{o\in\dom(H)\mid H(o)=\kw{object}\langle C,\sigma,\eta\rangle\},\\
 \mathsf{activationRecordIds}(H)
  &=\{a\in\dom(H)\mid
       H(a)=\kw{activation}\langle p,\rho,\ell,o,C,k,h,q_a,\gamma\rangle\}.
\end{aligned}
\end{equation}
Class-chain termination is the predicate
\begin{equation}\label{eq:reflective-class-chain-validity}
\begin{split}
 \mathsf{reachesTop}_{H}(C)\quad\Longleftrightarrow\quad{}
 &\exists k\in\mathbb N,C_0,\ldots,C_k.\;
 C_0=C\ \land\ C_k=\Top\ \land\\[-1mm]
 &\mathsf{pairwiseDistinct}(C_0,\ldots,C_k)\ \land\\
 &\forall i<k.\;C_i\in\mathsf{classRecordIds}(H)\ \land
                  \super_H(C_i)=C_{i+1}.
\end{split}
\end{equation}
Individual class-record consistency is
\begin{equation}\label{eq:reflective-class-record-validity}
\begin{split}
 \mathsf{classRecordOK}(P,H,C)\quad\Longleftrightarrow\quad{}
 &\mixin_H(C)\in\dom(P.\mathit{mixin})\ \land\\
 &H(C).C_m\in\mathsf{classRecordIds}(H)\ \land\\
 &H(C).eo\in\dom(H)\cup\{\nil\}\ \land
   (C=\Object\ \lor\ \super_H(C)\ne\Top).
\end{split}
\end{equation}
The complete class-graph check is
\begin{equation}\label{eq:reflective-class-graph-validity}
\begin{split}
 \mathsf{classGraphOK}(P,H)\quad\Longleftrightarrow\quad{}
 &\Object\in\mathsf{classRecordIds}(H)\ \land
   \super_H(\Object)=\Top\ \land\\
 &\forall C\in\mathsf{classRecordIds}(H).\;
      \mathsf{reachesTop}_{H}(C)\ \land\\
 &\forall C\in\mathsf{classRecordIds}(H)\setminus\{\Top\}.\;
      \mathsf{classRecordOK}(P,H,C).
\end{split}
\end{equation}
The declaration retained by class provenance and the complete ordinary-object
check are defined by
\begin{equation}\label{eq:class-origin-declaration}
 \mathsf{originDecl}(\kw{expression}(K))=K,
 \qquad \mathsf{originDecl}(\kw{actor}(K))=K.
\end{equation}
The cached class for one nested declaration is valid exactly when
\begin{equation}\label{eq:cached-nested-class-validity}
\begin{split}
 \mathsf{cachedNestedClassOK}(H,o,d)\quad\Longleftrightarrow\quad{}
 &H(o).\eta(d)\in\mathsf{classRecordIds}(H)\ \land\\
 &H(H(o).\eta(d)).eo=o\ \land\\
 &\mathsf{originDecl}(H(H(o).\eta(d)).\chi_C)=d.
\end{split}
\end{equation}
An ordinary object record is valid when
\begin{equation}\label{eq:reflective-object-record-validity}
\begin{aligned}
 \mathsf{objectRecordOK}(P,H,o)\quad\Longleftrightarrow\quad{}
 &\classof_H(o)\in\mathsf{classRecordIds}(H)\ \land\\
 &\dom(H(o).\sigma)=
   \mathsf{set}(\mathsf{layout}_{P,H}(\classof_H(o)))\ \land\\
 &\dom(H(o).\eta)\subseteq
   \mathsf{nestedKeys}_{P,H}(\classof_H(o))\ \land\\
 &\forall d\in\dom(H(o).\eta).\;\mathsf{cachedNestedClassOK}(H,o,d).
\end{aligned}
\end{equation}
The store-wide object check is
\begin{equation}\label{eq:reflective-object-validity}
 \mathsf{objectStoreOK}(P,H)\quad\Longleftrightarrow\quad
 \forall o\in\mathsf{ordinaryObjectIds}(H).\;
   \mathsf{objectRecordOK}(P,H,o).
\end{equation}

For activation validation, first extract the activation identities from each
actor's live stack:
\begin{equation}\label{eq:live-activation-order}
\begin{gathered}
 \mathsf{activationIds}(\epsilon)=\epsilon,\\
 \mathsf{activationIds}(\mathcal A\mathbin{\cdot}\langle a,\pi\rangle)
   =\mathsf{activationIds}(\mathcal A)\mathbin{\cdot}\langle a\rangle,\\
 \mathsf{stackActivationIds}_{\mathcal X}(A)=\epsilon
   \quad\text{if }\mathcal X(A)=\kw{idle},\\
 \mathsf{stackActivationIds}_{\mathcal X}(A)=\mathsf{activationIds}(\mathcal A),\\
 \hfill\text{if }\mathcal X(A)=\kw{runningTurn}\langle\mathcal A,t,p,u\rangle,\\
 \mathsf{stackActivationIds}_{\mathcal X}(A)=\mathsf{activationIds}(\mathcal A),\\
 \hfill\text{if }\mathcal X(A)=
   \kw{pausedTurn}\langle\mathcal A,t,p,u,r_s,d_s\rangle,\\
 \mathsf{liveActivationIds}(\mathcal X)
   =\bigcup_{A\in\dom(\mathcal X)}
       \mathsf{set}(\mathsf{stackActivationIds}_{\mathcal X}(A)),\\
 \mathsf{below}_{\mathcal X}(A,k,a)\Longleftrightarrow\exists i<j.\\
 \mathsf{stackActivationIds}_{\mathcal X}(A)[i]=k\ \land
 \mathsf{stackActivationIds}_{\mathcal X}(A)[j]=a.
\end{gathered}
\end{equation}
Definedness of every partial expression below uses the notation
\begin{equation}\label{eq:definedness-notation}
 e_p\downarrow\ \Longleftrightarrow\ \exists v.\;e_p=v,
 \qquad e_p\uparrow\ \Longleftrightarrow\ \neg(e_p\downarrow).
\end{equation}
Surviving method declarations determine the parameter domain of a live method
activation:
\begin{equation}\label{eq:method-formals}
 \mathsf{methodFormals}_{P}(i_f)=\bar x
 \Longleftrightarrow
 \exists M,s,\alpha,\ell,b,K.\;
 \methods_P(M)(s)=\langle i_f,s,\alpha,\bar x,\ell,b,K\rangle.
\end{equation}
The validity of one live activation is
\begin{equation}\label{eq:live-activation-record-validity}
\begin{aligned}
 \mathsf{liveActivationRecordOK}(P,H,\mathcal X,a)
 &\Longleftrightarrow
 \exists p,\rho,\ell,o,C,k,h,q_a,\gamma.\;\\[-1mm]
 &\quad H(a)=\kw{activation}\langle
       p,\rho,\ell,o,C,k,h,q_a,\gamma\rangle\ \land\\
 &\quad\bigl((p=\kw{top}\ \land\ \rho=\ell=\varnothing\ \land
              o=C=h=\nil)\ \lor\\
 &\qquad(p\ne\kw{top}\ \land\ o\in\dom(H)\ \land
          C\preceq_H\classof_H(o)\ \land\\
 &\hspace{39mm}h\in\mathsf{activationRecordIds}(H))\bigr)\ \land\\
 &\quad\bigl(k=\nil\ \lor
       \exists A.\;\mathsf{below}_{\mathcal X}(A,k,a)\bigr)\ \land\\
 &\quad\bigl(\mathsf{methodFormals}_{P}(p)\uparrow\ \lor\\
 &\hspace{22mm}\dom(\rho)=
       \mathsf{set}(\mathsf{methodFormals}_{P}(p))\bigr).
\end{aligned}
\end{equation}
The aggregate check is
\begin{equation}\label{eq:live-activation-validity}
\begin{split}
 \mathsf{liveActivationsOK}(P,H,\mathcal X)
 \quad\Longleftrightarrow\quad{}
 &\forall a\in\mathsf{liveActivationIds}(\mathcal X).\\[-1mm]
 &\mathsf{liveActivationRecordOK}(P,H,\mathcal X,a).
\end{split}
\end{equation}
Equation~\eqref{eq:live-activation-validity} therefore permits an already
materialized activation of a removed method, but requires the bound parameter
domain to agree with every surviving edited signature.

Let $H_0=\mathsf{vmStore}(\Omega)$ and let $H'=\mathsf{vmStore}(\Omega')$.
The command-specific activation conditions are the following total predicate:
\begin{equation}\label{eq:activation-parameter-edit-validity}
\begin{split}
 &\mathsf{activationEditOK}_{H_0,H',\mathcal X'}
   (\kw{writeActivationParameter}(g,a,x,v))\\
 &\qquad\Longleftrightarrow
   a\in\mathsf{activationRecordIds}(H_0)\ \land x\in\dom(H_0(a).\rho).
\end{split}
\end{equation}
\begin{equation}\label{eq:activation-local-edit-validity}
\begin{split}
 &\mathsf{activationEditOK}_{H_0,H',\mathcal X'}
   (\kw{writeActivationLocal}(g,a,j,v))\\
 &\qquad\Longleftrightarrow
   a\in\mathsf{activationRecordIds}(H_0)\ \land j\in\dom(H_0(a).\ell).
\end{split}
\end{equation}
\begin{equation}\label{eq:activation-class-edit-validity}
\begin{split}
 &\mathsf{activationEditOK}_{H_0,H',\mathcal X'}
   (\kw{changeActivationCurrentClass}(g,a,C))\\
 &\qquad\Longleftrightarrow
   a\in\mathsf{activationRecordIds}(H')\ \land
   C\preceq_{H'}\classof_{H'}(H'(a).o).
\end{split}
\end{equation}
\begin{equation}\label{eq:activation-continuation-edit-validity}
\begin{split}
 &\mathsf{activationEditOK}_{H_0,H',\mathcal X'}
   (\kw{changeActivationContinuation}(g,a,k))\\
 &\qquad\Longleftrightarrow
   a\in\mathsf{liveActivationIds}(\mathcal X')\ \land\\
 &\hspace{34mm}(k=\nil\ \lor\
      \exists A.\;\mathsf{below}_{\mathcal X'}(A,k,a)).
\end{split}
\end{equation}
\begin{equation}\label{eq:activation-uncontinuable-edit-validity}
\begin{split}
 &\mathsf{activationEditOK}_{H_0,H',\mathcal X'}
   (\kw{markActivationUncontinuable}(g,a))\\
 &\qquad\Longleftrightarrow a\in\mathsf{activationRecordIds}(H_0),\\
 &\mathsf{activationEditOK}_{H_0,H',\mathcal X'}(\delta)
   \Longleftrightarrow\kw{true}
   \quad\text{if }\delta\notin\mathit{ActivationCommand}.
\end{split}
\end{equation}
Continuation transfer is already partial in
Equation~\eqref{eq:reflective-continuation-transfer}; its transaction-level
cardinality and the remaining activation edits are checked by
\begin{equation}\label{eq:continuation-transfer-count}
 \mathsf{transferCount}(\Delta,B)=
 \left|\{\delta\in\Delta\mid\exists g,a,v.\;
       \delta=\kw{continueAtActivation}(g,B,a,v)\}\right|.
\end{equation}
\begin{equation}\label{eq:activation-transaction-validity}
\begin{aligned}
 \mathsf{activationEditsOK}(\Omega,\Delta,\Omega')
 \quad\Longleftrightarrow\quad{}
 &\bigl(\forall\delta\in\Delta.\;
   \mathsf{activationEditOK}_{H_0,H',\mathcal X'}(\delta)\bigr)\ \land\\
 &\forall B\in\dom(\mathcal X').\;
   \mathsf{transferCount}(\Delta,B)\le1.
\end{aligned}
\end{equation}

Let $\mathsf{refs}$ be the structural reference extractor, defined for every
runtime constructor, finite map, and sequence by
\begin{equation}\label{eq:structural-reference-extractor}
\begin{aligned}
 \mathsf{refs}(x)&=\{x\}&&x\in\mathit{Ref},\\
 \mathsf{refs}(z)&=\varnothing&&z\text{ is atomic and }z\notin\mathit{Ref},\\
 \mathsf{refs}(\mathsf{ctor}(z_1,\ldots,z_k))
   &=\bigcup_{i=1}^{k}\mathsf{refs}(z_i),\\
 \mathsf{refs}([z_i\mapsto z_i']_{i\in I})
   &=\bigcup_{i\in I}(\mathsf{refs}(z_i)\cup\mathsf{refs}(z_i')).
\end{aligned}
\end{equation}
The actor-residency check is then
\begin{equation}\label{eq:reflective-actor-isolation}
\begin{aligned}
 \mathsf{resident}_{\Omega}(A)
  &=\dom(V)\cup\dom(\mathcal H(A))\cup\dom(\Pi)\cup\dom(\Phi),\\
 \mathsf{actorIsolationOK}(\Omega)\quad\Longleftrightarrow\quad{}
 &\forall A\in\dom(\mathcal H).\;
   \mathsf{refs}(\mathcal X(A))\subseteq\mathsf{resident}_{\Omega}(A)\ \land\\
 &\quad\forall x\in\dom(\mathcal H(A)).\;
   \mathsf{refs}(\mathcal H(A)(x))\subseteq\mathsf{resident}_{\Omega}(A)\ \land\\
 &\forall x\in\dom(V).\;
   \mathsf{refs}(V(x))\subseteq\dom(V)\ \land
   \mathsf{value}_{P,\mathsf{vmStore}(\Omega)}(x).
\end{aligned}
\end{equation}

The finite validation-kind domain, its interpretation, and its fixed order are
{\small
\begin{equation}\label{eq:reflection-validation-kinds}
\begin{gathered}
 \kappa_r::=\kw{classGraphCheck}\mid\kw{objectStoreCheck}
  \mid\kw{liveActivationCheck}\mid\kw{activationEditCheck}
  \mid\kw{actorIsolationCheck},\\
 \kw{classGraphCheck}\mathrel{\triangleleft_{\mathsf{check}}}
 \kw{objectStoreCheck}\mathrel{\triangleleft_{\mathsf{check}}}
 \kw{liveActivationCheck},\\
 \kw{liveActivationCheck}\mathrel{\triangleleft_{\mathsf{check}}}
 \kw{activationEditCheck}\mathrel{\triangleleft_{\mathsf{check}}}
 \kw{actorIsolationCheck},\\
 \mathsf{validationHolds}(\kw{classGraphCheck},\Omega,\Delta,\Omega')
  \Longleftrightarrow\mathsf{classGraphOK}(P',H'),\\
 \mathsf{validationHolds}(\kw{objectStoreCheck},\Omega,\Delta,\Omega')
  \Longleftrightarrow\mathsf{objectStoreOK}(P',H'),\\
 \mathsf{validationHolds}(\kw{liveActivationCheck},\Omega,\Delta,\Omega')
  \Longleftrightarrow\mathsf{liveActivationsOK}(P',H',\mathcal X'),\\
 \mathsf{validationHolds}(\kw{activationEditCheck},\Omega,\Delta,\Omega')
  \Longleftrightarrow\mathsf{activationEditsOK}(\Omega,\Delta,\Omega'),\\
 \mathsf{validationHolds}(\kw{actorIsolationCheck},\Omega,\Delta,\Omega')
  \Longleftrightarrow\mathsf{actorIsolationOK}(\Omega').
\end{gathered}
\end{equation}
}
Define the failed validation set by
\begin{equation}\label{eq:failed-reflection-validations}
 \mathsf{failedValidations}(\Omega,\Delta,\Omega')=
 \{\kappa_r\mid
   \neg\mathsf{validationHolds}(\kappa_r,\Omega,\Delta,\Omega')\}.
\end{equation}
For $\Delta=\langle\delta_1,\ldots,\delta_m\rangle$, also define
\begin{equation}\label{eq:reflection-error-index-sets}
\begin{aligned}
 \mathsf{unauthorizedIndices}(\Omega,A,\Delta)
  &=\{i\mid1\le i\le m\ \land
       \neg\mathsf{authorized}_{\Omega,A}(\delta_i)\},\\
 \mathsf{conflictPairs}(\Delta)
  &=\{(i,j)\mid\mathsf{conflict}_{\Delta}(i,j)\}.
\end{aligned}
\end{equation}
The validity predicate now has no informal conjuncts:
\begin{equation}\label{eq:reflection-validity}
\begin{aligned}
 \mathsf{validReflection}(\Omega,A,\Delta,\Omega')
 \quad\Longleftrightarrow\quad{}
 &\mathsf{prepareReflection}(\Omega,A,\Delta)=\Omega'\ \land\\
 &\mathsf{unauthorizedIndices}(\Omega,A,\Delta)=\varnothing\ \land\\
 &\mathsf{conflictPairs}(\Delta)=\varnothing\ \land\\
 &\mathsf{failedValidations}(\Omega,\Delta,\Omega')=\varnothing.
\end{aligned}
\end{equation}

Successful preparation commits in one VM step:
\[
\dfrac{\mathsf{validReflection}(\Omega,A,\Delta,\Omega')}
 {\Omega\xrightarrow{\;\mathsf{commitReflection}(A,\Delta,n+1)\;}\Omega'}
\quad(\rulelabel{Reflection-Commit})
\]
The rejection diagnostic domain is
\begin{equation}\label{eq:reflection-failure-domain}
\begin{aligned}
 \xi_r::={}&\kw{unauthorizedCommand}(i)
 \mid\kw{commandConflict}(i,j)
 \mid\kw{sourcePatchFailure}\mid\kw{elaborationFailure}\\
 &\mid\kw{statePreparationFailure}\mid\kw{invariantFailure}(\kappa_r).
\end{aligned}
\end{equation}
For compactness, define the two partial candidates
\begin{equation}\label{eq:reflection-diagnostic-candidates}
\begin{aligned}
 \mathsf{reflectionSourceCandidate}(P,\Delta)
  &=\mathsf{applySourceCommands}(P,\mathsf{codeCommands}(\Delta)),\\
 \mathsf{reflectionProgramCandidate}(P,\Delta)
  &=\mathsf{reelaborateProgram}
      (P,\mathsf{reflectionSourceCandidate}(P,\Delta)).
\end{aligned}
\end{equation}
Writing
$\mathcal U_r=\mathsf{unauthorizedIndices}(\Omega,A,\Delta)$ and
$\mathcal C_r=\mathsf{conflictPairs}(\Delta)$, the previously implicit failure
function is the following partial function:
{\small
\begin{equation}\label{eq:reflection-failure}
\mathsf{reflectionFailure}(\Omega,A,\Delta)=
\begin{cases}
 \kw{unauthorizedCommand}(\min\mathcal U_r)
   &\mathcal U_r\ne\varnothing,\\
 \kw{commandConflict}(i,j)
   &\begin{gathered}\mathcal U_r=\varnothing\ \land\
       \mathcal C_r\ne\varnothing\ \land\\
       (i,j)=\min_{\rm lex}\mathcal C_r,
    \end{gathered}\\
 \kw{sourcePatchFailure}
   &\begin{gathered}\mathcal U_r=\mathcal C_r=\varnothing\ \land\\
    \mathsf{reflectionSourceCandidate}(P,\Delta)\uparrow,
    \end{gathered}\\
 \kw{elaborationFailure}
   &\begin{gathered}\mathcal U_r=\mathcal C_r=\varnothing\ \land\\
    \mathsf{reflectionSourceCandidate}(P,\Delta)\downarrow\ \land\\
    \mathsf{reflectionProgramCandidate}(P,\Delta)\uparrow,
    \end{gathered}\\
 \kw{statePreparationFailure}
   &\begin{gathered}\mathcal U_r=\mathcal C_r=\varnothing\ \land\\
    \mathsf{reflectionProgramCandidate}(P,\Delta)\downarrow\ \land\\
    \mathsf{prepareReflection}(\Omega,A,\Delta)\uparrow,
    \end{gathered}\\
 \kw{invariantFailure}(\min_{\triangleleft_{\mathsf{check}}}\mathcal B_r)
   &\begin{gathered}\mathcal U_r=\mathcal C_r=\varnothing\ \land\\
    \mathsf{prepareReflection}(\Omega,A,\Delta)=\Omega'\ \land\\
    \mathcal B_r=\mathsf{failedValidations}(\Omega,\Delta,\Omega')
       \ne\varnothing.
    \end{gathered}
\end{cases}
\end{equation}
}
It is undefined exactly when
$\mathsf{validReflection}(\Omega,A,\Delta,\Omega')$ holds for some $\Omega'$.
Failure leaves the state identical:
\[
\dfrac{\mathsf{reflectionFailure}(\Omega,A,\Delta)=\xi}
 {\Omega\xrightarrow{\;\mathsf{rejectReflection}(A,\Delta,\xi)\;}\Omega}
\quad(\rulelabel{Reflection-Reject})
\]
The case order in Equation~\eqref{eq:reflection-failure} affects only the
diagnostic; it cannot expose an intermediate state.  Except for a successful
Equation~\eqref{eq:replace-current-actor-stack}, whose requesting frame has
itself been replaced, the mirror library reifies the commit receipt or failure
as the return or exception of the primitive mirror message.  Current-stack
replacement instead continues directly from the installed top control.

\subsection{Causal connection, turns, and other VMs}

A reflective transition is a VM safepoint transition: it may occur between any
two actor, delivery, or sequential steps, but never inside one such step.  It
replaces the single $P$ in Equation~\eqref{eq:actor-configuration}; hence every
actor in the cohort observes the same version $n+1$ immediately.  The following
definition fixes the boundary between old and new behavior:
\begin{equation}\label{eq:reflection-version-observation}
\begin{array}{c|c}
\text{computation state at commit}&\text{program observed afterward}\\ \hline
\text{already materialized control }t\text{ and frames }\pi
  &\text{their existing terms, code image, and control map}\\
\text{future method resolution or invocation}&P'\\
\begin{gathered}\text{mailbox or network packet}\\[-1mm]
 \text{not yet dispatched}\end{gathered}&P'\text{ at its later turn}\\
\text{closure object already allocated}&\text{its captured body in }H(c)\\
\text{class or object retaining stable }M,C,i_f&
  \text{the updated records addressed by those identities.}
\end{array}
\end{equation}
Thus an install never splices a new expression into a currently executing
small step.  A method invocation that has already entered its body continues
that body; the next dispatch through the same stable method identity uses its
new definition.  In-flight asynchronous messages carry selectors and
arguments, not resolved methods, so they naturally dispatch against the
version current when Rule \rulelabel{Turn-Start} begins their turn.

Let a distributed system be
\begin{equation}\label{eq:vm-family}
 \mathfrak F:\mathit{VMId}\rightharpoonup\mathit{ActorConfiguration},
 \qquad \mathfrak F(w).w=w.
\end{equation}
Rule \rulelabel{Reflection-Commit} changes exactly one component
$\mathfrak F(w)$.  No rule copies $P'$ into $\mathfrak F(w')$ for $w'\ne w$.
Values in $V$ are shared by every actor of $w$, so their behavior changes with
the cohort's code; a value transferred to another VM is copied under
Equation~\eqref{eq:value-transfer-contract} and thereafter uses the recipient
VM's independently versioned program.  A remote VM may be changed only by a
separate, authorized mirror request delivered to that VM.  Its delivery is
ordered like any other eventual message; the resulting local commit is atomic
at the recipient safepoint.

The same rule governs remote debugging.  An actor mirror transferred to
another actor or VM remains a far capability.  Its \kw{actorStackQuery} result
is transported by the remote-representation judgment; object references in a
frame become values, far references, or authorized mirrors rather than raw
near references.  A remote controller first applies
Equation~\eqref{eq:pause-actor-effect}, reads the resulting paused stack
snapshot, and then names its pause token in
Equation~\eqref{eq:replace-paused-actor-stack},
\eqref{eq:resume-paused-actor-full-speed}, or
\eqref{eq:replace-resume-paused-actor-stack}.  Formally, a remote debugger step
has the component-wise shape
\begin{equation}\label{eq:remote-debugger-component-step}
 \frac{\mathfrak F(w),A,g\vdash\mathsf{debugRequest}(\delta)
          \Downarrow\mathfrak F'(w)}
      {\mathfrak F\longmapsto_{\mathsf{remoteDebug}}
       \mathfrak F[w\mapsto\mathfrak F'(w)]},
 \qquad
 \forall v\ne w.\ \mathfrak F'(v)=\mathfrak F(v).
\end{equation}
The request is delivered using the ordinary cross-VM transport and executes
the corresponding local reflective command at $w$'s safepoint.  There is no
distributed atomic transaction spanning several VMs: coordinated multi-actor
debugging is a protocol of individually atomic pause, inspect, install, and
resume operations.  Pause tokens reject delayed or replayed commands after
either an intervening edit or a resume.

Finally, VM primitives are ordinary messages to a mirror whose record has
target $\kw{vmTarget}(w)$ and right $\kw{primitive}$.  Their platform-specific state
transformers lie outside this language semantics, but the capability premise
is mandatory:
\[
\dfrac{\begin{gathered}
       \widehat H_A(g)=\kw{mirror}\langle C_g,w,\kw{vmTarget}(w),R\rangle
       \qquad\kw{primitive}\in R\\
       \mathsf{hostPrimitive}(\Omega,s,\bar o)=\langle\Omega',v\rangle
       \end{gathered}}
 {\Omega,A,g\vdash\mathsf{invokeVMPrimitive}(s,\bar o)
       \Downarrow\langle\Omega',v\rangle}
\quad(\rulelabel{VM-Primitive})
\]
There is no primitive-call syntax and no bypass around the mirror record.


\section{Optional garbage-collection layer}
\label{sec:garbage-collection}

This section is an optional conservative extension of the VM semantics.  It is
kept in this separate source file, and no preceding rule depends on one of its
auxiliary definitions.  Omitting the single
\verb|\section{Optional garbage-collection layer}
\label{sec:garbage-collection}

This section is an optional conservative extension of the VM semantics.  It is
kept in this separate source file, and no preceding rule depends on one of its
auxiliary definitions.  Omitting the single
\verb|\section{Optional garbage-collection layer}
\label{sec:garbage-collection}

This section is an optional conservative extension of the VM semantics.  It is
kept in this separate source file, and no preceding rule depends on one of its
auxiliary definitions.  Omitting the single
\verb|\input{gc-semantics}| line restores the monotonically growing heaps of
Sections~3--11.  Including it replaces that implementation abstraction with an
explicit, nondeterministic collector.  The collector introduces no Newspeak
syntax and never runs inside an actor step or reflective transaction.

The construction follows the explicit-heap rewrite method of Morrisett,
Felleisen, and Harper, \emph{Abstract Models of Memory Management} (1995), and
the treatment of observable collection in Donnelly, Hallett, and Kfoury,
\emph{Formal Semantics of Weak References} (2006).  The distinctive Newspeak
observations are Equation~\eqref{eq:all-instances-of}, whose implicit instance
registry is weak, and the platform classes \kw{WeakArray} and \kw{WeakMap}.
Collection can therefore change both later population snapshots and later
reads of weak containers.

\subsection{The collectable store}

Garbage collection applies to ordinary heap records and to the routing records
of promise and far-reference handles.  For
$\Omega=\langle P,V,\mathcal H,\mathcal X,\mathcal Q,\mathcal N,
\Pi,\Phi,\Gamma,\mathcal E,\mathcal D,w,n\rangle$, define the disjoint reference
store
\begin{equation}\label{eq:gc-reference-store}
 \mathsf{refStore}(\Omega)
   =\mathsf{vmStore}(\Omega)\uplus\Pi\uplus\Phi.
\end{equation}
The union is disjoint because its three domains are respectively subsets of
$\mathit{ObjRef}$, $\mathit{PromiseId}$, and $\mathit{FarRefId}$.  Put
$\mathcal S_\Omega=\mathsf{refStore}(\Omega)$ and expose a uniform payload
projection:
\begin{equation}\label{eq:gc-reference-payload}
 \mathsf{payload}_{\Omega}(x)=
 \begin{cases}
  \mathsf{vmStore}(\Omega)(x) & x\in\mathit{ObjRef},\\
  \Pi(x) & x\in\mathit{PromiseId},\\
  \Phi(x) & x\in\mathit{FarRefId}.
 \end{cases}
 \qquad x\in\dom(\mathcal S_\Omega).
\end{equation}
Thus a live promise strongly retains its settlement value and registered
waiters, and a live far-reference handle strongly retains its near target.
The tables $\Pi$ and $\Phi$ do not retain their own keys merely by containing
them; otherwise every handle and every remote target would live forever.

\subsection{WeakArray and WeakMap}
\label{sec:weak-containers}

\kw{WeakArray} and \kw{WeakMap} are mutable platform objects and therefore
belong to exactly one actor heap.  Their concrete VM layouts are abstracted by
two additional object-record forms:
\begin{equation}\label{eq:weak-container-records}
\begin{aligned}
 H_A(a_w)&=\kw{weakArray}\langle C,\langle y_1,\ldots,y_m\rangle\rangle,
   &a_w&\in\mathit{WeakArrayId},\\
 H_A(m_w)&=\kw{weakMap}\langle C,E_w\rangle,
   &m_w&\in\mathit{WeakMapId},
\end{aligned}
\end{equation}
where $\mathit{WeakArrayId}$ and $\mathit{WeakMapId}$ are disjoint subsets of
$\mathit{ObjRef}$, every $y_i\in\mathit{Ref}$, and
$E_w:\mathit{ObjRef}\rightharpoonup\mathit{Ref}$ is finite.  $C$ and any
ordinary fixed fields of the implementation are strong; the sequence elements
and map keys are weak.  A map value is conditionally strong exactly while its
key is strongly reachable.  This is the usual ephemeron interpretation of a
weak-key map: a value-to-key cycle cannot keep the key alive.

Every stored weak identity initially resolves; weakness controls retention, not
whether an arbitrary dangling identity may be written:
\begin{equation}\label{eq:weak-container-validity}
\begin{aligned}
 \mathsf{weakArrayStoreOK}(\Omega)
 &\Longleftrightarrow
   \forall A,a_w,C,\bar y.\
     \mathcal H(A)(a_w)=\kw{weakArray}\langle C,\bar y\rangle\\[-1mm]
 &\hspace{35mm}\Longrightarrow
     \mathsf{refs}(\bar y)\subseteq\dom(\mathcal S_\Omega),\\
 \mathsf{weakMapStoreOK}(\Omega)
 &\Longleftrightarrow
   \forall A,m_w,C,E_w.\
     \mathcal H(A)(m_w)=\kw{weakMap}\langle C,E_w\rangle\\[-1mm]
 &\hspace{35mm}\Longrightarrow
     \dom(E_w)\cup\mathsf{refs}(\mathsf{range}(E_w))
       \subseteq\dom(\mathcal S_\Omega),\\
 \mathsf{weakContainersOK}(\Omega)
 &\Longleftrightarrow
   \mathsf{weakArrayStoreOK}(\Omega)\land\mathsf{weakMapStoreOK}(\Omega).
\end{aligned}
\end{equation}

The platform method contract is given by distinct partial store functions
rather than overloaded notation:
\begin{equation}\label{eq:weak-array-operations}
\begin{aligned}
 \mathsf{weakArrayRead}(H_A,a_w,i)&=y_i,\\
 \mathsf{writeWeakArrayCell}(H_A,a_w,i,v)
   &=H_A[a_w\mapsto
       \kw{weakArray}\langle C,\bar y[i\mapsto v]\rangle],\\
 &\hspace{35mm}1\le i\le |\bar y|.
\end{aligned}
\end{equation}
\begin{equation}\label{eq:weak-map-operations}
\begin{aligned}
 \mathsf{weakMapLookup}(H_A,m_w,k)&=
   \begin{cases}E_w(k)&k\in\dom(E_w),\\ \nil&k\notin\dom(E_w),\end{cases}\\
 \mathsf{writeWeakMapEntry}(H_A,m_w,k,v)
   &=H_A[m_w\mapsto\kw{weakMap}\langle C,E_w[k\mapsto v]\rangle].
\end{aligned}
\end{equation}
Here $H_A(a_w)$ and $H_A(m_w)$ have the respective forms of
Equation~\eqref{eq:weak-container-records}.  Allocation creates an all-$\nil$
array of the requested nonnegative size or a map with $E_w=\varnothing$, using
the ordinary global-freshness obligation of Equation~\eqref{eq:gc-wrapper-state}.
The two writes are defined only when their new key and value references resolve
in $\mathcal S_\Omega$, and they preserve Equation~
\eqref{eq:weak-container-validity}.
Out-of-bounds array indices are undefined and are reified by the platform
exception boundary.  Storing $\nil$ is permitted.  Since lookup also returns
$\nil$ for an absent key, the public \kw{WeakMap} protocol deliberately does
not distinguish absence from a stored $\nil$ value.

The state outside the collectable store is
\begin{equation}\label{eq:gc-control-state}
 \mathsf{control}(\Omega)=
 \langle P,\mathcal X,\mathcal Q,\mathcal N,
          \Gamma,\mathcal E,\mathcal D,w,n\rangle.
\end{equation}
Let $\mathcal R_h\subseteq\dom(\mathcal S_\Omega)$ be the finite set of strong
references temporarily retained by the embedding host, such as arguments of a
platform primitive that is atomic with respect to collection.  A closed
Newspeak execution has $\mathcal R_h=\varnothing$.  First erase the weak
components of a payload while preserving its class and ordinary fields:
\begin{equation}\label{eq:gc-strong-payload}
\begin{aligned}
 \mathsf{strongPayload}(\kw{weakArray}\langle C,\bar y\rangle)
   &=\kw{weakArray}\langle C,\bar{\nil}\rangle,\\
 \mathsf{strongPayload}(\kw{weakMap}\langle C,E_w\rangle)
   &=\kw{weakMap}\langle C,\varnothing\rangle,\\
 \mathsf{strongPayload}(r)&=r
   \quad\text{for every other runtime record }r.
\end{aligned}
\end{equation}
The direct roots, unconditional strong successors, and map-value successors
enabled by a candidate live set $L$ are
\begin{equation}\label{eq:gc-roots-and-successors}
\begin{aligned}
 \mathsf{roots}_{\Omega}(\mathcal R_h)
   &=\bigl(\mathcal R_h\cup
       \mathsf{refs}(\mathsf{control}(\Omega))\bigr)
       \cap\dom(\mathcal S_\Omega),\\
 \mathsf{strongSucc}_{\Omega}(x)
   &=\mathsf{refs}(\mathsf{strongPayload}(
        \mathsf{payload}_{\Omega}(x)))
       \cap\dom(\mathcal S_\Omega),\\
 \mathsf{enabledMapValues}_{\Omega}(x,L)
   &=\begin{cases}
      \displaystyle
      \bigcup_{\substack{k\in\dom(E_w)\\ k\in L}}
        \mathsf{refs}(E_w(k))\cap\dom(\mathcal S_\Omega)
       & \mathsf{payload}_{\Omega}(x)=
           \kw{weakMap}\langle C,E_w\rangle,\\[3mm]
      \varnothing & \text{otherwise.}
     \end{cases}
\end{aligned}
\end{equation}
The structural extractor $\mathsf{refs}$ is Equation~
\eqref{eq:structural-reference-extractor}.  Static platform objects and atoms
are roots through $P$; running activations and their terms are roots through
$\mathcal X$; queued and in-flight arguments, reply promises, and settlement
values are roots through $\mathcal Q$ and $\mathcal N$.  Event identities and
actor identities that do not occur in $\dom(\mathcal S_\Omega)$ are discarded
by the intersection.

Strong reachability is the least fixed point
\begin{equation}\label{eq:gc-reachability}
\begin{aligned}
 \mathsf{reach}^{0}_{\Omega}(\mathcal R_h)
   &=\mathsf{roots}_{\Omega}(\mathcal R_h),\\
 \mathsf{reach}^{i+1}_{\Omega}(\mathcal R_h)
   &=\mathsf{reach}^{i}_{\Omega}(\mathcal R_h)
     \cup\bigcup_{x\in\mathsf{reach}^{i}_{\Omega}(\mathcal R_h)}
       \bigl(\mathsf{strongSucc}_{\Omega}(x)
        \cup\mathsf{enabledMapValues}_{\Omega}
          (x,\mathsf{reach}^{i}_{\Omega}(\mathcal R_h))\bigr),\\
 \mathsf{live}_{\Omega}(\mathcal R_h)
   &=\bigcup_{i\in\mathbb N}\mathsf{reach}^{i}_{\Omega}(\mathcal R_h),\\
 \mathsf{garbage}_{\Omega}(\mathcal R_h)
   &=\dom(\mathcal S_\Omega)
       \setminus\mathsf{live}_{\Omega}(\mathcal R_h).
\end{aligned}
\end{equation}
All stores are finite, so this sequence reaches a fixed point.  Cycles not
reachable from a root are garbage together; reference counting is neither
assumed nor implied.  Weak-array elements never enter the fixed point merely
because their array is live.  A weak-map value enters only after both its map
and its key have entered; consequently a value that points back to its otherwise
unreachable key does not activate its own entry.

\subsection{Admissible collection and identity history}

A collector need not reclaim all eligible records.  It may reclaim a nonempty
set $G$, provided no retained record would contain a dangling strong reference:
\begin{equation}\label{eq:admissible-garbage-set}
\begin{split}
 \mathsf{admissibleGC}_{\Omega,\mathcal R_h}(G)
 \quad\Longleftrightarrow\quad{}
 &\varnothing\ne G\subseteq
     \mathsf{garbage}_{\Omega}(\mathcal R_h)\ \land\\
 &\forall x\in\dom(\mathcal S_\Omega)\setminus G.\;
      \mathsf{strongSucc}_{\Omega}(x)\cap G=\varnothing.
\end{split}
\end{equation}
The second conjunct matters because $\kw{allInstancesOfQuery}$ can expose a
retained but previously unreachable object.  It rules out retaining such an
object while collecting another object to which one of its ordinary slots
still points.  Taking all of $\mathsf{garbage}_{\Omega}(\mathcal R_h)$ always
satisfies this condition; a generational or incremental collector may choose a
smaller closed subset.

Weak edges are intentionally absent from the closure premise.  If a selected
record is named by a retained weak array or by an inactive weak-map entry, the
collector clears that weak observation rather than rejecting the collection.
An enabled weak-map value cannot belong to $G$ because Equation~
\eqref{eq:gc-reachability} already makes it live.

Define weak-edge clearing on a retained payload by
\begin{equation}\label{eq:gc-clear-weak-payload}
\begin{aligned}
 \mathsf{clearWeak}_{G}
  (\kw{weakArray}\langle C,\langle y_1,\ldots,y_m\rangle\rangle)
  &=\kw{weakArray}\langle C,\langle y'_1,\ldots,y'_m\rangle\rangle,\\
 y'_i&=\begin{cases}
       \nil&\mathsf{refs}(y_i)\cap G\ne\varnothing,\\
       y_i&\mathsf{refs}(y_i)\cap G=\varnothing,
      \end{cases}\\
 \mathsf{clearWeak}_{G}(\kw{weakMap}\langle C,E_w\rangle)
  &=\kw{weakMap}\langle C,E'_w\rangle,\\
 E'_w&=E_w|_{\{k\mid k\notin G\ \land\
                    \mathsf{refs}(E_w(k))\cap G=\varnothing\}},\\
 \mathsf{clearWeak}_{G}(r)&=r
   \quad\text{for every other runtime record }r.
\end{aligned}
\end{equation}
The value-side test on $E'_w$ permits partial collection to remove an inactive
entry's value while leaving its still-unreachable key for a later collection;
no retained weak record then contains a dangling identity.

For $K=\dom(\mathcal S_\Omega)\setminus G$, let
\begin{equation}\label{eq:gc-pruned-actor-heap}
 \mathsf{prunedActorHeap}(H_A,G,K)=
 \{x\mapsto\mathsf{clearWeak}_{G}(H_A(x))
      \mid x\in\dom(H_A)\cap K\}.
\end{equation}
Store pruning is the total operation
\begin{equation}\label{eq:gc-prune-store}
\begin{split}
 \mathsf{prune}(\Omega,G)=\Omega[&
   V\mapsto V|_K,\quad
   \mathcal H(A)\mapsto
      \mathsf{prunedActorHeap}(\mathcal H(A),G,K)
      \text{ for every }A\in\dom(\mathcal H),\\
  &\Pi\mapsto\Pi|_K,\quad\Phi\mapsto\Phi|_K].
\end{split}
\end{equation}
Map restriction is by key.  Every other component, including causal and event
history, is unchanged.

Deleting a record must not make its identity fresh again.  Rather than adding
an allocation ledger to every existing configuration and rule, the optional
layer wraps the VM state:
\begin{equation}\label{eq:gc-wrapper-state}
 \Psi=\langle\Omega,\mathcal I\rangle,
 \qquad
 \dom(\mathcal S_\Omega)\subseteq\mathcal I\subseteq\mathit{Ref},
\end{equation}
where $\mathcal I$ is the finite set of reference identities allocated so far
in this execution.  Let the combined existing VM transition be the disjoint
union
\begin{equation}\label{eq:gc-underlying-transition}
\begin{aligned}
 \Omega\mathrel{\Longrightarrow^{\kw{actorStep}}}\Omega'
  &\Longleftrightarrow \Omega\longrightarrow_{\mathsf{act}}\Omega',\\
 \Omega\mathrel{\Longrightarrow^{\lambda_r}}\Omega'
  &\Longleftrightarrow \Omega\xrightarrow{\lambda_r}\Omega',
 \quad
 \lambda_r::=\kw{commitReflection}(A,\Delta,n)
       \mid\kw{rejectReflection}(A,\Delta,\xi_r).
\end{aligned}
\end{equation}
For one such transition define
\begin{equation}\label{eq:gc-new-identities}
 \mathsf{newRefs}(\Omega,\Omega')=
   \dom(\mathcal S_{\Omega'})\setminus\dom(\mathcal S_{\Omega}).
\end{equation}
The ordinary transition lifts only when all newly allocated identities have
never been used:
\[
\dfrac{
 \Omega\mathrel{\Longrightarrow^{\lambda}}\Omega'
 \qquad
 \mathsf{newRefs}(\Omega,\Omega')\cap\mathcal I=\varnothing
}{
 \mathcal R_h\vdash\langle\Omega,\mathcal I\rangle
   \mathrel{\Longrightarrow^{\lambda}_{\mathsf{gc}}}
 \langle\Omega',
   \mathcal I\cup\mathsf{newRefs}(\Omega,\Omega')\rangle
}
\quad(\rulelabel{GC-Lift})
\]
This gives every existing ``fresh'' premise its intended global meaning while
leaving the underlying rules unchanged.  An initial wrapped configuration has
$\mathcal I=\dom(\mathcal S_{\Omega_0})$.

Collection is an independently enabled VM-safepoint transition:
\[
\dfrac{
 \mathsf{admissibleGC}_{\Omega,\mathcal R_h}(G)
}{
 \mathcal R_h\vdash\langle\Omega,\mathcal I\rangle
   \longrightarrow_{\mathsf{gc}}
 \langle\mathsf{prune}(\Omega,G),\mathcal I\rangle
}
\quad(\rulelabel{GC-Collect})
\]
There is deliberately no fairness requirement for
\rulelabel{GC-Collect}: an eligible object may remain uncollected indefinitely.
The implementation chooses both when to collect and which admissible closed
garbage region to reclaim.

\subsection{Observable consequences}

Let $\Omega'=\mathsf{prune}(\Omega,G)$.  Equations~
\eqref{eq:gc-reachability}--\eqref{eq:gc-prune-store} immediately give the
safety obligations
\begin{equation}\label{eq:gc-safety}
\begin{aligned}
 \mathsf{live}_{\Omega}(\mathcal R_h)
   &\subseteq\dom(\mathcal S_{\Omega'}),\\
 \forall x\in\dom(\mathcal S_{\Omega'}).\quad
 \mathsf{refs}(\mathsf{payload}_{\Omega'}(x))
   \cap\mathit{Ref}
   &\subseteq\dom(\mathcal S_{\Omega'}).
\end{aligned}
\end{equation}
The first property says that no strongly reachable record is reclaimed.  The
second says that every retained encoded reference still resolves: selected
weak targets have already been replaced by $\nil$ or removed from their map.
Together
with Equation~\eqref{eq:gc-wrapper-state}, they state memory safety and
permanent identity freshness independently of any particular tracing
algorithm.

For exact-instance enumeration, collection has the intentionally observable
effect
\begin{equation}\label{eq:gc-all-instances-effect}
 \mathsf{allInstancesOf}_{\Omega'}(C)=
 \mathsf{filter}_{x\notin G}
   \bigl(\mathsf{allInstancesOf}_{\Omega}(C)\bigr).
\end{equation}
For mixin-application enumeration, Equation~\eqref{eq:mixin-applications}
similarly gives
\begin{equation}\label{eq:gc-mixin-applications-effect}
 \mathsf{mixinApplications}_{\Omega'}(M)=
 \mathsf{filter}_{C\notin G}
   \bigl(\mathsf{mixinApplications}_{\Omega}(M)\bigr).
\end{equation}
Let $H_A=\mathcal H(A)$ and
$H'_A=\mathsf{prunedActorHeap}(H_A,G,K)$.  For a retained weak array,
Equation~\eqref{eq:gc-clear-weak-payload} gives the observable read
\begin{equation}\label{eq:gc-weak-array-effect}
\begin{aligned}
 y&=\mathsf{weakArrayRead}(H_A,a_w,i),\qquad a_w\in K,\\
 \mathsf{weakArrayRead}(H'_A,a_w,i)
  &=\begin{cases}
     \nil & \mathsf{refs}(y)\cap G\ne\varnothing,\\
     y & \text{otherwise.}
    \end{cases}
\end{aligned}
\end{equation}
For a retained weak map, an entry disappears when collection removes its key
or any reference in its value:
\begin{equation}\label{eq:gc-weak-map-effect}
\begin{aligned}
 v&=\mathsf{weakMapLookup}(H_A,m_w,k),\qquad m_w\in K,\\
 \mathsf{weakMapLookup}(H'_A,m_w,k)
  &=\begin{cases}
     \nil & k\in G\ \lor\ \mathsf{refs}(v)\cap G\ne\varnothing,\\
     v & \text{otherwise.}
    \end{cases}
\end{aligned}
\end{equation}
These changes occur atomically with \rulelabel{GC-Collect}; no program step can
observe a cleared weak cell while the corresponding target record is still
being removed.  A strongly reachable key keeps its weak-map value live through
Equation~\eqref{eq:gc-reachability}, while neither the map nor its value keeps
an otherwise unreachable key live.
Consequently an unreachable instance or class application can appear in a
snapshot before a GC step and be absent after it.  Mirror inspection itself is
atomic with respect to \rulelabel{GC-Collect}; once a returned sequence is
installed in the caller's activation or heap, its elements are ordinary strong
references and are roots of the next collection through $\mathcal X$ or a heap
record.  Nothing can be resurrected after its record has been removed.

If programs have no GC-observing primitive, removing unreachable records is
expected to be a stuttering refinement of the uncollected semantics.  In this
semantics the population queries deliberately invalidate that premise, so the
timing of collection is part of the permitted observable nondeterminism.
Finalizers, if later added to the Newspeak platform, must extend the collection
rule here rather than modify the base evaluator.
| line restores the monotonically growing heaps of
Sections~3--11.  Including it replaces that implementation abstraction with an
explicit, nondeterministic collector.  The collector introduces no Newspeak
syntax and never runs inside an actor step or reflective transaction.

The construction follows the explicit-heap rewrite method of Morrisett,
Felleisen, and Harper, \emph{Abstract Models of Memory Management} (1995), and
the treatment of observable collection in Donnelly, Hallett, and Kfoury,
\emph{Formal Semantics of Weak References} (2006).  The distinctive Newspeak
observations are Equation~\eqref{eq:all-instances-of}, whose implicit instance
registry is weak, and the platform classes \kw{WeakArray} and \kw{WeakMap}.
Collection can therefore change both later population snapshots and later
reads of weak containers.

\subsection{The collectable store}

Garbage collection applies to ordinary heap records and to the routing records
of promise and far-reference handles.  For
$\Omega=\langle P,V,\mathcal H,\mathcal X,\mathcal Q,\mathcal N,
\Pi,\Phi,\Gamma,\mathcal E,\mathcal D,w,n\rangle$, define the disjoint reference
store
\begin{equation}\label{eq:gc-reference-store}
 \mathsf{refStore}(\Omega)
   =\mathsf{vmStore}(\Omega)\uplus\Pi\uplus\Phi.
\end{equation}
The union is disjoint because its three domains are respectively subsets of
$\mathit{ObjRef}$, $\mathit{PromiseId}$, and $\mathit{FarRefId}$.  Put
$\mathcal S_\Omega=\mathsf{refStore}(\Omega)$ and expose a uniform payload
projection:
\begin{equation}\label{eq:gc-reference-payload}
 \mathsf{payload}_{\Omega}(x)=
 \begin{cases}
  \mathsf{vmStore}(\Omega)(x) & x\in\mathit{ObjRef},\\
  \Pi(x) & x\in\mathit{PromiseId},\\
  \Phi(x) & x\in\mathit{FarRefId}.
 \end{cases}
 \qquad x\in\dom(\mathcal S_\Omega).
\end{equation}
Thus a live promise strongly retains its settlement value and registered
waiters, and a live far-reference handle strongly retains its near target.
The tables $\Pi$ and $\Phi$ do not retain their own keys merely by containing
them; otherwise every handle and every remote target would live forever.

\subsection{WeakArray and WeakMap}
\label{sec:weak-containers}

\kw{WeakArray} and \kw{WeakMap} are mutable platform objects and therefore
belong to exactly one actor heap.  Their concrete VM layouts are abstracted by
two additional object-record forms:
\begin{equation}\label{eq:weak-container-records}
\begin{aligned}
 H_A(a_w)&=\kw{weakArray}\langle C,\langle y_1,\ldots,y_m\rangle\rangle,
   &a_w&\in\mathit{WeakArrayId},\\
 H_A(m_w)&=\kw{weakMap}\langle C,E_w\rangle,
   &m_w&\in\mathit{WeakMapId},
\end{aligned}
\end{equation}
where $\mathit{WeakArrayId}$ and $\mathit{WeakMapId}$ are disjoint subsets of
$\mathit{ObjRef}$, every $y_i\in\mathit{Ref}$, and
$E_w:\mathit{ObjRef}\rightharpoonup\mathit{Ref}$ is finite.  $C$ and any
ordinary fixed fields of the implementation are strong; the sequence elements
and map keys are weak.  A map value is conditionally strong exactly while its
key is strongly reachable.  This is the usual ephemeron interpretation of a
weak-key map: a value-to-key cycle cannot keep the key alive.

Every stored weak identity initially resolves; weakness controls retention, not
whether an arbitrary dangling identity may be written:
\begin{equation}\label{eq:weak-container-validity}
\begin{aligned}
 \mathsf{weakArrayStoreOK}(\Omega)
 &\Longleftrightarrow
   \forall A,a_w,C,\bar y.\
     \mathcal H(A)(a_w)=\kw{weakArray}\langle C,\bar y\rangle\\[-1mm]
 &\hspace{35mm}\Longrightarrow
     \mathsf{refs}(\bar y)\subseteq\dom(\mathcal S_\Omega),\\
 \mathsf{weakMapStoreOK}(\Omega)
 &\Longleftrightarrow
   \forall A,m_w,C,E_w.\
     \mathcal H(A)(m_w)=\kw{weakMap}\langle C,E_w\rangle\\[-1mm]
 &\hspace{35mm}\Longrightarrow
     \dom(E_w)\cup\mathsf{refs}(\mathsf{range}(E_w))
       \subseteq\dom(\mathcal S_\Omega),\\
 \mathsf{weakContainersOK}(\Omega)
 &\Longleftrightarrow
   \mathsf{weakArrayStoreOK}(\Omega)\land\mathsf{weakMapStoreOK}(\Omega).
\end{aligned}
\end{equation}

The platform method contract is given by distinct partial store functions
rather than overloaded notation:
\begin{equation}\label{eq:weak-array-operations}
\begin{aligned}
 \mathsf{weakArrayRead}(H_A,a_w,i)&=y_i,\\
 \mathsf{writeWeakArrayCell}(H_A,a_w,i,v)
   &=H_A[a_w\mapsto
       \kw{weakArray}\langle C,\bar y[i\mapsto v]\rangle],\\
 &\hspace{35mm}1\le i\le |\bar y|.
\end{aligned}
\end{equation}
\begin{equation}\label{eq:weak-map-operations}
\begin{aligned}
 \mathsf{weakMapLookup}(H_A,m_w,k)&=
   \begin{cases}E_w(k)&k\in\dom(E_w),\\ \nil&k\notin\dom(E_w),\end{cases}\\
 \mathsf{writeWeakMapEntry}(H_A,m_w,k,v)
   &=H_A[m_w\mapsto\kw{weakMap}\langle C,E_w[k\mapsto v]\rangle].
\end{aligned}
\end{equation}
Here $H_A(a_w)$ and $H_A(m_w)$ have the respective forms of
Equation~\eqref{eq:weak-container-records}.  Allocation creates an all-$\nil$
array of the requested nonnegative size or a map with $E_w=\varnothing$, using
the ordinary global-freshness obligation of Equation~\eqref{eq:gc-wrapper-state}.
The two writes are defined only when their new key and value references resolve
in $\mathcal S_\Omega$, and they preserve Equation~
\eqref{eq:weak-container-validity}.
Out-of-bounds array indices are undefined and are reified by the platform
exception boundary.  Storing $\nil$ is permitted.  Since lookup also returns
$\nil$ for an absent key, the public \kw{WeakMap} protocol deliberately does
not distinguish absence from a stored $\nil$ value.

The state outside the collectable store is
\begin{equation}\label{eq:gc-control-state}
 \mathsf{control}(\Omega)=
 \langle P,\mathcal X,\mathcal Q,\mathcal N,
          \Gamma,\mathcal E,\mathcal D,w,n\rangle.
\end{equation}
Let $\mathcal R_h\subseteq\dom(\mathcal S_\Omega)$ be the finite set of strong
references temporarily retained by the embedding host, such as arguments of a
platform primitive that is atomic with respect to collection.  A closed
Newspeak execution has $\mathcal R_h=\varnothing$.  First erase the weak
components of a payload while preserving its class and ordinary fields:
\begin{equation}\label{eq:gc-strong-payload}
\begin{aligned}
 \mathsf{strongPayload}(\kw{weakArray}\langle C,\bar y\rangle)
   &=\kw{weakArray}\langle C,\bar{\nil}\rangle,\\
 \mathsf{strongPayload}(\kw{weakMap}\langle C,E_w\rangle)
   &=\kw{weakMap}\langle C,\varnothing\rangle,\\
 \mathsf{strongPayload}(r)&=r
   \quad\text{for every other runtime record }r.
\end{aligned}
\end{equation}
The direct roots, unconditional strong successors, and map-value successors
enabled by a candidate live set $L$ are
\begin{equation}\label{eq:gc-roots-and-successors}
\begin{aligned}
 \mathsf{roots}_{\Omega}(\mathcal R_h)
   &=\bigl(\mathcal R_h\cup
       \mathsf{refs}(\mathsf{control}(\Omega))\bigr)
       \cap\dom(\mathcal S_\Omega),\\
 \mathsf{strongSucc}_{\Omega}(x)
   &=\mathsf{refs}(\mathsf{strongPayload}(
        \mathsf{payload}_{\Omega}(x)))
       \cap\dom(\mathcal S_\Omega),\\
 \mathsf{enabledMapValues}_{\Omega}(x,L)
   &=\begin{cases}
      \displaystyle
      \bigcup_{\substack{k\in\dom(E_w)\\ k\in L}}
        \mathsf{refs}(E_w(k))\cap\dom(\mathcal S_\Omega)
       & \mathsf{payload}_{\Omega}(x)=
           \kw{weakMap}\langle C,E_w\rangle,\\[3mm]
      \varnothing & \text{otherwise.}
     \end{cases}
\end{aligned}
\end{equation}
The structural extractor $\mathsf{refs}$ is Equation~
\eqref{eq:structural-reference-extractor}.  Static platform objects and atoms
are roots through $P$; running activations and their terms are roots through
$\mathcal X$; queued and in-flight arguments, reply promises, and settlement
values are roots through $\mathcal Q$ and $\mathcal N$.  Event identities and
actor identities that do not occur in $\dom(\mathcal S_\Omega)$ are discarded
by the intersection.

Strong reachability is the least fixed point
\begin{equation}\label{eq:gc-reachability}
\begin{aligned}
 \mathsf{reach}^{0}_{\Omega}(\mathcal R_h)
   &=\mathsf{roots}_{\Omega}(\mathcal R_h),\\
 \mathsf{reach}^{i+1}_{\Omega}(\mathcal R_h)
   &=\mathsf{reach}^{i}_{\Omega}(\mathcal R_h)
     \cup\bigcup_{x\in\mathsf{reach}^{i}_{\Omega}(\mathcal R_h)}
       \bigl(\mathsf{strongSucc}_{\Omega}(x)
        \cup\mathsf{enabledMapValues}_{\Omega}
          (x,\mathsf{reach}^{i}_{\Omega}(\mathcal R_h))\bigr),\\
 \mathsf{live}_{\Omega}(\mathcal R_h)
   &=\bigcup_{i\in\mathbb N}\mathsf{reach}^{i}_{\Omega}(\mathcal R_h),\\
 \mathsf{garbage}_{\Omega}(\mathcal R_h)
   &=\dom(\mathcal S_\Omega)
       \setminus\mathsf{live}_{\Omega}(\mathcal R_h).
\end{aligned}
\end{equation}
All stores are finite, so this sequence reaches a fixed point.  Cycles not
reachable from a root are garbage together; reference counting is neither
assumed nor implied.  Weak-array elements never enter the fixed point merely
because their array is live.  A weak-map value enters only after both its map
and its key have entered; consequently a value that points back to its otherwise
unreachable key does not activate its own entry.

\subsection{Admissible collection and identity history}

A collector need not reclaim all eligible records.  It may reclaim a nonempty
set $G$, provided no retained record would contain a dangling strong reference:
\begin{equation}\label{eq:admissible-garbage-set}
\begin{split}
 \mathsf{admissibleGC}_{\Omega,\mathcal R_h}(G)
 \quad\Longleftrightarrow\quad{}
 &\varnothing\ne G\subseteq
     \mathsf{garbage}_{\Omega}(\mathcal R_h)\ \land\\
 &\forall x\in\dom(\mathcal S_\Omega)\setminus G.\;
      \mathsf{strongSucc}_{\Omega}(x)\cap G=\varnothing.
\end{split}
\end{equation}
The second conjunct matters because $\kw{allInstancesOfQuery}$ can expose a
retained but previously unreachable object.  It rules out retaining such an
object while collecting another object to which one of its ordinary slots
still points.  Taking all of $\mathsf{garbage}_{\Omega}(\mathcal R_h)$ always
satisfies this condition; a generational or incremental collector may choose a
smaller closed subset.

Weak edges are intentionally absent from the closure premise.  If a selected
record is named by a retained weak array or by an inactive weak-map entry, the
collector clears that weak observation rather than rejecting the collection.
An enabled weak-map value cannot belong to $G$ because Equation~
\eqref{eq:gc-reachability} already makes it live.

Define weak-edge clearing on a retained payload by
\begin{equation}\label{eq:gc-clear-weak-payload}
\begin{aligned}
 \mathsf{clearWeak}_{G}
  (\kw{weakArray}\langle C,\langle y_1,\ldots,y_m\rangle\rangle)
  &=\kw{weakArray}\langle C,\langle y'_1,\ldots,y'_m\rangle\rangle,\\
 y'_i&=\begin{cases}
       \nil&\mathsf{refs}(y_i)\cap G\ne\varnothing,\\
       y_i&\mathsf{refs}(y_i)\cap G=\varnothing,
      \end{cases}\\
 \mathsf{clearWeak}_{G}(\kw{weakMap}\langle C,E_w\rangle)
  &=\kw{weakMap}\langle C,E'_w\rangle,\\
 E'_w&=E_w|_{\{k\mid k\notin G\ \land\
                    \mathsf{refs}(E_w(k))\cap G=\varnothing\}},\\
 \mathsf{clearWeak}_{G}(r)&=r
   \quad\text{for every other runtime record }r.
\end{aligned}
\end{equation}
The value-side test on $E'_w$ permits partial collection to remove an inactive
entry's value while leaving its still-unreachable key for a later collection;
no retained weak record then contains a dangling identity.

For $K=\dom(\mathcal S_\Omega)\setminus G$, let
\begin{equation}\label{eq:gc-pruned-actor-heap}
 \mathsf{prunedActorHeap}(H_A,G,K)=
 \{x\mapsto\mathsf{clearWeak}_{G}(H_A(x))
      \mid x\in\dom(H_A)\cap K\}.
\end{equation}
Store pruning is the total operation
\begin{equation}\label{eq:gc-prune-store}
\begin{split}
 \mathsf{prune}(\Omega,G)=\Omega[&
   V\mapsto V|_K,\quad
   \mathcal H(A)\mapsto
      \mathsf{prunedActorHeap}(\mathcal H(A),G,K)
      \text{ for every }A\in\dom(\mathcal H),\\
  &\Pi\mapsto\Pi|_K,\quad\Phi\mapsto\Phi|_K].
\end{split}
\end{equation}
Map restriction is by key.  Every other component, including causal and event
history, is unchanged.

Deleting a record must not make its identity fresh again.  Rather than adding
an allocation ledger to every existing configuration and rule, the optional
layer wraps the VM state:
\begin{equation}\label{eq:gc-wrapper-state}
 \Psi=\langle\Omega,\mathcal I\rangle,
 \qquad
 \dom(\mathcal S_\Omega)\subseteq\mathcal I\subseteq\mathit{Ref},
\end{equation}
where $\mathcal I$ is the finite set of reference identities allocated so far
in this execution.  Let the combined existing VM transition be the disjoint
union
\begin{equation}\label{eq:gc-underlying-transition}
\begin{aligned}
 \Omega\mathrel{\Longrightarrow^{\kw{actorStep}}}\Omega'
  &\Longleftrightarrow \Omega\longrightarrow_{\mathsf{act}}\Omega',\\
 \Omega\mathrel{\Longrightarrow^{\lambda_r}}\Omega'
  &\Longleftrightarrow \Omega\xrightarrow{\lambda_r}\Omega',
 \quad
 \lambda_r::=\kw{commitReflection}(A,\Delta,n)
       \mid\kw{rejectReflection}(A,\Delta,\xi_r).
\end{aligned}
\end{equation}
For one such transition define
\begin{equation}\label{eq:gc-new-identities}
 \mathsf{newRefs}(\Omega,\Omega')=
   \dom(\mathcal S_{\Omega'})\setminus\dom(\mathcal S_{\Omega}).
\end{equation}
The ordinary transition lifts only when all newly allocated identities have
never been used:
\[
\dfrac{
 \Omega\mathrel{\Longrightarrow^{\lambda}}\Omega'
 \qquad
 \mathsf{newRefs}(\Omega,\Omega')\cap\mathcal I=\varnothing
}{
 \mathcal R_h\vdash\langle\Omega,\mathcal I\rangle
   \mathrel{\Longrightarrow^{\lambda}_{\mathsf{gc}}}
 \langle\Omega',
   \mathcal I\cup\mathsf{newRefs}(\Omega,\Omega')\rangle
}
\quad(\rulelabel{GC-Lift})
\]
This gives every existing ``fresh'' premise its intended global meaning while
leaving the underlying rules unchanged.  An initial wrapped configuration has
$\mathcal I=\dom(\mathcal S_{\Omega_0})$.

Collection is an independently enabled VM-safepoint transition:
\[
\dfrac{
 \mathsf{admissibleGC}_{\Omega,\mathcal R_h}(G)
}{
 \mathcal R_h\vdash\langle\Omega,\mathcal I\rangle
   \longrightarrow_{\mathsf{gc}}
 \langle\mathsf{prune}(\Omega,G),\mathcal I\rangle
}
\quad(\rulelabel{GC-Collect})
\]
There is deliberately no fairness requirement for
\rulelabel{GC-Collect}: an eligible object may remain uncollected indefinitely.
The implementation chooses both when to collect and which admissible closed
garbage region to reclaim.

\subsection{Observable consequences}

Let $\Omega'=\mathsf{prune}(\Omega,G)$.  Equations~
\eqref{eq:gc-reachability}--\eqref{eq:gc-prune-store} immediately give the
safety obligations
\begin{equation}\label{eq:gc-safety}
\begin{aligned}
 \mathsf{live}_{\Omega}(\mathcal R_h)
   &\subseteq\dom(\mathcal S_{\Omega'}),\\
 \forall x\in\dom(\mathcal S_{\Omega'}).\quad
 \mathsf{refs}(\mathsf{payload}_{\Omega'}(x))
   \cap\mathit{Ref}
   &\subseteq\dom(\mathcal S_{\Omega'}).
\end{aligned}
\end{equation}
The first property says that no strongly reachable record is reclaimed.  The
second says that every retained encoded reference still resolves: selected
weak targets have already been replaced by $\nil$ or removed from their map.
Together
with Equation~\eqref{eq:gc-wrapper-state}, they state memory safety and
permanent identity freshness independently of any particular tracing
algorithm.

For exact-instance enumeration, collection has the intentionally observable
effect
\begin{equation}\label{eq:gc-all-instances-effect}
 \mathsf{allInstancesOf}_{\Omega'}(C)=
 \mathsf{filter}_{x\notin G}
   \bigl(\mathsf{allInstancesOf}_{\Omega}(C)\bigr).
\end{equation}
For mixin-application enumeration, Equation~\eqref{eq:mixin-applications}
similarly gives
\begin{equation}\label{eq:gc-mixin-applications-effect}
 \mathsf{mixinApplications}_{\Omega'}(M)=
 \mathsf{filter}_{C\notin G}
   \bigl(\mathsf{mixinApplications}_{\Omega}(M)\bigr).
\end{equation}
Let $H_A=\mathcal H(A)$ and
$H'_A=\mathsf{prunedActorHeap}(H_A,G,K)$.  For a retained weak array,
Equation~\eqref{eq:gc-clear-weak-payload} gives the observable read
\begin{equation}\label{eq:gc-weak-array-effect}
\begin{aligned}
 y&=\mathsf{weakArrayRead}(H_A,a_w,i),\qquad a_w\in K,\\
 \mathsf{weakArrayRead}(H'_A,a_w,i)
  &=\begin{cases}
     \nil & \mathsf{refs}(y)\cap G\ne\varnothing,\\
     y & \text{otherwise.}
    \end{cases}
\end{aligned}
\end{equation}
For a retained weak map, an entry disappears when collection removes its key
or any reference in its value:
\begin{equation}\label{eq:gc-weak-map-effect}
\begin{aligned}
 v&=\mathsf{weakMapLookup}(H_A,m_w,k),\qquad m_w\in K,\\
 \mathsf{weakMapLookup}(H'_A,m_w,k)
  &=\begin{cases}
     \nil & k\in G\ \lor\ \mathsf{refs}(v)\cap G\ne\varnothing,\\
     v & \text{otherwise.}
    \end{cases}
\end{aligned}
\end{equation}
These changes occur atomically with \rulelabel{GC-Collect}; no program step can
observe a cleared weak cell while the corresponding target record is still
being removed.  A strongly reachable key keeps its weak-map value live through
Equation~\eqref{eq:gc-reachability}, while neither the map nor its value keeps
an otherwise unreachable key live.
Consequently an unreachable instance or class application can appear in a
snapshot before a GC step and be absent after it.  Mirror inspection itself is
atomic with respect to \rulelabel{GC-Collect}; once a returned sequence is
installed in the caller's activation or heap, its elements are ordinary strong
references and are roots of the next collection through $\mathcal X$ or a heap
record.  Nothing can be resurrected after its record has been removed.

If programs have no GC-observing primitive, removing unreachable records is
expected to be a stuttering refinement of the uncollected semantics.  In this
semantics the population queries deliberately invalidate that premise, so the
timing of collection is part of the permitted observable nondeterminism.
Finalizers, if later added to the Newspeak platform, must extend the collection
rule here rather than modify the base evaluator.
| line restores the monotonically growing heaps of
Sections~3--11.  Including it replaces that implementation abstraction with an
explicit, nondeterministic collector.  The collector introduces no Newspeak
syntax and never runs inside an actor step or reflective transaction.

The construction follows the explicit-heap rewrite method of Morrisett,
Felleisen, and Harper, \emph{Abstract Models of Memory Management} (1995), and
the treatment of observable collection in Donnelly, Hallett, and Kfoury,
\emph{Formal Semantics of Weak References} (2006).  The distinctive Newspeak
observations are Equation~\eqref{eq:all-instances-of}, whose implicit instance
registry is weak, and the platform classes \kw{WeakArray} and \kw{WeakMap}.
Collection can therefore change both later population snapshots and later
reads of weak containers.

\subsection{The collectable store}

Garbage collection applies to ordinary heap records and to the routing records
of promise and far-reference handles.  For
$\Omega=\langle P,V,\mathcal H,\mathcal X,\mathcal Q,\mathcal N,
\Pi,\Phi,\Gamma,\mathcal E,\mathcal D,w,n\rangle$, define the disjoint reference
store
\begin{equation}\label{eq:gc-reference-store}
 \mathsf{refStore}(\Omega)
   =\mathsf{vmStore}(\Omega)\uplus\Pi\uplus\Phi.
\end{equation}
The union is disjoint because its three domains are respectively subsets of
$\mathit{ObjRef}$, $\mathit{PromiseId}$, and $\mathit{FarRefId}$.  Put
$\mathcal S_\Omega=\mathsf{refStore}(\Omega)$ and expose a uniform payload
projection:
\begin{equation}\label{eq:gc-reference-payload}
 \mathsf{payload}_{\Omega}(x)=
 \begin{cases}
  \mathsf{vmStore}(\Omega)(x) & x\in\mathit{ObjRef},\\
  \Pi(x) & x\in\mathit{PromiseId},\\
  \Phi(x) & x\in\mathit{FarRefId}.
 \end{cases}
 \qquad x\in\dom(\mathcal S_\Omega).
\end{equation}
Thus a live promise strongly retains its settlement value and registered
waiters, and a live far-reference handle strongly retains its near target.
The tables $\Pi$ and $\Phi$ do not retain their own keys merely by containing
them; otherwise every handle and every remote target would live forever.

\subsection{WeakArray and WeakMap}
\label{sec:weak-containers}

\kw{WeakArray} and \kw{WeakMap} are mutable platform objects and therefore
belong to exactly one actor heap.  Their concrete VM layouts are abstracted by
two additional object-record forms:
\begin{equation}\label{eq:weak-container-records}
\begin{aligned}
 H_A(a_w)&=\kw{weakArray}\langle C,\langle y_1,\ldots,y_m\rangle\rangle,
   &a_w&\in\mathit{WeakArrayId},\\
 H_A(m_w)&=\kw{weakMap}\langle C,E_w\rangle,
   &m_w&\in\mathit{WeakMapId},
\end{aligned}
\end{equation}
where $\mathit{WeakArrayId}$ and $\mathit{WeakMapId}$ are disjoint subsets of
$\mathit{ObjRef}$, every $y_i\in\mathit{Ref}$, and
$E_w:\mathit{ObjRef}\rightharpoonup\mathit{Ref}$ is finite.  $C$ and any
ordinary fixed fields of the implementation are strong; the sequence elements
and map keys are weak.  A map value is conditionally strong exactly while its
key is strongly reachable.  This is the usual ephemeron interpretation of a
weak-key map: a value-to-key cycle cannot keep the key alive.

Every stored weak identity initially resolves; weakness controls retention, not
whether an arbitrary dangling identity may be written:
\begin{equation}\label{eq:weak-container-validity}
\begin{aligned}
 \mathsf{weakArrayStoreOK}(\Omega)
 &\Longleftrightarrow
   \forall A,a_w,C,\bar y.\
     \mathcal H(A)(a_w)=\kw{weakArray}\langle C,\bar y\rangle\\[-1mm]
 &\hspace{35mm}\Longrightarrow
     \mathsf{refs}(\bar y)\subseteq\dom(\mathcal S_\Omega),\\
 \mathsf{weakMapStoreOK}(\Omega)
 &\Longleftrightarrow
   \forall A,m_w,C,E_w.\
     \mathcal H(A)(m_w)=\kw{weakMap}\langle C,E_w\rangle\\[-1mm]
 &\hspace{35mm}\Longrightarrow
     \dom(E_w)\cup\mathsf{refs}(\mathsf{range}(E_w))
       \subseteq\dom(\mathcal S_\Omega),\\
 \mathsf{weakContainersOK}(\Omega)
 &\Longleftrightarrow
   \mathsf{weakArrayStoreOK}(\Omega)\land\mathsf{weakMapStoreOK}(\Omega).
\end{aligned}
\end{equation}

The platform method contract is given by distinct partial store functions
rather than overloaded notation:
\begin{equation}\label{eq:weak-array-operations}
\begin{aligned}
 \mathsf{weakArrayRead}(H_A,a_w,i)&=y_i,\\
 \mathsf{writeWeakArrayCell}(H_A,a_w,i,v)
   &=H_A[a_w\mapsto
       \kw{weakArray}\langle C,\bar y[i\mapsto v]\rangle],\\
 &\hspace{35mm}1\le i\le |\bar y|.
\end{aligned}
\end{equation}
\begin{equation}\label{eq:weak-map-operations}
\begin{aligned}
 \mathsf{weakMapLookup}(H_A,m_w,k)&=
   \begin{cases}E_w(k)&k\in\dom(E_w),\\ \nil&k\notin\dom(E_w),\end{cases}\\
 \mathsf{writeWeakMapEntry}(H_A,m_w,k,v)
   &=H_A[m_w\mapsto\kw{weakMap}\langle C,E_w[k\mapsto v]\rangle].
\end{aligned}
\end{equation}
Here $H_A(a_w)$ and $H_A(m_w)$ have the respective forms of
Equation~\eqref{eq:weak-container-records}.  Allocation creates an all-$\nil$
array of the requested nonnegative size or a map with $E_w=\varnothing$, using
the ordinary global-freshness obligation of Equation~\eqref{eq:gc-wrapper-state}.
The two writes are defined only when their new key and value references resolve
in $\mathcal S_\Omega$, and they preserve Equation~
\eqref{eq:weak-container-validity}.
Out-of-bounds array indices are undefined and are reified by the platform
exception boundary.  Storing $\nil$ is permitted.  Since lookup also returns
$\nil$ for an absent key, the public \kw{WeakMap} protocol deliberately does
not distinguish absence from a stored $\nil$ value.

The state outside the collectable store is
\begin{equation}\label{eq:gc-control-state}
 \mathsf{control}(\Omega)=
 \langle P,\mathcal X,\mathcal Q,\mathcal N,
          \Gamma,\mathcal E,\mathcal D,w,n\rangle.
\end{equation}
Let $\mathcal R_h\subseteq\dom(\mathcal S_\Omega)$ be the finite set of strong
references temporarily retained by the embedding host, such as arguments of a
platform primitive that is atomic with respect to collection.  A closed
Newspeak execution has $\mathcal R_h=\varnothing$.  First erase the weak
components of a payload while preserving its class and ordinary fields:
\begin{equation}\label{eq:gc-strong-payload}
\begin{aligned}
 \mathsf{strongPayload}(\kw{weakArray}\langle C,\bar y\rangle)
   &=\kw{weakArray}\langle C,\bar{\nil}\rangle,\\
 \mathsf{strongPayload}(\kw{weakMap}\langle C,E_w\rangle)
   &=\kw{weakMap}\langle C,\varnothing\rangle,\\
 \mathsf{strongPayload}(r)&=r
   \quad\text{for every other runtime record }r.
\end{aligned}
\end{equation}
The direct roots, unconditional strong successors, and map-value successors
enabled by a candidate live set $L$ are
\begin{equation}\label{eq:gc-roots-and-successors}
\begin{aligned}
 \mathsf{roots}_{\Omega}(\mathcal R_h)
   &=\bigl(\mathcal R_h\cup
       \mathsf{refs}(\mathsf{control}(\Omega))\bigr)
       \cap\dom(\mathcal S_\Omega),\\
 \mathsf{strongSucc}_{\Omega}(x)
   &=\mathsf{refs}(\mathsf{strongPayload}(
        \mathsf{payload}_{\Omega}(x)))
       \cap\dom(\mathcal S_\Omega),\\
 \mathsf{enabledMapValues}_{\Omega}(x,L)
   &=\begin{cases}
      \displaystyle
      \bigcup_{\substack{k\in\dom(E_w)\\ k\in L}}
        \mathsf{refs}(E_w(k))\cap\dom(\mathcal S_\Omega)
       & \mathsf{payload}_{\Omega}(x)=
           \kw{weakMap}\langle C,E_w\rangle,\\[3mm]
      \varnothing & \text{otherwise.}
     \end{cases}
\end{aligned}
\end{equation}
The structural extractor $\mathsf{refs}$ is Equation~
\eqref{eq:structural-reference-extractor}.  Static platform objects and atoms
are roots through $P$; running activations and their terms are roots through
$\mathcal X$; queued and in-flight arguments, reply promises, and settlement
values are roots through $\mathcal Q$ and $\mathcal N$.  Event identities and
actor identities that do not occur in $\dom(\mathcal S_\Omega)$ are discarded
by the intersection.

Strong reachability is the least fixed point
\begin{equation}\label{eq:gc-reachability}
\begin{aligned}
 \mathsf{reach}^{0}_{\Omega}(\mathcal R_h)
   &=\mathsf{roots}_{\Omega}(\mathcal R_h),\\
 \mathsf{reach}^{i+1}_{\Omega}(\mathcal R_h)
   &=\mathsf{reach}^{i}_{\Omega}(\mathcal R_h)
     \cup\bigcup_{x\in\mathsf{reach}^{i}_{\Omega}(\mathcal R_h)}
       \bigl(\mathsf{strongSucc}_{\Omega}(x)
        \cup\mathsf{enabledMapValues}_{\Omega}
          (x,\mathsf{reach}^{i}_{\Omega}(\mathcal R_h))\bigr),\\
 \mathsf{live}_{\Omega}(\mathcal R_h)
   &=\bigcup_{i\in\mathbb N}\mathsf{reach}^{i}_{\Omega}(\mathcal R_h),\\
 \mathsf{garbage}_{\Omega}(\mathcal R_h)
   &=\dom(\mathcal S_\Omega)
       \setminus\mathsf{live}_{\Omega}(\mathcal R_h).
\end{aligned}
\end{equation}
All stores are finite, so this sequence reaches a fixed point.  Cycles not
reachable from a root are garbage together; reference counting is neither
assumed nor implied.  Weak-array elements never enter the fixed point merely
because their array is live.  A weak-map value enters only after both its map
and its key have entered; consequently a value that points back to its otherwise
unreachable key does not activate its own entry.

\subsection{Admissible collection and identity history}

A collector need not reclaim all eligible records.  It may reclaim a nonempty
set $G$, provided no retained record would contain a dangling strong reference:
\begin{equation}\label{eq:admissible-garbage-set}
\begin{split}
 \mathsf{admissibleGC}_{\Omega,\mathcal R_h}(G)
 \quad\Longleftrightarrow\quad{}
 &\varnothing\ne G\subseteq
     \mathsf{garbage}_{\Omega}(\mathcal R_h)\ \land\\
 &\forall x\in\dom(\mathcal S_\Omega)\setminus G.\;
      \mathsf{strongSucc}_{\Omega}(x)\cap G=\varnothing.
\end{split}
\end{equation}
The second conjunct matters because $\kw{allInstancesOfQuery}$ can expose a
retained but previously unreachable object.  It rules out retaining such an
object while collecting another object to which one of its ordinary slots
still points.  Taking all of $\mathsf{garbage}_{\Omega}(\mathcal R_h)$ always
satisfies this condition; a generational or incremental collector may choose a
smaller closed subset.

Weak edges are intentionally absent from the closure premise.  If a selected
record is named by a retained weak array or by an inactive weak-map entry, the
collector clears that weak observation rather than rejecting the collection.
An enabled weak-map value cannot belong to $G$ because Equation~
\eqref{eq:gc-reachability} already makes it live.

Define weak-edge clearing on a retained payload by
\begin{equation}\label{eq:gc-clear-weak-payload}
\begin{aligned}
 \mathsf{clearWeak}_{G}
  (\kw{weakArray}\langle C,\langle y_1,\ldots,y_m\rangle\rangle)
  &=\kw{weakArray}\langle C,\langle y'_1,\ldots,y'_m\rangle\rangle,\\
 y'_i&=\begin{cases}
       \nil&\mathsf{refs}(y_i)\cap G\ne\varnothing,\\
       y_i&\mathsf{refs}(y_i)\cap G=\varnothing,
      \end{cases}\\
 \mathsf{clearWeak}_{G}(\kw{weakMap}\langle C,E_w\rangle)
  &=\kw{weakMap}\langle C,E'_w\rangle,\\
 E'_w&=E_w|_{\{k\mid k\notin G\ \land\
                    \mathsf{refs}(E_w(k))\cap G=\varnothing\}},\\
 \mathsf{clearWeak}_{G}(r)&=r
   \quad\text{for every other runtime record }r.
\end{aligned}
\end{equation}
The value-side test on $E'_w$ permits partial collection to remove an inactive
entry's value while leaving its still-unreachable key for a later collection;
no retained weak record then contains a dangling identity.

For $K=\dom(\mathcal S_\Omega)\setminus G$, let
\begin{equation}\label{eq:gc-pruned-actor-heap}
 \mathsf{prunedActorHeap}(H_A,G,K)=
 \{x\mapsto\mathsf{clearWeak}_{G}(H_A(x))
      \mid x\in\dom(H_A)\cap K\}.
\end{equation}
Store pruning is the total operation
\begin{equation}\label{eq:gc-prune-store}
\begin{split}
 \mathsf{prune}(\Omega,G)=\Omega[&
   V\mapsto V|_K,\quad
   \mathcal H(A)\mapsto
      \mathsf{prunedActorHeap}(\mathcal H(A),G,K)
      \text{ for every }A\in\dom(\mathcal H),\\
  &\Pi\mapsto\Pi|_K,\quad\Phi\mapsto\Phi|_K].
\end{split}
\end{equation}
Map restriction is by key.  Every other component, including causal and event
history, is unchanged.

Deleting a record must not make its identity fresh again.  Rather than adding
an allocation ledger to every existing configuration and rule, the optional
layer wraps the VM state:
\begin{equation}\label{eq:gc-wrapper-state}
 \Psi=\langle\Omega,\mathcal I\rangle,
 \qquad
 \dom(\mathcal S_\Omega)\subseteq\mathcal I\subseteq\mathit{Ref},
\end{equation}
where $\mathcal I$ is the finite set of reference identities allocated so far
in this execution.  Let the combined existing VM transition be the disjoint
union
\begin{equation}\label{eq:gc-underlying-transition}
\begin{aligned}
 \Omega\mathrel{\Longrightarrow^{\kw{actorStep}}}\Omega'
  &\Longleftrightarrow \Omega\longrightarrow_{\mathsf{act}}\Omega',\\
 \Omega\mathrel{\Longrightarrow^{\lambda_r}}\Omega'
  &\Longleftrightarrow \Omega\xrightarrow{\lambda_r}\Omega',
 \quad
 \lambda_r::=\kw{commitReflection}(A,\Delta,n)
       \mid\kw{rejectReflection}(A,\Delta,\xi_r).
\end{aligned}
\end{equation}
For one such transition define
\begin{equation}\label{eq:gc-new-identities}
 \mathsf{newRefs}(\Omega,\Omega')=
   \dom(\mathcal S_{\Omega'})\setminus\dom(\mathcal S_{\Omega}).
\end{equation}
The ordinary transition lifts only when all newly allocated identities have
never been used:
\[
\dfrac{
 \Omega\mathrel{\Longrightarrow^{\lambda}}\Omega'
 \qquad
 \mathsf{newRefs}(\Omega,\Omega')\cap\mathcal I=\varnothing
}{
 \mathcal R_h\vdash\langle\Omega,\mathcal I\rangle
   \mathrel{\Longrightarrow^{\lambda}_{\mathsf{gc}}}
 \langle\Omega',
   \mathcal I\cup\mathsf{newRefs}(\Omega,\Omega')\rangle
}
\quad(\rulelabel{GC-Lift})
\]
This gives every existing ``fresh'' premise its intended global meaning while
leaving the underlying rules unchanged.  An initial wrapped configuration has
$\mathcal I=\dom(\mathcal S_{\Omega_0})$.

Collection is an independently enabled VM-safepoint transition:
\[
\dfrac{
 \mathsf{admissibleGC}_{\Omega,\mathcal R_h}(G)
}{
 \mathcal R_h\vdash\langle\Omega,\mathcal I\rangle
   \longrightarrow_{\mathsf{gc}}
 \langle\mathsf{prune}(\Omega,G),\mathcal I\rangle
}
\quad(\rulelabel{GC-Collect})
\]
There is deliberately no fairness requirement for
\rulelabel{GC-Collect}: an eligible object may remain uncollected indefinitely.
The implementation chooses both when to collect and which admissible closed
garbage region to reclaim.

\subsection{Observable consequences}

Let $\Omega'=\mathsf{prune}(\Omega,G)$.  Equations~
\eqref{eq:gc-reachability}--\eqref{eq:gc-prune-store} immediately give the
safety obligations
\begin{equation}\label{eq:gc-safety}
\begin{aligned}
 \mathsf{live}_{\Omega}(\mathcal R_h)
   &\subseteq\dom(\mathcal S_{\Omega'}),\\
 \forall x\in\dom(\mathcal S_{\Omega'}).\quad
 \mathsf{refs}(\mathsf{payload}_{\Omega'}(x))
   \cap\mathit{Ref}
   &\subseteq\dom(\mathcal S_{\Omega'}).
\end{aligned}
\end{equation}
The first property says that no strongly reachable record is reclaimed.  The
second says that every retained encoded reference still resolves: selected
weak targets have already been replaced by $\nil$ or removed from their map.
Together
with Equation~\eqref{eq:gc-wrapper-state}, they state memory safety and
permanent identity freshness independently of any particular tracing
algorithm.

For exact-instance enumeration, collection has the intentionally observable
effect
\begin{equation}\label{eq:gc-all-instances-effect}
 \mathsf{allInstancesOf}_{\Omega'}(C)=
 \mathsf{filter}_{x\notin G}
   \bigl(\mathsf{allInstancesOf}_{\Omega}(C)\bigr).
\end{equation}
For mixin-application enumeration, Equation~\eqref{eq:mixin-applications}
similarly gives
\begin{equation}\label{eq:gc-mixin-applications-effect}
 \mathsf{mixinApplications}_{\Omega'}(M)=
 \mathsf{filter}_{C\notin G}
   \bigl(\mathsf{mixinApplications}_{\Omega}(M)\bigr).
\end{equation}
Let $H_A=\mathcal H(A)$ and
$H'_A=\mathsf{prunedActorHeap}(H_A,G,K)$.  For a retained weak array,
Equation~\eqref{eq:gc-clear-weak-payload} gives the observable read
\begin{equation}\label{eq:gc-weak-array-effect}
\begin{aligned}
 y&=\mathsf{weakArrayRead}(H_A,a_w,i),\qquad a_w\in K,\\
 \mathsf{weakArrayRead}(H'_A,a_w,i)
  &=\begin{cases}
     \nil & \mathsf{refs}(y)\cap G\ne\varnothing,\\
     y & \text{otherwise.}
    \end{cases}
\end{aligned}
\end{equation}
For a retained weak map, an entry disappears when collection removes its key
or any reference in its value:
\begin{equation}\label{eq:gc-weak-map-effect}
\begin{aligned}
 v&=\mathsf{weakMapLookup}(H_A,m_w,k),\qquad m_w\in K,\\
 \mathsf{weakMapLookup}(H'_A,m_w,k)
  &=\begin{cases}
     \nil & k\in G\ \lor\ \mathsf{refs}(v)\cap G\ne\varnothing,\\
     v & \text{otherwise.}
    \end{cases}
\end{aligned}
\end{equation}
These changes occur atomically with \rulelabel{GC-Collect}; no program step can
observe a cleared weak cell while the corresponding target record is still
being removed.  A strongly reachable key keeps its weak-map value live through
Equation~\eqref{eq:gc-reachability}, while neither the map nor its value keeps
an otherwise unreachable key live.
Consequently an unreachable instance or class application can appear in a
snapshot before a GC step and be absent after it.  Mirror inspection itself is
atomic with respect to \rulelabel{GC-Collect}; once a returned sequence is
installed in the caller's activation or heap, its elements are ordinary strong
references and are roots of the next collection through $\mathcal X$ or a heap
record.  Nothing can be resurrected after its record has been removed.

If programs have no GC-observing primitive, removing unreachable records is
expected to be a stuttering refinement of the uncollected semantics.  In this
semantics the population queries deliberately invalidate that premise, so the
timing of collection is part of the permitted observable nondeterminism.
Finalizers, if later added to the Newspeak platform, must extend the collection
rule here rather than modify the base evaluator.


\section{Source front ends and metacircular compilation}
\label{sec:front-end}

This section closes the boundary left abstract in
Equation~\eqref{eq:compilation-unit}.  It has two separate purposes.  First,
it defines the relation from source text through PEG recognition, surface AST
construction, stable source identities, and elaboration to the annotated core
used by the preceding rules.  Second, it makes the implementations of those
stages ordinary Newspeak objects.  The latter introduces no new reflective
command: changing a parser, elaborator, or compiler is an ordinary change to
its methods or object graph under Section~\ref{sec:reflection}.

\subsection{Source text and PEG recognition}
\label{sec:peg-recognition}

Let $u\in\mathit{SourceText}=\mathit{UnicodeScalar}^{*}$, let
$0\le i\le |u|$ be a source offset, let $N\in\mathit{Nonterminal}$, and let
$X\subseteq\mathit{UnicodeScalar}$.  Character ranges, optional expressions,
positive closure, and the Specification's prefix \texttt{\textasciitilde}
notation are abbreviations for the following PEG core:
\begin{equation}\label{eq:peg-expression-domain}
 r::=\epsilon\mid\kw{char}(c)\mid\kw{set}(X)\mid\kw{any}
   \mid N\mid r_1r_2\mid r_1/r_2\mid r^{*}
   \mid {!r}\mid {\&r}\mid\kw{capture}(n,r).
\end{equation}
A grammar is $G=\langle N_0,\rho_G\rangle$, where $N_0$ is its start symbol
and $\rho_G:\mathit{Nonterminal}\rightharpoonup\mathit{PEGExpr}$ is finite.
A concrete tree is
\begin{equation}\label{eq:concrete-tree-domain}
 t::=\kw{leaf}\langle i,j\rangle
   \mid\kw{node}\langle n,i,j,\bar t\rangle,
 \qquad \mathsf{span}(t)=[i,j).
\end{equation}
Uncaptured recognition returns an ordered forest.  Write
$\mathsf{ok}(j,\bar t)$ for success and $\mathsf{fail}(i)$ for failure at
offset $i$.  The deterministic function
$\mathsf{peg}_{G}(r,u,i)$ is defined by the following exhaustive equations.
The failure offset is diagnostic only; ordered choice always restarts its
right alternative at the original offset.
\begin{equation}\label{eq:peg-atomic-recognition}
\begin{aligned}
 \mathsf{peg}_{G}(\epsilon,u,i)&=\mathsf{ok}(i,\epsilon),\\
 \mathsf{peg}_{G}(\kw{char}(c),u,i)&=
  \begin{cases}\mathsf{ok}(i+1,\langle\kw{leaf}\langle i,i+1\rangle\rangle)
       &i<|u|\land u_i=c,\\
    \mathsf{fail}(i)&\text{otherwise,}\end{cases}\\
 \mathsf{peg}_{G}(\kw{set}(X),u,i)&=
  \begin{cases}\mathsf{ok}(i+1,\langle\kw{leaf}\langle i,i+1\rangle\rangle)
       &i<|u|\land u_i\in X,\\
    \mathsf{fail}(i)&\text{otherwise,}\end{cases}\\
 \mathsf{peg}_{G}(\kw{any},u,i)&=
  \begin{cases}\mathsf{ok}(i+1,\langle\kw{leaf}\langle i,i+1\rangle\rangle)
       &i<|u|,\\
    \mathsf{fail}(i)&i=|u|.\end{cases}
\end{aligned}
\end{equation}
\begin{equation}\label{eq:peg-sequence-choice}
\begin{aligned}
 \mathsf{peg}_{G}(r_1r_2,u,i)&=
  \begin{cases}
   \mathsf{ok}(k,\bar t_1\bar t_2)&
     \mathsf{peg}_{G}(r_1,u,i)=\mathsf{ok}(j,\bar t_1)\land
     \mathsf{peg}_{G}(r_2,u,j)=\mathsf{ok}(k,\bar t_2),\\
   \mathsf{fail}(h)&\mathsf{peg}_{G}(r_1,u,i)=\mathsf{fail}(h),\\
   \mathsf{fail}(h)&\mathsf{peg}_{G}(r_1,u,i)=\mathsf{ok}(j,\bar t_1)\land
                      \mathsf{peg}_{G}(r_2,u,j)=\mathsf{fail}(h),
  \end{cases}\\
 \mathsf{peg}_{G}(r_1/r_2,u,i)&=
  \begin{cases}
   q&\mathsf{peg}_{G}(r_1,u,i)=q=\mathsf{ok}(j,\bar t),\\
   \mathsf{peg}_{G}(r_2,u,i)&
      \mathsf{peg}_{G}(r_1,u,i)=\mathsf{fail}(h).
  \end{cases}
\end{aligned}
\end{equation}
Greedy repetition and lookahead are
\begin{equation}\label{eq:peg-repetition-lookahead}
\begin{aligned}
 \mathsf{peg}_{G}(r^{*},u,i)&=
  \begin{cases}
   \mathsf{ok}(i,\epsilon)&\mathsf{peg}_{G}(r,u,i)=\mathsf{fail}(h),\\
   \mathsf{ok}(k,\bar t\bar t')&
     \mathsf{peg}_{G}(r,u,i)=\mathsf{ok}(j,\bar t)\land j>i\land
     \mathsf{peg}_{G}(r^{*},u,j)=\mathsf{ok}(k,\bar t'),
  \end{cases}\\
 \mathsf{peg}_{G}(!r,u,i)&=
  \begin{cases}\mathsf{ok}(i,\epsilon)&
          \mathsf{peg}_{G}(r,u,i)=\mathsf{fail}(h),\\
        \mathsf{fail}(i)&\mathsf{peg}_{G}(r,u,i)=\mathsf{ok}(j,\bar t),
  \end{cases}\\
 \mathsf{peg}_{G}(\&r,u,i)&=
  \begin{cases}\mathsf{ok}(i,\epsilon)&
          \mathsf{peg}_{G}(r,u,i)=\mathsf{ok}(j,\bar t),\\
        \mathsf{fail}(i)&\mathsf{peg}_{G}(r,u,i)=\mathsf{fail}(h).
  \end{cases}
\end{aligned}
\end{equation}
Nonterminal invocation and capture are
\begin{equation}\label{eq:peg-nonterminal-capture}
\begin{aligned}
 \mathsf{peg}_{G}(N,u,i)&=\mathsf{peg}_{G}(\rho_G(N),u,i),\\
 \mathsf{peg}_{G}(\kw{capture}(n,r),u,i)&=
  \begin{cases}
   \mathsf{ok}(j,\langle\kw{node}\langle n,i,j,\bar t\rangle\rangle)&
      \mathsf{peg}_{G}(r,u,i)=\mathsf{ok}(j,\bar t),\\
   \mathsf{fail}(h)&\mathsf{peg}_{G}(r,u,i)=\mathsf{fail}(h).
  \end{cases}
\end{aligned}
\end{equation}
Undefined nonterminals are failures.  The grammar admissibility predicate is
\begin{equation}\label{eq:peg-grammar-validity}
\begin{split}
 \mathsf{pegOK}(G)\Longleftrightarrow{}&
 N_0\in\dom(\rho_G)\ \land\
 \mathsf{refsNT}(\mathsf{range}(\rho_G))\subseteq\dom(\rho_G)\ \land\\
 &\neg\exists N.\;N\mathrel{\mathsf{lead}_{G}^{+}}N\ \land\
 \forall r^{*}\preceq\mathsf{range}(\rho_G).\;\neg\mathsf{nullable}_{G}(r).
\end{split}
\end{equation}
Here $\mathsf{nullable}_{G}$ is the least predicate obtained by applying
Equations~\eqref{eq:peg-atomic-recognition}--\eqref{eq:peg-nonterminal-capture}
to empty input, and $N\mathrel{\mathsf{lead}_{G}}N'$ when $N'$ can be invoked
from $\rho_G(N)$ without consuming a character.  Thus
Equation~\eqref{eq:peg-grammar-validity} excludes left recursion and a nullable
operand of repetition; the recursive calls above consequently terminate.
The mechanization checks this relational predicate by a finite certificate.
Let $\nu_G:N\to\mathbb B$ be the computed nullable upper approximation and
let $r_G:N\to\mathbb N$ be the computed rank of leading nonterminal calls.
Write $\mathsf{nullable}^{+}_{\nu}(q)$ for the structural Boolean
approximation that interprets a nonterminal $N$ as $\nu(N)$, and write
$\mathsf{lead}^{+}_{\nu}(q)$ for the finite set of nonterminals reachable at
the head of $q$ under the same approximation.  The executable certificate is

\begin{equation}\label{eq:peg-admissibility-certificate}
\begin{split}
 \mathsf{pegCertOK}(G)\Longleftrightarrow{}&
 N_0\in\dom(\rho_G)\ \land\\
 &\forall N\in\dom(\rho_G).\ \rho_G(N)=b\Rightarrow\\
 &\quad \mathsf{refsNT}(b)\subseteq\dom(\rho_G)\ \land\\
 &\quad(\mathsf{nullable}^{+}_{\nu_G}(b)\Rightarrow\nu_G(N))\ \land\\
 &\quad\forall q^*\preceq b.\ \neg\mathsf{nullable}^{+}_{\nu_G}(q)\ \land\\
 &\quad\forall N'\in\mathsf{lead}^{+}_{\nu_G}(b).\
               r_G(N')<r_G(N).
\end{split}
\end{equation}
In the fourth line $q$ is the operand of that particular occurrence $q^*$.
Character-set
predicates are erased only while computing this structural certificate: they
cannot recognize empty input and therefore cannot affect any of its four
analyses.  The certificate checker and its semantic consequence are

\begin{equation}\label{eq:peg-certificate-soundness}
 \mathsf{pegCertOK}(G)=\kw{true}\Rightarrow\mathsf{pegOK}(G),
 \qquad
 \mathsf{pegCertOK}(G_{\NS})=\kw{true}.
\end{equation}
Thus the normalized grammar below is not merely assumed admissible; its finite
certificate is evaluated and the checker is proved sound for
Equation~\eqref{eq:peg-grammar-validity}.
Complete recognition is
\begin{equation}\label{eq:peg-complete-parse}
 \mathsf{parse}_{G}(u)=t
 \quad\Longleftrightarrow\quad
 \mathsf{pegOK}(G)\land
 \mathsf{peg}_{G}(N_0,u,0)=\mathsf{ok}(|u|,\langle t\rangle).
\end{equation}
The remainder of this subsection gives the distinguished grammar explicitly.
For compactness, the following are definitions, not additional PEG operators:
\begin{equation}\label{eq:peg-derived-notation}
 \begin{gathered}
 r?\equiv r/\epsilon,\qquad r^+\equiv rr^*,\qquad
 \lit{c_1\cdots c_n}\equiv
   \kw{char}(c_1)\cdots\kw{char}(c_n),\\
 \mathsf{triv}\equiv(\mathsf{white}/\mathsf{comment})^*,\qquad
 \mathsf{tok}(w)\equiv\mathsf{triv}\,\lit w,\qquad
 \mathsf{word}(w)\equiv\mathsf{triv}\,\lit w\,!\mathsf{idRest}.
 \end{gathered}
\end{equation}
$\mathsf{white}$ recognizes one Unicode \texttt{White\_Space} scalar.
The finite punctuation family is
\begin{equation}\label{eq:newspeak-punctuation-peg}
\begin{gathered}
 \rho_{\NS}(T_p)=\mathsf{tok}(p),\qquad p\in\mathcal P,\\
 \mathcal P=\{\texttt{:},\texttt{,},\texttt{.},\texttt{||},\texttt{=},
 \texttt{\char94},\texttt{[},\{,\texttt{(},\texttt{<},\texttt{<:},\\
 \texttt{\#},\texttt{>},\texttt{]},\},\texttt{)},\texttt{;},
 \texttt{/},\texttt{|},\texttt{::=},\texttt{<-:}\}.
\end{gathered}
\end{equation}
The longest member is selected before any of its prefixes.  This resolves the
otherwise implicit lexical priority of \texttt{||}, \texttt{::=},
\texttt{<:}, and \texttt{<-:}.

Let $L,U,D$ be the ASCII lower-case, upper-case, and decimal-digit sets,
$A=L\cup U$, and let $Q$ be the Specification's finite
\emph{specialCharacter} set.  The lexical part of $G_{\NS}$ is
\begin{equation}\label{eq:newspeak-lexical-peg}
\begin{aligned}
 \rho_{\NS}(\mathsf{idRest})&=\kw{set}(A\cup D)/\kw{char}(\texttt{\_}),\\
 \rho_{\NS}(\mathsf{rawId})&=(\kw{set}(A)/\kw{char}(\texttt{\_}))
   \mathsf{idRest}^*,\\
 \rho_{\NS}(\mathsf{id})&=\mathsf{triv}\,!\mathsf{reservedWord}
      \mathsf{rawId},\\
 \rho_{\NS}(\mathsf{rawKeyword})&=\mathsf{rawId}\lit{:},\\
 \rho_{\NS}(\mathsf{keyword})&=\mathsf{triv}\,\mathsf{rawKeyword},\\
 \rho_{\NS}(\mathsf{setterKeyword})&=
   \mathsf{triv}\,\mathsf{rawKeyword}\lit{:},\\
 \rho_{\NS}(\mathsf{keywords})&=\mathsf{rawKeyword}^+,\\
 \rho_{\NS}(\mathsf{rawBinary})&=(\kw{set}(Q)/\lit{-})\kw{set}(Q)^*,\\
 \rho_{\NS}(\mathsf{binarySelector})&=
   \mathsf{triv}\,\mathsf{rawBinary},\\
 \rho_{\NS}(\mathsf{metadataTag})&=\lit{:}\mathsf{rawId}\lit{:},\\
 \rho_{\NS}(\mathsf{commentItem})&=
   (!\lit{*)}\,\kw{any})/\mathsf{comment},\\
 \rho_{\NS}(\mathsf{comment})&=\lit{(*}\,\mathsf{metadataTag}?\\[-1mm]
   &\quad\mathsf{commentItem}^*\lit{*)}.
\end{aligned}
\end{equation}
Reserved words are recognized by the dedicated tokens
$\mathsf{selfToken}$, $\mathsf{superToken}$, $\mathsf{outerToken}$,
$\mathsf{trueToken}$, $\mathsf{falseToken}$, and $\mathsf{nilToken}$;
$\mathsf{id}$ additionally has negative lookahead for those complete words.
Thus the prose reservation is part of the grammar rather than a later
implementation convention.
The dedicated-word productions are
\begin{equation}\label{eq:newspeak-word-tokens}
\begin{aligned}
 \rho_{\NS}(\mathsf{reservedWord})&=
   \mathsf{word}(\texttt{self})/\mathsf{word}(\texttt{super})/
   \mathsf{word}(\texttt{outer})/\\[-1mm]
 &\quad\mathsf{word}(\texttt{true})/\mathsf{word}(\texttt{false})/
   \mathsf{word}(\texttt{nil}),\\
 \rho_{\NS}(\mathsf{selfToken})&=\mathsf{word}(\texttt{self}),\\
 \rho_{\NS}(\mathsf{superToken})&=\mathsf{word}(\texttt{super}),\\
 \rho_{\NS}(\mathsf{outerToken})&=\mathsf{word}(\texttt{outer}),\\
 \rho_{\NS}(\mathsf{trueToken})&=\mathsf{word}(\texttt{true}),\\
 \rho_{\NS}(\mathsf{falseToken})&=\mathsf{word}(\texttt{false}),\\
 \rho_{\NS}(\mathsf{nilToken})&=\mathsf{word}(\texttt{nil}),\\
 \rho_{\NS}(\mathsf{classToken})&=\mathsf{word}(\texttt{class}),\\
 \rho_{\NS}(\mathsf{lazyModifier})&=\mathsf{word}(\texttt{lazy}).
\end{aligned}
\end{equation}

Numeric, string, and symbol lexemes are defined by
\begin{equation}\label{eq:newspeak-literal-peg}
\begin{aligned}
 \mathsf{digits}&=\kw{set}(D)^+,\\
 \mathsf{extendedDigits}&=\kw{set}(D\cup U)^+,\\
 \mathsf{fraction}&=T_{\texttt{.}}\mathsf{digits},\\
 \mathsf{extendedFraction}&=T_{\texttt{.}}\mathsf{extendedDigits},\\
 \mathsf{exponent}&=\lit e\lit{-}?\mathsf{digits},\\
 \mathsf{decimalNum}&=\lit{-}?\mathsf{digits}
       \mathsf{fraction}?\mathsf{exponent}?,\\
 \mathsf{radixNum}&=\mathsf{digits}\lit r\lit{-}?
       \mathsf{extendedDigits}\mathsf{extendedFraction}?
       \mathsf{exponent}?,\\
 \rho_{\NS}(\mathsf{number})&=\mathsf{triv}
       (\mathsf{radixNum}/\mathsf{decimalNum}),\\
 \rho_{\NS}(\mathsf{singleString})&=
   \lit{' }((\lit{''})/(!\lit{' }\kw{any}))^*\lit{' },\\
 \rho_{\NS}(\mathsf{doubleString})&=
   \lit{"}((\lit{""})/(!\lit{"}\kw{any}))^*\lit{"},\\
 \rho_{\NS}(\mathsf{string})&=\mathsf{triv}
   (\mathsf{singleString}/\mathsf{doubleString}),\\
 \rho_{\NS}(\mathsf{symbol})&=\mathsf{triv}\,T_{\texttt{\#}}\\[-1mm]
  &\quad(\mathsf{string}/\mathsf{keywords}/
          \mathsf{rawBinary}/\mathsf{rawId}).
\end{aligned}
\end{equation}
The action in Equation~\eqref{eq:newspeak-literal-actions} below performs radix
and escape decoding and NFC normalization; recognition itself loses no source
spelling.

Optional type annotations affect parsing, tooling, and reflection but not the
untyped dynamic kernel.  Let
\begin{equation}\label{eq:newspeak-type-word-set}
\begin{split}
 \mathcal W_T=\{&\texttt{arg},\texttt{for},\texttt{generic},
 \texttt{inheritedTypeOf},\texttt{is},\texttt{message},\texttt{of},\\[-1mm]
 &\texttt{receiverType},\texttt{subtypeOf},\texttt{typeArg},\texttt{where}\}.
\end{split}
\end{equation}
Their complete PEG is
\begin{equation}\label{eq:newspeak-type-peg}
\begin{aligned}
 \rho_{\NS}(T_w)&=\mathsf{word}(w)\quad(w\in\mathcal W_T),\\
 \rho_{\NS}(\mathsf{returnType})&=T_{\texttt{\char94}}\mathsf{type},\\
 \rho_{\NS}(\mathsf{type})&=T_{\texttt{<}}\mathsf{typeExpr}T_{\texttt{>}},\\
 \rho_{\NS}(\mathsf{typePrimary})&=\mathsf{id}\mathsf{typeArguments}?,\\
 \rho_{\NS}(\mathsf{typeFactor})&=\mathsf{typePrimary}/\mathsf{blockType}/\\[-1mm]
 &\quad\mathsf{tupleType}/\mathsf{parenthesizedTypeExpression},\\
 \rho_{\NS}(\mathsf{parenthesizedTypeExpression})&=
      T_{\texttt{(}}\mathsf{typeExpr}T_{\texttt{)}},\\
 \rho_{\NS}(\mathsf{typeTerm})&=\mathsf{typeFactor}\mathsf{id}^*,\\
 \rho_{\NS}(\mathsf{typeExpr})&=\mathsf{typeTerm}\\[-1mm]
 &\quad((T_{\texttt{|}}/T_{\texttt{;}}/T_{\texttt{/}})
      \mathsf{typeExpr})?,\\
 \rho_{\NS}(\mathsf{typeArguments})&=T_{\texttt{[}}\mathsf{typeExpr}
      (T_{\texttt{,}}\mathsf{typeExpr})^*T_{\texttt{]}},\\
 \rho_{\NS}(\mathsf{tupleType})&=T_{\{}(
      \mathsf{typeExpr}(T_{\texttt{.}}\mathsf{typeExpr})^*)?T_{\}},\\
 \rho_{\NS}(\mathsf{blockArgType})&=T_{\texttt{:}}\mathsf{typeTerm},\\
 \rho_{\NS}(\mathsf{blockReturnType})&=\mathsf{typeExpr},\\
 \rho_{\NS}(\mathsf{nonEmptyBlockArgList})&=
      \mathsf{blockArgType}^+(T_{\texttt{|}}\mathsf{blockReturnType})?,\\
 \rho_{\NS}(\mathsf{blockType})&=T_{\texttt{[}}
      (\mathsf{nonEmptyBlockArgList}/\mathsf{blockReturnType}?)T_{\texttt{]}}.
\end{aligned}
\end{equation}
Type parameters and inference patterns are
\begin{equation}\label{eq:newspeak-type-pattern-peg}
\begin{aligned}
 \rho_{\NS}(\mathsf{typeParameterDecl})&=T_{\texttt{<}}T_{\texttt{[}}
      \mathsf{id}(T_{\texttt{,}}\mathsf{id})^*T_{\texttt{]}}T_{\texttt{>}},\\
 \rho_{\NS}(\mathsf{typePattern})&=T_{\texttt{<}}\mathsf{typeFormal}
      (T_{\texttt{;}}\mathsf{typeFormal})^*T_{\texttt{>}},\\
 \rho_{\NS}(\mathsf{typeFormal})&=T_{\texttt{where}}\mathsf{id}
      \mathsf{typeParamConstraint}?T_{\texttt{is}}\mathsf{inferenceClause},\\
 \rho_{\NS}(\mathsf{typeParamConstraint})&=T_{\texttt{<}}
      \mathsf{typeBoundQualifier}?\mathsf{typeExpr}T_{\texttt{>}},\\
 \rho_{\NS}(\mathsf{typeBoundQualifier})&=
      T_{\texttt{subtypeOf}}/T_{\texttt{inheritedTypeOf}},\\
 \rho_{\NS}(\mathsf{inferenceClause})&=T_{\texttt{receiverType}}/
      (\mathsf{returnType}\mathsf{returnTypeInferenceClause})/
      \mathsf{typeArgInferenceClause}/\\[-1mm]
 &\quad(T_{\texttt{arg}}\mathsf{number}
      (T_{\texttt{of}}\mathsf{msgSelector})?),\\
 \rho_{\NS}(\mathsf{returnTypeInferenceClause})&=
      T_{\texttt{of}}\mathsf{msgSelector},\\
 \rho_{\NS}(\mathsf{msgSelector})&=
      \mathsf{symbol}T_{\texttt{message}}T_{\texttt{of}}
      \mathsf{inferenceClause},\\
 \rho_{\NS}(\mathsf{typeArgInferenceClause})&=T_{\texttt{typeArg}}
      \mathsf{number}T_{\texttt{for}}T_{\texttt{generic}}\\[-1mm]
 &\quad\mathsf{symbol}T_{\texttt{of}}\mathsf{inferenceClause}.
\end{aligned}
\end{equation}

Let $\mathcal T_I$ contain every nonterminal on the left of
Equation~\eqref{eq:newspeak-type-peg} except $\mathsf{type}$, together with
every left-hand nonterminal of Equation~\eqref{eq:newspeak-type-pattern-peg}
except $\mathsf{typeParameterDecl}$ and $\mathsf{typePattern}$.  Type syntax is
retained as source metadata by
\begin{equation}\label{eq:newspeak-type-actions}
\begin{aligned}
 \mathsf{staticTypeToken}_{z}(n,w)&=
   \langle\kw{staticType},\langle n,w\rangle,z\rangle,\\
 \mathcal J_{\NS}(n,\bar v,w)&=\kw{erase}\\[-1mm]
 &\quad(n\in\mathcal T_I\cup\{T_w\mid w\in\mathcal W_T\}),\\
 \mathcal J_{\NS}(n,\bar v,w)&=
   \kw{metadata}(\mathsf{staticTypeToken}_{z}(n,w))
   \quad(n\in\{\mathsf{type},\mathsf{typeParameterDecl},
                 \mathsf{typePattern}\}).
\end{aligned}
\end{equation}

The message and precedence productions are the following PEG.  They include
the asynchronous-send token required by Specification \S5.3.7, which the
printed 0.109 grammar describes in prose but omits from its displayed
productions.
\begin{equation}\label{eq:newspeak-message-peg}
\begin{aligned}
 \rho_{\NS}(\mathsf{unarySelector})&=\mathsf{id},\\
 \rho_{\NS}(\mathsf{binaryMessage})&=
      \mathsf{binarySelector}\,\mathsf{unaryExpression},\\
 \rho_{\NS}(\mathsf{keywordMessage})&=
      (\mathsf{keyword}\,\mathsf{binaryExpression})^+,\\
 \rho_{\NS}(\mathsf{message})&=
      \mathsf{keywordMessage}/\mathsf{binaryMessage}/\mathsf{unarySelector},\\
 \rho_{\NS}(\mathsf{receiverPrimary})&=
      \mathsf{outerReceiver}/\mathsf{superToken}/\mathsf{primary},\\
 \rho_{\NS}(\mathsf{unaryExpression})&=
      \mathsf{receiverPrimary}\,\mathsf{unarySelector}^*,\\
 \rho_{\NS}(\mathsf{binaryExpression})&=
      \mathsf{unaryExpression}\,\mathsf{binaryMessage}^*,\\
 \rho_{\NS}(\mathsf{keywordExpression})&=
      \mathsf{binaryExpression}\,\mathsf{keywordMessage}?,\\
 \rho_{\NS}(\mathsf{eventualExpression})&=
      \mathsf{keywordExpression}\,T_{\texttt{<-:}}\,\mathsf{message},\\
 \rho_{\NS}(\mathsf{cascadeMessage})&=
      T_{\texttt{;}}(\mathsf{keywordMessage}/\mathsf{binaryMessage}/
                       \mathsf{unarySelector}),\\
 \rho_{\NS}(\mathsf{implicitKeywordSend})&=\mathsf{keywordMessage},\\
 \rho_{\NS}(\mathsf{nontrivialUnaryMessages})&=
      \mathsf{unarySelector}^+\mathsf{binaryMessage}^*
      \mathsf{keywordMessage}?,\\
 \rho_{\NS}(\mathsf{nontrivialBinaryMessages})&=
      \mathsf{binaryMessage}^+\mathsf{keywordMessage}?,\\
 \rho_{\NS}(\mathsf{keywordMessages})&=\mathsf{keywordMessage},\\
 \rho_{\NS}(\mathsf{nonEmptyMessages})&=
      \mathsf{nontrivialUnaryMessages}/\\[-1mm]
 &\quad\mathsf{nontrivialBinaryMessages}/\mathsf{keywordMessages},\\
 \rho_{\NS}(\mathsf{messageCascade})&=
      \mathsf{nonEmptyMessages}\,\mathsf{cascadeMessage}^*,\\
 \rho_{\NS}(\mathsf{cascadedExpression})&=
      \mathsf{primary}\,\mathsf{messageCascade}?,\\
 \rho_{\NS}(\mathsf{sendExpression})&=
      \mathsf{eventualExpression}/\mathsf{implicitKeywordSend}/
      \mathsf{cascadedExpression},\\
 \rho_{\NS}(\mathsf{expression})&=
      \mathsf{setterKeyword}?\,\mathsf{sendExpression}.
\end{aligned}
\end{equation}
The first selector of a receiver-less unary or keyword expression is retained
as an implicit send by the actions below.  Binary sends always have an
explicit left operand.  A bare \kw{outer} or \kw{super} is rejected by
Equation~\eqref{eq:surface-context-validity}; permitting it temporarily in
$\mathsf{receiverPrimary}$ keeps the precedence grammar non-left-recursive.

The literal, block, pattern, statement, and primary productions are
\begin{equation}\label{eq:newspeak-expression-peg}
\begin{aligned}
 \rho_{\NS}(\mathsf{literal})&=\mathsf{pattern}/\mathsf{number}/
   \mathsf{symbol}/\mathsf{string}/\\[-1mm]
 &\quad\mathsf{tuple}/\mathsf{trueToken}/\mathsf{falseToken}/
   \mathsf{nilToken},\\
 \rho_{\NS}(\mathsf{tuple})&=T_{\{}
   (\mathsf{expression}(T_{\texttt{.}}\mathsf{expression})^*
      T_{\texttt{.}}?)?T_{\}},\\
 \rho_{\NS}(\mathsf{blockParameter})&=T_{\texttt{:}}\mathsf{slotDecl},\\
 \rho_{\NS}(\mathsf{blockParameters})&=
      \mathsf{blockParameter}^+T_{\texttt{|}},\\
 \rho_{\NS}(\mathsf{block})&=T_{\texttt{[}}
      \mathsf{blockParameters}?\mathsf{codeBody}T_{\texttt{]}},\\
 \rho_{\NS}(\mathsf{pattern})&=T_{\texttt{<}}\mathsf{patternLiteral}T_{\texttt{>}},\\
 \rho_{\NS}(\mathsf{patternLiteral})&=\mathsf{wildcardPattern}/
      \mathsf{literalPattern}/\mathsf{keywordPattern},\\
 \rho_{\NS}(\mathsf{wildcardPattern})&=\mathsf{tok}(\texttt{\_}),\\
 \rho_{\NS}(\mathsf{literalPattern})&=\mathsf{number}/\mathsf{symbol}/
      \mathsf{string}/\mathsf{tuple},\\
 \rho_{\NS}(\mathsf{keywordPattern})&=\mathsf{keywordPatternPair}^+,\\
 \rho_{\NS}(\mathsf{keywordPatternPair})&=
      \mathsf{keyword}\,\mathsf{keywordPatternValue}?,\\
 \rho_{\NS}(\mathsf{keywordPatternValue})&=\mathsf{wildcardPattern}/
      \mathsf{literalPattern}/\\[-1mm]
 &\quad\mathsf{variablePattern}/\mathsf{nestedPattern},\\
 \rho_{\NS}(\mathsf{variablePattern})&=
      \mathsf{triv}\lit{?}\mathsf{rawId},\\
 \rho_{\NS}(\mathsf{nestedPattern})&=\mathsf{triv}\mathsf{pattern},\\
 \rho_{\NS}(\mathsf{parenthesized})&=
      T_{\texttt{(}}\mathsf{expression}T_{\texttt{)}},\\
 \rho_{\NS}(\mathsf{primary})&=\mathsf{literal}/\mathsf{block}/
      \mathsf{classDeclaration}/\\[-1mm]
 &\quad\mathsf{objectLiteral}/\mathsf{parenthesized}/\\[-1mm]
 &\quad\mathsf{selfToken}/\mathsf{unarySelector},\\
 \rho_{\NS}(\mathsf{returnStatement})&=T_{\texttt{\char94}}
      \mathsf{expression}T_{\texttt{.}}?,\\
 \rho_{\NS}(\mathsf{statements})&=\mathsf{returnStatement}/\\[-1mm]
 &\quad(\mathsf{expression}(T_{\texttt{.}}\mathsf{statements})?)/\epsilon,\\
 \rho_{\NS}(\mathsf{codeBody})&=\mathsf{slotDecls}?\mathsf{statements}.
\end{aligned}
\end{equation}

Finally, class, member, and compilation-unit syntax is
\begin{equation}\label{eq:newspeak-declaration-peg}
\begin{aligned}
 \rho_{\NS}(\mathsf{messagePattern})&=(\mathsf{keywordPatternDecl}/
      \mathsf{binaryPatternDecl}/\mathsf{unarySelector})\mathsf{returnType}?,\\
 \rho_{\NS}(\mathsf{keywordPatternDecl})&=
      (\mathsf{keyword}\mathsf{slotDecl})^+,\\
 \rho_{\NS}(\mathsf{binaryPatternDecl})&=
      \mathsf{binarySelector}\mathsf{slotDecl},\\
 \rho_{\NS}(\mathsf{accessModifier})&=
      \mathsf{tok}(\texttt{private})/\mathsf{tok}(\texttt{public})/\\[-1mm]
 &\quad\mathsf{tok}(\texttt{protected}),\\
 \rho_{\NS}(\mathsf{methodDecl})&=\mathsf{accessModifier}?
      \mathsf{messagePattern}T_{\texttt{=}}T_{\texttt{(}}
      \mathsf{codeBody}T_{\texttt{)}},\\
 \rho_{\NS}(\mathsf{slotDecl})&=\mathsf{id}\mathsf{type}?,\\
 \rho_{\NS}(\mathsf{slotDef})&=\mathsf{accessModifier}?\mathsf{slotDecl}
   (((T_{\texttt{=}}/T_{\texttt{::=}})\mathsf{expression}T_{\texttt{.}})/
      T_{\texttt{.}})?,\\
 \rho_{\NS}(\mathsf{sequentialSlots})&=T_{\texttt{|}}
      \mathsf{slotDef}^*T_{\texttt{|}},\\
 \rho_{\NS}(\mathsf{simultaneousSlots})&=T_{\texttt{||}}
      \mathsf{slotDef}^*T_{\texttt{||}},\\
 \rho_{\NS}(\mathsf{slotDecls})&=
      \mathsf{sequentialSlots}/\mathsf{simultaneousSlots},\\
 \rho_{\NS}(\mathsf{lazySlotDecl})&=
   ((\mathsf{lazyModifier}\mathsf{accessModifier}?)/\\[-1mm]
 &\quad(\mathsf{accessModifier}\mathsf{lazyModifier}))
   \mathsf{slotDecl}(T_{\texttt{=}}/T_{\texttt{::=}})\\[-1mm]
 &\quad\mathsf{expression}T_{\texttt{.}},\\
 \rho_{\NS}(\mathsf{classHeader})&=T_{\texttt{(}}\mathsf{comment}?
      \mathsf{slotDecls}?\mathsf{statements}T_{\texttt{)}},\\
 \rho_{\NS}(\mathsf{nestedClassDecl})&=
      \mathsf{accessModifier}?\mathsf{classDeclaration},\\
 \rho_{\NS}(\mathsf{instanceSide})&=T_{\texttt{(}}\\[-1mm]
 &\quad(\mathsf{nestedClassDecl}/\mathsf{lazySlotDecl}/
      \mathsf{methodDecl})^*T_{\texttt{)}},\\
 \rho_{\NS}(\mathsf{classSide})&=T_{\texttt{:}}T_{\texttt{(}}
      \mathsf{methodDecl}^*T_{\texttt{)}},\\
 \rho_{\NS}(\mathsf{classBody})&=\mathsf{classHeader}
      \mathsf{instanceSide}\mathsf{classSide}?,\\
 \rho_{\NS}(\mathsf{outerReceiver})&=\mathsf{outerToken}\mathsf{id},\\
 \rho_{\NS}(\mathsf{inheritancePrefix})&=\mathsf{outerReceiver}/
      \mathsf{selfToken}/\mathsf{superToken},\\
 \rho_{\NS}(\mathsf{inheritanceClause})&=\mathsf{inheritancePrefix}?
      \mathsf{unarySelector}\mathsf{message}?,\\
 \rho_{\NS}(\mathsf{mixinSuffix})&=
      (T_{\texttt{<:}}\mathsf{inheritanceClause})^+
      (T_{\texttt{.}}/\mathsf{classBody}),\\
 \rho_{\NS}(\mathsf{inheritanceAndBody})&=\mathsf{classBody}/\\[-1mm]
 &\quad(\mathsf{inheritanceClause}
        (\mathsf{classBody}/\mathsf{mixinSuffix})),\\
 \rho_{\NS}(\mathsf{classDeclaration})&=\mathsf{classToken}\mathsf{id}
      \mathsf{typeParameterDecl}?\\[-1mm]
 &\quad\mathsf{messagePattern}?T_{\texttt{=}}
      \mathsf{inheritanceAndBody},\\
 \rho_{\NS}(\mathsf{objectLiteral})&=
      (\mathsf{id}\mathsf{keywordMessage}?)?\mathsf{classBody},\\
 \rho_{\NS}(\mathsf{compilationUnit})&=\mathsf{id}\mathsf{string}?
      \mathsf{classDeclaration}\,\mathsf{triv}.
\end{aligned}
\end{equation}
$G_{\NS}=\langle\mathsf{compilationUnit},\rho_{\NS}\rangle$ with every
syntactic right-hand side in
Equations~\eqref{eq:newspeak-message-peg}--\eqref{eq:newspeak-declaration-peg}
wrapped in a same-named capture.  Equations~\eqref{eq:peg-derived-notation}--
\eqref{eq:newspeak-declaration-peg} are the normalized 0.109 grammar.  They
make explicit four facts left to prose or typography in the printed grammar:
Unicode whitespace, reserved-word exclusion, asynchronous sends, and the use
of class declarations, object literals, and \kw{self} as primaries.  No
unimplemented chain production is included.

\subsection{Concrete trees, metadata, and stable AST identities}
\label{sec:ast-construction}

Let $z=[i,j)$ be the span of a captured node.  Message clauses and patterns are
not expressions; their complete surface domains are
\begin{equation}\label{eq:surface-message-pattern-domain}
\begin{aligned}
 c&::=\kw{clause}_{z}(s,\bar R),\\
 p&::=\kw{wild}_{z}\mid\kw{literalPat}_{z}(R)
   \mid\kw{variablePat}_{z}(x)\mid\kw{nestedPat}_{z}(p)
   \mid\kw{keywordPat}_{z}(\langle k_i,p_i?\rangle_{i=1}^{n}).
\end{aligned}
\end{equation}
The complete executable surface AST is
\begin{equation}\label{eq:surface-expression-domain}
\begin{aligned}
 R::={}&\kw{atom}_{z}(\omega)\mid\kw{identifier}_{z}(x)
  \mid\kw{self}_{z}\mid\kw{outer}_{z}(x)\mid\kw{super}_{z}
  \mid\kw{paren}_{z}(R)\\
 &\mid\kw{implicit}_{z}(s,\bar R)
  \mid\kw{message}_{z}(R,s,\bar R)
  \mid\kw{eventual}_{z}(R,s,\bar R)
  \mid\kw{cascade}_{z}(R,\bar c_0,\bar c)\\
 &\mid\kw{setter}_{z}(x,R)\mid\kw{tuple}_{z}(\bar R)
  \mid\kw{pattern}_{z}(p)\mid\kw{closure}_{z}(\bar d,\bar g,\bar s)\\
 &\mid\kw{object}_{z}(o_h)\mid\kw{classExpr}_{z}(d)
  \mid\kw{sequence}_{z}(\bar s),\\
 s_R::={}&\kw{exprStmt}_{z}(R)\mid\kw{returnStmt}_{z}(R).
\end{aligned}
\end{equation}
Declarations and class structure are
\begin{equation}\label{eq:surface-declaration-domain}
\begin{aligned}
 d::={}&\kw{formal}_{z}(x)\mid
  \kw{slot}_{z}(\alpha?,x,m,e?)\mid
  \kw{lazySlot}_{z}(\alpha?,x,m,e)\\
 &\mid\kw{method}_{z}(\alpha?,s,\bar d,\bar g,\bar s)
  \mid\kw{nestedClass}_{z}(\alpha?,d)
  \mid\kw{classDecl}_{z}(x,s_f,\bar d,b),\\
 g::={}&\kw{sequentialGroup}_{z}(\bar d)\mid
         \kw{simultaneousGroup}_{z}(\bar d),\\
 h::={}&\kw{classStructure}_{z}(g?,\bar s,\bar d,\bar d_c),\\
 o_h::={}&\kw{implicitObjectHeader}_{z}(h)\mid
       \kw{explicitObjectHeader}_{z}(R,c,h),\\
 b::={}&\kw{defaultInheritance}_{z}(h)\mid
  \kw{explicitInheritance}_{z}(R,c,h)\mid
  \kw{mixinChain}_{z}(R,c,\langle c_i\rangle_{i=1}^{n},h?),\\
 U::={}&\kw{unit}_{z}(\lambda,\gamma?,d).
\end{aligned}
\end{equation}
Here $m\in\{\kw{immutable},\kw{mutable}\}$ records \texttt{=} versus
\texttt{::=}; $\gamma$ is the optional category string; and absent access
remains absent until elaboration supplies the protected default.  Every
declaration and expression constructor carries its captured source span and,
after Equation~\eqref{eq:stable-ast-identification}, its stable identity.

The grammar carries a finite semantic-action map
\begin{equation}\label{eq:peg-action-map}
\begin{aligned}
 \mathcal J_G:
   \mathit{CaptureName}\times\mathit{ActionValue}^{*}\times
   \mathit{SourceSlice}&\rightharpoonup\mathit{ActionResult},\\
 a&::=\kw{emit}(v)\mid\kw{erase}\mid\kw{metadata}(m).
\end{aligned}
\end{equation}
Projection returns an ordered sequence because punctuation contributes no
semantic child.  Metadata likewise leaves the semantic-child sequence but is
accumulated separately.  Put
\begin{equation}\label{eq:ast-tree-projection}
\begin{aligned}
 \mathsf{project}_{G,u}(\kw{leaf}\langle i,j\rangle)
   &=\langle\langle\kw{token}(u[i:j],[i,j))\rangle,\epsilon\rangle,\\
 \bar v&=\bar v_1\cdots\bar v_r,
 &\bar m&=\bar m_1\cdots\bar m_r,\\[-1mm]
 \mathsf{project}_{G,u}(t_k)&=\langle\bar v_k,\bar m_k\rangle
   \quad(1\le k\le r),\\
 \mathsf{project}_{G,u}(\kw{node}\langle n,i,j,\bar t\rangle)
   &=\mathsf{emitResult}
      (\mathcal J_G(n,\bar v,u[i:j]),\bar m),\\
 \mathsf{emitResult}(\kw{emit}(v),\bar m)&=\langle\langle v\rangle,\bar m\rangle,\\
 \mathsf{emitResult}(\kw{erase},\bar m)&=\langle\epsilon,\bar m\rangle,\\
 \mathsf{emitResult}(\kw{metadata}(m),\bar m)&=
      \langle\epsilon,\bar m\cdot\langle m\rangle\rangle,\\
 \mathsf{ast}_{G,u}(t)&=\langle R,\bar m\rangle
   \quad\Longleftrightarrow\quad
   \mathsf{project}_{G,u}(t)=\langle\langle R\rangle,\bar m\rangle.
\end{aligned}
\end{equation}
The metadata sequence $\bar m$ is attached by
Equation~\eqref{eq:metadata-association}; it is not an executable AST child.
Below, an equation $\mathcal J_G(\cdots)=v$ abbreviates
$\mathcal J_G(\cdots)=\kw{emit}(v)$.
To define $\mathcal J_{G_{\NS}}$ without a production-sized prose gap, first
define the message folds.  If $c_i=\kw{clause}(s_i,\bar R_i)$, then
\begin{equation}\label{eq:surface-message-folds}
\begin{aligned}
 \mathsf{asClause}_{z}(\kw{selector}(s))&=\kw{clause}_{z}(s,\epsilon),\\
 \mathsf{asClause}_{z}(c)&=c
   \quad(c=\kw{clause}_{z'}(s,\bar R)),\\
 \mathsf{flattenClauses}_{z}(\epsilon)&=\epsilon,\\
 \mathsf{flattenClauses}_{z}(v\cdot\bar v)&=
   \mathsf{asClause}_{z}(v)\cdot\mathsf{flattenClauses}_{z}(\bar v)
   \quad(v\notin\mathit{Sequence}),\\
 \mathsf{flattenClauses}_{z}(\bar c\cdot\bar v)&=
   \bar c\cdot\mathsf{flattenClauses}_{z}(\bar v)
   \quad(\bar c\in\mathit{Clause}^{*}),\\
 \mathsf{foldExplicit}_{z}(R,\epsilon)&=R,\\
 \mathsf{foldExplicit}_{z}(R,v_1\cdot\bar v)&=
   \mathsf{foldExplicit}_{z}
     (\kw{message}_{z}(R,s_1,\bar R_1),\bar v),\\[-1mm]
 &\qquad\mathsf{asClause}_{z}(v_1)=\kw{clause}(s_1,\bar R_1),\\
 \mathsf{foldImplicit}_{z}(v_1\cdot\bar v)&=
   \mathsf{foldExplicit}_{z}
     (\kw{implicit}_{z}(s_1,\bar R_1),\bar v),\\[-1mm]
 &\qquad\mathsf{asClause}_{z}(v_1)=\kw{clause}(s_1,\bar R_1),\\
 \mathsf{makeCascade}_{z}(R,\bar c_0,\epsilon)&=
   \mathsf{foldExplicit}_{z}(R,\bar c_0),\\
 \mathsf{makeCascade}_{z}(R,\bar c_0,c_1\cdot\bar c)&=
   \kw{cascade}_{z}(R,\bar c_0,c_1\cdot\bar c).
\end{aligned}
\end{equation}
The first two equations preserve unary-before-binary-before-keyword order
already imposed by Equation~\eqref{eq:newspeak-message-peg}; the third retains
one receiver and every cascade clause.

Lexical actions are
\begin{equation}\label{eq:newspeak-literal-actions}
\begin{aligned}
 \mathcal J_{\NS}(\mathsf{id},\bar v,w)&=
   \kw{name}(\mathsf{decodeId}(w)),\\
 \mathcal J_{\NS}(\mathsf{unarySelector},\langle\kw{name}(x)\rangle,w)
   &=\kw{selector}_{z}(\#x),\\
 \mathcal J_{\NS}(\mathsf{binarySelector},\bar v,w)&=
   \kw{selector}_{z}(\mathsf{decodeBinary}(w)),\\
 \mathcal J_{\NS}(\mathsf{keyword},\bar v,w)&=
   \kw{keyword}_{z}(\mathsf{decodeKeyword}(w)),\\
 \mathcal J_{\NS}(\mathsf{setterKeyword},\bar v,w)&=
   \kw{name}(\mathsf{decodeSetterName}(w)),\\
 \mathcal J_{\NS}(\mathsf{selfToken},\bar v,w)&=\kw{self}_{z},\\
 \mathcal J_{\NS}(\mathsf{superToken},\bar v,w)&=\kw{super}_{z},\\
 \mathcal J_{\NS}(\mathsf{outerReceiver},
      \langle\kw{name}(x)\rangle,w)&=\kw{outer}_{z}(x),\\
 \mathcal J_{\NS}(\mathsf{comment},\bar v,w)&=
   \kw{metadata}(\mathsf{metadataToken}_{z}(\mathsf{tag}(\bar v),
      \mathsf{payload}(w))),\\
 \mathcal J_{\NS}(\mathsf{number},\bar v,w)&=
   \kw{atom}_{z}(\mathsf{decodeNumber}(w)),\\
 \mathcal J_{\NS}(\mathsf{string},\bar v,w)&=
   \kw{atom}_{z}(\mathsf{NFC}(\mathsf{decodeString}(w))),\\
 \mathcal J_{\NS}(\mathsf{symbol},\bar v,w)&=
   \kw{atom}_{z}(\mathsf{decodeSymbol}(w)),\\
 \mathcal J_{\NS}(\mathsf{trueToken},\bar v,w)&=\kw{atom}_{z}(\kw{true}),\\
 \mathcal J_{\NS}(\mathsf{falseToken},\bar v,w)&=\kw{atom}_{z}(\kw{false}),\\
 \mathcal J_{\NS}(\mathsf{nilToken},\bar v,w)&=\kw{atom}_{z}(\kw{nil}).
\end{aligned}
\end{equation}
$\mathsf{decodeNumber}$ is the positional interpretation in Specification
\S5.1.1; it rejects a digit outside the stated radix.

Punctuation and lexical scaffolding are erased only after their enclosing
capture has used its source slice.  Pure wrapper productions pass their unique
semantic child through:
\begin{equation}\label{eq:newspeak-structural-actions}
\begin{aligned}
 \mathcal E_{\NS}={}&\{T_p\mid p\in\mathcal P\}\cup\{\mathsf{white},
   \mathsf{reservedWord},\mathsf{idRest},\mathsf{rawId},\\[-1mm]
 &\mathsf{rawKeyword},\mathsf{keywords},\mathsf{rawBinary},
   \mathsf{metadataTag},\mathsf{commentItem},\\[-1mm]
 &\mathsf{classToken},\mathsf{lazyModifier},\mathsf{outerToken}\},\\
 \mathcal I_{\NS}={}&\{\mathsf{literal},\mathsf{receiverPrimary},
   \mathsf{sendExpression},\\[-1mm]
 &\mathsf{slotDecls},\mathsf{blockParameter},\mathsf{inheritancePrefix},
   \mathsf{patternLiteral},\\[-1mm]
 &\mathsf{keywordPatternValue}\},\\
 \mathcal J_{\NS}(n,\bar v,w)&=\kw{erase}
      \quad(n\in\mathcal E_{\NS}),\\
 \mathcal J_{\NS}(n,\langle v\rangle,w)&=v
      \quad(n\in\mathcal I_{\NS}).
\end{aligned}
\end{equation}

Sequence normalization used by the expression actions is
\begin{equation}\label{eq:newspeak-action-sequence-helpers}
\begin{aligned}
 \mathsf{orEmpty}(\kw{none})&=\epsilon,\\
 \mathsf{orEmpty}(v)&=\langle v\rangle\quad(v\ne\kw{none}),\\
 \mathsf{orEmptySeq}(\kw{none})&=\epsilon,\\
 \mathsf{orEmptySeq}(\bar v)&=\bar v\quad(\bar v\ne\kw{none}),\\
 \mathsf{asStatement}(R)&=\kw{exprStmt}(R)
      \quad(R\in\mathit{SurfaceExpression}),\\
 \mathsf{asStatement}(s_R)&=s_R
      \quad(s_R\in\mathit{SurfaceStatement}),\\
 \mathsf{flattenStatements}(\epsilon)&=\epsilon,\\
 \mathsf{flattenStatements}(v\cdot\bar v)&=
      \mathsf{asStatement}(v)\cdot\mathsf{flattenStatements}(\bar v)
      \quad(v\notin\mathit{Sequence}),\\
 \mathsf{flattenStatements}(\kw{sequence}(\bar s)\cdot\bar v)&=
      \bar s\cdot\mathsf{flattenStatements}(\bar v),\\
\end{aligned}
\end{equation}
The message and expression actions are exhaustive:
\begin{equation}\label{eq:newspeak-expression-actions}
\begin{aligned}
 \mathcal J_{\NS}(\mathsf{binaryMessage},
      \langle\kw{selector}(s),R\rangle,w)
   &=\kw{clause}_{z}(s,\langle R\rangle),\\
 \mathcal J_{\NS}(\mathsf{keywordMessage},
      \langle\kw{keyword}(k_i),R_i\rangle_{i=1}^{n},w)
   &=\kw{clause}_{z}(k_1\cdots k_n,\langle R_i\rangle_{i=1}^{n}),\\
 \mathcal J_{\NS}(\mathsf{message},
      \langle\kw{selector}(s)\rangle,w)&=\kw{clause}_{z}(s,\epsilon),\\
 \mathcal J_{\NS}(\mathsf{message},\langle c\rangle,w)&=c
      \quad(c=\kw{clause}_{z'}(s,\bar R)),\\
 \mathcal J_{\NS}(\mathsf{cascadeMessage},
      \langle\kw{selector}(s)\rangle,w)&=\kw{clause}_{z}(s,\epsilon),\\
 \mathcal J_{\NS}(\mathsf{cascadeMessage},\langle c\rangle,w)&=c,\\
\end{aligned}
\end{equation}
The productions that normalize to an ordered sequence of message clauses are
\begin{equation}\label{eq:newspeak-message-sequence-names}
 \mathcal N_M=\{\mathsf{nontrivialUnaryMessages},
   \mathsf{nontrivialBinaryMessages},
   \mathsf{keywordMessages},\mathsf{nonEmptyMessages}\}.
\end{equation}
\begin{equation}\label{eq:newspeak-send-fold-actions}
\begin{aligned}
 \mathcal J_{\NS}(\mathsf{primary},
      \langle\kw{selector}_{z'}(\#x)\rangle,w)&=\kw{identifier}_{z}(x),\\
 \mathcal J_{\NS}(\mathsf{primary},
      \langle\kw{classDecl}_{z'}(\bar v)\rangle,w)&=
      \kw{classExpr}_{z}(\kw{classDecl}_{z'}(\bar v)),\\
 \mathcal J_{\NS}(\mathsf{primary},\langle R\rangle,w)&=R
      \quad\text{for every other primary},\\
 \mathcal J_{\NS}(\mathsf{unaryExpression},\langle R,\bar c\rangle,w)
   &=\mathsf{foldExplicit}_{z}(R,\bar c),\\
 \mathcal J_{\NS}(\mathsf{binaryExpression},\langle R,\bar c\rangle,w)
   &=\mathsf{foldExplicit}_{z}(R,\bar c),\\
 \mathcal J_{\NS}(\mathsf{keywordExpression},\langle R,c?\rangle,w)
   &=\mathsf{foldExplicit}_{z}(R,c?),\\
 \mathcal J_{\NS}(\mathsf{eventualExpression},\langle R,c\rangle,w)
   &=\kw{eventual}_{z}(R,\selector(c),\mathsf{arguments}(c)),\\
 \mathcal J_{\NS}(\mathsf{implicitKeywordSend},\langle c\rangle,w)
   &=\kw{implicit}_{z}(\selector(c),\mathsf{arguments}(c)),\\
 \mathcal J_{\NS}(n,\bar v,w)&=
      \mathsf{flattenClauses}_{z}(\bar v)
      \quad(n\in\mathcal N_M),\\
 \mathcal J_{\NS}(\mathsf{messageCascade},
      \langle\bar c_0,\bar c\rangle,w)
   &=\kw{cascadeParts}_{z}(\bar c_0,\bar c),\\
 \mathcal J_{\NS}(\mathsf{cascadedExpression},\langle R\rangle,w)&=R,\\
 \mathcal J_{\NS}(\mathsf{cascadedExpression},
      \langle R,\kw{cascadeParts}(\bar c_0,\bar c)\rangle,w)
   &=\mathsf{makeCascade}_{z}(R,\bar c_0,\bar c),\\
\end{aligned}
\end{equation}
\begin{equation}\label{eq:newspeak-compound-expression-actions}
\begin{aligned}
 \mathcal J_{\NS}(\mathsf{expression},
      \langle\kw{name}(x),R\rangle,w)&=\kw{setter}_{z}(x,R),\\
 \mathcal J_{\NS}(\mathsf{expression},\langle R\rangle,w)&=R,\\
 \mathcal J_{\NS}(\mathsf{parenthesized},\langle R\rangle,w)&=\kw{paren}_{z}(R),\\
 \mathcal J_{\NS}(\mathsf{tuple},\bar R,w)&=\kw{tuple}_{z}(\bar R),\\
 \mathcal J_{\NS}(\mathsf{blockParameters},\bar d,w)&=\bar d,\\
 \mathcal J_{\NS}(\mathsf{codeBody},
      \langle g?,\kw{sequence}(\bar s)\rangle,w)
   &=\kw{body}_{z}(g?,\bar s),\\
 \mathcal J_{\NS}(\mathsf{block},
      \langle\bar d?,\kw{body}(g?,\bar s)\rangle,w)
   &=\kw{closure}_{z}(\mathsf{orEmptySeq}(\bar d?),
       \mathsf{orEmpty}(g?),\bar s),\\
 \mathcal J_{\NS}(\mathsf{returnStatement},\langle R\rangle,w)
   &=\kw{returnStmt}_{z}(R),\\
 \mathcal J_{\NS}(\mathsf{statements},\bar v,w)&=
      \kw{sequence}_{z}(\mathsf{flattenStatements}(\bar v)).
\end{aligned}
\end{equation}
A unary selector in primary position becomes $\kw{identifier}_{z}(x)$;
Equation~\eqref{eq:expression-elaboration} later makes it an implicit nullary
send.  An implicit keyword message instead takes the explicit
$\mathsf{implicitKeywordSend}$ branch.  Thus neither form introduces a
variable-read node.

Pattern actions are
\begin{equation}\label{eq:newspeak-pattern-actions}
\begin{aligned}
 \mathcal J_{\NS}(\mathsf{wildcardPattern},\epsilon,w)&=\kw{wild}_{z},\\
 \mathcal J_{\NS}(\mathsf{literalPattern},\langle R\rangle,w)&=
      \kw{literalPat}_{z}(R),\\
 \mathcal J_{\NS}(\mathsf{variablePattern},\langle\kw{name}(x)\rangle,w)&=
      \kw{variablePat}_{z}(x),\\
 \mathcal J_{\NS}(\mathsf{nestedPattern},\langle p\rangle,w)&=
      \kw{nestedPat}_{z}(p),\\
 \mathcal J_{\NS}(\mathsf{keywordPatternPair},
      \langle\kw{keyword}(k),p?\rangle,w)&=\langle k,p?\rangle,\\
 \mathcal J_{\NS}(\mathsf{keywordPattern},\bar v,w)&=
      \kw{keywordPat}_{z}(\bar v),
      \quad\bar v=\langle k_i,p_i?\rangle_{i=1}^{n},\\
 \mathcal J_{\NS}(\mathsf{pattern},\langle p\rangle,w)&=\kw{pattern}_{z}(p).
\end{aligned}
\end{equation}

For an inheritance clause define
\begin{equation}\label{eq:newspeak-inheritance-actions}
\begin{aligned}
 \mathsf{inheritanceReceiver}_{z}(\kw{none},\kw{selector}(s))
   &=\kw{implicit}_{z}(s,\epsilon),\\
 \mathsf{inheritanceReceiver}_{z}(R,\kw{selector}(s))
   &=\kw{message}_{z}(R,s,\epsilon),\\
 \mathsf{initializerClause}_{z}(\kw{none})&=
   \kw{clause}_{z}(\#\kw{new},\epsilon),\\
 \mathsf{initializerClause}_{z}(c)&=c,\\
 R_I&=\mathsf{inheritanceReceiver}_{z}(R?,\kw{selector}(s)),\\
 c_I&=\mathsf{initializerClause}_{z}(c?),\\
 \mathcal J_{\NS}(\mathsf{inheritanceClause},
      \langle R?,\kw{selector}(s),c?\rangle,w)
   &=\kw{inheritanceHead}_{z}(R_I,c_I).
\end{aligned}
\end{equation}
For compactness in the remaining action equations, define the non-overloaded
abbreviation
\begin{equation}\label{eq:newspeak-action-abbreviation}
 \mathsf{act}_{[i,j)}(n,\bar v,u[i:j])
   =\mathcal J_{\NS}(n,\bar v,u[i:j]).
\end{equation}
The complete declaration actions are
\begin{equation}\label{eq:newspeak-declaration-actions}
\begin{aligned}
 N_I&=\mathsf{inheritanceAndBody},\\
 \mathsf{sel}(\kw{patDecl}(s,\bar d))&=s,\\
 \mathsf{formals}(\kw{patDecl}(s,\bar d))&=\bar d,\\
 \mathsf{factorySel}(\kw{none})&=\#\kw{new},\\
 \mathsf{factorySel}(\kw{patDecl}(s,\bar d))&=s,\\
 \mathsf{factoryFormals}(\kw{none})&=\epsilon,\\
 \mathsf{factoryFormals}(\kw{patDecl}(s,\bar d))&=\bar d,\\
 \mathsf{act}_{z}(\mathsf{keywordPatternDecl},
      \langle\kw{keyword}(k_i),d_i\rangle_{i=1}^{n},w)&=\\[-1mm]
 &\quad\kw{patDecl}_{z}(k_1\cdots k_n,\langle d_i\rangle_{i=1}^{n}),\\
 \mathsf{act}_{z}(\mathsf{binaryPatternDecl},
      \langle\kw{selector}(s),d\rangle,w)&=
      \kw{patDecl}_{z}(s,\langle d\rangle),\\
 \mathsf{act}_{z}(\mathsf{messagePattern},
      \langle\kw{selector}(s)\rangle,w)&=\kw{patDecl}_{z}(s,\epsilon),\\
 \mathsf{act}_{z}(\mathsf{messagePattern},\langle p\rangle,w)&=p
      \quad(p=\kw{patDecl}_{z'}(s,\bar d)),\\
 \mathsf{act}_{z}(\mathsf{accessModifier},\bar v,w)&=
      \mathsf{decodeAccess}(w),\\
 \mathsf{act}_{z}(\mathsf{slotDecl},\langle\kw{name}(x)\rangle,w)&=
      \kw{formal}_{z}(x),\\
 \mathsf{act}_{z}(\mathsf{slotDef},\langle\alpha?,d,R?\rangle,w)&=\\[-1mm]
 &\quad\kw{slot}_{z}(\alpha?,\mathsf{name}(d),
      \mathsf{mutability}(w),R?),\\
 \mathsf{act}_{z}(\mathsf{sequentialSlots},\bar d,w)&=
      \kw{sequentialGroup}_{z}(\bar d),\\
 \mathsf{act}_{z}(\mathsf{simultaneousSlots},\bar d,w)&=
      \kw{simultaneousGroup}_{z}(\bar d),\\
 \mathsf{act}_{z}(\mathsf{lazySlotDecl},
      \langle\alpha?,d,R\rangle,w)&=\\[-1mm]
 &\quad\kw{lazySlot}_{z}(\alpha?,\mathsf{name}(d),
      \mathsf{mutability}(w),R),\\
 \mathsf{act}_{z}(\mathsf{methodDecl},
      \langle\alpha?,p,\kw{body}(g?,\bar s)\rangle,w)&=\\[-1mm]
 &\quad\kw{method}_{z}(\alpha?,\mathsf{sel}(p),
      \mathsf{formals}(p),\mathsf{orEmpty}(g?),\bar s),\\
 \mathsf{act}_{z}(\mathsf{nestedClassDecl},\langle\alpha?,d\rangle,w)&=
      \kw{nestedClass}_{z}(\alpha?,d),\\
 \mathsf{act}_{z}(\mathsf{classHeader},
      \langle g?,\kw{sequence}(\bar s)\rangle,w)&=
      \kw{header}_{z}(g?,\bar s),\\
 \mathsf{act}_{z}(\mathsf{instanceSide},\bar d,w)&=\bar d,\\
 \mathsf{act}_{z}(\mathsf{classSide},\bar d_c,w)&=\bar d_c,\\
 \mathsf{act}_{z}(\mathsf{classBody},
      \langle\kw{header}(g?,\bar s),\bar d,\bar d_c?\rangle,w)&=\\[-1mm]
 &\quad\kw{classStructure}_{z}
      (g?,\bar s,\bar d,\mathsf{orEmptySeq}(\bar d_c?)),\\
 \mathsf{act}_{z}(\mathsf{mixinSuffix},\langle\bar h,h?\rangle,w)&=
      \kw{mixinTail}_{z}(\bar h,h?),\\
 \mathsf{act}_{z}(N_I,\langle h\rangle,w)&=
      \kw{defaultInheritance}_{z}(h),\\
 \mathsf{act}_{z}(N_I,
      \langle\kw{inheritanceHead}(R,c),h\rangle,w)&=
      \kw{explicitInheritance}_{z}(R,c,h),\\
 h_I&=\kw{inheritanceHead}(R,c),\\
 t_I&=\kw{mixinTail}(\bar c,h?),\\
 \mathsf{act}_{z}(N_I,\langle h_I,t_I\rangle,w)&=\\[-1mm]
 &\quad\kw{mixinChain}_{z}(R,c,\bar c,h?),\\
 \mathsf{act}_{z}(\mathsf{classDeclaration},
      \langle\kw{name}(x),p?,b\rangle,w)&=
      \kw{classDecl}_{z}(x,\mathsf{factorySel}(p?),
         \mathsf{factoryFormals}(p?),b),\\
 \mathsf{act}_{z}(\mathsf{objectLiteral},\langle h\rangle,w)&=
      \kw{object}_{z}(\kw{implicitObjectHeader}_{z}(h)),\\
 R_o&=\kw{identifier}_{z}(x),\\
 c_o&=\mathsf{initializerClause}_{z}(c?),\\
 o_h&=\kw{explicitObjectHeader}_{z}(R_o,c_o,h),\\
 \mathsf{act}_{z}(\mathsf{objectLiteral},
      \langle\kw{name}(x),c?,h\rangle,w)&=\\[-1mm]
 &\quad\kw{object}_{z}(o_h),\\
 \mathsf{act}_{z}(\mathsf{compilationUnit},
      \langle\lambda,\gamma?,d\rangle,w)&=\\[-1mm]
 &\quad\kw{unit}_{z}(\lambda,\gamma?,d).
\end{aligned}
\end{equation}
$\mathsf{mutability}(w)$ is immutable exactly when the slot definition's
source slice $w$ contains \texttt{=} rather than \texttt{::=}; either
assignment operator makes a lazy slot mutable or immutable in the same way,
and omission is mutable.
The inheritance actions construct exactly the three $b$ alternatives in
Equation~\eqref{eq:surface-declaration-domain}, in source order.  All
punctuation, trivia, and wrapper productions have either
the discard action or the unique-child action.  Equations
\eqref{eq:newspeak-literal-actions}--\eqref{eq:newspeak-declaration-actions},
together with those two structural actions, define $\mathcal J_{G_{\NS}}$ on
every successful captured production; there is no residual action case.

A metadata token is $m=\langle\tau,w,[i,j)\rangle$, where $\tau$ is its tag
and $w$ its uninterpreted payload; an ordinary comment has tag $\kw{comment}$.
Let $\mathsf{triviaOnly}_{u}(i,j)$ mean that $u[i:j]$ contains only whitespace
and comments, and define
\begin{equation}\label{eq:metadata-trivia-gap}
 \mathsf{triviaBetween}_{u}(X,Y)\Longleftrightarrow
 \mathsf{end}(X)\le\mathsf{start}(Y)\land
 \mathsf{triviaOnly}_{u}(\mathsf{end}(X),\mathsf{start}(Y)).
\end{equation}
The association relation is
\begin{equation}\label{eq:metadata-association}
\begin{aligned}
 m\mathrel{\mathsf{before}_{u}}R
 &\Longleftrightarrow
   \mathsf{triviaBetween}_{u}(m,R)\land\\[-1mm]
 &\qquad R\in\mathit{Expression}\uplus\mathit{FormalDeclaration}
                    \uplus\mathit{SlotDeclaration},\\
 m\mathrel{\mathsf{afterPattern}_{u}}R
 &\Longleftrightarrow
   \mathsf{triviaBetween}_{u}(\mathsf{pattern}(R),m)\land\\[-1mm]
 &\qquad R\in\mathit{MethodDeclaration}
                   \uplus\mathit{FactoryDeclaration},\\
 m\mathrel{\mathsf{afterHeader}_{u}}R
 &\Longleftrightarrow
   \mathsf{triviaBetween}_{u}(\mathsf{header}(R),m)\land
   R\in\mathit{ClassDeclaration},\\
 m\mathrel{\mathsf{staticTypeOf}_{u}}R
 &\Longleftrightarrow
   \mathsf{tag}(m)=\kw{staticType}\land
   \mathsf{span}(m)\subseteq\mathsf{span}(R)\land\\[-1mm]
 &\qquad R=\min_{\subseteq}\{R'\mid
      \mathsf{span}(m)\subseteq\mathsf{span}(R')\land
      R'\in\mathit{TypedDeclaration}\}.
\end{aligned}
\end{equation}
The applicable relation is selected by node kind in displayed order.  Metadata
attachment and its insertion into the tree are the functions
\begin{equation}\label{eq:metadata-attachment}
\begin{aligned}
 \mathsf{related}_{u}(m,R)&=
  \begin{cases}
   m\mathrel{\mathsf{staticTypeOf}_{u}}R&\mathsf{tag}(m)=\kw{staticType},\\
   m\mathrel{\mathsf{afterHeader}_{u}}R
      &R\in\mathit{ClassDeclaration},\\
   m\mathrel{\mathsf{afterPattern}_{u}}R
      &R\in\mathit{MethodDeclaration}\uplus\mathit{FactoryDeclaration},\\
   m\mathrel{\mathsf{before}_{u}}R&\text{otherwise,}
  \end{cases}\\
 \mathsf{metadata}_{u,\bar m}(R)&=
   \mathsf{sourceOrder}(\langle m\in\bar m\mid
      \mathsf{related}_{u}(m,R)\rangle),\\
 \mathsf{attachMetadata}_{u}(R,\bar m)&=
   R[\mathsf{meta}(n)\mapsto\mathsf{metadata}_{u,\bar m}(n)
       \mid n\in\mathsf{nodes}(R)].
\end{aligned}
\end{equation}
Punctuation, whitespace, and metadata tokens therefore do not become
executable AST children.

Parsed nodes next receive identities.  Let
$\kappa:\mathit{NodeKey}\rightharpoonup
 (\mathit{DeclId}\uplus\mathit{SiteId})$ be an optional identity-retention map,
where a key contains a source-unit identity, structural path, constructor kind,
and declared name when one exists.  Define
\begin{equation}\label{eq:stable-ast-identification}
 \mathsf{identify}_{\kappa}(R)=\langle S,\kappa'\rangle
 \quad\Longleftrightarrow\quad
 \begin{cases}
  \mathsf{id}_{S}(n)=\kappa(\mathsf{key}_{R}(n))
    &\mathsf{key}_{R}(n)\in\dom(\kappa),\\
  \mathsf{id}_{S}(n)\text{ is fresh of the required sort}
    &\mathsf{key}_{R}(n)\notin\dom(\kappa),
 \end{cases}
\end{equation}
for every declaration or expression node $n$, with $\kappa'$ extending
$\kappa$ by the chosen keys.  The result must satisfy
\begin{equation}\label{eq:stable-ast-identity-validity}
 \mathsf{idsOK}(S)\Longleftrightarrow
 \mathsf{id}_{S}\text{ is injective}\ \land\
 \mathsf{id}_{S}(n)\in
 \begin{cases}\mathit{DeclId}&n\in\mathit{DeclarationNode}(S),\\
               \mathit{SiteId}&n\in\mathit{ExpressionNode}(S).
 \end{cases}
\end{equation}
An initial compilation uses $\kappa=\varnothing$.  A source-preserving editor
may supply a retained map; mirror changes that directly edit identified AST
nodes preserve identities according to Equation~\eqref{eq:source-patch} and do
not reparse text.

\subsection{Binding, elaboration, and static well-formedness}
\label{sec:static-elaboration}

An elaboration context is
\begin{equation}\label{eq:elaboration-context}
 \Gamma=\langle K,q,D_0,\bar B,r_t,h_t\rangle,
\end{equation}
where $K$ is the enclosing class-body declaration or $\kw{top}$, $q$ is the
enclosing activation declaration or $\kw{none}$, $D_0$ is the immediately
enclosing named class or $\kw{none}$, $\bar B$ is the innermost-first sequence
of lexical binder maps, and $r_t,h_t$ record whether local and non-local return
are permitted.  Lexical binding is the partial function
\begin{equation}\label{eq:static-binding-function}
 \mathsf{bind}_{\Gamma}(x)=d
 \quad\Longleftrightarrow\quad
 j=\min\{i\mid x\in\dom(B_i)\}\land d=B_j(x),
 \qquad \bar B=\langle B_0,\ldots,B_m\rangle.
\end{equation}
An absent set on the right makes the function undefined.  Class-name binding
uses the same definition restricted to class declarations; write
$\mathsf{classBind}_{\Gamma}(x)=D_t$.
The total annotation used by an ordinary implicit send is
\begin{equation}\label{eq:implicit-binding-annotation}
 \mathsf{bindingAnnotation}_{\Gamma}(x)=
 \begin{cases}
  \mathsf{bind}_{\Gamma}(x)&\mathsf{bind}_{\Gamma}(x)\ne\undef,\\
  \kw{none}&\mathsf{bind}_{\Gamma}(x)=\undef.
 \end{cases}
\end{equation}

Parentheses do not erase receiver authority.  Define
\begin{equation}\label{eq:surface-receiver-kind}
\begin{aligned}
 \mathsf{receiverKind}_{\Gamma}((R))&=\mathsf{receiverKind}_{\Gamma}(R),\\
 \mathsf{receiverKind}_{\Gamma}(\kw{self})&=\kw{selfReceiver}(D_0),\\
 \mathsf{receiverKind}_{\Gamma}(\kw{outer}(x))
   &=\kw{outerReceiver}(\mathsf{classBind}_{\Gamma}(x),D_0),\\
 \mathsf{receiverKind}_{\Gamma}(\kw{super})&=\kw{superReceiver},\\
 \mathsf{receiverKind}_{\Gamma}(R)&=\kw{ordinaryReceiver}
   \quad\text{otherwise.}
\end{aligned}
\end{equation}
For already elaborated arguments $\bar e$, put
$r_k=\mathsf{receiverKind}_{\Gamma}(R)$.  Send construction is the distinct
partial function
\begin{equation}\label{eq:elaborated-send-construction}
\begin{aligned}
 \mathsf{makeSend}_{\Gamma}(R,s,\bar e)&=
 \begin{cases}
  \kw{selfSend}(s,\bar e,D_0)&
    r_k=\kw{selfReceiver}(D_0),\\
  \kw{outerSend}(s,\bar e,D_t,D_0)&
    r_k=\kw{outerReceiver}(D_t,D_0),\\
  \kw{superSend}(s,\bar e)&
    r_k=\kw{superReceiver},\\
  \kw{send}(e_R,s,\bar e)&
    r_k=\kw{ordinaryReceiver}\land
      \mathsf{elabExpr}_{\Gamma}(R)=e_R.
 \end{cases}
\end{aligned}
\end{equation}
It is undefined if $D_0=\kw{none}$ in a self or outer case, if the named outer
class does not bind, or if super occurs outside a method activation.  Notice
again that the super case receives no declaration annotation.

Use distinct functions for a single expression, an expression list, a
statement sequence, a declaration, and a complete unit:
\begin{equation}\label{eq:elaboration-function-signatures}
\begin{aligned}
 \mathsf{elabExpr}_{\Gamma}&:\mathit{SurfaceExpression}\rightharpoonup
     \mathit{CoreExpression},\\
 \mathsf{elabExprList}_{\Gamma}&:\mathit{SurfaceExpression}^{*}
     \rightharpoonup\mathit{CoreExpression}^{*},\\
 \mathsf{elabStatements}_{\Gamma}&:\mathit{SurfaceStatement}^{*}
     \rightharpoonup\mathit{CoreStatement}^{*},\\
 \mathsf{elabDecl}_{\Gamma}&:\mathit{SurfaceDeclaration}
     \rightharpoonup\mathit{DerivedDeclaration},\\
 \mathsf{deriveProgram}&:\mathit{IdentifiedUnit}\rightharpoonup
     (\mathit{Program}\times\mathit{CoreExpression}).
\end{aligned}
\end{equation}
The list and optional functions, selector annotations, and activation contexts
are defined by
\begin{equation}\label{eq:elaboration-helper-functions}
\begin{aligned}
 \mathsf{elabExprList}_{\Gamma}(\epsilon)&=\epsilon,\\
 \mathsf{elabExprList}_{\Gamma}(R\cdot\bar R)&=
   \mathsf{elabExpr}_{\Gamma}(R)\cdot
   \mathsf{elabExprList}_{\Gamma}(\bar R),\\
 \mathsf{elabOptionalExpr}_{\Gamma}(\kw{none})&=\kw{none},\\
 \mathsf{elabOptionalExpr}_{\Gamma}(R)&=\mathsf{elabExpr}_{\Gamma}(R),\\
 \mathsf{elabOptionalPattern}_{\Gamma}(\kw{none})&=\_,\\
 \mathsf{elabOptionalPattern}_{\Gamma}(p)&=\mathsf{elabPattern}_{\Gamma}(p),\\
 \mathsf{selectorBindingAnnotation}_{\Gamma}(s)&=
  \begin{cases}
   \mathsf{bindingAnnotation}_{\Gamma}(x)&s=\#x,\\
   \kw{none}&\mathsf{arity}(s)>0.
  \end{cases}
\end{aligned}
\end{equation}
For a declaration sequence, let
\begin{equation}\label{eq:elaboration-context-extension}
\begin{aligned}
 \mathsf{groupDecls}(\epsilon)&=\epsilon,\\
 \mathsf{groupDecls}(g\cdot\bar g)&=
     \mathsf{members}(g)\cdot\mathsf{groupDecls}(\bar g),\\
 \mathsf{binder}(\bar d)&=
     \{\mathsf{name}(d)\mapsto\mathsf{id}(d)\mid d\in\bar d\},\\
 \mathsf{returnFlags}_{\Gamma}(q)&=
  \begin{cases}
   \langle\mathsf{true},\mathsf{false}\rangle&q\in\mathit{MethodDeclId},\\
   \langle\mathsf{false},\Gamma.r_t\lor\Gamma.h_t\rangle
       &q\in\mathit{ClosureDeclId},
  \end{cases}\\
 \mathsf{activationContext}_{\Gamma}(q,\bar d,\bar g)&=
   \langle\Gamma.K,q,\Gamma.D_0,
     \mathsf{binder}(\bar d\cdot\mathsf{groupDecls}(\bar g))
       \cdot\Gamma.\bar B,r_q,h_q\rangle,\\[-1mm]
 &\qquad\langle r_q,h_q\rangle=\mathsf{returnFlags}_{\Gamma}(q).
\end{aligned}
\end{equation}
Thus $\Gamma_q=\mathsf{activationContext}_{\Gamma}(q,\bar d,\bar g)$ in every
closure or method case below.  The atomic and send cases are
\begin{equation}\label{eq:expression-elaboration}
\begin{aligned}
 \mathsf{elabExpr}_{\Gamma}(\kw{identifier}_{z}(x))
   &=\kw{implicit}(\#x,\epsilon,
      \mathsf{bindingAnnotation}_{\Gamma}(x)),\\
 \mathsf{elabExpr}_{\Gamma}(\kw{implicit}_{z}(s,\bar R))
   &=\kw{implicit}(s,\mathsf{elabExprList}_{\Gamma}(\bar R),
      \mathsf{selectorBindingAnnotation}_{\Gamma}(s)),\\
 \mathsf{elabExpr}_{\Gamma}(\kw{message}_{z}(R,s,\bar R'))
   &=\mathsf{makeSend}_{\Gamma}
       (R,s,\mathsf{elabExprList}_{\Gamma}(\bar R')),\\
 \mathsf{elabExpr}_{\Gamma}(\kw{eventual}_{z}(R,s,\bar R'))
   &=\kw{asyncSend}(\mathsf{elabExpr}_{\Gamma}(R),s,
                    \mathsf{elabExprList}_{\Gamma}(\bar R')),\\
 \mathsf{elabExpr}_{\Gamma}(\kw{atom}_{z}(\omega))&=\kw{atom}(\omega),\\
 \mathsf{elabExpr}_{\Gamma}(\kw{self}_{z})&=\kw{self},\\
 \mathsf{elabExpr}_{\Gamma}(\kw{paren}_{z}(R))&=\mathsf{elabExpr}_{\Gamma}(R),\\
 \mathsf{elabExpr}_{\Gamma}(\kw{tuple}_{z}(\bar R))
   &=\kw{tuple}(\mathsf{elabExprList}_{\Gamma}(\bar R)),\\
 \mathsf{elabExpr}_{\Gamma}(\kw{pattern}_{z}(p))
   &=\kw{pattern}(\mathsf{elabPattern}_{\Gamma}(p),
       \mathsf{bindingAnnotation}_{\Gamma}(\kw{Pattern})),\\
 \mathsf{elabExpr}_{\Gamma}(\kw{setter}_{z}(x,R))
   &=\mathsf{elabExpr}_{\Gamma}(\mathsf{setterExpansion}_{z}(x,R)),\\
 \mathsf{elabExpr}_{\Gamma}(\kw{cascade}_{z}(R,\bar c_0,\bar c))
   &=\mathsf{elabExpr}_{\Gamma}
      (\mathsf{cascadeExpansion}_{z}(R,\bar c_0,\bar c)).
\end{aligned}
\end{equation}
Here $\mathsf{selectorBindingAnnotation}_{\Gamma}$ is
Equation~\eqref{eq:implicit-binding-annotation} for a unary activation member
and is $\kw{none}$ otherwise.  Tuple and pattern nodes are retained because
their dynamic expansions are Equations~\eqref{eq:tuple-expansion} and
\eqref{eq:pattern-components}--\eqref{eq:variable-pattern-parameter}.

Pattern elaboration recursively elaborates only embedded expressions:
\begin{equation}\label{eq:pattern-elaboration}
\begin{aligned}
 \mathsf{elabPattern}_{\Gamma}(\kw{wild}_{z})&=\_,\\
 \mathsf{elabPattern}_{\Gamma}(\kw{literalPat}_{z}(R))&=
      \mathsf{elabExpr}_{\Gamma}(R),\\
 \mathsf{elabPattern}_{\Gamma}(\kw{variablePat}_{z}(x))&=?x_z,\\
 \mathsf{elabPattern}_{\Gamma}(\kw{nestedPat}_{z}(p))&=
      \langle\mathsf{elabPattern}_{\Gamma}(p)\rangle,\\
 \mathsf{elabPattern}_{\Gamma}
    (\kw{keywordPat}_{z}(\langle k_i,p_i?\rangle_i))&=
      \langle k_i,
        \mathsf{elabOptionalPattern}_{\Gamma}(p_i?)\rangle_i.
\end{aligned}
\end{equation}
An absent keyword component denotes a wildcard.  The resulting forms are
exactly the arguments of Equations~\eqref{eq:pattern-components}--
\eqref{eq:keyword-pattern-expansion}.

The setter expansion, with fresh $p\ne x$, is
\begin{equation}\label{eq:setter-expansion}
\begin{aligned}
 \mathsf{setterBody}_{z}(x,p)&=
   \langle\kw{exprStmt}(\kw{implicit}(\#x{:},
      \langle\kw{identifier}(p)\rangle)),
      \kw{exprStmt}(\kw{identifier}(p))\rangle,\\
 \mathsf{setterExpansion}_{z}(x,R)&=
   \kw{message}_{z}\bigl(\kw{closure}_{z}
      (\langle\kw{formal}(p)\rangle,\epsilon,
       \mathsf{setterBody}_{z}(x,p)),\\[-1mm]
 &\hspace{38mm}\selsym{value:},\langle R\rangle\bigr).
\end{aligned}
\end{equation}
It therefore evaluates $R$ once, sends the setter, and returns the argument.

A cascade receiver is evaluated once.  For fresh synthetic closure identity
$q_c$, formal declaration $d_c$, and identifier $x_c$, define
\begin{equation}\label{eq:cascade-expansion}
\begin{aligned}
 &\mathsf{cascadeExpansion}_{z}
   (R,\bar c_0,
      \langle\kw{clause}(s_1,\bar R_1),\ldots,
             \kw{clause}(s_m,\bar R_m)\rangle)\\
 &\quad=\kw{message}_{z}\bigl(
   \kw{closure}_{q_c}(\langle\kw{formal}_{d_c}(x_c)\rangle,\epsilon,\\[-1mm]
 &\hspace{31mm}\langle
   \kw{exprStmt}(\mathsf{foldExplicit}_{z}
      (\kw{identifier}(x_c),\bar c_0))\rangle\cdot\\[-1mm]
 &\hspace{31mm}\langle
   \kw{exprStmt}(\kw{message}(\kw{identifier}(x_c),s_i,\bar R_i))
   \mid1\le i\le m\rangle),\\[-1mm]
 &\hspace{31mm}\selsym{value:},\langle R\rangle\bigr).
\end{aligned}
\end{equation}
The binding map inside the synthesized closure maps $x_c$ to $d_c$.

The remaining expression cases and statement elaboration are
preceded by normalization of local declarations:
\begin{equation}\label{eq:local-group-elaboration}
\begin{aligned}
 \mathsf{elabLocalDecl}_{\Gamma}
  (\kw{slot}_{d}(\alpha?,x,\kw{immutable},R))
   &=\kw{immutable}(d,\mathsf{elabExpr}_{\Gamma}(R)),\\
 \mathsf{elabLocalDecl}_{\Gamma}
  (\kw{slot}_{d}(\alpha?,x,\kw{mutable},R?))
   &=\kw{mutable}(d,\mathsf{elabOptionalExpr}_{\Gamma}(R?)),\\
 \mathsf{elabLocalGroup}_{\Gamma}(\kw{sequentialGroup}(\bar d))
   &=\langle\mathsf{elabLocalDecl}_{\Gamma}(d)\mid d\in\bar d\rangle,\\
 \mathsf{elabLocalGroup}_{\Gamma}(\kw{simultaneousGroup}(\bar d))
   &=\langle\kw{simultaneous}
      (\langle\mathsf{elabLocalDecl}_{\Gamma}(d)\mid d\in\bar d\rangle)\rangle,\\
 \mathsf{elabLocalGroups}_{\Gamma}(\epsilon)&=\epsilon,\\
 \mathsf{elabLocalGroups}_{\Gamma}(g\cdot\bar g)&=
   \mathsf{elabLocalGroup}_{\Gamma}(g)\cdot
   \mathsf{elabLocalGroups}_{\Gamma}(\bar g).
\end{aligned}
\end{equation}
\begin{equation}\label{eq:compound-expression-elaboration}
\begin{aligned}
 \Gamma_q&=\mathsf{activationContext}_{\Gamma}(q,\bar d,\bar g),\\
 \ell_q&=\mathsf{elabLocalGroups}_{\Gamma_q}(\bar g),\\
 b_q&=\mathsf{elabStatements}_{\Gamma_q}(\bar s),\\
 \mathsf{elabExpr}_{\Gamma}(\kw{closure}_{q}(\bar d,\bar g,\bar s))
   &=\kw{closureLiteral}(q,\mathsf{names}(\bar d),\ell_q,b_q),\\
 \mathsf{elabExpr}_{\Gamma}(\kw{object}_{L}(o_h))
   &=\kw{objectLiteral}(\mathsf{objectDescriptor}_{\Gamma}(L,o_h),
       \mathsf{objectSuperclass}_{\Gamma}(o_h)),\\
 \mathsf{elabExpr}_{\Gamma}(\kw{classExpr}_{D}(d))
   &=\mathsf{elabClass}_{\Gamma}(D,d),\\
 \mathsf{elabStatements}_{\Gamma}
   (\langle\kw{exprStmt}(R_i)\rangle_i)&=\\[-1mm]
 &\quad\langle\kw{expr}(\mathsf{elabExpr}_{\Gamma}(R_i))\rangle_i,\\
 \mathsf{elabStatements}_{\Gamma}
   (\bar s\cdot\langle\kw{returnStmt}(R)\rangle)
   &=\mathsf{elabStatements}_{\Gamma}(\bar s)\cdot\langle e_r\rangle,\\
 e_r&=\kw{return}(\mathsf{elabExpr}_{\Gamma}(R)).
\end{aligned}
\end{equation}
$\Gamma_q$ prepends the formal and local binder maps, sets the activation to
$q$, and enables a non-local return exactly when $q$ is a closure with an
enclosing method.  A return can only be the final surface statement by
Equation~\eqref{eq:newspeak-expression-peg}.

Access defaulting and the declaration cases are explicit:
\begin{equation}\label{eq:declaration-elaboration}
\begin{aligned}
 \mathsf{accessOf}(\kw{none})&=\Protected,\\
 \mathsf{accessOf}(\kw{public})&=\Public,\\
 \mathsf{accessOf}(\kw{protected})&=\Protected,\\
 \mathsf{accessOf}(\kw{private})&=\Private,\\
 \Gamma_f&=\mathsf{activationContext}_{\Gamma}(i_f,\bar d,\bar g),\\
 \ell_f&=\mathsf{elabLocalGroups}_{\Gamma_f}(\bar g),\\
 b_f&=\mathsf{elabStatements}_{\Gamma_f}(\bar s),\\
 a_f&=\mathsf{accessOf}(\alpha?),\\
 d_f&=\kw{method}_{i_f}(\alpha?,s,\bar d,\bar g,\bar s),\\
 d_s&=\kw{slot}_{d}(\alpha?,x,m,R?),\\
 d_l&=\kw{lazySlot}_{d}(\alpha?,x,m,R),\\
 d_n&=\kw{nestedClass}_{d}(\alpha?,D),\\
 \mathsf{elabDecl}_{\Gamma}(d_f)
   &=\langle i_f,s,a_f,\mathsf{names}(\bar d),
       \ell_f,b_f,\Gamma.K\rangle,\\
 \mathsf{elabDecl}_{\Gamma}(d_s)&=\\[-1mm]
 &\quad\mathsf{ordinarySlotDef}(d,x,a_f,m,
       \mathsf{elabOptionalExpr}_{\Gamma}(R?)),\\
 \mathsf{elabDecl}_{\Gamma}(d_l)&=\\[-1mm]
 &\quad\mathsf{lazySlotDef}(d,x,a_f,m,
       \mathsf{elabExpr}_{\Gamma}(R)),\\
 \mathsf{elabDecl}_{\Gamma}(d_n)&=\mathsf{nestedClassDef}(d,a_f,
       \mathsf{elabClass}_{\Gamma}(D,D)).
\end{aligned}
\end{equation}
$\Gamma_f$ prepends the method's formal and local declarations and enables
local return.  The ordinary- and lazy-slot constructors produce the physical
slot and the exact getter/setter bodies of
Section~\ref{sec:synthesized-methods}.  The nested-class constructor produces
its memoizing getter.  These constructors have different domains and names.

Class and object elaboration uses the message-template domain in
Equation~\eqref{eq:message-template-domain}.  Define
\begin{equation}\label{eq:front-end-message-template}
 \mathsf{elabMessageTemplate}_{\Gamma}
   (\kw{clause}(s,\bar R))=
   \langle s,\mathsf{elabExprList}_{\Gamma}(\bar R)\rangle.
\end{equation}
For a named class declaration $D$ with structure $h$, let $M_D$ and $M_D^c$
be the stable instance- and class-side mixins induced by $h$.  Its descriptor
constructor is
\begin{equation}\label{eq:front-end-body-descriptor}
\begin{aligned}
 \mathsf{bodyDescriptor}_{\Gamma}
   (D,x,u,s_f,\bar d,h,c)
   =\langle D,x,u,s_f,\mathsf{names}(\bar d),
      M_D,M_D^c,\mathsf{elabMessageTemplate}_{\Gamma}(c)\rangle.
\end{aligned}
\end{equation}
This is exactly the tuple $\chi_b$ of Equation~\eqref{eq:body-descriptor}.
Omitted inheritance is
\begin{equation}\label{eq:front-end-default-superclass}
\begin{aligned}
 \mathsf{defaultSuperclass}_{\Gamma}&=
  \begin{cases}
   \Object&\Gamma.K=\kw{top},\\
   \kw{implicit}(\#\kw{Object},\epsilon,
      \mathsf{bindingAnnotation}_{\Gamma}(\kw{Object}))&\text{otherwise,}
  \end{cases}\\
 \mathsf{defaultInitializer}&=\kw{clause}(\#\kw{new},\epsilon).
\end{aligned}
\end{equation}
The two non-application class cases are therefore
\begin{equation}\label{eq:front-end-class-body-elaboration}
\begin{aligned}
 &\mathsf{elabClass}_{\Gamma}
  (D,\kw{classDecl}(x,s_f,\bar d,\kw{defaultInheritance}(h)))\\
 &\quad=\kw{classBody}(\chi_b,\mathsf{defaultSuperclass}_{\Gamma}),\\[-1mm]
 &\qquad\chi_b=\mathsf{bodyDescriptor}_{\Gamma}
      (D,x,\kw{implicitSuperclass},s_f,\bar d,h,
       \mathsf{defaultInitializer}),\\
 &\mathsf{elabClass}_{\Gamma}
  (D,\kw{classDecl}(x,s_f,\bar d,
       \kw{explicitInheritance}(R,c,h)))\\
 &\quad=\kw{classBody}
   (\chi_b,\mathsf{elabExpr}_{\Gamma}(R)),\\[-1mm]
 &\qquad\chi_b=\mathsf{bodyDescriptor}_{\Gamma}
      (D,x,\kw{explicitSuperclass},s_f,\bar d,h,c).
\end{aligned}
\end{equation}

For a mixin chain, write its first inheritance head as
$h_0=\kw{inheritanceHead}(R_0,c_0)$ and its nonempty suffix as
$\bar h=\langle h_1,\ldots,h_k\rangle$, where
$h_j=\kw{inheritanceHead}(R_j,c_j)$.  Each application uses $c_{j-1}$ to
initialize its superclass and $c_j$ to initialize the newly applied mixin.
The synthetic identities and names required for an intermediate application
are deterministic injections of the enclosing declaration identity and its
one-based position:
\begin{equation}\label{eq:front-end-application-identities}
\begin{aligned}
 D_{A,j}&=\mathsf{applicationDeclId}(D,j),&
 i_{A,j}&=\mathsf{applicationInitId}(D,j),\\
 n_{A,j}&=\mathsf{applicationName}(D,j),&
 \bar x_{A,j}&=\mathsf{applicationFormals}(D,j,\selector(c_j)).
\end{aligned}
\end{equation}
The names in $\bar x_{A,j}$ are pairwise fresh, and their number is the arity
of $\selector(c_j)$.  Let $M^c_{A,j}$ be the stable class-side mixin containing
only the synthesized public primary factory for that signature.  The complete
descriptor of an intermediate application is
\begin{equation}\label{eq:front-end-anonymous-application-descriptor}
\begin{aligned}
 \mathsf{anonAppDesc}_{\Gamma}(D,j,c_{j-1},c_j)
   &=\langle D_{A,j},n_{A,j},\selector(c_j),\bar x_{A,j},\\[-1mm]
   &\quad M^c_{A,j},i_{A,j},m_{j-1},m_j\rangle,\\
 m_{j-1}&=\mathsf{elabMessageTemplate}_{\Gamma}(c_{j-1}),\\
 m_j&=\mathsf{elabMessageTemplate}_{\Gamma}(c_j).
\end{aligned}
\end{equation}
For the last application of a body-less named class, the declaration's own
identity, name, factory formals, and class-side mixin replace those synthetic
components:
\begin{equation}\label{eq:front-end-named-application-descriptor}
\begin{aligned}
 &\mathsf{namedAppDesc}_{\Gamma}
   (D,x,s_f,\bar d,c_{k-1},c_k)={}\\
 &\quad\langle D,x,s_f,\mathsf{names}(\bar d),M_D^c,i_{A,k},
   \mathsf{elabMessageTemplate}_{\Gamma}(c_{k-1}),
   \mathsf{elabMessageTemplate}_{\Gamma}(c_k)\rangle.
\end{aligned}
\end{equation}
The anonymous prefix fold returns both the resulting class expression and the
last initializer clause; retaining the latter makes the next superclass
initializer explicit:
\begin{equation}\label{eq:front-end-mixin-fold}
\begin{aligned}
 \mathsf{foldAnonymous}_{\Gamma,D}(E,c_0,\epsilon)&=\langle E,c_0\rangle,\\
 \mathsf{foldAnonymous}_{\Gamma,D}
  (E,c_{j-1},\langle h_j\rangle\cdot\bar h)
   &=\mathsf{foldAnonymous}_{\Gamma,D}(E_j,c_j,\bar h),\\
 E_j&=\kw{mixinApply}(\chi_{A,j},E,
       \mathsf{elabExpr}_{\Gamma}(R_j)),\\
 \chi_{A,j}&=\mathsf{anonAppDesc}_{\Gamma}(D,j,c_{j-1},c_j).
\end{aligned}
\end{equation}
A body-less chain reserves its last application for the named descriptor.  Put
\begin{equation}\label{eq:front-end-bodyless-prefix}
 \mathsf{foldAnonymous}_{\Gamma,D}
  (\mathsf{elabExpr}_{\Gamma}(R_0),c_0,
   \langle h_1,\ldots,h_{k-1}\rangle)
   =\langle E_{k-1},c_{k-1}\rangle.
\end{equation}
Then
\begin{equation}\label{eq:front-end-bodyless-mixin-elaboration}
\begin{aligned}
 &\mathsf{elabClass}_{\Gamma}
  (D,\kw{classDecl}(x,s_f,\bar d,
     \kw{mixinChain}(R_0,c_0,\bar h,\kw{none})))={}\\
 &\quad\kw{mixinApply}(
   \mathsf{namedAppDesc}_{\Gamma}
      (D,x,s_f,\bar d,c_{k-1},c_k),
   E_{k-1},\mathsf{elabExpr}_{\Gamma}(R_k)).
\end{aligned}
\end{equation}
This equation is defined only when $s_f=\selector(c_k)$ and $k\ge1$, as the
grammar requires.  If the chain ends in a body, put
\begin{equation}\label{eq:front-end-bodied-prefix}
 \mathsf{foldAnonymous}_{\Gamma,D}
  (\mathsf{elabExpr}_{\Gamma}(R_0),c_0,\bar h)=\langle E_k,c_k\rangle.
\end{equation}
Then
\begin{equation}\label{eq:front-end-bodied-mixin-elaboration}
\begin{aligned}
 &\mathsf{elabClass}_{\Gamma}
  (D,\kw{classDecl}(x,s_f,\bar d,
     \kw{mixinChain}(R_0,c_0,\bar h,h)))={}\\
 &\quad\kw{classBody}
   (\mathsf{bodyDescriptor}_{\Gamma}
      (D,x,\kw{explicitSuperclass},s_f,\bar d,h,c_k),E_k).
\end{aligned}
\end{equation}
Equations~\eqref{eq:front-end-anonymous-application-descriptor} and
\eqref{eq:front-end-named-application-descriptor} are instances of the complete
application tuple $\chi_a$ in Equation~\eqref{eq:application-descriptor}; no
descriptor component is supplied informally.

An object literal uses the same body derivation with an object-literal identity
$L$, the public nullary synthetic class-side factory, and the optional header's
inheritance receiver and initializer clause:
\begin{equation}\label{eq:front-end-object-elaboration}
\begin{aligned}
 m_0&=\mathsf{elabMessageTemplate}_{\Gamma}
         (\mathsf{defaultInitializer}),\\
 m_1&=\mathsf{elabMessageTemplate}_{\Gamma}(c),\\
 \mathsf{objectDescriptor}_{\Gamma}
  (L,\kw{implicitObjectHeader}(h))
   &=\langle L,\kw{implicitSuperclass},M_{L,h},M^c_{L,h},m_0\rangle,\\
 \mathsf{objectDescriptor}_{\Gamma}
  (L,\kw{explicitObjectHeader}(R,c,h))
   &=\langle L,\kw{explicitSuperclass},M_{L,h},M^c_{L,h},m_1\rangle,\\
 \mathsf{objectSuperclass}_{\Gamma}(\kw{implicitObjectHeader}(h))
   &=\mathsf{defaultSuperclass}_{\Gamma},\\
 \mathsf{objectSuperclass}_{\Gamma}(\kw{explicitObjectHeader}(R,c,h))
   &=\mathsf{elabExpr}_{\Gamma}(R).
\end{aligned}
\end{equation}
Equation~\eqref{eq:front-end-object-elaboration} is exactly the descriptor
$\chi_o$ in Equation~\eqref{eq:object-literal-descriptor}.

For an identified unit $S$, the maps consumed by $P$ are now defined directly.
Let $\mathsf{ancestors}_{S}(n)$ be the strict lexical-ancestor sequence of
node $n$, nearest first.  Let $M_K$ be the stable mixin identity derived from
class-body identity $K$.  Then
\begin{equation}\label{eq:derived-program-indexes}
\begin{aligned}
 \mathsf{sourceMap}_{S}(d)&=\mathsf{subtree}_{S}(d),\\
 \mathsf{declMixinMap}_{S}(K)&=M_K,\\
 \mathsf{lexParentMap}_{S}(D)&=
   \mathsf{firstClass}(\mathsf{ancestors}_{S}(D)),\\
 \mathsf{bodyMap}_{S}(i_f)&=
   \mathsf{elabStatements}_{\Gamma_{i_f}}(\mathsf{surfaceBody}_{S}(i_f)),\\
 \mathsf{scopeChainMap}_{S}(z)&=
   \langle d\in\mathsf{ancestors}_{S}(z)\mid
      d\in\mathit{ClassBodyDeclId}\uplus\mathit{ActDeclId}\rangle,\\
 \mathsf{activationMemberMap}_{S}(q)&=
   \{\#x\mid x\text{ is a formal or local declared directly by }q\},\\
 \mathsf{ownerMap}_{S}(K)&=
  \begin{cases}
   \kw{activationOwner}&\mathsf{firstOwner}_{S}(K)\in\mathit{ActDeclId},\\
   \kw{classOwner}&\mathsf{firstOwner}_{S}(K)\in\mathit{ClassDeclId},\\
   \kw{objectLiteralOwner}&\mathsf{firstOwner}_{S}(K)\in\mathit{ObjLitDeclId},\\
   \kw{topOwner}&\mathsf{firstOwner}_{S}(K)=\kw{none}.
  \end{cases}
\end{aligned}
\end{equation}
The mixin entry itself is
\begin{equation}\label{eq:derived-mixin-entry}
\begin{aligned}
 \phi_K&=\mathsf{explicitMethods}_{S}(K)
          \uplus\mathsf{synthMethods}_{S}(K),\\
 \delta_K&=\mathsf{physicalSlots}_{S}(K),
 &\nu_K&=\mathsf{nestedDecls}_{S}(K),\\
 \iota_K&=\mathsf{initializer}_{S}(K),
 &\mathsf{mixinMap}_{S}(M_K)&=
   \langle K,\phi_K,\delta_K,\nu_K,\iota_K\rangle.
\end{aligned}
\end{equation}
The disjoint union is defined only when selectors are unique.
$\mathsf{initializer}_{S}(K)$ is the tuple in
Equation~\eqref{eq:mixin-initializer}: its factory selector and formals come
from the class header, its source-ordered groups are transformed by
Equation~\eqref{eq:declaration-elaboration}, and its body by
Equation~\eqref{eq:compound-expression-elaboration}.

Top-level actor seeds and the one underspecified pattern form are
\begin{equation}\label{eq:derived-auxiliary-indexes}
\begin{aligned}
 \mathsf{actorSeedMap}_{S}(M_K)&=
   \langle K,M_K,\mathsf{factoryDescriptor}_{S}(K)\rangle
   &&\mathsf{ownerMap}_{S}(K)=\kw{topOwner},\\
 \mathsf{patternVariableMap}_{S}(z)&=
   \mathsf{patternVariableParameter}_{S}(z)
   &&\kw{variablePat}_{z}(x)\in S.
\end{aligned}
\end{equation}
The latter is precisely the explicit parameter isolated by
Equation~\eqref{eq:variable-pattern-parameter}; it is not used for any other
surface form.

Collecting Equations~\eqref{eq:derived-program-indexes}--
\eqref{eq:derived-auxiliary-indexes}, define
\begin{equation}\label{eq:declaration-derivation}
\begin{aligned}
 \mathsf{deriveProgram}(S)=\langle P,e\rangle
 \Longleftrightarrow{}&
 &P.\mathit{unit}=S\land\\
 &P.\mathit{source}=\mathsf{sourceMap}_{S}\land\\
 &P.\mathit{declMixin}=\mathsf{declMixinMap}_{S}\land\\
 &P.\mathit{mixin}=\mathsf{mixinMap}_{S}\land\\
 &P.\mathit{lexParent}=\mathsf{lexParentMap}_{S}\land\\
 &P.\mathit{body}=\mathsf{bodyMap}_{S}\land\\
 &P.\mathit{scopeChain}=\mathsf{scopeChainMap}_{S}\land\\
 &P.\mathit{activationMembers}=\mathsf{activationMemberMap}_{S}\land\\
 &P.\mathit{owner}=\mathsf{ownerMap}_{S}\land\\
 &P.\mathit{actorSeed}=\mathsf{actorSeedMap}_{S}\land\\
 &P.\mathit{patternVariable}=\mathsf{patternVariableMap}_{S}\land\\
 &e=\mathsf{elabTop}_{S}(S).
\end{aligned}
\end{equation}
The unchanged atom and platform maps are supplied as parameters to
$\mathsf{deriveProgram}$ when a program is installed.

The static predicates are defined without an informal residual clause:
\begin{equation}\label{eq:static-well-formedness-components}
\begin{aligned}
 \mathsf{bindingsOK}(S)&\Longleftrightarrow
   \bigl(\forall\Gamma,B_i\in\Gamma.\bar B.\;
      \mathsf{declaredNames}(B_i)\text{ are pairwise distinct}\bigr)\land\\[-1mm]
 &\quad\bigl(\forall\kw{outer}_{z}(x)\in S.\;
      \mathsf{classBind}_{\Gamma_z}(x)\ne\undef\land
      \Gamma_z.D_0\ne\kw{none}\bigr)\land\\[-1mm]
 &\quad\bigl(\forall\kw{self}_{z}\text{ used as receiver in }S.\;
      \Gamma_z.D_0\ne\kw{none}\bigr),\\
 \mathsf{membersOK}(S)&\Longleftrightarrow
   \forall K,s.\;|\{f\mid f\text{ is induced directly in }K\land
                         \selector(f)=s\}|\le1,\\
 \mathsf{arityOK}(S)&\Longleftrightarrow
   \forall f.\;|f.\bar x|=\mathsf{arity}(\selector(f)),\\
 \mathsf{returnsOK}(S)&\Longleftrightarrow
   \forall R\in\mathit{ReturnNode}(S).\;
   (\mathsf{local}(R)\Rightarrow r_t(R))\land
   (\mathsf{nonlocal}(R)\Rightarrow h_t(R)),\\
 \mathsf{contextsOK}(S)&\Longleftrightarrow
   \mathsf{bindingsOK}(S)\land\mathsf{returnsOK}(S)\land
   \forall R\in\mathit{SuperSendNode}(S).\;q(R)\ne\kw{none},\\
 \mathsf{surfaceOK}(S)&\Longleftrightarrow
   \mathsf{idsOK}(S)\land\mathsf{membersOK}(S)\land
   \mathsf{arityOK}(S)\land\mathsf{contextsOK}(S)\land
   \mathsf{shapesOK}(S).
\end{aligned}
\end{equation}
The surface-shape condition is
\begin{equation}\label{eq:surface-context-validity}
\begin{split}
 \mathsf{lastClause}
  (\langle h_1,\ldots,h_k\rangle)&=c_k
   \quad\text{when }h_k=\kw{inheritanceHead}(R_k,c_k),\\
 \mathsf{shapesOK}(S)\Longleftrightarrow{}&
  \text{every }\kw{outer}_{z}(x)\text{ or }\kw{super}_{z}
     \text{ is the receiver of a message}\ \land\\
 &\text{every return is final in its statement sequence}\ \land\\
 &\text{every immutable slot has an initializer}\ \land\\
 &\text{every lazy slot has an initializer}\ \land\\
 &\forall\kw{classDecl}(x,s_f,\bar d,
       \kw{mixinChain}(R_0,c_0,\bar h,\kw{none}))\in S.\;
       s_f=\selector(\mathsf{lastClause}(\bar h)).
\end{split}
\end{equation}
Ordinary unbound implicit sends are permitted by
Equation~\eqref{eq:implicit-binding-annotation}; named outer targets and
duplicate declarations in one binder map are rejected by the first predicate.

The source-to-core judgment is
\begin{equation}\label{eq:source-to-annotated-core}
 \begin{gathered}
 G,\kappa,u\vdash\mathsf{frontEnd}\Downarrow
   \langle P,e,\kappa'\rangle\\[-1mm]
 \Longleftrightarrow\\[-1mm]
  \mathsf{parse}_{G}(u)=t,\quad
  \mathsf{ast}_{G,u}(t)=\langle R_0,\bar m\rangle,\quad
  R=\mathsf{attachMetadata}_{u}(R_0,\bar m),\quad
  \mathsf{identify}_{\kappa}(R)=\langle S,\kappa'\rangle,\\
  \mathsf{surfaceOK}(S),\quad
  \mathsf{deriveProgram}(S)=\langle P,e\rangle,\quad
  \mathsf{TL}_{P}(e).
 \end{gathered}
\end{equation}
Core validity must cover expressions stored throughout $P$, not only the
top-level expression.  Let $\mathsf{coreRoots}(P,e)$ be the least finite set
containing $e$ and every expression occurring in method and closure bodies,
their local initializers, mixin initializer groups and bodies, pattern-variable
expansions, and the deferred superclass or mixin messages stored in class,
mixin-application, object-literal, and actor-seed descriptors:
\begin{equation}\label{eq:annotated-core-roots}
\begin{split}
 \mathsf{coreRoots}(P,e)=\mathsf{sub}\bigl(&\{e\}
   \cup P.\mathit{methodBodies}\cup P.\mathit{methodLocals}\\
 &\cup P.\mathit{closureBodies}\cup P.\mathit{closureLocals}\\
 &\cup P.\mathit{mixinInitializers}\cup P.\mathit{patternExpansions}\\
 &\cup P.\mathit{classMessages}\cup P.\mathit{mixinMessages}\\
 &\cup P.\mathit{objectMessages}\cup P.\mathit{actorMessages}\bigr).
\end{split}
\end{equation}
Here $\mathsf{sub}$ recursively includes tuple, pattern, cascade, send, lazy,
class, mixin-application, object-literal, and nested-class subexpressions.  For
Equation~\eqref{eq:annotated-core-roots}, $P.\mathit{mixinInitializers}$
projects the groups and bodies of all mixin initializers,
$P.\mathit{patternExpansions}$ is the range of the finite pattern-variable
table, and the four $\mathit{Messages}$ projections collect the message
templates in the corresponding finite descriptor tables.  These are
abbreviations for existing Program fields, not new semantic stores.  For
fuel $k$, $\mathsf{annOK}^{k}_{P}(r,b)$ checks every declaration and scope
annotation in $r$, where $b$ says whether a super send is permitted.  In
particular, it checks that named class and closure declarations occur in the
appropriate finite Program table, that an annotated implicit send is declared
by its recorded scope, and that a super send occurs only in installed class
code.  The executable gate is
\begin{equation}\label{eq:annotated-core-checker}
\begin{split}
 \mathsf{coreCheck}_{k}(P,e)=\kw{true}\Longleftrightarrow{}&
 P.\mathit{declMixin}\text{ is injective}\ \land\\
 &\forall r\in\mathsf{coreRoots}(P,e)\setminus\{e\}.\;
       \mathsf{annOK}^{k}_{P}(r,\kw{true})\ \land\\
 &\mathsf{annOK}^{k}_{P}(e,\kw{false})\ \land\ \mathsf{TL}_{P}(e).
\end{split}
\end{equation}
Core validity, used by loaders and reflective validation, is therefore the
independently checkable predicate
\begin{equation}\label{eq:annotated-core-validity}
\begin{split}
 \mathsf{coreOK}(P,e)\Longleftrightarrow{}&
 \exists k>0.\ \mathsf{coreCheck}_{k}(P,e)=\kw{true}\ \land\\
 &\mathsf{deriveProgram}(P.\mathit{unit})=\langle P,e\rangle\ \land
   \mathsf{surfaceOK}(P.\mathit{unit}).
\end{split}
\end{equation}
For reflective validation, which replaces a program but does not itself start
its top-level expression, define the projection
\begin{equation}\label{eq:annotated-program-validity}
 \mathsf{coreProgramOK}(P)
 \Longleftrightarrow
 \exists e.\;\mathsf{deriveProgram}(P.\mathit{unit})=\langle P,e\rangle
              \land\mathsf{coreOK}(P,e).
\end{equation}
Equations~\eqref{eq:source-to-annotated-core} and
\eqref{eq:annotated-core-validity} are the promised static connection from the
Specification's PEG AST to every declaration and send annotation consumed by
the dynamic semantics.  Moreover, determinism of PEG recognition and the
functional equations above give the front-end uniqueness property
\begin{equation}\label{eq:front-end-uniqueness}
 \frac{G_{\NS},\kappa,u\vdash\mathsf{frontEnd}\Downarrow X\qquad
       G_{\NS},\kappa,u\vdash\mathsf{frontEnd}\Downarrow Y}{X=Y}.
\end{equation}
Totality holds on exactly the inputs that parse completely and satisfy
$\mathsf{surfaceOK}$; failure at any earlier partial function is a static
error, never a stuck core term.

\subsection{Reified language services}
\label{sec:reified-front-ends}

The host registry already introduced in Equation~\eqref{eq:compilation-unit}
now has the concrete type
\begin{equation}\label{eq:language-registry}
 \mathcal L:\mathit{LanguageId}\rightharpoonup\mathit{ObjRef}.
\end{equation}
$\mathcal L(\lambda)$ is an ordinary Newspeak language object.  It answers the
public message \texttt{frontEndImage} with an immutable value
\begin{equation}\label{eq:front-end-image}
 I_L=\kw{frontEndImage}\langle o_p,o_e,o_c,o_s,G,v_L\rangle,
\end{equation}
whose components are parser, elaborator, compiler, and Simulator objects, the
declarative PEG value, and a language-defined revision value.  This record is
a protocol value, not a new privileged heap-record kind.  Its component
objects have ordinary classes, methods, enclosing objects, and slots.

For a caller actor $A$, write
\begin{equation}\label{eq:ordinary-service-call}
 \Omega,A\vdash\mathsf{call}(o,s,\bar v)
    \Downarrow\langle\Omega',v\rangle
\end{equation}
iff the ordinary public send of $\langle s,\bar v\rangle$ to $o$ starts in
$A$, follows the actor and sequential transition relations already defined,
and returns $v$ to its caller.  A signaled exception or broken reply produces
the corresponding failure rather than a successful judgment.  Thus
Equation~\eqref{eq:ordinary-service-call} is notation for the reflexive
transitive closure of existing rules, not a second evaluator.

Surface ASTs and annotated-core images are ordinary immutable Newspeak value
graphs.  Executable artifacts are immutable values as well.  The VM boundary
supplies injective partial
decoders
\begin{equation}\label{eq:front-end-value-decoders}
\begin{aligned}
 \mathsf{decodeSurface}&:\mathit{ObjRef}\rightharpoonup\mathit{SurfaceAST},\\
 \mathsf{decodeCore}&:\mathit{ObjRef}\rightharpoonup
        (\mathit{Program}\times\mathit{CoreExpression}),\\
 \mathsf{decodeArtifact}&:\mathit{ObjRef}\rightharpoonup\mathit{Artifact}.
\end{aligned}
\end{equation}
They inspect representation and allocate no object, invoke no user method, and
accept only finite immutable graphs.  Their inverses are ordinary constructors
provided by the language implementation.

Admission is attached to the values obtained at a use boundary, not to the
permanent identity of a component object.  A front-end image is admitted when
all four component references are live ordinary Newspeak objects, its grammar
passes Equation~\eqref{eq:peg-admissibility-certificate}, and successful
elaboration satisfies the complete checker:
\begin{equation}\label{eq:front-end-artifact-admission}
\begin{split}
 \mathsf{frontEndAdmitted}_{\Omega}(I_L)\Longleftrightarrow{}&
 I_L=\kw{frontEndImage}\langle o_p,o_e,o_c,o_s,G,v_L\rangle\ \land\\
 &\bigwedge_{o\in\{o_p,o_e,o_c,o_s\}}\mathsf{liveOrdinary}_{\Omega}(o)
 \ \land\ \mathsf{pegCertOK}(G)=\kw{true}\ \land\\
 &\forall S,P,e.\;
   \mathsf{surfaceOK}(S)\land
   \mathsf{deriveProgram}(S)=\langle P,e\rangle\\
 &\hspace{24mm}\Rightarrow\exists k>0.\;
       \mathsf{coreCheck}_{k}(P,e)=\kw{true}.
\end{split}
\end{equation}
Here $\mathsf{liveOrdinary}_{\Omega}(o)$ means that $o$ is present in the
ordinary-object component of the appropriate heap in $\Omega$.
The compiler's product is separately admitted by
Equation~\eqref{eq:artifact-validity}; the Simulator's code/PC representation
is separately admitted by the certificate in
Equation~\eqref{eq:simulator-step-correspondence}.  Replacing a component
therefore grants no authority to bypass the fixed core boundary.

Capture one front-end image and run its three stages by ordinary sends:
\begin{equation}\label{eq:metacircular-compilation-pipeline}
\begin{split}
 &\Omega_0,A,\lambda,u\vdash\mathsf{compileUnit}
      \Downarrow\langle\Omega_4,B,P,e,I_L\rangle
 \quad\Longleftrightarrow\\
 &\quad o_L=\mathcal L(\lambda)\ \land\
 \Omega_0,A\vdash\mathsf{call}(o_L,\selsym{frontEndImage},\epsilon)
      \Downarrow\langle\Omega_1,I_L\rangle\ \land\\
 &\quad I_L=\kw{frontEndImage}\langle o_p,o_e,o_c,o_s,G,v_L\rangle\ \land\\
 &\quad \mathsf{frontEndAdmitted}_{\Omega_1}(I_L)\ \land\
 \Omega_1,A\vdash\mathsf{call}(o_p,\selsym{parse:},\langle u\rangle)
      \Downarrow\langle\Omega_2,v_s\rangle\ \land\\
 &\quad \mathsf{decodeSurface}(v_s)=S\ \land\mathsf{idsOK}(S)\ \land\
 \Omega_2,A\vdash\mathsf{call}(o_e,\selsym{elaborate:},\langle v_s\rangle)
      \Downarrow\langle\Omega_3,v_k\rangle\ \land\\
 &\quad \mathsf{decodeCore}(v_k)=\langle P,e\rangle\ \land
       \mathsf{coreOK}(P,e)\ \land\
 \Omega_3,A\vdash\mathsf{call}(o_c,\selsym{compile:},\langle v_k\rangle)
      \Downarrow\langle\Omega_4,v_b\rangle\ \land\\
 &\quad \mathsf{decodeArtifact}(v_b)=B\ \land
       \mathsf{artifactOK}(B,P,e).
\end{split}
\end{equation}
The source string $u$ is represented by the platform's immutable String value;
the notation passes its contents rather than granting the compiler a raw host
pointer.  A language implementation may combine stages internally, but its
observable result must refine Equation~\eqref{eq:metacircular-compilation-pipeline}.

An artifact contains the core image, concrete code images, relocation data,
source/control maps, and the captured language revision:
\begin{equation}\label{eq:artifact-domain}
 B=\langle P,e,\mathcal B,\mathsf{reloc},\mathsf{controlMap},v_L\rangle.
\end{equation}
Its validity is
\begin{equation}\label{eq:artifact-validity}
\begin{split}
 \mathsf{artifactOK}(B,P,e)\Longleftrightarrow{}&
 \mathsf{coreOK}(P,e)\land B.P=P\land B.e=e\land\\
 &\dom(B.\mathsf{controlMap})=\mathsf{pcs}(B.\mathcal B)\land\\
 &\mathsf{targetsOK}(B.\mathcal B,B.\mathsf{reloc})\land
 \mathsf{codeRefinesCore}(B.\mathcal B,P,e).
\end{split}
\end{equation}
The final conjunct of Equation~\eqref{eq:artifact-validity} is the
compiler-correctness obligation.  Concrete execution, observed at the
source/control map, must weakly simulate the sequential and actor semantics.
A changed compiler may accept new surface forms or omit old conveniences by
changing parsing and elaboration, but an artifact admitted by the VM must
still satisfy the fixed core contract.  Changing the instruction set or
primitive meanings requires a VM refinement, not merely a language front-end
change.

The earlier compilation-unit judgment is therefore the projection
\begin{equation}\label{eq:compilation-unit-projection}
 \mathcal L,\Omega,A\vdash\langle\lambda,u\rangle
      \mathsf{compilesTo}\langle P,e,B\rangle
 \quad\Longleftrightarrow\quad
 \exists\Omega',I_L.\;
 \Omega,A,\lambda,u\vdash\mathsf{compileUnit}
      \Downarrow\langle\Omega',B,P,e,I_L\rangle.
\end{equation}

\subsection{Ordinary change, phase behavior, and bootstrap}
\label{sec:metacircular-change}

Let $\mathsf{components}_{\Omega}(\lambda)$ be the value obtained by the first
call in Equation~\eqref{eq:metacircular-compilation-pipeline}.  A successful
existing reflective commit may change methods of $o_L,o_p,o_e,o_c,o_s$, or may
change ordinary slots from which \texttt{frontEndImage} constructs its result.
Its effect on the next compilation is simply
\begin{equation}\label{eq:front-end-reflective-change}
 \Omega\xrightarrow{\kw{commitReflection}(A,\Delta,n+1)}\Omega'
 \ \land\
 \mathsf{components}_{\Omega}(\lambda)\ne
 \mathsf{components}_{\Omega'}(\lambda).
\end{equation}
Equation~\eqref{eq:front-end-reflective-change} is a consequence that an
ordinary reflection transition may have; it is not a new transition rule and
there is no $\mathsf{installParser}$ or $\mathsf{installCompiler}$ command.
Likewise, assigning a new component object through the language object's
ordinary mutable protocol needs no reflection at all when the caller already
has that authority.

The phase behavior follows from existing activation and dispatch rules:
\begin{equation}\label{eq:front-end-phase-behavior}
\begin{aligned}
 \mathsf{capturedImage}(c)&=I_L
   &&\text{for a compilation }c\text{ already past its first call},\\
 \mathsf{currentTerm}_{\Omega'}(a)&=\mathsf{currentTerm}_{\Omega}(a)
   &&\text{for every materialized stage activation }a,\\
 \mathsf{dispatchProgram}(a')&=P'
   &&\text{for a later send dispatched after the commit.}
\end{aligned}
\end{equation}
Thus self-replacement needs no circular premise.  Source defining a replacement
parser is parsed and compiled by the currently executing language objects; an
ordinary commit changes their code or references; subsequent compilation sees
the replacement.  A replacement not expressible in the old accepted language
must enter as an artifact already satisfying Equation~\eqref{eq:artifact-validity}.

Finally, metacircularity has a finite base.  An initial image supplies
\begin{equation}\label{eq:bootstrap-basis}
 \mathsf{bootstrap}=\langle
   \Omega_0,\mathcal L_0,
   \mathsf{decodeSurface},\mathsf{decodeCore},\mathsf{decodeArtifact},
   \mathsf{loadCode},\mathsf{primitiveStep}\rangle.
\end{equation}
After bootstrap, every object reachable through $\mathcal L_0$ is an ordinary
Newspeak object and can be changed by the mechanisms above.  The decoders,
artifact loader, concrete instruction dispatch, allocation, and host I/O remain
VM parameters unless a particular implementation supplies a Newspeak
refinement of them.  This boundary states what is not metacircular without
preventing a more completely bootstrapped VM from discharging those parameters.


\section{Established invariants and mechanization status}

The following are semantic invariants, admissibility conditions for reflective
transactions, and the principal claims tracked by the machine-checked
companion:
\begin{enumerate}
\item every ordinary object has exactly one live class;
\item every base class chain is finite and ends in \Top;
\item each class retains the identity of its applied mixin and enclosing object;
\item each class retains a tagged creation origin and its inducing source
      declaration;
\item each created class has a fresh metaclass, while repeated applications may
      deliberately share an instance-side mixin;
\item each mixin retains its inducing declaration identity;
\item an object's physical layout contains exactly the disjoint root-to-leaf
      slot layouts of its class chain, all initialized to $\nil$ before any
      instance initializer executes;
\item a nested-class cache is per receiver and contains at most one class for
      each inducing declaration identity;
\item equal atomic payloads denote the same canonical literal object;
\item each object-literal evaluation creates a fresh object, class, and
      metaclass but reuses the literal declaration's mixin;
\item a module definition is precisely an instantiable class created by a
      top-level class expression and is a value object;
\item direct method selectors are unique within a mixin after elaboration;
\item a dispatch result determines both the method and the activation's current
      class---method bodies never recompute the latter from the receiver alone;
\item private methods participate only in lexical private selection and
      unrestricted system lookup, never inherited public/protected lookup;
\item failed user dispatch invokes DNU exactly once for the original message;
      failure inside the chosen DNU body is an ordinary new computation;
\item mutable objects and active computations belong to exactly one actor;
      cross-actor arguments contain only values, far references, or promise
      identities;
\item no actor observes a far reference into its own heap, because transfer
      back to the target actor normalizes it to the near referent;
\item every actor executes at most one turn, takes mailbox entries in FIFO
      order, and receives causally related packets in E-order;
\item each promise is settled at most once, by fulfillment or breaking, and
      every dependent eventual send is resumed only after settlement;
\item reflective source-structure changes and all required re-elaboration are
      committed atomically.
\item every reflective read or change is authorized by an unforgeable mirror
      whose target and rights cover that exact operation;
\item a successful reflective transaction advances exactly one VM program
      version, and a failed transaction changes no program, heap, activation,
      continuation, or version component;
\item reflective layout migration preserves retained physical slot identities,
      initializes new slots to $\nil$, and makes removed slots unreachable;
\item every actor in one VM cohort observes the same committed program version,
      while a distinct VM changes only through its own authorized transaction;
\item already materialized control is not rewritten by a code install, whereas
      later dispatch---including dispatch of an in-flight message---uses the
      version current when that dispatch occurs;
\item every running or paused nonempty actor stack has exactly one active
      control image at its top; each lower activation has an explicit suspended
      term and resumable control location;
\item every copied live frame has a fresh activation identity, while retained
      frames preserve identity and stack-structural links internal to a copied
      segment are consistently renamed;
\item a paused-stack mutation succeeds only for its current unforgeable pause
      token, and every successful mutation invalidates that token;
\item stack replacement is atomic, makes discarded activations
      uncontinuable, and exposes no partially materialized frame;
\item full-speed resumption continues the ordinary sequential relation from
      the installed control image rather than introducing a second execution
      semantics.
\item a weak array does not retain its elements; collection atomically replaces
      every cell whose target is reclaimed by $\nil$;
\item a weak map does not retain its keys, while the value of an entry is
      conditionally strong exactly when both the map and its key are strongly
      reachable; a value-to-key cycle cannot retain that key;
\item every retained weak-container reference resolves after collection,
      because entries or cells naming a selected identity are cleared in the
      same safepoint transition.
\item every accepted compilation unit has a successful PEG parse, injective
      source identities, a statically well-formed surface AST, and an
      independently checkable annotated-core image;
\item parser, elaborator, compiler, and Simulator components are ordinary
      Newspeak objects: existing reflection and mutation change them, while
      every executable artifact admitted by the VM refines the fixed core
      semantics.
\end{enumerate}

\subsection{Machine-checked status}

The companion development in \texttt{formal-semantics/lean} formalizes the
domains and transition relations above in Lean.  Its
\texttt{Newspeak.Properties} module exports only established results, and its
\texttt{THEOREMS.md} ledger records every item below as checked.  The
development contains no \texttt{sorry}, \texttt{admit}, or additional axiom.
At the language-semantics boundary it establishes:
\begin{itemize}
\item uniqueness of class chains; determinism of unrestricted, public, and
      protected lookup; equivalence of the executable and named-rule lookup
      formulations; and both the static and dynamic forms of the
      elaborated/ECOOP-style implicit-binding equivalence;
\item determinism and well-formedness preservation of the complete sequential
      relation through evaluation order, DNU, invocation, closures, local and
      non-local return, slots, class and instance construction, initialization,
      object literals, nested classes, exception reification, tuples, patterns,
      and top-level execution;
\item exact finite heap and environment domains, exact duplicate-free
      \kw{allInstancesOf} and mixin-application enumeration, unique physical
      layouts, stable provenance, and the object/module invariants listed
      above;
\item deterministic selected actor turns, isolation and transfer
      normalization, FIFO dequeue, causal delivery constraints, terminal
      promise persistence, at-most-once settlement, and notification only
      after terminal commitment, while leaving actor and network scheduling
      relational;
\item authorization, atomic success/failure, exact version advancement,
      validation preservation, materialized-control stability, debugger stack
      imaging and replacement, full-speed resumption, and commutation of
      pairwise independent reflective commands;
\item exact strong and conditional reachability, stabilization and least
      fixed-point results for weak-map ephemerons, atomic WeakArray and WeakMap
      clearing, post-collection validity, well-formedness preservation, and a
      finite stuttering refinement of the base semantics; and
\item determinism of PEG recognition, projection, stable identification, and
      elaboration; soundness of the finite grammar certificate; complete
      derivation and validation of the annotated core; and finite-execution
      refinement from an admitted compiled artifact to annotated-core
      execution.
\end{itemize}

\paragraph{Remaining implementation verification.}\leavevmode\par
Equation~\eqref{eq:simulator-step-correspondence} is proved generically from a
mandatory compiler/Simulator admission certificate.  The currently shipped
V5 compiler and Simulator must still be verified against that boundary:
concrete bytecode artifacts, PCs, operand stacks, and frame objects must supply
the certificate and representation relation.  A replacement implementation
has the same obligation.  This is not a missing rule of the language
semantics; it is the implementation-refinement project that instantiates the
proved abstract theorem.  A useful further strengthening, also outside the
language definition, is a provenance theorem tracing every dequeued wake
packet to its unique earlier promise-settlement event.

\bibliographystyle{alpha}
\bibliography{references}

\end{document}
